% This document is compiled using pdfLaTeX
% You can switch XeLaTeX/pdfLaTeX/LaTeX/LuaLaTeX in Settings

\documentclass{article}
\usepackage{ctex}
\usepackage{amsmath} % 确保引入了此宏包以支持 aligned

\title{因子日历2026}

\date{\today}

\begin{document}

\maketitle

\section{资本利得突出量(行为金融因子-处置效应)}

\subsection{因子定义}
使用过去一段时间的换手率加权价格作为参考价格，用于衡量股票的综合持仓成本，进而衡量大部分投资者在某只股票上是盈利还是亏损：
\begin{equation}
    \mathrm{RP}_t = \frac{1}{k} \sum_{n=0}^{T-1} \left( V_{t-n} \prod_{s=1}^{n-1} (1 - V_{t-n+s}) \right) P_{t-n}
\end{equation}
计算第 \( t \) 日股票 \( i \) 的持仓者相对参考价格数学 \( \mathrm{RP}_i \) 的平均盈亏情况即为资本利得突出量：
\begin{equation}
    \mathrm{CGO}_{i,t} = \frac{\mathrm{Close}_{i,t} - \mathrm{RP}_t}{\mathrm{Close}_{i,t}}
\end{equation}

\subsection{变量定义}
\begin{itemize}
    \item \( V_t \)、\( \mathrm{Close} \)、\( P_t \) 分别为为 \( t \) 日的换手率、收盘价、VWAP 均价；
    \item \( k \) 为权重系数，用于保证计算 \( \mathrm{RP}_t \) 时的权重和等于 1，即为 \( P_t \) 前面权重系数之和。
\end{itemize}

\subsection{说明}
处置效应指投资者“售盈持亏”的情感倾向；CGO 因子值越大，浮盈越大；历史股价上涨，高 CGO 意味着大量投资者盈利，抛压大，卖盘压力增加，股价上涨受阻；历史股价下跌，低或负 CGO 意味着投资者普遍被套，抛售意愿弱，卖压小，一旦有利好，股价更容易反弹；CGO 因子对应反转效应，与未来收益负相关。

\subsection{参考文献}
\begin{enumerate}
    \item 陈开锐, 姚紫薇. 2025, 处置效应与V型处置效应对在量化选股中的应用. 中信建投.
    \item Grinblatt, M., \& Han, B. (2005). \textit{Prospect theory, mental accounting, and the disposition effect}. Journal of Financial Economics, 78(2), 311–339.
\end{enumerate}

\section{残差反转}

\subsection{因子定义（高频因子-动量反转类）}

计算小单资金流强度：分子为（小单买额-小单卖额）之和，分母为其绝对值；计算大单资金流强度：分子为（小单买额-小单卖额）之和。
\begin{equation}
    \mathrm{S}_t = \frac{\sum_{t-T}^t (\mathrm{buy}_t - \mathrm{sell}_t)}{\sum_{t-T}^t |\mathrm{buy}_t - \mathrm{sell}_t|}
\end{equation}
将反转因子（过去20日涨跌幅）与大单资金强度和小单资金强度各自做回归，得到的残差即为得到的大单和小单的残差反转因子：
\begin{equation}
    \mathrm{Ret20}_t = a + b * \mathrm{S}_t + \varepsilon_t
\end{equation}

\subsection{说明}
残差反转因子剥离了资金流强度对传统反转因子的影响，选股效果得到进一步提升。

\subsection{参考文献}
\begin{enumerate}
    \item 张建超, 高鹏. 2021. 大单与小单资金流的 alpha 能力. 开源证券.
\end{enumerate}

\section{个人投资者交易热度（行为金融因子-投资者注意力）}

\subsection{因子定义}

\begin{equation}
    \mathrm{TURN\_RETAIL} = \frac{1}{m} \sum_{t=1}^{t-m+1} \frac{\mathrm{INFLOW\_RETAIL}_{i,t} + \mathrm{OUTFLOW\_RETAIL}_{i,t}}{\mathrm{CAP}_{i,t}}
\end{equation}

\subsection{变量定义}
\begin{itemize}
    \item ① \(\mathrm{INFLOW\_RETAIL}_{i,t}\) 和 \(\mathrm{OUTFLOW\_RETAIL}_{i,t}\) 分别为股票 \(i\) 在交易日 \(t\) 的个人投资者买入金额和卖出金额；
    \item ② 将单笔成交额低于 4 万元的投资者作为个人投资者；
    \item ③ \(\mathrm{CAP}_{i,t}\) 为股票 \(i\) 在交易日 \(t\) 的流通市值；
    \item ④ \(m\) 为交易日窗口，一般取为月度。
\end{itemize}

\subsection{说明}
个人投资者交易热度属于股票交易活跃度类因子，通过衡量股票交易活跃度来反映投资者关注度。股票交易越活跃，投资者对该股票的关注度越高；而投资者有限注意力会推动关注度高的股票更易出现非理性净买入溢价，进而导致股价后续反转。

\subsection{参考文献}
\begin{enumerate}
    \item 陈升锐. 2024. 投资者有限关注及注意力捕捉与溢出. 中信建投.
\end{enumerate}

\section{换手率分布均匀度(高频因子-成交分布类)}

\subsection{因子定义}


每月月底，回溯所有股票过去20个交易日，计算当日分钟换手率的标准差：
\begin{equation}
    \mathrm{TurnVolDaily} = \mathrm{std}(\text{分钟换手率})
\end{equation}
其中，
\begin{equation}
    \text{分钟换手率} = \frac{\text{分钟成交量}}{\text{当日流通股本}}
\end{equation}
基于这20个交易日的分钟换手率序列，计算标准差和均值，然后进行标准化和市值中性化处理得到换手率分布均匀度 \(\mathrm{UTD}\)：
\begin{equation}
    \mathrm{UTD} = \frac{\mathrm{std}(\mathrm{TurnVolDaily})}{\mathrm{mean}(\mathrm{TurnVolDaily})}
\end{equation}

\subsection{说明}
换手率分布均匀度因子衡量了股票换手率在日内不同交易时段的分布均匀程度，以及该均匀程度在日间是否平稳，股票换手率在日内及日间越平稳，即因子值越小越好。

\subsection{参考文献}
\begin{enumerate}
    \item 高子剑. 2021. 换手率分布均匀度：基于分钟成交量的选股因子. 东吴证券.
\end{enumerate}

\section{筹码分布成本重心(行为金融因子-筹码分布)}

\subsection{因子定义}

\begin{equation}
    \mathrm{CKDP}_t = \frac{\mathrm{mean}_t - \mathrm{Chip\_Lowest}_t}{\mathrm{Chip\_Highest}_t - \mathrm{Chip\_Lowest}_t}
\end{equation}

\subsection{变量定义}
\begin{itemize}
    \item ① 假设某只股票在 \( t \) 日的筹码分布共有 \( n \) 个价位（\( i = 1, 2, ..., n \)），\( \mathrm{vol}_{i,t} \) 为第 \( t \) 日 \( i \) 价位上的筹码峰高度，\( \mathrm{price}_{i,t} \) 为该筹码峰对应的价格，\( \mathrm{Chip\_Highest}_t \)、\( \mathrm{Chip\_Lowest}_t \) 分别对应筹码最高价和筹码最低价；
    \item ② \(\displaystyle \mathrm{mean}_t = \sum (\mathrm{price}_{i,t} * p_{i,t})\) 为筹码分布加权平均成本。
\end{itemize}

\subsection{说明}
筹码分布成本重心衡量了筹码平均成本在筹码分布中的相对位置，可用于识别建仓或出货阶段。成本重心接近 0，说明平均成本靠近成本区间下沿位置，成本偏低，可能处于低位建仓阶段；成本重心接近 1，说明平均成本靠近区间上沿位置，成本偏高，可能处于高位派发阶段；成本重心接近 0.5，成本分布对称，市场处于平衡状态。经测试，筹码分布成本宽带与未来收益负相关。

\subsection{参考文献}
\begin{enumerate}
    \item 陈升锐, 姚紫薇. 2025. 筹码分布因子系统构建. 中信建投.
\end{enumerate}

\section{下行跳跃波动率(高级因子-波动跳跃类)}

\subsection{因子定义}

\begin{equation}
    \mathrm{RVJN}_t = \max\left(\mathrm{RS}_t^- - \frac{1}{2} \hat{IV}_t, 0\right)
\end{equation}
\begin{equation}
    \mathrm{RS}_t^- = \sum_{i=1}^n r_{t,i}^2 \mathrm{I}_{r_{t,i} > 0} \rightarrow_u \frac{1}{2} \int_0^t \sigma_s^2 ds + \sum_{s \leq t} \Delta X_s^2 \mathrm{I}_{\Delta X_s < 0}
\end{equation}
\begin{equation}
    \hat{IV}_t = \mu^{-3}_{\frac{2}{3}} \sum_{i=k}^n |r_{t_i}|^{\frac{2}{3}} |r_{t_{i-1}}|^{\frac{2}{3}} \cdots |r_{t_{i-k+1}}|^{\frac{2}{3}}
\end{equation}

\subsection{变量定义}
\begin{itemize}
    \item ① \( r_{t_i} \) 为某股票交易日 \( t \) 内的分钟对数收益率序列，如 5 分钟对数收益率序列；
    \item ② \( k=3 \)，\( n \) 表示该股票在交易日 \( t \) 中特定频率的分钟级别数据个数，如 5 分钟频率通常有 48 个数据点；
    \item ③ \( \mu_q \) 为标准正态分布的第q阶矩对短；
    \item ④ 下行跳跃波动率 \(\mathrm{RVJN}\) 整体取值为正，所以需要用 \(\max\) 对负差值做修正。
\end{itemize}

\subsection{说明}
按收益涨跌情况，股价波动可分为上行波动和下行波动，跳跃波动也可据此做划分。下行跳跃波动就是下行波动中的跳跃部分；短期内，下行跳跃波动通常有正的风险溢价。

\subsection{参考文献}
\begin{enumerate}
    \item Bo Yu, Bruce Mizrach, Norman R. Swanson. (2020). \textit{New Evidence of the Marginal Predictive Content of Small and Large Jumps in the Cross-Section}. Econometrics, MDPI, 8(2), 1-52.
\end{enumerate}

\section{下游传导动量(图谱网络-动量溢出)}

\subsection{因子定义}
计算上游公司 A 与下游公司 B 间每一对上下游产品 \(i \to j\) 之间的关联度：
\begin{equation}
    \mathrm{wgt}_{A,B,t} = \frac{\sum_i \mathrm{prdIncome}_{A,i,t} \times \mathbf{1}\{i \to j, i \in A\} + \sum_j \mathrm{prdIncome}_{B,j,t} \times \mathbf{1}\{i \to j, j \in B\}}{\sum_i \mathrm{prdIncome}_{A,i,t} + \sum_j \mathrm{prdIncome}_{B,j,t}}
\end{equation}
对上游公司 A，以产业链关联度为权重加权其所有下游公司的月度收益率，即为下游传导动量因子：
\begin{equation}
    \mathrm{upstream\_Moment}_{A,t} = \frac{\sum_{A \neq B} \mathrm{wgt}_{A,B,t} \times \mathrm{rtn}_{B,t}}{\sum_{A \neq B} \mathrm{wgt}_{A,B,t}}
\end{equation}

\subsection{变量定义}
\begin{itemize}
    \item ① \(\mathrm{prdIncome}\) 为公司 A 在产品 i 上的营业收入；
    \item ② \(\sum_i \mathrm{prdIncome}_{A,i,t}\) 为公司 A 的当期总营收；
    \item ③ \(\mathbf{1}\{i \to j\}\) 表示产品 i 为产品 j 的上游；
    \item ④ \(\mathrm{rtn}_{B,t}\) 为公司 B 在 t 时刻的月度收益率；
    \item ⑤ 可以将下游传导动量因子对月度收益率做正交化，剥离月度反转效应。
\end{itemize}

\subsection{说明}
下游传导动量因子值越大，表示公司 A 作为上游企业时，位于其产业链下游的高关联度高集群公司在考察月份涨幅相对较高，公司 A 接下来也会有上涨的可能。

\subsection{参考文献}
\begin{enumerate}
    \item 郑兆磊.2023.产业链视角下的Alpha传导研究.兴业证券
\end{enumerate}

\section{超大单和小单的同步相关性（高频因子-资金流类）}

\subsection{因子定义}

\begin{equation}
    \mathrm{RankCorr}(\mathrm{EL}_t, \mathrm{S}_t)
\end{equation}

\subsection{变量定义}
\begin{itemize}
    \item ① \( \mathrm{EL}_t \) 为超大单（>100万元）过去20个交易日的资金净流入序列；
    \item ② \( \mathrm{S}_t \) 为小单（<4万元）过去20个交易日同期的资金净流入序列。
\end{itemize}

\subsection{说明}
在资金流同步相关性因子中，上述超大单和小单的同步相关性因子表现最好；但资金流同步相关性因子是“伪动力学”因子，在恒等式“超大单+大单+中单+小单=0”的约束下，该因子的 alpha，可以被两者流入量级所解释，其底层 alpha 来源仍是资金流强度。

\subsection{参考文献}
\begin{enumerate}
    \item 魏建榕, 盛少成. 2022. 新型因子：资金流动力学与散户羊群效应. 开源证券.
\end{enumerate}

\section{买入浮盈系数修正后的积极灵活筹码（行为金融因子-筹码分布）}

\subsection{因子定义}
计算筹码当日收盘价相对筹码买入价格的涨跌幅作为该筹码的买入浮盈系数：
\begin{equation}
    a_i = \frac{\mathrm{Close}_t}{\mathrm{TradePrice}_i}
\end{equation}
将买入浮盈系数与筹码买入量相乘即可得到修正后的积极灵活筹码值：
\begin{equation}
    \mathrm{AFH\_Close} = \frac{\sum_{i=1}^N a_i * \mathrm{Vol\_AFH}_i}{\sum_{t=T-20}^T \mathrm{Volume}_t}
\end{equation}

\subsection{变量定义}
\begin{itemize}
    \item ① \(i\) 为筹码购入时的订单号，统计对象为过去 20 日的订单号；
    \item ② \( \mathrm{Vol\_AFH}_i \) 为主动买入时订单金额小于 5 万元的订单成交量；
    \item ③ \( \mathrm{TradePrice}_i \) 为该订单的买入价格；
    \item ④ \( \mathrm{Close}_t \)、\( \mathrm{Volume}_t \) 分别为订单买入当日的开盘价和当日总成交量。
\end{itemize}

\subsection{说明}
投资者会由自我归因偏差引起过度自信，这将加强投资者的处置效应。当投资者买入相应个股之后，若当日即获得浮盈，将使投资者形成“买对了”的心理行为，这将提升投资者对股票判断的自信，此时的买入筹码也更具处置效应，因此可以计算筹码当日收盘价相对筹码买入价格的涨跌幅作为该筹码的买入浮盈系数，用于修正原始的积极灵活筹码，提升因子表现。

\subsection{参考文献}
\begin{enumerate}
    \item 任瞳, 许继宏. 2024. 基于逐笔数据的 AFH 改进因子. 招商证券.
\end{enumerate}

\section{非流动性（高频因子-流动性类）}

\subsection{因子定义}

不考虑日内价格变动方向，直接进行无向叠加，求得股价变动的最短路径：
\begin{equation}
    \mathrm{shortcut} = 2 * (\mathrm{high} - \mathrm{low}) - |\mathrm{close} - \mathrm{open}|
\end{equation}
利用高频分钟 K 线，每日的 K 线最短路径非流动性因子：
\begin{equation}
    \mathrm{ILLIQ}_t = \frac{1}{d} \sum_{i=1}^d \left( \sum_{j=1}^p \frac{\mathrm{ShortCut}_j}{\mathrm{Value}_j} \right)_{t-i}
\end{equation}

\subsection{变量定义}
\begin{itemize}
    \item ① \(\mathrm{ShortCut}_j\) 表示日内频率下第 \(j\) 根 K 线的最短路径；
    \item ② \(\mathrm{Value}_j\) 表示日内频率下第 \(j\) 根 K 线内的成交额；
    \item ③ \(p\) 表示日内频率分段个数；
    \item ④ \(d\) 表示日间移动平均的周期参数。
\end{itemize}

\subsection{说明}
通过使用更高频的 K 线数据提升 K 线最短路径对股票交易时市场冲击的代理精度，大大提升了 K 线最短路径非流动性因子的有效性与预测能力。同时随着数据频率的提高，K 线最短路径定义对经典定义下的非流动性因子提升程度愈加显著。比起经典的经典定义方式，K 线最短路径能更充分地利用高频数据新引入的信息。

\subsection{参考文献}
\begin{enumerate}
    \item 刘均伟. 2017. 基于 K 线最短路径构造的非流动性因子. 光大证券.
\end{enumerate}

\section{基于PB的前景价值TK（行为金融因子-前景理论）}

\subsection{因子定义}

① 以个股过去N期的估值变化率为衡量公司价值的代理变量：
\begin{equation}
    (r_{pb,-m},\frac{1}{N}; r_{pb,-1},\frac{1}{N}; r_{pb,0},\frac{1}{N}}; r_{pb,1},\frac{1}{N}; \ldots; r_{pb,m},\frac{1}{N})
\end{equation}

② 利用价值函数、权重函数计算前景价值 TK：
\begin{equation}
    \mathrm{TK} = \sum_{i=-m}^{-1} v(r_{pb,i}) \left[ w^-\left(\frac{i + m + 1}{N}\right) - w^-\left(\frac{i + m}{N}\right) \right] + \sum_{i=1}^{n} v(r_{pb,i}) \left[ w^+\left(\frac{n - i + 1}{N}\right) - w^+\left(\frac{n - i}{N}\right) \right]
\end{equation}

③ 价值函数、权重函数都保持不变，函数参数取值为：\(\alpha=0.88\)，\(\lambda=2.25\)；\(\gamma=0.61\)，\(\delta=0.69\)。

\subsection{说明}
前景价值代表了投资者给每支个股价值的“打分”，而用于衡量公司价值的参考指标不一定为历史收益率，此处选用了PB估值变化率来衡量公司价值：投资者更倾向于持有低估的股票，卖出高估的股票，即高价值量的股票（高估值股票）由于卖出压力反而会被低估，而低价值量的股票（低估股票）反而会被高估，此时的前景价值与股票未来收益呈正相关。

\subsection{参考文献}
\begin{enumerate}
    \item 丁鲁明, 胡一江. 2020. 前景理论因子的选股能力. 中信建投.
\end{enumerate}

\section{加权收盘价比因子(高频因子-量价相关性类)}

\subsection{因子定义}

\begin{equation}
    \text{加权收盘价比} = \frac{\sum_{t=1}^T \frac{\mathrm{vol}_t}{VOL}\mathrm{close}_t }{\frac{1}{T} \sum_{t=1}^T \mathrm{close}_t}
\end{equation}

\subsection{变量定义}
\begin{itemize}
    \item ① \(\mathrm{close}_t\) 和 \(\mathrm{vol}_t\) 分别为 \(t\) 时间段的收盘价和成交量；
    \item ② \(\mathrm{Vol} = \sum^{T}_{t=1} \mathrm{vol}_t\) 为整个时间段总成交量；
    \item ③ \(T\) 为划分时间段的个数，如日内5分钟频率的 \(T=48\)。
\end{itemize}

\subsection{说明}
加权收盘价比因子为成交量加权的收盘价求和同收盘价等权求和的比值，刻画了个股在高位成交密集水平；与未来收益负相关，因子值越高，个股在高位成交越密集，因交易中羊群效应的影响，推动股价高估，未来收益会回调。

\subsection{参考文献}
\begin{enumerate}
    \item 川桃, 郑起. 2020. 高频因子（八）：高位成交因子——从量价匹配说起. 长江证券.
\end{enumerate}

\section{行业动量溢出效应(图谱网络-动量溢出)}

\subsection{因子定义}
\begin{equation}
    \mathrm{RET}^{IND}_{it} =  \mathrm{Industrial\_weight}_{ij,t} \times \mathrm{ReT}_{jt}
\end{equation}
\begin{equation}
    \mathrm{Industrial\_weight}_{ij,t} = \frac{\mathrm{IMV}_{j,t}}{\sum_{j=1 and j \neq i}^N \mathrm{IMV}_{j,t}}
\end{equation}

\subsection{变量定义}
\begin{itemize}
    \item ① \(\mathrm{Ret}_{jt}\) 为股票 \(j\) 在 \(t\) 时刻的月度收益率；
    \item ② \(\mathrm{IMV}_{j,t}\) 为 \(t\) 时期同行业内公司 \(j\) 的市值。
\end{itemize}

\subsection{说明}
行业动量溢出效应因子刻画了与单只股票同属于相同行业下其余股票涨跌对该股票涨跌的联动效应（领先滞后效应）。

\subsection{参考文献}
\begin{enumerate}
    \item Usman Ali, David Hirshleifer. (2020). \textit{Shared analyst coverage: Unifying momentum spillover effects}. Journal of Financial Economics, 6(136), 649-675.
\end{enumerate}

\section{成交量占比峰度(高频因子-成交分布类)}

\subsection{因子定义}

\begin{equation}
    \mathrm{kurt}\left( \frac{\mathrm{volume}_i}{\sum \mathrm{volume}_i} \right)
\end{equation}

\subsection{变量定义}
\begin{itemize}
    \item ① \(\mathrm{volume}_i\) 为日内分钟成交量，如10分钟等；成交量占比是成交量在个股内纵向可比指标，主要从成交的分布上给出衡量；
    \item ② 此外也可将成交量占比替换为对数成交量、换手率（全市场个股间横向可比指标，从流动性上给出衡量），计算成交量峰度类因子。
\end{itemize}

\subsection{说明}
成交量峰度类因子均衡量了个股过去一段时间的交易极端行为，与未来收益负相关；个股成交呈尖峰时，大部分时间成交在合理范围波动，但存在特定区间交易异常活跃或几乎无成交的极端情况，此时往往多空博弈激烈，个股收益不确定性强。

\subsection{参考文献}
\begin{enumerate}
    \item 川桃, 郑起. 2019. 高频因子（四）：高阶矩高频因子. 长江证券.
\end{enumerate}

\section{空间网络相对中心度(图谱网络-网络结构)}

\subsection{因子定义}
计算股票 \( i \) 和样本股票池中其他股票 \( j \) 的近 20 个交易日收益率序列之间的 Pearson 相关系数并求均值：
\begin{equation}
    \bar{p}_{i,t} = \frac{1}{N-1} \sum_{j=1, j \neq i}^{N-1} \mathrm{corr}(r_{i,t}, r_{j,t})
\end{equation}
基于相关系数计算股票 \( i \) 和其他股票间的平均距离：
\begin{equation}
    \bar{d}_{i,t} = \sqrt{2 \times (1 - \bar{p}_{i,t})}
\end{equation}
计算空间维度的网络相对中心度：
\begin{equation}
    \mathrm{SCC}_{i,t} = \frac{1}{d_{i,t}^2}
\end{equation}

\subsection{说明}
SCC 因子衡量了股票在网络中与其余股票的相对距离（关联程度），因子值越高，说明股票在网络中与其余股票的距离越近，越处于网络相对中心的位置，对其他股票的相对影响程度也越高。

\subsection{参考文献}
\begin{enumerate}
    \item 张自力, 闫红蕾, 张楠. 2020. 股票网络、系统性风险与股票定价. 经济学, 1, 329-350.
    \item 曹春晓, 杨国平. 2021. 股票网络与网络中心度因子研究. 华西证券.
\end{enumerate}

\section{理想振幅因子(高频因子-波动跳跃类)}

\subsection{因子定义}


对单只股票，回溯取其最近 \(N\) 个交易日（默认 \(N=20\)）的数据，计算每日的振幅（最高价/最低价-1）；

选择收盘价较高的 \(\lambda\)（默认 25\%）有效交易日，计算振幅均值得到高价振幅因子 \(\mathrm{V}_{\mathrm{high}}(\lambda)\)；选择收盘价较低的 \(\lambda\)（默认 25\%）有效交易日，计算振幅均值得到低价振幅因子 \(\mathrm{V}_{\mathrm{low}}(\lambda)\)；

计算理想振幅因子：
\begin{equation}
    V(\lambda) = \mathrm{V}_{\mathrm{high}}(\lambda) - \mathrm{V}_{\mathrm{low}}(\lambda)
\end{equation}

\subsection{说明}
基于股价将振幅因子进行切割得到的理想振幅因子，考虑了不同价格区间的振幅分布信息差异（高价振幅因子具有更强的负向选股能力），选股能力更强。

\subsection{参考文献}
\begin{enumerate}
    \item 魏建榕. 2020. 振幅因子的隐藏结构. 开源证券.
\end{enumerate}

\section{“随波逐流”因子（高级因子-量价相关性类）}

\subsection{因子定义}

1. \textbf{对股票 \(i\)，计算第 \(t\) 日“合理收益率”：}
   \begin{equation}
       \text{合理收益率} = \frac{1}{20} \sum_{n=0}^{19} \text{日内收益率}_{t-n} = \frac{1}{20} \sum_{n=0}^{19} \left( \frac{\text{收盘价}_{t-n}}{\text{开盘价}_{t-n}} - 1 \right)
   \end{equation}

2. \textbf{对股票 \(i\)，取第 \(t\) 日内 1 分钟价格序列（剔除集合竞价），计算该日第 \(m\) 分钟“相对开盘收益率”：}
   \begin{equation}
       \text{相对开盘收益率}_{m,t} = \frac{\text{收盘价}_{m,t}}{\text{开盘价}_t} - 1
   \end{equation}

3. \textbf{计算第 \(t\) 日的“高位成交额”和“低位成交额”：}
   \begin{equation}
       \text{高位成交额}_t = \sum_{m=1}^{240} \text{分钟成交额}_m \cdot \mathrm{I}\left( \text{相对开盘收益率}_{m,t} > \text{合理收益率}_t \right)
   \end{equation}
   \begin{equation}
       \text{低位成交额}_t = \sum_{m=1}^{240} \text{分钟成交额}_m \cdot \mathrm{I}\left( \text{相对开盘收益率}_{m,t} < \text{合理收益率}_t \right)
   \end{equation}

4. \textbf{计算股票 \(i\) 在 \(t\) 日价格处于高位时的相对成交额——“高低额差”：}
   \begin{equation}
       \text{高低额差}_t = \frac{\text{高位成交额}_t - \text{低位成交额}_t}{\text{收盘时流通市值}_t}
   \end{equation}

5. \textbf{对股票 \(i\)，每月末，取所有股票过去 20 日的“高低额差”序列，计算股票 \(i\) 与其余所有股票“高低额差”序列间的 spearman 相关系数，对相关系数取绝对值并求均值，即得到股票 \(i\) 的“随波逐流”因子。}

\subsection{说明}
“随波逐流”因子衡量了个股股价处于相对高位时，其成交额与市场趋势的关联性，与未来收益正相关；当个股股价处于相对高位时，相关系数越大的股票，其成交额更多由市场趋势带动，投资者对其股价的分歧度越小，预示个股未来收益会相对更高。

\subsection{参考文献}
\begin{enumerate}
    \item 曹春晓. 个股成交额的市场跟随性与“水中行舟”因子, 2023, 方正证券.
\end{enumerate}

\section{产业链地位优势因子（图谱网络-网络结构）}

\subsection{因子定义}
方法1：按单一维度方式计算公司 B 与公司 A 间每一对上下游产品之间的关联度，并取最大值作为两个公司最终的关联度：
\begin{equation}
    \mathrm{wgt}_{(B,j),(A,i),t} = \frac{(\mathrm{prdIncome}_{B,j,t} + \mathrm{prdIncome}_{A,i,t}) \times \mathbf{1}\{j \to i\}}{\sum_i \mathrm{prdIncome}_{A,i,t} + \sum_j \mathrm{prdIncome}_{B,j,t}}
\end{equation}

方法2：按整体维度方式计算公司 B 与公司 A 间每一对上下游产品之间的关联度：
\begin{equation}
    \mathrm{wgt}_{A,B,t} = \frac{\sum_i \mathrm{prdIncome}_{A,i,t} \times \mathbf{1}\{i \to j, i \in A\} + \sum_j \mathrm{prdIncome}_{B,j,t} \times \mathbf{1}\{i \to j, j \in B\}}{\sum_i \mathrm{prdIncome}_{A,i,t} + \sum_j \mathrm{prdIncome}_{B,j,t}}
\end{equation}

\subsection{变量定义}
\begin{itemize}
    \item ① \( \mathrm{prdIncome}_{A,i,t} \) 为公司 A 在产品 i 上的营业收入；
    \item ② \( \sum_i \mathrm{prdIncome}_{A,i,t} \) 为公司 A 的当期总营收；
    \item ③ \( \mathbf{1}\{i \to j\} \) 表示产品 j 为产品 i 的上游；
    \item ④ 若公司 A 为上游公司，统计其与下游公司的关联度，得到的为下游关联优势地位因子；反之，为上游关联优势地位因子；若同时考虑公司 A 的上游和下游关联度，则为上下游关联度地位优势因子；
    \item ⑤ 可以将产业链地位优势因子进行市值行业中心化处理。
\end{itemize}

\subsection{说明}
产业链地位优势因子刻画了上市公司在产业链中的相对优势；产业链地位优势因子值越高，越处于产业链的中心地位，而位于中心行业的股票越容易受到其他行业冲击的影响，市场风险更大，进而能获得更高的股票收益。

\subsection{参考文献}
\begin{enumerate}
    \item 郑兆磊.2023.产业链视角下的Alpha传导研究.兴业证券
\end{enumerate}

\section{高频已实现偏度（高频因子-收益分布类）}

\subsection{因子定义}
\begin{equation}
    \mathrm{RSkew}_i = \frac{\sqrt{N} \sum_{j=1}^N (r_{ij} - \bar{r}_i)^3}{\mathrm{RVar}_i^{3/2}}
\end{equation}
\begin{equation}
    \mathrm{RVar}_i = \sum_{j=1}^N (r_{ij} - \bar{r}_i)^2
\end{equation}

\subsection{变量定义}
\begin{itemize}
    \item ① \( r_{ij} \) 为股票 \( i \) 的 1 分钟对数收益率序列或 5 分钟对数收益率序列；
    \item ② \( N \) 表示个股在交易日 \( t \) 中特定频率的分钟级别数据个数；
    \item ③ 在任意选股时刻，因子值为前几日指标的均值，如月度选股，计算因子值时使用的是股票过去 20 个交易日的均值。
\end{itemize}

\subsection{说明}
高频已实现偏度因子刻画了股价日内快速拉升或下跌的特征，与股票未来收益负相关；短期大幅拉升的股票，未来更有可能出现反转；而日内快速下跌的股票，往往具有更高的下行风险溢价。

\subsection{参考文献}
\begin{enumerate}
    \item 冯佳睿, 袁林青. 2017. 高频因子之股票收益分布特征. 海通证券.
    \item Amaya, D., Christoffersen, P., Jacobs, K., and Vasquez, A. (2016). \textit{Does Realized Skewness Predict the Cross-Section of Equity Returns?}. Journal of Financial Economics, 118, 135–167.
\end{enumerate}

\section{尾盘半小时收益率（高频因子-动量反转类）}

\subsection{因子定义}
\begin{equation}
    \mathrm{ret}_{14:30-15:00,i} = \frac{\mathrm{close}_{15:00,i}}{\mathrm{close}_{14:30,i}} - 1
\end{equation}

\subsection{变量定义}
\begin{itemize}
    \item ① \( \mathrm{close}_{14:30,i} \) 和 \( \mathrm{close}_{15:00,i} \) 分别为尾盘 14:30 分钟收盘价和当日收盘价；
    \item ② 基于相同的逻辑，也可以用开盘价和 10:00 的收盘价计算开盘半小时收益率。
\end{itemize}

\subsection{说明}
通常股票在开盘后半小时和收盘前半小时的成交相对更活跃，多空博弈激烈，蕴含的信息相对更多，所以可以单独统计这些时间段的量价因子。

\subsection{参考文献}
\begin{enumerate}
    \item 艾玛钧, 安宁宁, 罗军. 2021. 高频量化数据因子化方法. 广发证券.
\end{enumerate}

\section{相似反转（量价因子改进）}

\subsection{因子定义}
1. 滚动计算历史收盘序列（长度为 120 个交易日）中任意时刻 \( T \) 起连续 \( \mathrm{RW}=6 \) 日序列同当前时刻 \( t \) 的股票收盘价序列的皮尔逊相关系数：
\begin{equation}
    \mathrm{Correlation}_{i,T} = \rho_{\mathrm{Sequence}_{i,t}, \mathrm{Temp  History}_{i,T}}
\end{equation}
\begin{equation}
    \mathrm{Sequence}_{i},t = [\mathrm{Close}_{i,t-\mathrm{RW}+1}, \mathrm{Close}_{i,t-\mathrm{RW}+2}, \dots, \mathrm{Close}_{i,t}]
\end{equation}
\begin{equation}
    \mathrm{Temp  History}_{i,T} = [\mathrm{Close}_{i,T}, \mathrm{Close}_{i,T+1}, \dots, \mathrm{Close}_{i,T+\mathrm{RW}-1}]
\end{equation}

2. 设定相关系数的阈值 \( \mathrm{Threshold}=0.4 \) 筛选出走势相似度较高的时刻：
\begin{equation}
    \{ \tau \mid |\mathrm{Correlation}_{i,\tau}| \geq \mathrm{Threshold} \}
\end{equation}

3. 通过指数衰减法，计算历史出现相似走势后既定持仓时间内股票累计超额收益率的加权平均值即为相似反转因子：
\begin{equation}
    \mathrm{SimilarReverse}_{i,t} = -\sum_{k=1}^n \omega_{i,k} * \mathrm{ER}_{i,\tau_k}
\end{equation}
\begin{equation}
    \mathrm{ER}_{i,\tau} = \prod_{k=\tau+\mathrm{RW}}^{\tau+\mathrm{RW}+\mathrm{HoldingTime}-1} (1 + \mathrm{ER}_{i,k}) - 1, \quad 1 \leq \mathrm{HoldingTime} \leq \mathrm{RW}
\end{equation}
\begin{equation}
    \omega_{i,k} = \frac{e^{-\lambda k}}{\sum_{k=1}^n e^{-\lambda k}}, \quad \lambda = \frac{\ln(2)}{H}
\end{equation}

\subsection{说明}
相似反转因子即为多次相似走势后在既定持仓时间内股票累计超额收益的平均值，历史中出现相似走势后股票一段时间内的平均超额收益率越高，未来持有同样时间的超额收益率就越低。

\subsection{参考文献}
\begin{enumerate}
    \item 郑琳琳, 王天业. 2024. 基于历史相似走势的因子选股研究. 西南证券.
\end{enumerate}

\section{价格弹性（高频因子-流动性类）}

\subsection{因子定义}
\begin{equation}
    \mathrm{resiliency} = \frac{\mathrm{high} - \mathrm{low}}{\mathrm{turnover}}
\end{equation}

\subsection{变量定义}
\begin{itemize}
    \item ① \(\mathrm{high}\)、\(\mathrm{low}\)、\(\mathrm{turnover}\) 分别 tick 数据中的最高价、最低价和成交额；
    \item ② 可对上述因子日内序列求均值得到日度因子值，并计算近期的变化趋势：（20 个交易日指数移动平均 - 120 个交易日移动平均）/ 120 个交易日的标准差。
\end{itemize}

\subsection{说明}
价格弹性反映了单位成交额下股价的波动幅度，与未来收益正相关，弹性越大说明单位成交额对价格冲击越大，市场流动性越弱，而流动性较低的个股未来收益表现相对较好。

\subsection{参考文献}
\begin{enumerate}
    \item 胡骥驼等. 2021. 高频视角下的微观流动性与波动性. 中金公司.
\end{enumerate}

\section{主力波动率（高频因子-波动跳跃类）}

\subsection{因子定义}

1. 按照日内分钟频收益率方向变化的时刻（用前一分钟收盘价与本分钟收盘价判断），将日内收益率序列合并为上涨或下跌趋势段 \([t_i, t_j]\)；

2. 对每一分钟判断当前时刻成交量 \(V_1\) 与上一个趋势段的最后一分钟成交量的相对大小 \(V_0\)：若 \(V_0 < V_1\)，则为成交量放量；\(V_0 \geq V_1\)，则为成交量缩量；结合收益率涨跌情况，可以将日内每分钟行情划分为：放量上涨、缩量上涨、放量下跌、缩量下跌；

3. 对每个趋势段 \([t_i, t_j]\) 内的每一分钟 \(t_k (i < k \leq j)\)，比较其对应成交量 \(V_{t_k}\) 与区间 \([t_i, t_j]\) 内最大成交量 \(V_{\max}\) 和最小成交量 \(V_{\min}\)：若 \(V_{t_k} < V_{\min}\)，则为成交量缩量；若 \(V_{t_k} > V_{\min}\)，则为成交量放量；结合收益率涨跌情况，可以将日内每分钟行情进一步划分为：放量持续上涨、缩量持续上涨、放量持续下跌、缩量持续下跌；

4. 对于“放量上涨”、“放量持续上涨”、“放量下跌”、“放量持续下跌”4种技术形态，计算每种技术形态对应的时段收益率作为该形态下的日内收益率，对过去20个日度收益率先做截面标准化，取其绝对值，再做截面标准化，然后计算绝对值调整后的过去20日收益率的标准差，由此得到4种技术形态下的波动率因子，再将这4个因子等权合成即得到主力波动率因子。

\subsection{说明}
上述“放量上涨”和“放量持续上涨”状态可捕捉主力拉升行为，“放量下跌”和“放量持续下跌”状态捕捉主力出货行为，由此构建的主力波动率刻画了主力资金行为对股价影响的波动程度，波动程度越大，越有可能导致股价过度反应，未来收益率更低。

\subsection{参考文献}
\begin{enumerate}
    \item 叶尔乐.2024.基于分钟K线的“主力波动率”构造及应用.民生证券
\end{enumerate}

\section{跌幅最大时刻笔均成交金额（高频因子-成交分布类）}

\subsection{因子定义}

1. 将日内 240 个 1 分钟数据集合标记为 \( T = \{t_1, t_2, \dots, t_{240}\} \)，判断股票 i 的每分钟涨跌幅，将跌幅最大的 10\% 时刻标记为 \( P = \{t_{i1}, t_{i2}, \dots, t_{ik}\} \)；

2. 对跌幅最大时刻进行汇总，根据其成交额 \( \mathrm{Amt} \) 之和除以成交笔数 \( \mathrm{DealNum} \) 之和，构造跌幅最大时刻的笔均成交金额 \( \mathrm{ATD}_{\mathrm{DownTop10\%}} \)：
\begin{equation}
    \mathrm{ATD}_{\mathrm{DownTop10\%}} = \frac{\sum_{t \in P} \mathrm{Amt}_t}{\sum_{t \in P} \mathrm{DealNum}_t}
\end{equation}

3. 同理，将全天的成交金额除以全天的成交笔数，得到全天的笔均成交金额 \( \mathrm{ATD}_T \)：
\begin{equation}
    \mathrm{ATD}_T = \frac{\sum_{t \in T} \mathrm{Amt}_t}{\sum_{t \in T} \mathrm{DealNum}_t}
\end{equation}

4. 将跌幅最大时刻的笔均成交金额除以全天的笔均成交金额，即得到去量纲化后的跌幅最大时刻标准化笔均成交金额因子 \( \mathrm{SATD}_{\mathrm{DownTop10\%}} \)：
\begin{equation}
    \mathrm{SATD}_{\mathrm{DownTop10\%}} = \frac{\mathrm{ATD}_{\mathrm{DownTop10\%}}}{\mathrm{ATD}_T}
\end{equation}

\subsection{说明}
笔均成交额类因子通过汇总“特殊的、具有较多信息含量”时刻的成交金额情况来刻画主力资金的行为动向；笔均成交金额越大，说明该特殊时间内资金优势越明显，其属于主力资金的可能性越大。“跌幅最大时刻”除了考虑股价涨跌方向，还考虑了涨跌程度，能进一步追踪主力资金抄底意愿强度（跌幅越大，抄底意愿越强），比只考虑下跌时刻的相关性因子表现更优。

\subsection{参考文献}
\begin{enumerate}
    \item 张欣鹏, 张宇. 2025. 日内特殊时刻蕴含的主力资金 Alpha 信息. 国信证券.
\end{enumerate}

\section{知情主卖占比(高频因子-资金流类)}

\subsection{因子定义}
使用股票过去 \( T=20 \) 个交易日的分钟逐笔成交数据进行如下时序回归，得到回归残差 \(\varepsilon_{ij}\) 即为股票预期外收益：
\begin{equation}
    R_{i,j,n} = \gamma_0 + \sum_{k=1}^4 \gamma_{1,k} D_{k,i,j,n}^{weekday} + \sum_{k=1}^3 \gamma_{2,k} D_{k,i,j,n}^{period} + \gamma_{3,1} R_{i,j-1,n} + \varepsilon_{i,j,n}
\end{equation}

将预期收益为正的分钟时刻的投资者主动卖出行为定义为知情主卖，对这些分钟时刻的主卖成交额求和得到知情主卖额，计算知情主卖占比：
\begin{equation}
    \text{知情主卖占比} = \frac{1}{T} \sum_{n=t}^{n=t-T+1} \frac { \sum_{j=9:30}^{14:56} \text{知情主卖成交额}_{i,j,n} } {\left( \sum_{j=9:30}^{14:56} \text{成交额}_{i,j,n} )}
\end{equation}

\subsection{变量定义}
\begin{itemize}
    \item ① \( R_{i,j,n} \) 为股票 \( i \) 在 \( n \) 日第 \( j \) 分钟的收益率，\( R_{i,j-1,n} \) 为收盘滞后项；
    \item ② \( D_{i,j,n}^{\mathrm{weekday}} \)，\( k=1,2,3,4 \) 为虚拟变量，表示星期一至星期四；
    \item ③ \( D_{i,j,n}^{\mathrm{Period}} \)，\( k=1,2,3 \) 为时间段虚拟变量，表示开盘后 30 分钟、盘中时段、收盘前 30 分钟；
    \item ④ 基于逐笔成交数据中的 BS 标志将逐笔成交数据合成为分钟级别的主动买卖金额，\( B \) 为主动买入，\( S \) 为主动卖出。
\end{itemize}

\subsection{说明}
可以利用对知情交易的判断，将主卖进行过滤，得到知情主卖；测试发现，知情主卖占比越高，股票未来收益越差，因为具有信息优势的知情交易者不看好股票未来表现。

\subsection{参考文献}
\begin{enumerate}
    \item 冯佳睿, 袁林青. 2020. 知情交易与主卖主卖. 海通证券.
\end{enumerate}

\section{量峰加权价格分位点(高频因子-量价相关性类)}

\subsection{因子定义}
\begin{enumerate}
    \item 对过去 20 日同时点成交量按照 1 倍标准差的标准划分，高于 1 倍标准差划分为“喷发成交量”，低于 1 倍标准差划分为“温和成交量”。
    \item 对“喷发成交量”进一步按照连续与否划分：若当前分钟为“喷发成交量”，且前后 1 分钟均为“温和成交量”，则划为“孤立喷发成交量”，称为“量峰”；若当前分钟为“喷发成交量”，且前后 1 分钟也存在“喷发成交量”，则划为“连续喷发成交量”，称为“量岭”；将“温和成交量”称为“量谷”。
    \item 将“日内最高价、日内最低价、昨日收盘价”三者的最高值与最低值作为【昨收、今收】区间的价格高低位，并计算每日量峰成交量加权价格的相对分位点，20 日的分位点均值即为量峰加权价格分位点因子；
    \item 需对量峰加权价格分位点因子做 20 日反转因子的中性化处理。
\end{enumerate}

\subsection{说明}
量峰加权价格分位点因子衡量了知情交易价格的相对水平，因子值越高，知情交易者情绪越乐观，未来收益也可能越高。

\subsection{参考文献}
\begin{enumerate}
    \item 樊建裕, 王志强. 2025. 高频成交量的峰、岭、谷信息. 开源证券.
\end{enumerate}

\section{三小将-撤单率（高频因子-流动性类）}

\subsection{因子定义}
基于逐笔委托数据，将早盘集合竞价阶段（只取9:15至9:20委托订单可撤销的阶段）的撤单数据划分为成交、全撤、部撤和废单四类；只对卖方撤单数据计算全撤、部撤和废单三类撤单率指标，将其等权合成得到三小将-撤单率因子：
\begin{equation}
    \text{撤单率} = \frac{\text{撤单量}}{\text{自由流通股本}}
\end{equation}

\subsection{变量定义}
\begin{itemize}
    \item 可对计算的撤单率指标计算20日均值进行低频化处理。
\end{itemize}

\subsection{说明}
撤单率因子取值越高，股票当天的撤单越多，交易越不活跃（成交不迅速、买卖价差大），而市场会给予低流动性标的更高的收益补偿，是一个正向因子。

\subsection{参考文献}
\begin{enumerate}
    \item 魏建榕, 苏良. 2024. 订单流系列：撤单行为规律初探. 开源证券.
\end{enumerate}

\section{跌幅时间重心偏离（高频因子-收益分布类）}

\subsection{因子定义}
基于个股上涨和下跌的 1 分钟收益率序列，逐日统计幅度和时间信息，记为：
\begin{itemize}
    \item \textbf{上涨幅度}: \( R_u = [r_1^u, r_2^u, \cdots, r_n^u] \)
    \item \textbf{上涨时间}: \( U = [u_1, u_2, \cdots, u_n] \)
    \item \textbf{下跌幅度}: \( R_d = [r_1^d, r_2^d, \cdots, r_m^d] \)
    \item \textbf{下跌时间}: \( D = [d_1, d_2, \cdots, d_m] \)
\end{itemize}

分别计算日频指标：
\begin{equation}
    \text{涨幅时间重心}: G_u = \frac{U R_u^T}{\|R_u\|_1}
\end{equation}
\begin{equation}
    \text{跌幅时间重心}: G_d = \frac{D R_d^T}{\|R_d\|_1}
\end{equation}

通过横截面回归，构造“时间差”指标，并取 \( N=20 \) 日均值作为因子：
\begin{equation}
    G_d = \alpha + \beta G_u + \varepsilon
\end{equation}
\begin{equation}
    \mathrm{TGD} = \frac{1}{N} \sum_{i=1}^N \varepsilon
\end{equation}

\subsection{说明}
涨、跌时间重心通过分钟涨跌幅对修正的时间截加权平均得到，反映了股价在日内交易中上涨和下跌的位置，捕捉了股票的交易行为特征，两者的相对位置（时间差）是一个有效的 Alpha 因子，与未来收益正相关；“时间差 Alpha”是收益率结构和“低波效应”的综合。

\subsection{参考文献}
\begin{enumerate}
    \item 魏建榕, 苏良, 盛少成. 2022. 日内分钟收益率的时序特征：逻辑讨论与因子增强. 开源证券.
\end{enumerate}

\section{每笔成交量筛选的局部反转（高频因子-动量反转类）}

\subsection{因子定义}
\begin{equation}
    \mathrm{sum}\left( \left\{ \mathrm{ret}_i \middle| q_j(\mathrm{pvol}_i) < \mathrm{pvol}_i < q_{j+1}(\mathrm{pvol}_i) \right\} \right)
\end{equation}

\subsection{变量定义}
\begin{itemize}
    \item ① \(\mathrm{ret}_i\) 和 \(\mathrm{pvol}_i\) 分别为每个时间段的收益率和每笔成交量=成交量/成交笔数，时间段为分钟频率；
    \item ② \(q()\) 为分位数函数，将每笔成交量进行排列，取每笔成交量最大组（位于80\%-100\%区间）对应时间段的收益率计算反转因子。
\end{itemize}

\subsection{说明}
股价变动在每笔成交量较小的时间段表现出动量效应，而每笔成交量较大的时间段表现出反转效应，所以可以利用每笔成交量对反转因子进行“提纯”，即对于过去20个交易日的所有时间段，将每笔成交量较小、表现出动量效应的时间段提取出来，剩余的每笔成交量最大的时间段内反转效应最为纯净，收益能力也较“提纯”前的传统反转因子大幅增强。每笔成交量实际同时衡量了成交活跃程度以及价格的变动，而成交活跃度和价格的变动共同反映了交易异常活跃的程度，交易异常活跃的程度是市场过度反应导致价格变动产生反转效应的来源。

\subsection{参考文献}
\begin{enumerate}
    \item 覃川桃, 郑起. 2021. 高频因子（十）：量价关系中的反转微观结果. 长江证券.
\end{enumerate}

\section{量岭间隔偏度（高频因子-成交分布类）}

\subsection{因子定义}


① 对过去 20 日同时点成交量按照 1 倍标准差的标准划分，高于 1 倍标准差划分为“喷发成交量”，低于 1 倍标准差划分为“温和成交量”；

② 对“喷发成交量”进一步按照连续与否划分：若当前分钟为“喷发成交量”，且前后 1 分钟均为“温和成交量”，则划为“孤立喷发成交量”，称为“量峰”；若当前分钟为“喷发成交量”，且前后 1 分钟也存在“喷发成交量”，则划为“连续喷发成交量”，称为“量岭”；将“温和成交量”称为“量谷”。

③ 对过去 20 日同日前后两个量岭之间的时间间隔分布，计算其偏度，即为量岭间隔偏度因子。

\subsection{变量定义}
\begin{itemize}
    \item 此外，还可统计两个量岭之间的时间间隔分布的标准差、峰度等指标，即可得到量岭间隔标准差因子、量岭间隔峰度因子。
\end{itemize}

\subsection{说明}
量岭间隔偏度因子衡量了个人投资者交易的时间间隔的分布情况，个人投资者通常是跟风交易，交易频率高，间隔时间短，易过度交易，所以时间间隔分布的异常值相对较少，且主要集中在左侧的短间隔区域，与未来收益负相关。（报告中测试发现，量岭间隔峰度因子也为负向因子，但量岭间隔标准差因子为正向因子。）

\subsection{参考文献}
\begin{enumerate}
    \item 魏建榕, 王志强. 2025. 高频成交量的峰、岭、谷信息. 开源证券.
\end{enumerate}

\section{开盘后净主买强度（高频因子-资金流类）}

\subsection{因子定义}
\begin{equation}
    \frac{1}{T} \sum_{n=t}^{t-T+1} \frac{\mathrm{mean}(\text{净主买成交额}_{i,j,n})}{\mathrm{std}(\text{净主买成交额}_{i,j,n})}
\end{equation}
\begin{equation}
    \text{净主买成交额}_{i,j,n} = \text{主动买入成交额}_{i,j,n} - \text{主动卖出成交额}_{i,j,n}
\end{equation}
\subsection{变量定义}
\begin{itemize}
    \item ① 基于逐笔成交数据中的 BS 标志将逐笔成交数据合成为分钟级别的主动买卖金额，B 为主动买入（卖出方先挂单，买入方主动触碰卖单并成交）；S 为主动卖出（买入方先挂单，卖出方主动触碰买单并成交）；
    \item ② 剔除了处于涨跌停分位上的主动买入金额和主动卖出金额数据；
    \item ③ \(i, j, n\) 分别表示第 \(i\) 只股票在第 \(n\) 个交易日内的第 \(j\) 分钟的成交额；
    \item ④ 开盘后的买入意愿占比用的是 9:30-10:00 范围内的数据；
    \item ⑤ 月度选股下，\(T=20\) 个交易日；周度选股下，\(T=5\) 个交易日。
\end{itemize}

\subsection{说明}
同时刻画了投资者在开盘后 30 分钟内主动买入行为的强度和稳健性，开盘后净主买强度越高，投资者的主动买入行为越稳健。

\subsection{参考文献}
\begin{enumerate}
    \item 冯佳睿, 袁林青. 2020. 基于主动买入行为的选股因子. 海通证券.
    \item 冯佳睿等. 2020. 高频因子的发现与思考. 海通证券.
\end{enumerate}

\section{单笔成交金额分位数(高频因子-成交分布类)}

\subsection{因子定义}
\begin{equation}
    S = \frac{A_{0.1} - A_{\min}}{A_{\max} - A_{\min}}
\end{equation}

\subsection{变量定义}
\begin{itemize}
    \item ① \(A\) 为某只股票某个交易内的分钟单笔成交金额序列，下标 0.1、max、min 分别表示该序列的 10\% 分位数、最大值、最小值；
    \item ② 需要将 \(A\) 序列从小到大排序，剔除最大的 10 个样本值（极端值）；
    \item ③ 滚动 20 个交易日计算 \(S\) 指标的均值，即可得到 \(\mathrm{QUA}\) 因子。
\end{itemize}

\subsection{说明}
单笔成交金额分位数 \(\mathrm{QUA}\) 本质上刻画了大单相对小单的成交金额偏离程度，\(\mathrm{QUA}\) 取值越小，双方偏离程度越高，而大单相对小单的成交金额偏离越大，主力（“相对大单”）对于该股票的关注度越高，主力资金参与交易的意愿越强，未来股价表现会越好；单笔成交金额分布形态除了计算分位数外，还可以计算标准差 \(\mathrm{STD}\)、峰度 \(\mathrm{KURT}\)、偏度 \(\mathrm{SKEW}\) 等。

\subsection{参考文献}
\begin{enumerate}
    \item 魏建榕, 苏良. 2022. 高频因子：分钟单笔金额序列中的主力行为刻画. 开源证券.
\end{enumerate}

\section{供应商行业集中度（图谱网络-网络结构）}

\subsection{因子定义}


供应商行业集中度 \(\mathrm{SH}\) 的计算公式如下：
\begin{equation}
    \mathrm{SH} = \sum_{i=1}^{I} w_{ji} \times H_j
\end{equation}

其中 \(H_j\) 为供应商公司所属行业的赫芬达尔指数：
\begin{equation}
    H_j = \sum_{i=1}^{I} s_{ij}^2
\end{equation}

\subsection{变量定义}
\begin{itemize}
    \item ① \(H_j\) 为供应商公司j所属行业的赫芬达尔指数；
    \item ② \(s_{ij}\) 为属于行业 \(j\) 的公司在该行业的市场份额，市场份额可以通过营业收入或市值来计算；
    \item ③ \(w_{ij}\) 为客户 \(i\) 对供应商 \(j\) 的采购额占自身总采购额的占比。
\end{itemize}

\subsection{说明}
供应商行业集中度为公司所有供应商所属行业集中度的加权平均，反映了公司在其供应商行业中的竞争地位；若供应商行业集中度过高，意味着公司依赖少数几家供应商，更易面临供应中断的风险；若供应商行业集中度过低，意味着公司可以从众多供应商中选择，具有较高的灵活性和议价能力。

\subsection{参考文献}
\begin{enumerate}
    \item Jing Wu, John R. Birge. (2014). \textit{Supply Chain Network Structure and Firm Returns}. SSRN Electronic Journal.
\end{enumerate}

\section{单和小单的错位相关性(高额因子-资金流类)}

\subsection{因子定义}

\begin{equation}
    \mathrm{RankCorr}(S_t, S_{t+1})
\end{equation}

\subsection{变量定义}
\begin{itemize}
    \item ① \( S_t \) 和 \( S_{t+1} \) 都为小单（<4万元）过去20个交易日同期的资金净流入序列，但是双方日期错开1个交易日。
\end{itemize}

\subsection{说明}
在资金流错位相关性因子中，上述小单和小单的错误相关性因子表现最好；该因子的alpha来源是“散户追涨杀跌的羊群效应”，其中，散户羊群效应可以通过股票“过去N日收益率”与“小单未来N日净流入”的秩相关系数数学公式：
\begin{equation}
    \mathrm{RankCorr}(R_t, S_{t+1})
\end{equation}
来衡量。

\subsection{参考文献}
\begin{enumerate}
    \item 魏建榕, 盛少成. 2022. 新型因子：资金流动力学与散户羊群效应. 开源证券.
\end{enumerate}

\section{卖单非流动性（高频因子-流动性类）}

\subsection{因子定义}
\begin{equation}
    r_{i,t} = \alpha + \beta_1 * S_{i,t} + \beta_2 * B_{i,t} + \epsilon_{it}
\end{equation}

\subsection{变量定义}
\begin{itemize}
    \item ① \( \beta_1 \) 为卖出非流动性系数；
    \item ② \( \beta_2 \) 为买入非流动性系数；
    \item ③ \( S_{i,t} \) 为股票 \( i \) 在 \( t \) 时间区间内的主动卖出金额；
    \item ④ \( B_{i,t} \) 为股票 \( i \) 在 \( t \) 时间区间内主动买入金额。
\end{itemize}

\subsection{说明}
卖单非流动性衡量了高频数据下主动卖出的交易金额对于股票价格变动的影响；卖单非流动性在控制风险后的 Fama-MacBeth 截面回归对收益率显著，且预测效果要好于买单非流动性，主要是由于投资者存在亏损厌恶的心理。

\subsection{参考文献}
\begin{enumerate}
    \item Brennan, M. J., Chordia, T., Subrahmanyam, A., \& Tong, Q. (2012). \textit{Sell-order liquidity and the cross-section of expected stock returns}. Journal of Financial Economics, 105(3), 523-541.
    \item 朱剑涛. 2017. 技术类新 Alpha 因子的批量测试. 东方证券.
\end{enumerate}

\section{地理动量溢出效应（图谱网络-动量溢出）}

\subsection{因子定义}
\begin{equation}
    \mathrm{RET}_{it}^{\mathrm{GEO}} = \sum_{j \neq i} \mathrm{geographic\_weight}_{ij,t} \times \mathrm{Re T}_{jt}
\end{equation}
\begin{equation}
    \mathrm{geographic\_weight}_{ij,t} = \frac{\mathrm{GMV}_{j,t}}{\sum_{j=1 and j \neq i}^N \mathrm{GMV}_{j,t}}
\end{equation}

\subsection{变量定义}
\begin{itemize}
    \item ① \( \mathrm{ReT}_{jt} \) 为股票 \( j \) 在 \( t \) 时刻的月度收益率；
    \item ② \( \mathrm{GMV}_{j,t} \) 为 \( t \) 时期同地理区域内公司 \( j \) 的市值。
\end{itemize}

\subsection{说明}
地理动量溢出效应因子刻画了与单只股票同位于相同地理其余股票涨跌对该股票涨跌的联动效应（领先滞后效应）。

\subsection{参考文献}
\begin{enumerate}
    \item Usman Ali, David Hirshleifer. (2020). \textit{Shared analyst coverage: Unifying momentum spillover effects}. Journal of Financial Economics, 6(136), 649-675.
\end{enumerate}

\section{上下行跳跃波动的不对称性（高频因子-波动跳跃类）}

\subsection{因子定义}

\begin{equation}
    \mathrm{SRVJ}_t = \mathrm{RVJP}_t - \mathrm{RVJN}_t
\end{equation}
\begin{equation}
    \mathrm{RVJP}_t = \max\left(\mathrm{RS}^+_t - \frac{1}{2}\hat{IV}_t, 0\right)
\end{equation}
\begin{equation}
    \mathrm{RVJN}_t = \max\left(\mathrm{RS}^-_t - \frac{1}{2}\hat{IV}_t, 0\right)
\end{equation}

\subsection{变量定义}
\begin{itemize}
    \item ① \( \mathrm{RVJP}_t \) 和 \( \mathrm{RVJN}_t \) 分别为交易日 \( t \) 的上行跳跃波动率和下行跳跃波动率，计算细节可参考相关日历页。
\end{itemize}

\subsection{说明}
将上行跳跃波动率减下行跳跃波动率即可得到上下行跳跃波动的不对称性；短期内，跳跃波动的不对称性通常有负的风险溢价，股价短期大幅上涨，未来更易补跌，股价短期大幅下跌，未来更易补涨。

\subsection{参考文献}
\begin{enumerate}
    \item Bo Yu, Bruce Mizrach, Norman R. Swanson. (2020). \textit{New Evidence of the Marginal Predictive Content of Small and Large Jumps in the Cross-Section}. Econometrics, MDPI, 8(2), 1-52.
\end{enumerate}

\section{成交量分桶熵(高频因子-成交分布类)}

\subsection{因子定义}
\begin{equation}
    \mathrm{vol\_entropy}(V) = -\sum_{k=0}^{\min(\mathrm{maxbin}, \mathrm{len}(V))} p_k \ln(p_k) \cdot \mathbf{1}_{(p_k > 0)}
\end{equation}

\subsection{变量定义}
\begin{itemize}
    \item ① 对日内分钟成交量 \( V \) 序列基于 \([\min(V), \max(V)]\) 区间进行等距分桶，并计算成交量 \( V \) 序列取值落在第 \( k \) 个桶的概率 \( p_k \)；
    \item ② \( \mathrm{maxbin} \) 为桶的个数，\( \mathrm{len}(V) = T \) 为日内分钟成交量 \( V \) 序列的长度；
    \item ③ 计算个股分桶熵数值近 20 日标准差以衡量个股分桶熵在时序上的离散程度。
\end{itemize}

\subsection{说明}
成交量分桶熵（标准差）衡量了成交量不稳定性的分散程度，该因子值越大，说明该股在近一段时间内成交量不稳定性的分散程度较大，股价曾经或正在受到知情交易者的带动，此时个股不稳定性极高，整体下行风险较大；日度成交量分桶熵取值越大，日内成交量分布越均匀。

\subsection{参考文献}
\begin{enumerate}
    \item 郑兆森. 2022. 成交量分布中的 Alpha. 兴业证券.
\end{enumerate}

\section{机构交易热度(行为金融因子-投资者注意力)}

\subsection{因子定义}
\begin{equation}
    \mathrm{TURN\_INST} = \frac{1}{m}\frac{ \sum_{t=1}^{t-m+1} \left( \mathrm{INFLOW\_INST}_{i,t} + \mathrm{OUTFLOW\_INST}_{i,t} \right)}{\mathrm{CAP}_{i,t}}
\end{equation}

\subsection{变量定义}
\begin{itemize}
    \item ① \(\mathrm{INFLOW\_INST}_{i,t}\) 和 \(\mathrm{OUTFLOW\_INST}_{i,t}\) 分别为股票 \(i\) 在交易日 \(t\) 的机构买入金额和卖出金额；
    \item ② 将单笔成交额高于 100 万元的投资者作为机构投资者；
    \item ③ \(\mathrm{CAP}_{i,t}\) 为股票 \(i\) 在交易日 \(t\) 的流通市值；
    \item ④ \(m\) 为交易日窗口，一般取为月度。
\end{itemize}

\subsection{说明}
机构交易热度属于股票交易活跃度类因子，通过衡量股票交易活跃度来反映投资者关注度。股票交易越活跃，投资者对该股票的关注度越高；而投资者有限注意力会推动关注度高的股票更易出现非理性净买入溢价，进而导致股价后续反转。

\subsection{参考文献}
\begin{enumerate}
    \item 陈升锐. 2024. 投资者有限关注及注意力捕捉与溢出. 中信建投.
\end{enumerate}

\section{净委买变化率(高频因子-资金流类)}

\subsection{因子定义}
\begin{equation}
    \text{净委买变化率}^{T}_{k,t} = \frac{\text{净委买变化量}^{T}_{k,t}}{\text{流通股本}_{T}}
\end{equation}
\begin{equation}
    \text{净委买变化量}^{T}_{k,t} = \sum_{j=1}^{k}  \text{委买变化量}^{T}_{j,t} - \sum_{j=1}^{k}\text{委卖变化量}_{j,t}^T 
\end{equation}

\subsection{变量定义}
\begin{itemize}
    \item ① \(\text{委买变化量}^{T}_{j,t}\) 和 \(\text{委卖变化量}^{T}_{j,t}\) 分别为 \(T\) 日内 \(t\) 至 \(t+1\) 时刻的第 \(j\) 档委买的变化量和第 \(j\) 档委卖的变化量；
    \item ② \(\text{流通股本}_{T}\) 为 \(T\) 日股票的流通股本；
    \item ③ 一般用前 1 档位数据即可；使用的档位数量越多，因子的选股能力越弱；
    \item ④ 此外还可计算 \(T\) 日内净委买变化率序列的均值、标准差和偏度等分布特征；也可分时间区间计算因子，如使用 9:30-10:00 范围内的数据计算开盘后 30 分钟净委买变化率；
    \item ⑤ 可将过去 5 日、20 日上述因子的均值和中位数作为周度和月度的因子值。
\end{itemize}

\subsection{说明}
刻画了投资者的买入意愿，委买量的增加代表了投资者买入意愿的增强，委卖量的增加代表了投资者卖出意愿的增强；开盘后 30 分钟的数据更是包含了投资者对于前一天收盘后股票信息的集中反馈，集中反馈越正面，体现出的买入意愿越强，股票未来的超额收益表现越好。

\subsection{参考文献}
\begin{enumerate}
    \item 冯佳睿, 袁林青. 2019. 捕捉投资者的交易意愿. 海通证券.
\end{enumerate}

\section{日内黄金分割反转(高频因子-动量反转类)}

\subsection{因子定义}

\begin{equation}
    \mathrm{Ret}_{10:00\text{-to-close}} = \sum_{t=1}^{20} \ln \left( \frac{\mathrm{Close}_t}{\mathrm{Price}_t^{10:00}} \right)
\end{equation}

\subsection{变量定义}
\begin{itemize}
    \item ① \( \mathrm{Close}_t \) 和 \( \mathrm{Price}_t^{10:00} \) 分别为 \( t \) 交易日的收盘价和早盘 10 点的价格。
\end{itemize}

\subsection{说明}
10:00 的成交价更能反映出当日开盘时投资者的心理预期和价格水平，是理想的黄金分割时点，由此计算的反转因子表现要优于日内收益和日度收益。

\subsection{参考文献}
\begin{enumerate}
    \item 朱定豪. 2020. 昼夜分离：隔夜跳空与日内反转选股因子. 华安证券.
\end{enumerate}

\section{时间网络相对中心度(图谱网络-网络结构)}

\subsection{因子定义}


计算股票 \( i \) 在 \( t \) 时刻截面上相对于样本股票池的平均收益的偏离程度：
\begin{equation}
    z_{i,t} = \frac{r_{i,t} - \hat{r}_{m,t}}{\sigma_{m,t}}
\end{equation}

计算过去 20 个交易日股票 \( i \) 对其他股票收益的偏离程度的平均值：
\begin{equation}
    \bar{z}_{i,t} = \sqrt{\frac{1}{\Delta t} \sum_{\tau=1}^{\Delta t} z_{i,t-\tau}^2}
\end{equation}

计算时间维度的网络相对中心度：
\begin{equation}
    \mathrm{TCC}_{i,t} = \frac{1}{\bar{z}_{i,t}^2}
\end{equation}

\subsection{说明}
TCC 因子衡量了股票的网络位置随着时间推移的稳定性，因子值越高，说明过去时间段内股票 \( i \) 对网络中其它股票的收益偏离程度越小，股票在网络中的位置越稳定。

\subsection{参考文献}
\begin{enumerate}
    \item 张自力, 闫红蕾, 张楠. 2020. 股票网络、系统性风险与股票定价. 经济学, 1, 329-350.
    \item 曹春晓, 杨国平. 2021. 股票网络与网络中心度因子研究. 华西证券.
\end{enumerate}

\section{“孤雁出群”因子（高频因子-量价相关性类）}

\subsection{因子定义}


1. 对于交易日 \( t \)，计算每只股票的分钟收益率：\( t \) 分钟收盘价 / \( t-1 \) 分钟收盘价 - 1；

2. 计算“分钟市场分化度”，即每一分钟所有股票的分钟收益率的标准差；

3. 计算这一天所有“分钟市场分化度”的均值，找到那些“分钟市场分化度”小于均值的时刻，记为 \( t \) 日的“不分化时刻”；

4. 取市场上所有股票在当日“不分化时刻”的成交额序列，然后分别计算每只股票的分钟成交额序列与其余股票分钟成交额序列的 Pearson 相关系数；

5. 对上述的相关系数取绝对值，然后分别计算每只股票与其余股票的相关系数绝对值的均值，即为“日孤雁出群”因子；

6. 每月末，分别计算过去 20 个交易日“日孤雁出群”因子的均值和标准差，并将二者等权合成，即为“孤雁出群”因子。

\subsection{说明}
“孤雁出群”因子衡量了个股在市场分化不明显时，其成交额与市场趋势的关联性，与未来收益负相关；当市场分化不明显缺乏热点时，相关系数越小的股票，其成交额走出了独立趋势，更有可能在酝酿新的热点。

\subsection{参考文献}
\begin{enumerate}
    \item 曹春晓. 个股成交额的市场跟随性与“水中行舟”因子, 2023, 方正证券.
\end{enumerate}

\section{卖盘深度（高频因子-流动性类）}

\subsection{因子定义}
\begin{equation}
    \mathrm{ask\_depth} = \sum \left| \frac{\mathrm{av}_i \times \mathrm{mid_{price}}}{\mathrm{a}_i - \mathrm{last_{mid}}+ \epsilon} \right| 
\end{equation}

\subsection{变量定义}
\begin{itemize}
    \item ① 使用五档卖单与当下时刻中间价的价差的倒数对卖单的数量进行加权得到卖盘深度；
    \item ② 当单边所有卖单档位为空时，令市场深度为0；
    \item ③ 可对上述因子日内序列求均值得到日度因子值，再进行20个交易日指数移动平均得到月度因子值。
\end{itemize}

\subsection{说明}
卖盘深度统计了卖单的挂单量情况，反映了卖单盘口深度，五档卖单中，对报价离中间价越远的卖单的卖量给予越小的权重；与未来收益负相关，因子取值越大，挂单量越多，盘口深度越大，流动性越高。也可对日内该因子序列求标准差，衡量其在当日均匀程度，与流动性负相关。

\subsection{参考文献}
\begin{enumerate}
    \item 胡骏等. 2021. 高频视角下的微观流动性与波动性. 中金公司.
\end{enumerate}

\section{CSSD 模型(行为金融因子-羊群效应)}

\subsection{因子定义}
利用个股收益 \( r_i \) 和市场收益 \( r_m \)，计算股票收益的横截面标准差：
\begin{equation}
    \mathrm{CSSD} = \sqrt{\frac{\sum_{i=1}^N (r_i - r_m)^2}{N-1}}
\end{equation}

以 \(\mathrm{CSSD}\) 为 \( y \)，以 \( D_t^u \)（市场收益高于计算收益阈值时取值为 1，否则取值为 0）和 \( D_t^d \)（市场收益低于计算收益阈值时取值为 1，否则取值为 0）为 \( x \)，进行如下时序回归：
\begin{equation}
    \mathrm{CSSD} = \alpha + \beta_1 D_t^u + \beta_2 D_t^d + \varepsilon
\end{equation}

\subsection{说明}
CSSD 模型对羊群效应的检验主要基于时序回归上市场极端收益项前回归系数正负及显著性；当 \( \beta_1 \) 或 \( \beta_2 \) 显著小于 0 时，说明市场极端上涨或下跌时，CSSD 取值明显下降，投资者的投资行为越来越一致，市场上羊群效应越明显。

\subsection{参考文献}
\begin{enumerate}
    \item Christie, W. G., and Huang, R. D. (1995). \textit{Following the Pied Piper: Do Individual Returns Herd around the Market?}. Financial Analysts Journal, 51(4), 31–37.
\end{enumerate}

\section{修正的振幅（高频因子-波动跳跃类）}

\subsection{因子定义}


对股票 \( i \)，计算 \( n \) 日的日跳跌度和日度振幅：
\begin{equation}
    t \text{ 分钟跳跌度} = (t \text{ 分钟单复利差}) \times 2 - (t \text{ 分钟连续复利收益率})^2
\end{equation}
\begin{equation}
    \text{日度振幅} = \frac{\text{当日最高价} - \text{当日最低价}}{\text{上一日收盘价}}
\end{equation}

每日计算日跳跌度因子的横截面均值，将日跳跌度小于横截面均值的股票的振幅做修正：即将其这一天的振幅取相反数，其余股票振幅取值不变，由此得到“翻转振幅”因子。

\subsection{变量定义}
\begin{itemize}
    \item ① 日跳跌度因子为日内所有分钟跳跌度序列的均值，具体计算逻辑见相关日历页；
    \item ② 每月末对过去 20 天的“翻转振幅”求均值，即得到“修正振幅因子”；
    \item ③ 也可用日频量价数据计算日间跳跌度，利用日间跳跌度对振幅做相同逻辑的修正。
\end{itemize}

\subsection{说明}
修正振幅因子利用股票日跳跌度对传统振幅的变化过程做了区分：那些未发生跳跃或跳跌程度较小的股票，可能并未吸引到博彩偏好型投资者，股价也未过度反应，可对这些时刻的振幅做反向修正，增强其低波动异象的属性。

\subsection{参考文献}
\begin{enumerate}
    \item 曹春晓. 个股股价跳跃及其对振幅因子的改进, 2022, 方正证券.
    \item Jiang, G. J., \& Oomen, R. C. A. (2008). \textit{Testing for jumps when asset prices are observed with noise - a "swap variance" approach}. Journal of Econometrics, 144(2), 352-370.
\end{enumerate}

\section{高频已实现峰度（高频因子-收益分布类）}

\subsection{因子定义}
\begin{equation}
    \mathrm{Rkurt}_i = \frac{N \sum_{j=1}^N (r_{ij} - \bar{r}_i)^4}{\mathrm{RVar}_i^2}
\end{equation}
\begin{equation}
    \mathrm{RVar}_i = \sum_{j=1}^N (r_{ij} - \bar{r}_i)^2
\end{equation}

\subsection{变量定义}
\begin{itemize}
    \item ① \( r_{ij} \) 为股票 \( i \) 的 1 分钟对数收益率序列或 5 分钟对数收益率序列；
    \item ② \( N \) 表示个股在交易日 \( t \) 中特定频率的分钟级别数据个数；
    \item ③ 在任意选股时刻，因子值为前几日指标的均值，如月度选股，计算因子值时使用的是股票过去 20 个交易日的均值。
\end{itemize}

\subsection{说明}
股票日内的价格形态分布特征对于股票未来收益具有一定预测作用，高频收益峰度就是刻画日内价格形态的特征之一；高频峰度因子通常与股票未来收益负相关。

\subsection{参考文献}
\begin{enumerate}
    \item 冯佳睿, 袁林青. 2017. 高频因子之股票收益分布特征. 海通证券.
    \item Amaya, D., Christoffersen, P., Jacobs, K., and Vasquez, A. (2016). \textit{Does Realized Skewness Predict the Cross-Section of Equity Returns?}. Journal of Financial Economics, 118, 135-167.
\end{enumerate}

\section{上游传导动量（图谱网络-动量溢出）}

\subsection{因子定义}
计算上游公司 B 与下游公司 A 间每一对上下游产品 \( j \to i \) 间的关联度，并取最大值作为两个公司最终的关联度：
\begin{equation}
    \mathrm{wgt}_{(B,j),(A,i),t} = \frac{(\mathrm{prdIncome}_{B,j,t} + \mathrm{prdIncome}_{A,i,t}) \times \mathbf{1}\{j \to i\}}{\sum_i \mathrm{prdIncome}_{A,i,t} + \sum_j \mathrm{prdIncome}_{B,j,t}}
\end{equation}

对下游公司 A，以产业链关联度为权重加权其所有上游公司的月度收益率，即为上游传导动量因子：
\begin{equation}
    \mathrm{upstream\_Moment}_{A,t} = \frac{\sum_{A \neq B} \mathrm{wgt}_{A,B,t} \times \mathrm{rtn}_{B,t}}{\sum_{A \neq B} \mathrm{wgt}_{A,B,t}}
\end{equation}

\subsection{变量定义}
\begin{itemize}
    \item ① \( \mathrm{prdIncome}_{A,i,t} \) 为公司 A 在产品 i 上的营业收入；
    \item ② \( \sum_i \mathrm{prdIncome}_{A,i,t} \) 为公司 A 的当期总营收；
    \item ③ \( \mathbf{1}\{j \to i\} \) 表示产品 j 为产品 i 的上游；
    \item ④ \( \mathrm{rtn}_{B,t} \) 为公司 B 在 t 时刻的月度收益率；
    \item ⑤ 可以将上游传导动量因子对月度收益率做正交化，剥离月度反转效应。
\end{itemize}

\subsection{说明}
上游传导动量因子值越大，表示公司 A 作为下游企业时，位于其产业链上游的高关联度高集群公司在考察月份涨幅相对较高，公司 A 接下来也会有上涨的可能。

\subsection{参考文献}
\begin{enumerate}
    \item 郑兆磊.2023.产业链视角下的Alpha传导研究.兴业证券
\end{enumerate}

\section{知情主买占比（高级因子-资金流类）}

\subsection{因子定义}


使用股票过去20个交易日的分钟逐笔成交数据进行如下时序回归，得到回归残差 \( \varepsilon_{i,j} \) 即为股票预期外收益：
\begin{equation}
    R_{i,T,j} = \gamma_0 + \sum_{k=1}^4 \gamma_{1,k} D_{T,k,j}^{\mathrm{weekday}} + \sum_{k=1}^3 \gamma_{2,k} D_{T,k,j}^{\mathrm{Period}} + \gamma_{3,1} R_{i,T,j-1} + \varepsilon_{i,j}
\end{equation}

将预期收益为正的分钟时刻的投资者主动卖出行为定义为知情主卖，对这些分钟时刻的主卖成交额求和得到知情主卖额，计算知情主卖占比：
\begin{equation}
    \text{知情主买占比} = \frac{1}{T} \sum_{n=t}^{t-T+1}  \frac{\sum_{j=9:30}^{14:56} \text{知情主买成交额}_{i,j,n} }{ \sum_{j=9:30}^{14:56} \text{成交额}_{i,j,n} }
\end{equation}

\subsection{变量定义}
\begin{itemize}
    \item ① \( R_{i,T,j} \) 为股票 \( i \) 在 \( T \) 日第 \( j \) 分钟的收益率，\( R_{i,T,j-1} \) 为收益滞后项；
    \item ② \( D_{T,k,j}^{\mathrm{weekday}}, k=1,2,3,4 \) 为虚拟变量，表示星期一至星期四；
    \item ③ \( D_{T,k,j}^{\mathrm{Period}}, k=1,2,3 \) 为时间段虚拟变量，表示开盘后 30 分钟、盘中时段、收盘前 30 分钟；
    \item ④ 基于逐笔成交数据中的 BS 标志将逐笔成交数据合成为分钟级别的主动买卖金额，B 为主动买入，S 为主动卖出。
\end{itemize}

\subsection{说明}
可以利用对知情交易的判断，将主买进行过滤，得到知情主买；测试发现，知情主买占比越高，股票未来收益越差，可能是因为知情交易下也存在过度反应的倾向。

\subsection{参考文献}
\begin{enumerate}
    \item 冯佳睿, 袁林奇. 2020. 知情交易与主买主卖. 海通证券.
\end{enumerate}

\section{量涌波动率（高级因子-波动跳跌类）}

\subsection{因子定义}

a 采用二分法的思想，循环对每只股票按照成交量峰值时刻对时段进行划分（剔除开盘的10分钟）；若时段内包含至少3分钟，根据1分钟成交量的最大时刻（寻找最大时刻的时候不包括时段内第一分钟和最后一分钟）将时段划分为两个时段（最大时刻归为第一个时段）；若时段包含2分钟，划分为两个长度为1分钟的时段；若时段仅包含1分钟，不再进行划分；

b 对每只股票计算步骤a中划分得到的每个小时段的收益率，并计算这些区间收益率的标准差，作为该股票当日因子值；

c 高频因子低频化：计算步骤b中得到的日频因子标准化后，计算过去20天的标准差，记为“量涌波动率”因子。

\subsection{说明}
“量涌波动率”刻画不同交易情绪投资者对股票收益率的影响程度，波动率越小，说明不同情绪投资者对股票收益率的影响程度基本一致，更不容易出现过度反应现象，股票的“区段收益稳定性”更强，股票未来收益率更高；反之，股票的“区段收益稳定性”更弱，股票未来收益率更低。

\subsection{参考文献}
\begin{enumerate}
    \item 叶尔乐. 2025. 日内“动量脉冲”与股价过度反应的精细刻画. 民生证券.
\end{enumerate}

\section{筹码集中度(行为金融因子-筹码分布)}

\subsection{因子定义}
\begin{equation}
    \text{方式1: } \mathrm{chip\_distri} = \frac{\mathrm{price95}_t - \mathrm{price05}_t}{\mathrm{Chip\_Highest}}
\end{equation}
\begin{equation}
    \text{方式2: } \mathrm{chip\_distri} = \frac{2 \times (\mathrm{price95}_t - \mathrm{price05}_t)}{\mathrm{price95}_t + \mathrm{price05}_t}
\end{equation}

\subsection{变量定义}
\begin{itemize}
    \item ① 假设某只股票在 \( t \) 日的筹码分布共有 \( n \) 个价位（\( i = 1, 2, \ldots, n \)），\( \mathrm{vol}_{i,t} \) 为第 \( t \) 日 \( i \) 价位上的筹码峰高度，\( \mathrm{price}_{i,t} \) 为该筹码峰对应的价格，\( \mathrm{Chip\_Highest}_t \)、\( \mathrm{Chip\_Lowest}_t \) 分别对应筹码最高价和筹码最低价，\( \mathrm{price95}_t \)、\( \mathrm{price05}_t \) 分别为筹码分布 95\% 和 5\% 分位数价格；
    \item ② 上述计算的是 90\% 筹码集中度因子；此外，还可利用筹码分布 85\% 和 15\% 分位数价格计算 70\% 的筹码集中度因子；利用筹码分布 75\% 和 25\% 分位数价格计算 50\% 的筹码集中度因子等。
\end{itemize}

\subsection{说明}
筹码集中度可用于追踪股票吸筹、拉升、出货和下跌的阶段表现；吸筹阶段表现为低位集中筹码建立底仓，但该阶段通常比较隐蔽，换手率不会显著放大，筹码集中度变化不大；拉升/下跌阶段表现为股价脱离筹码集中区域，向上/向下打开筹码价格区间，使集中的筹码向上/向下转移发散，表现为筹码集中度降低；出货阶段对应的是筹码高位充分换手，某一价格区间的成交量放大，换手率显著放大，筹码集中度显著提升。因子值越小，筹码集中度越高，越有可能处于出货阶段，未来表现可能不佳。

\subsection{参考文献}
\begin{enumerate}
    \item 陈升锐, 姚紫薇. 2025. 筹码分布因子系统构建. 中信出版社.
\end{enumerate}

\section{谷岭加权价格比(高频因子-量价相关性类)}

\subsection{因子定义}


1. 对过去 20 日同时点成交量按照 1 倍标准差的标准划分，高于 1 倍标准差划分为“喷发成交量”，低于 1 倍标准差划分为“温和成交量”；

2. 对“喷发成交量”进一步按照连续与否划分：若当前分钟为“喷发成交量”，且前后 1 分钟均为“温和成交量”，则划为“孤立喷发成交量”，称为“量峰”；若当前分钟为“喷发成交量”，且前后 1 分钟也存在“喷发成交量”，则划为“连续喷发成交量”，称为“量岭”；将“温和成交量”称为“量谷”。

3. 计算如下因子的 20 日价格比均值即为谷岭加权价格比因子：
\begin{equation}
    \text{谷岭加权价格比} = \frac{\text{每日量谷成交量加权价格}}{\text{同日量岭成交量加权价格}}
\end{equation}

\subsection{变量定义}
\begin{itemize}
    \item ① 类似的，也可将每日量峰成交量加权价格与量岭成交量加权价格做比较，得到峰岭加权价格比因子，该因子也是正向因子。
\end{itemize}

\subsection{说明}
量谷时点，交易低迷，价格过度反应概率较低；量岭时点，个人投资者交易使得价格过度反应概率较高；因子值越高，个人投资者过度反应使得价格相对交易低迷时期价格向下偏离程度越高，个股未来表现越好。

\subsection{参考文献}
\begin{enumerate}
    \item 魏建榕, 王志豪. 2025. 高频成交量的峰、岭、谷信息. 开源证券.
\end{enumerate}

\section{毒流动性(高频因子-流动性类)}

\subsection{因子定义}

基于逐笔委托中的撤单数据，计算如下毒流动性因子：
\begin{equation}
    \text{毒流动性} = \frac{5\text{秒内撤单数量}}{30\text{秒内撤单数量}}
\end{equation}

\subsection{变量定义}
\begin{itemize}
    \item ① 可对计算的因子计算20日均值进行低频化处理。
\end{itemize}

\subsection{说明}
毒流动性因子从“高频撤单”角度来追踪机构交易行为（散户投资者通常不容易接触到程序化交易和高速交易通道，在订单簿中挂单时间越短的撤单，同样是机构行为的概率也越大）；高频撤单的相对比例越高，说明在盘口集中的撤单中，部分机构在股票中更多地起到了做市商的作用，然而这些流动性并非是可以获得的简单筹码，当趋势呈现反向时他们会迅速地改变原有交易方向，因此这些流动性也是“有毒的流动性”，高频流动性越多，实则股票交易的流动性越差。

\subsection{参考文献}
\begin{enumerate}
    \item 魏建榕, 苏良. 2024. 订单流系列：撤单行为规律初探. 开源证券.
\end{enumerate}

\section{相似低波(量价因子改进)}

\subsection{因子定义}
① 滚动计算历史收盘序列（长度为 120 个交易日）中任意时刻 \( T \) 起连续 \( \mathrm{RW}=6 \) 日序列同当前时刻 \( t \) 的股票收盘价序列的皮尔逊相关系数：
\begin{equation}
    \mathrm{Correlation}_{i,T} = \rho_{\mathrm{Sequence}_{i,t}, \mathrm{TempHistory}_{i,T}}
\end{equation}
\begin{equation}
    \mathrm{Sequence}_{i},t = [\mathrm{Close}_{i,t-\mathrm{RW}+1}, \mathrm{Close}_{i,t-\mathrm{RW}+2}, \dots, \mathrm{Close}_{i,t}]
\end{equation}
\begin{equation}
    \mathrm{TempHistory}_{i,T} = [\mathrm{Close}_{i,T}, \mathrm{Close}_{i,T+1}, \dots, \mathrm{Close}_{i,T+\mathrm{RW}-1}]
\end{equation}

② 设定相关系数的阈值 \( \mathrm{Threshold}=0.4 \) 筛选出走势相似度较高的时刻：
\begin{equation}
    \{ \tau \mid |\mathrm{Correlation}_{i,\tau}| \geq \mathrm{Threshold} \}
\end{equation}

③ 计算历史出现相似走势后既定持仓时间内股票累计超额收益率标准差的倒数即为相似低波因子：
\begin{equation}
    \mathrm{SimilarLowVolatility}_{i,t} = \frac{1}{\mathrm{std}(\mathrm{ER}_{i,\tau})}
\end{equation}
\begin{equation}
    \mathrm{ER}_{i,\tau} = \prod_{k=\tau+\mathrm{RW}}^{\tau+\mathrm{RW}+\mathrm{HoldingTime}-1} (1 + \mathrm{ER}_{i,k}) - 1, \quad 1 < \mathrm{HoldingTime} < \mathrm{RW}
\end{equation}
\begin{equation}
    \omega_{i,k} = \frac{e^{-\lambda k}}{\sum_{k=1}^{n} e^{-\lambda k}}, \quad \lambda = \frac{\ln(2)}{H}
\end{equation}

\subsection{说明}
同一只股票的历史收盘价序列中往往会出现多个走势相似的时刻，每次出现相似走势后，可以计算得到出现相似走势后股票一段时间内的累计超额收益率，多次相似走势后得到的累计超额收益率样本的波动率越低，未来继续持有该股票获得的超额收益率的期望越高。

\subsection{参考文献}
\begin{enumerate}
    \item 郑琳琳, 王天业. 2024. 基于历史相似走势的因子选股研究. 西南证券.
\end{enumerate}

\section{大单推动涨幅(高频因子-动量反转类)}

\subsection{因子定义}


\begin{equation}
    \prod_{n=t}^{t-T+1} \left( prod\left( 1 + r_{i,j,n} \cdot I_{\{j\in IdxSet\}}\right) \right) - 1
\end{equation}

\begin{equation}
    \text{平均单笔成交金额} = \frac{\sum_{j=1}^{N} \mathrm{Amt}_{i,j,n}}{\sum_{j=1}^{N} \mathrm{TrdNum}_{i,j,n}}
\end{equation}

\subsection{变量定义}
\begin{itemize}
    \item ① \( r_{i,j,n} \) 为股票 \( i \) 在第 \( n \) 个交易日内第 \( j \) 分钟的收益率序列；
    \item ② \( \mathrm{Amt}_{i,j,n} \) 为股票 \( i \) 在第 \( n \) 个交易日内第 \( j \) 分钟的成交额序列；
    \item ③ \( \mathrm{TrdNum}_{i,j,n} \) 为股票 \( i \) 在第 \( n \) 个交易日内第 \( j \) 分钟的成交笔数序列；
    \item ④ \( \mathrm{IdxSet} \) 为第 \( n \) 个交易日内平均单笔成交金额最大的 30\% 的分钟 \( K \) 线的序号；
    \item ⑤ 月度选股下，\( T=20 \) 个交易日；周度选股下，\( T=5 \) 个交易日。
\end{itemize}

\subsection{说明}
平均单笔成交金额较大的 \( K \) 线多空博弈激烈，未来的反转效应更强；该因子与股票未来收益负相关。

\subsection{参考文献}
\begin{enumerate}
    \item 冯佳睿，姚石.2019.日内分时交易中的玄机.海通证券
\end{enumerate}

\section{跌幅最大时刻-主卖笔均成交金额（高频因子-成交分布类）}

\subsection{因子定义}


1. 将日内 240 个 1 分钟数据集合标记为 \( T = \{t_1, t_2, \dots, t_{240}\} \)，判断股票 i 的每分钟涨跌幅，将跌幅最大的 10\% 时刻标记为 \( P = \{t_{i1}, t_{i2}, \dots, t_{ik}\} \)；

2. 利用逐笔成交数据，进一步对 P 中时刻的所有成交记录划分为主动买入成交 \( P\_Buy \) 和主动卖出成交 \( P\_Sell \) 两类；

3. 对跌幅最大时刻下的主动卖出成交额和成交笔数进行汇总，根据其成交额 \( \mathrm{Amt}_t \) 之和除以成交笔数 \( \mathrm{DealNum}_t \) 之和，构造跌幅最大时刻且主卖的笔均成交金额 \( \mathrm{ATD}_{\mathrm{DownTop10\%\_Sell}} \)：
\begin{equation}
    \mathrm{ATD}_{\mathrm{DownTop10\%\_Sell}} = \frac{\sum_{t \in P\_\mathrm{Sell}} \mathrm{Amt}_t}{\sum_{t \in P\_\mathrm{Sell}} \mathrm{DealNum}_t}
\end{equation}

4. 同理，将全天的成交金额除以全天的成交笔数，得到全天的笔均成交金额 \( \mathrm{ATD}_T \)：
\begin{equation}
    \mathrm{ATD}_T = \frac{\sum_{t \in T} \mathrm{Amt}_t}{\sum_{t \in T} \mathrm{DealNum}_t}
\end{equation}

5. 将跌幅最大时刻且主卖的笔均成交金额除以全天的笔均成交金额，即得到去量纲化后的标准化因子 \( \mathrm{SATD}_{\mathrm{DownTop10\%\_Sell}} \)：
\begin{equation}
    \mathrm{SATD}_{\mathrm{DownTop10\%\_Sell}} = \frac{\mathrm{ATD}_{\mathrm{DownTop10\%\_Sell}}}{\mathrm{ATD}_T}
\end{equation}

\subsection{说明}
笔均成交额类因子通过汇总“特殊的、具有较多信息含量”时刻的成交金额情况来刻画主力资金的行为动向；笔均成交金额越大，说明该特殊时间内资金优势越明显，其属于主力资金的可能性越大。“跌幅最大时刻”通过涨跌幅度来追踪主力资金抄底意愿强度，叠加逐笔订单主卖意愿做更细切分能进一步提升因子表现（卖方主导下，跌幅越大，主力抄底意愿越强，未来收益越高）。

\subsection{参考文献}
\begin{enumerate}
    \item 张欣悦, 张宇. 2025. 日内特殊时刻蕴含的主力资金 Alpha 信息. 国信证券.
\end{enumerate}

\section{改进大单交易占比（高级因子-资金流类）}

\subsection{因子定义}
\begin{equation}
    \text{改进大单交易占比} = \frac{(-1) \times \text{大买} \& \text{非大卖} + (-1) \times \text{非大买}\&\text{大卖} + \text{大买}\&\text{大卖}}{\text{总成交量}}
\end{equation}

\subsection{变量定义}
\begin{itemize}
    \item ① 大单和非大单的判断标准为：先将同一委买 ID（或委卖 ID）对应的实际成交量进行加总，再将全部实际成交的委买 ID（或委卖 ID）成交量进行降序排列，最后将成交量大于前 10\% 分位点的委买单（或委卖单）记为“大买单”（或“大卖单”），要剔除开盘集合竞价的成交记录。
\end{itemize}

\subsection{说明}
经测试发现对于“委买单为大单、委卖单为非大单”和“委买单为非大单、委卖单为大单”两种情况下的交易占比因子与未来股价负相关，只有“委买单及委卖单均为大单”的交易占比因子才与未来股价正相关，所以在计算大单交易占比因子时，对前面两种情况下做方向上的调整（通过乘 -1 将对未来股价的负向影响转为正向影响）有助于提升原始大单交易占比因子的表现。

\subsection{参考文献}
\begin{enumerate}
    \item 张欣鹏, 张宇. 2024. 高频订单成交数据蕴含的 Alpha 信息. 国信证券.
\end{enumerate}

\section{时间重心偏离（高频因子-收益分布类）}

\subsection{因子定义}
将涨、跌幅的时间重心单独剥离干扰因子（截面回归中心化处理）：
\begin{equation}
    G_u = f_u(\bar{R}_u, R_1, R_2, R_{\mathrm{overnight}}) + \varepsilon_u
\end{equation}
\begin{equation}
    G_d = f_d(\bar{R}_d, R_1, R_2, R_{\mathrm{overnight}}) + \varepsilon_d
\end{equation}

通过横截面回归，构造“时间差”指标，并取 \( N=20 \) 日均值作为因子：
\begin{equation}
    \varepsilon_d = \alpha + \beta \varepsilon_u + \varepsilon
\end{equation}
\begin{equation}
    \mathrm{TGD} = \frac{1}{N} \sum_{i=1}^N \varepsilon
\end{equation}

\subsection{变量定义}
\begin{itemize}
    \item ① \( G_u \) 和 \( G_d \) 分别为涨幅时间重心和跌幅时间重心，具体计算方式见相关日历页；
    \item ② \( \bar{R}_u \) 和 \( \bar{R}_d \) 分别表示涨、跌幅较大的 17 根分钟 Bar 的平均涨幅和平均跌幅；
    \item ③ \( R_1, R_2, R_{\mathrm{overnight}} \) 分别为 09:31-10:00 区间涨跌幅、10:01-10:30 区间涨跌幅、隔夜涨跌幅。
\end{itemize}

\subsection{说明}
涨、跌时间重心通过分钟涨跌幅对修正的时间戳加权平均得到，反映了股价在日内交易中上涨和下跌的位置，捕捉了股票的交易行为特征，两者的相对位置是一个有效的 Alpha 因子，与未来收益正相关；剥离收益率结构和极端收益反转效应等干扰因子，可以增强因子表现。

\subsection{参考文献}
\begin{enumerate}
    \item 魏建榕, 苏良, 张少成. 2022. 日内分钟收益率的时序特征：逻辑讨论与因子增强. 开源证券.
\end{enumerate}

\section{峰岭成交比（高频因子-成交分布类）}

\subsection{因子定义}

1. 对过去 20 日同时点成交量按照 1 倍标准差的标准划分，高于 1 倍标准差划分为“喷发成交量”，低于 1 倍标准差划分为“温和成交量”；

2. 对“喷发成交量”进一步按照连续与否划分：若当前分钟为“喷发成交量”，且前后 1 分钟均为“温和成交量”，则划为“孤立喷发成交量”，称为“量峰”；若当前分钟为“喷发成交量”，且前后 1 分钟也存在“喷发成交量”，则划为“连续喷发成交量”，称为“量岭”。将“温和成交量”称为“量谷”。

3. 计算 20 日量峰总成交额与量谷总成交额，两者做比即为峰岭成交比因子：
\begin{equation}
    \text{峰岭成交比} = \frac{\text{20日量峰总成交额}}{\text{20日量谷总成交额}}
\end{equation}

\subsection{说明}
峰岭成交比因子衡量了知情交易相对个人投资者交易的相对参与程度，因子值越高，未来收益越高。

\subsection{参考文献}
\begin{enumerate}
    \item 魏建榕, 王志豪. 2025. 高频成交量的峰、谷信息. 开源证券.
\end{enumerate}

\section{小买单主动成交度（高频因子-资金流类）}

\subsection{因子定义}
\begin{equation}
    \frac{1}{T} \sum_{n=t}^{t-T+1} \frac{\sum_{j=1}^{N} \text{小买单主动成交额}_{i,j,n}}{\sum_{j=1}^{N} \text{小买单成交额}_{i,j,n}}
\end{equation}

\subsection{变量定义}
\begin{itemize}
    \item ① 基于逐笔成交数据中的叫买与叫卖单号将逐笔成交数据合成为买卖单数据；
    \item ② 根据买卖单分布界定大小单，将小于多日买卖单成交单数据对数调整后的均值作为小单筛选阈值；
    \item ③ \( i, j, n \) 分别表示第 \( i \) 只股票在第 \( n \) 个交易日内的第 \( j \) 分钟的成交额；
    \item ④ 月度选股下，\( T=20 \) 个交易日；
    \item ⑤ 可将买单替换成卖单，即可计算得到小卖单主动成交度。
\end{itemize}

\subsection{说明}
主动成交度类因子刻画了逐笔买卖单中主动成交的程度；可以对大中小的买卖单都计算成交度，其中小单的相关因子表现最优，无论是小买单还是小卖单，主动成交度越高，股票未来收益表现更优。

\subsection{参考文献}
\begin{enumerate}
    \item 冯佳睿, 袁林青. 2021. 买卖单主动成交中的隐藏信息. 海通证券.
\end{enumerate}

\section{关注事件(行为金融因子-投资者注意力)}

\subsection{因子定义}
\begin{equation}
    \text{涨跌停事件数量} = \text{每月个股发生涨停、跌停的次数}
\end{equation}
\begin{equation}
    \text{龙虎榜事件数量} = \text{每月个股登上龙虎榜的次数}
\end{equation}

\subsection{变量定义}
\begin{itemize}
    \item ① 可以每期买入上述涨跌停事件或龙虎榜事件排名前 50\% 的股票构建组合，检验关注事件的选股能力。
\end{itemize}

\subsection{说明}
涨跌停股票、龙虎榜股票通常容易引起市场的广泛关注度，这些关注事件也可以作为投资者关注度的代理指标。

\subsection{参考文献}
\begin{enumerate}
    \item 陈升锐. 2024. 投资者有限关注及注意力捕捉与溢出. 中信建投.
\end{enumerate}

\section{成交额占比熵(高频因子-量价相关性类)}

\subsection{因子定义}
\begin{equation}
    \text{成交额占比熵} = H(\frac{amount_1}{AMOUNT},\frac{amount_2}{AMOUNT},\cdots,\frac{amount_N}{AMOUNT})
\end{equation}
\begin{equation}
    H(p_1, p_2, \dots, p_n) = -\sum_{i=1}^N p_i \ln(p_i)
\end{equation}

\subsection{变量定义}
\begin{itemize}
    \item ① \( \mathrm{amount}_t \) 为 \( t \) 时间段的成交额；
    \item ② \(\displaystyle \mathrm{AMOUNT} = \sum_{t=1}^T \mathrm{amount}_t\) 为整个时间段总成交额；
    \item ③ \( T \) 为划分时间段的个数，如日内 5 分钟频率的 \( T=48 \)。
\end{itemize}

\subsection{说明}
以成交额占比作为每一个时间段状态的发生概率，以熵的定义刻画成交的混乱程度，即得到成交额占比熵因子；当因子值越小时，体系越混乱，成交额占比彼此之间存在较大分歧，成交在价格高位较为密集，易造成股价的高估。

\subsection{参考文献}
\begin{enumerate}
    \item 覃川桃, 郑超. 2020. 高频因子（八）：高位成交因子——从量价匹配说起. 长江证券.
\end{enumerate}

\section{主力交易强度(高频因子-成交分布类)}

\subsection{因子定义}
\begin{equation}
    \mathrm{TS} = \mathrm{RankCorr}(A_{\mathrm{order}}, A_{\mathrm{minute}})
\end{equation}

\subsection{变量定义}
\begin{itemize}
    \item ① \( A_{\mathrm{order}} \) 为某只股票某个交易内的分钟单笔成交金额序列；
    \item ② \( A_{\mathrm{minute}} \) 为某只股票某个交易内的分钟成交金额序列；
    \item ③ 滚动 20 个交易日计算 \( \mathrm{TS} \) 指标的均值，即主力交易强度因子 \( \mathrm{MTS} \)。
\end{itemize}

\subsection{说明}
单笔成交金额与成交额之间的相关性强弱，描述的是代表主力的“相对大单”对分钟成交额的影响。从资金行为学的角度来看，单笔成交金额与成交额的相关性越强，主力资金主导成交节奏的能力也越强。

\subsection{参考文献}
\begin{enumerate}
    \item 姚建榕, 苏良. 2022. 高频因子：分钟单笔金额序列中的主力行为刻画. 开源证券.
\end{enumerate}

\section{复杂公司动量（行为金融因子-筹码分布）}

\subsection{因子定义}
\begin{equation}
    \mathrm{COMRET}_{i,t} = \frac{\sum_{j \in \mathrm{Ind}_i} W_{i,j,t} \mathrm{RET}_{j,t}}{\sum_{j \in \mathrm{Ind}_i} W_{i,j,t}}
\end{equation}

\subsection{变量定义}
\begin{itemize}
    \item ① \( W_{i,j,t} = \frac{\mathrm{Sales}_{i,j,t}}{\mathrm{Sales}_{i,t}} \) 为 \( j \) 行业营业收入在 \( i \) 公司营业收入中的占比；
    \item ② \( \mathrm{RET}_{j,t} \) 为只从事 \( j \) 行业的关联公司（有相同业务）在 \( t \) 时刻的月度收益。
\end{itemize}

\subsection{说明}
多元业务公司的信息处理复杂性相对较高，投资者可能无法及时对这些信息做出反应，这些复杂信息也就不容易及时反映在股价上；可以利用业务单一公司的股价表现来预测复杂公司的股价表现。

\subsection{参考文献}
\begin{enumerate}
    \item Cohen, L., \& Lou, D. (2012). \textit{Complicated firms}. Journal of Financial Economics, 104(2), 383-400.
\end{enumerate}

\section{高频上行波动占比（高频因子-波动跳跃类）}

\subsection{因子定义}
\begin{equation}
    \text{高频上行波动占比} = \frac{\sum_t(r_i^tI_{\{r_i^t>0\}})^2}{\sum_t(t_i^t)^2}
\end{equation}

\subsection{变量定义}
\begin{itemize}
    \item ① \( r_{i}^t \) 为分钟频率下的股票收益序列，如1分钟、5分钟、10分钟等；
    \item ② 在任意选股时刻，因子值为前 \( N \) 日指标的均值，如月度选股，计算因子值时使用的是股票过去20个交易日的均值。
\end{itemize}

\subsection{说明}
股票高频收益的上行波动衡量了股票价格拉升的特征。假设有两只股票在过去一段时间有着相同的涨幅，其中一只股票的涨幅由持续稳定的小幅上涨累计带动，而另一只股票的上涨源自于股票短期的大幅拉升，那么后者更有可能在收益上出现反转，而后者在因子值上也会体现出较高的上行波动率。

\subsection{参考文献}
\begin{enumerate}
    \item 冯佳睿, 袁林青. 2017. 高频因子之已实现波动分解. 海通证券.
\end{enumerate}

\section{开盘后净委买增额占比(高频因子-资金流类)}

\subsection{因子定义}
\begin{equation}
    \text{开盘后净委买增额占比}=\frac{1}{T} \sum_{n=t}^{t-T+1} \frac{\sum_{j \in 9:30-10:00} \text{净委买增额}_{i,j,n}}{ \text{成交额}_{i,j,n}}
\end{equation}

\subsection{变量定义}
\begin{itemize}
    \item ① 盘口委托快照数据包含买一至买十的委买量、卖一至卖十的委卖量等指标；
    \item ② 委托买单增加量为日内 \( t \) 至 \( t+1 \) 时刻各档的委买变化量之和，委托卖单增加量同理；
    \item ③ 一般用前 1 档位数据即可；使用的档位数量越多，因子的选股能力越弱；
    \item ④ \( i, j, n \) 分别表示第 \( i \) 只股票在第 \( n \) 个交易日内的第 \( j \) 分钟的数据；
    \item ⑤ 月度选股下，\( T=20 \) 个交易日；周度选股下，\( T=5 \) 个交易日。
\end{itemize}

\subsection{说明}
刻画了投资者的买入意愿，委买量的增加代表了投资者买入意愿的增强，委卖量的增加代表了投资者卖出意愿的增强；开盘后 30 分钟的数据更是包含了投资者对于前一天收盘后股票信息的集中反馈，集中反馈越正面，体现出的买入意愿越强，股票未来的超额收益表现越好。

\subsection{参考文献}
\begin{enumerate}
    \item 冯佳睿等. 2020. 高频因子的现实与幻想. 海通证券.
\end{enumerate}

\section{CSAD 模型(行为金融因子-羊群效应)}

\subsection{因子定义}
利用个股收益 \( r_i \) 和市场收益 \( R_m \)，计算股票收益的横截面绝对偏差：
\begin{equation}
    \mathrm{CSAD} = \frac{\sum_{i=1}^N |r_i - R_m|}{N}
\end{equation}

以 \( \mathrm{CSAD} \) 为 \( y \)，以市场收益绝对值和市场收益平方为 \( x \)，进行如下时序回归：
\begin{equation}
    \mathrm{CSAD} = \alpha + \beta_1 |R_{mt}| + \beta_2 |R_{mt}|^2 + \varepsilon
\end{equation}

\subsection{说明}
CSAD 模型对羊群效应的检验主要基于时序回归上市场收益平方项前回归系数正负及显著性；当 \( \beta_2 \) 显著小于 0 时，说明市场极端上涨或下跌时，CSAD 取值明显下降，投资者的投资行为越来越一致，市场上羊群效应越明显；当 \( \beta_1 \) 和 \( \beta_2 \) 都显著小于 0 时，说明市场羊群效应非常显著。

\subsection{参考文献}
\begin{enumerate}
    \item Chang, E. C., Cheng, J., \& Khorana, A. (2000). \textit{An Examination of Herd Behavior in Equity Markets: An International Perspective}. Journal of Banking \& Finance, 24(10), 1651-1679.
\end{enumerate}

\section{买盘深度（高频因子-流动性类）}

\subsection{因子定义}

\begin{equation}
    \mathrm{bid\_depth} = \sum \left| \frac{\mathrm{bv}_i \times \mathrm{mid}_{price}}{\mathrm{b}_i - \mathrm{last}_{mid} + \epsilon} \right|
\end{equation}

\subsection{变量定义}
\begin{itemize}
    \item ① 使用五档买单与当下时刻中间价的价差的倒数对买单的数量进行加权得到买盘深度；
    \item ② 当单边所有买单档位为空时，令市场深度为 0；
    \item ③ 可对上述因子日内序列求均值得到日度因子值，再进行 20 个交易日指数移动平均得到月度因子值。
\end{itemize}

\subsection{说明}
买盘深度统计了买单的挂单量情况，反映了买单盘口深度，五档买单中，对报价离中间价越远的买单的买量给予越小的权重；与未来收益负相关，因子取值越大，挂单量越多，盘口深度越大，流动性越高。也可对日内该因子序列求标准差，衡量其在当日均匀程度，与流动性负相关。

\subsection{参考文献}
\begin{enumerate}
    \item 胡骏骅等. 2021. 高频视角下的微观流动性与波动性. 中金公司.
\end{enumerate}

\section{隔夜跳空(高频因子-动量反转类)}

\subsection{因子定义}

\begin{equation}
    \mathrm{absRet}_{\mathrm{night}} = \sum_{t=1}^{20} \left| ln (\frac{\mathrm{Open}_t}{\mathrm{Close}_{t-1}}  )\right|
\end{equation}

\subsection{变量定义}
\begin{itemize}
    \item ① \( \mathrm{Open}_t \) 和 \( \mathrm{Close}_{t-1} \) 分别为 \( t \) 交易日的开盘价和上一日收盘价。
\end{itemize}

\subsection{说明}
隔夜跳空因子是隔夜收益率绝对值的累加（传统隔夜收益有弱动量效应，分组收益呈一定抛物线型），代表过去一段时间内隔夜累计跳空的幅度，与未来收益负相关。隔夜累计跳空幅度越大，未来收益越差，隔夜平开的股票未来收益越好。

\subsection{参考文献}
\begin{enumerate}
    \item 朱定豪. 2020. 昼夜分离：隔夜跳空与日内反转选股因子. 华安证券.
\end{enumerate}

\section{股票网络中心度(图谱网络-网络结构)}

\subsection{因子定义}
\begin{equation}
    \mathrm{CC}_{i,t} = \mathrm{SCC}_{i,t} \times 0.5 + \mathrm{TCC}_{i,t} \times 0.5
\end{equation}

\subsection{变量定义}
\begin{itemize}
    \item ① \( \mathrm{SCC}_{i,t} \) 为空间网络相对中心度，计算逻辑见相关日历页；
    \item ② \( \mathrm{TCC}_{i,t} \) 为时间网络相对中心度，计算逻辑见相关日历页。
\end{itemize}

\subsection{说明}
\( \mathrm{CC} \) 因子同时衡量了股票在网络中所处位置的中心程度和位置随时间推移的稳定程度，因子值越大，股票的网络中心度越高。

\subsection{参考文献}
\begin{enumerate}
    \item 张自力, 闫红蕾, 张楠. 2020. 股票网络、系统性风险与股票定价. 经济学, 1, 329-350.
    \item 曹春晓, 杨国平. 2021. 股票网络与网络中心度因子研究. 华西证券.
\end{enumerate}

\section{加权偏度(高频因子-收益分布类)}

\subsection{因子定义}
\begin{equation}
    \text{加权偏度} = \frac{ \sum_{t=1}^T w_t (\mathrm{close}_t - \overline{\mathrm{close}})^3}{\mathrm{close}_{\sigma}^3}
\end{equation}

\subsection{变量定义}
\begin{itemize}
    \item ① \( w_t = \frac{\mathrm{vol}_t}{VOL} \) 为正比于成交量的权重；
    \item ② \( \overline{\mathrm{close}} \) 为个股收盘价均值；
    \item ③ \( \mathrm{close}_{\sigma} \) 为个股收盘价标准差；
    \item ④ \( T \) 为划分时间段的个数，如日内 5 分钟频率的 \( T=48 \)。
\end{itemize}

\subsection{说明}
上述加权偏度是以成交量加权计算的价格分布的偏离程度，在计算偏度时，给成交量较高的时间段更大的权重；与未来收益负相关，负偏态越明显，个股成交集中在价格相对高位的水平，容易造成股价的高估。

\subsection{参考文献}
\begin{enumerate}
    \item 覃川桃, 郑起. 2020. 高频因子（八）：高位成交因子——从量价匹配说起. 长江证券.
\end{enumerate}

\section{“朝没晨雾”因子(高频因子-量价相关性类)}

\subsection{因子定义}

1. 对股票 \( i \)，用每日 1 分钟收盘价计算得到第 \( t \) 分钟收益率序列 \( \mathrm{ret}_t \)；用每日 1 分钟成交量计算得到 1 分钟增量成交量 \( \mathrm{voldiff}_t \)：
   \begin{equation}
       \mathrm{ret}_t = \frac{\mathrm{close}_t}{\mathrm{close}_{t-1}} - 1
   \end{equation}
   \begin{equation}
       \mathrm{voldiff}_t = \mathrm{volume}_t - \mathrm{volume}_{t-1}
   \end{equation}

2. 对每天第 6 分钟至第 240 分钟的上述数据进行带截距项的最小二乘回归：
   \begin{equation}
       \mathrm{ret}_t = \alpha_0 + \alpha_1 \mathrm{voldiff}_t + \alpha_2 \mathrm{voldiff}_{t-1} + \cdots + \alpha_6 \mathrm{voldiff}_{t-5} + \varepsilon_t
   \end{equation}

3. 记上述回归得到的第 \( t-1 \)、\( t-2 \)、\( t-3 \)、\( t-4 \)、\( t-5 \) 分钟的增量成交量的回归系数的 \( t \) 值分别为 \( t_1 \)、\( t_2 \)、\( t_3 \)、\( t_4 \)、\( t_5 \)，计算这 5 个 \( t \) 值的标准差，即为股票 \( i \) 在 \( T \) 日的“日朝没晨雾”因子；

4. 每月末，计算过去 20 个交易日的“日朝没晨雾”因子均值，即得到“朝没晨雾”因子。

\subsection{说明}
“朝没晨雾”因子衡量了每分钟的前 5 分钟的信息到来的平稳程度，与未来收益负相关；因子取值越小，表示该股票短时间内流入的信息越平稳，越没有突然到来的信息，表明该股票相对冷门，投资者对其的关注较少并且相对理智，因此其当前价格可能相对被低估，未来可能会产生较高收益。

\subsection{参考文献}
\begin{enumerate}
    \item 曹春晓. 推动个股价格变化的因素分解与“花隆林间”因子, 2023, 方正证券.
\end{enumerate}

\section{早盘主动净流入率稳定性(量价因子改进)}

\subsection{因子定义}
剔除当月早盘振幅最高的 \(1/5\) 标的，对剩余标的计算如下因子：
\begin{equation}
    \text{早盘主动净流入率稳定性因子}=\frac{\mathrm{mean}(\text{当月每日早盘主动净流入率}_{i,t})}{\mathrm{std}(\text{当月每日早盘主动净流入率}_{i,t})}
\end{equation}
其中，
\begin{equation}
    \text{早盘主动净流入率}_{i,t} = \frac{\text{开盘10:00前净主动买入额}_{i,t}}{\text{10:00前总成交额}_{i,t}}
\end{equation}

\subsection{说明}
该因子衡量了早盘主动资金流入的稳定性，可用于筛选主动资金近期具有稳定偏好的标的；使用早盘振幅是为了剔除非稳定性资金干扰较高的标的。

\subsection{参考文献}
\begin{enumerate}
    \item 王琦. 2023. JASON's alpha: 基本面+量价复合策略. 东北证券.
\end{enumerate}

\section{成交量支撑区域下限与收盘价差异（高频因子-成交分布类）}

\subsection{因子定义}

1. 计算同价成交量：将日内相同分钟收盘价的成交量累加至一起，得到在当前价格上的成交量总和；

2. 将成交量累计最大的价格定义为该股当日的成交量支撑点 \( VolumeSupportPrice \)，包含该支撑点的附近区域定义为成交量支撑区域 \( VolumeSupportArea \)；

3. 逐步计算 \( \mathrm{VSP} \) 周围的成交量累计值（价格由近及远，成交量由大及小），该累计值与全天成交量总和的比值超 50\% 的最小区域，即为成交量支撑区域 \( VSA \)，该区域上限价格为 \( \mathrm{VSA\_High} \)，下限价格为 \( \mathrm{VSA\_Low} \)；

4. 计算下限价格 \( \mathrm{VSA\_Low} \) 与当日收盘价之间的差异即得到 \( \mathrm{vsa\_ratio} \)：

\subsection{说明}
成交量支撑区域是判断个股日内公允价值区域的重要依据，过高/过低于该区域的价格将难以维持；若下限价格与当日收盘价之间差异较大，说明该股收盘后处于异常推动状态，当日大幅跌穿公允价值区域，其未来预计会回到公允价值区域，出现回弹上涨。

\subsection{参考文献}
\begin{enumerate}
    \item 郑兆磊. 2022. 成交量分布中的 Alpha. 兴业证券.
\end{enumerate}

\section{大的上行跳跃波动率(高频因子-波动跳跃类)}

\subsection{因子定义}
\begin{equation}
    \mathrm{RVLJP}_{\gamma,t} = \min\left(\mathrm{RVJP}_t, \sum_{i=1}^n r_{i,t}^2 \mathrm{I}_{\{r_{i,t} > \gamma\}}\right)
\end{equation}
\begin{equation}
    \mathrm{RVJP}_t = \max\left(\mathrm{RS}_t^+ - \frac{1}{2} \hat{IV}_t, 0\right)
\end{equation}
\begin{equation}
    \gamma = \alpha \Delta_n^{0.49} \sqrt{\hat{IV}_t}
\end{equation}

\subsection{变量定义}
\begin{itemize}
    \item ① \( \mathrm{RVJP}_t \) 为交易日 \( t \) 的上行跳跃波动率，计算细节可参考相关日历页；
    \item ② \( \gamma \) 为区分大、小跳跃的阈值，\( \Delta_n \) 为日内分钟数据个数的倒数，如 5 分钟频率的 \( \Delta_n = 1/48 \)，经验参数 \( \alpha = 4 \)；
    \item ③ \( r_{i,t} \) 为某股票交易日 \( t \) 内的分钟对数收益率序列，如 5 分钟对数收益率序列。
\end{itemize}

\subsection{说明}
跳跃波动率可以根据阈值进一步分解为大、小两个部分，即用于区分信息冲击导致的大跳跃波动和无法用信息冲击解释的小跳跃波动。上行波动率也可以做上述划分；由于大信息冲击更有可能改变投资者对股票基本面的预期，所以大跳跃波动通常比小跳跃波动有更强的收益预测能力。

\subsection{参考文献}
\begin{enumerate}
    \item Duong, D., \& Swanson, N. R. (2015). \textit{Empirical evidence on the importance of aggregation, asymmetry, and jumps for volatility prediction}. Journal of Econometrics, 187, 606-621.
    \item Bo Yu, Bruce Mizrach, Norman R. Swanson. (2020). \textit{New Evidence of the Marginal Predictive Content of Small and Large Jumps in the Cross-Section}. Econometrics, MDPI, 8(2), 1-52.
\end{enumerate}

\section{亏损卖出倾向(行为金融因子-处置效应)}

\subsection{因子定义}

1. 计算 \( t \) 时刻收盘价 \( \mathrm{Close}_t \) 与前 \( n \) 天 VWAP 均价 \( P_{t-n} \) 之间的涨跌幅，取涨跌幅为负的部分作为亏损收益率：
\begin{equation}
    \mathrm{loss}_t = \frac{\mathrm{Close}_t - P_{t-n}}{\mathrm{Close}_t} \cdot \mathbf{1}_{\{\mathrm{Close}_t < P_{t-n}\}}
\end{equation}

2. 以换手率 \( V_t \) 为权重加权亏损收益率得到资本亏损量即为亏损卖出倾向因子 \( \mathrm{Loss} \)：
\begin{equation}
    \mathrm{Loss}_t = \sum_{n=0}^{N-1} w_{t-n} \cdot \mathrm{loss}_{t-n}
\end{equation}
\begin{equation}
    w_{t-n} = \frac{1}{k} V_{t-n} \prod_{s=1}^{n-1} (1 - V_{t-n+s})
\end{equation}
\begin{equation}
    k = \sum_{n=1}^{T} V_{t-n} \prod_{s=1}^{n-1} (1 - V_{t-n+s})
\end{equation}

\subsection{说明}
亏损卖出倾向因子衡量了大部分投资者“未实现”的亏损，因子值越小，投资者未实现亏损越大，卖出意愿越强，股票出现超卖，股价偏低，未来股价更有可能上升，与未来收益负相关。

\subsection{参考文献}
\begin{enumerate}
    \item 陈升锐, 姚紫薇. 2025. 处置效应与 V 型处置效应在量化选股中的应用. 中信建投.
    \item An, L. (2016). \textit{Asset Pricing When Traders Sell Extreme Winners and Losers}. The Review of Financial Studies, 29(3), 823-861.
\end{enumerate}

\section{散户羊群效应(高毅因子-资金流类)}

\subsection{因子定义}
\begin{equation}
    \mathrm{RankCorr}(R_t, S_{t+1})
\end{equation}

\subsection{变量定义}
\begin{itemize}
    \item ① \(R_t\) 为过去 20 个交易日的日度收益率序列，用于表示市场涨跌；
    \item ② \(S_{t+1}\) 为小单（<4万元）过去 20 个交易日周期的资金净流入序列（前置 1 期），用于表示散户行为；
    \item ③ 使用 \(\frac{\mathrm{close}_t}{\mathrm{open}_t} - 1\) 日内收益，使用小单非主动净流入（开盘至 10 点之间），可以提升原始散户羊群效应的表现。
\end{itemize}

\subsection{说明}
衡量了散户追涨杀跌的程度，与股票未来收益负相关，即股票的散户羊群效应越高，未来预期收益越差。

\subsection{参考文献}
\begin{enumerate}
    \item 魏建榕, 盛少成. 2022. 新型因子：资金流动力学与散户羊群效应. 开源证券.
    \item 魏建榕, 盛少成. 2022. 大小单资金流 alpha 研究 2.0: 变量精筛与高频测算. 开源证券.
\end{enumerate}

\section{买单非流动性(高频因子-流动性类)}

\subsection{因子定义}
\begin{equation}
    r_{i,t} = \alpha + \beta_1 * S_{i,t} + \beta_2 * B_{i,t} + \epsilon_{it}
\end{equation}

\subsection{变量定义}
\begin{itemize}
    \item ① \( \beta_1 \) 为卖出非流动性系数；
    \item ② \( \beta_2 \) 为买入非流动性系数；
    \item ③ \( S_{i,t} \) 为股票 \( i \) 在 \( t \) 时间区间内的主动卖出金额；
    \item ④ \( B_{i,t} \) 为股票 \( i \) 在 \( t \) 时间区间内的主动性买入金额。
\end{itemize}

\subsection{说明}
买单非流动性衡量了高频数据下主动买入的交易金额对于股票价格变动的影响；买单非流动性对收益率的预测效果不及卖单非流动性，主要是由于投资者存在亏损厌恶的心理。

\subsection{参考文献}
\begin{enumerate}
    \item Brennan, M. J., Chordia, T., Subrahmanyam, A., \& Tong, Q. (2012). \textit{Sell-order liquidity and the cross-section of expected stock returns}. Journal of Financial Economics, 105(3), 523–541.
    \item 朱剑涛. 2017. 技术类新 Alpha 因子的批量测试. 东方证券.
\end{enumerate}

\section{年度搜索比 (downstream)(图谱网络-动量溢出)}

\subsection{因子定义}
\begin{equation}
    f_{ij}^t = \frac{\sum_{d=1}^{365} (\text{Daily Searches for } j \text{ after } i)_d}{\sum_{d=1}^{365} (\text{Daily Searches for } i \text{ and then any firm } j \neq i)_d}
\end{equation}
\begin{equation}
    R_{i,t} = \frac{\sum_{j=1}^N f_{ij}^t \times \mathrm{Ret}_j}{\sum_{j=1}^N f_{ij}^t}
\end{equation}

\subsection{变量定义}
\begin{itemize}
    \item ① \( f_{ij}^t \) 为年度搜索比、\( \mathrm{Ret}_j \) 为月度涨跌幅；
    \item ② 原始数据为从美国证券交易委员会（SEC）的 EDGAR 网站上获取的用户对公司信息的搜索流量；
    \item ③ 只保留同一天内至少搜索了两个不同公司（基于 CIK 编号识别的不同公司）的用户记录；
    \item ④ 剔除同一天内，用户对同一公司的连续搜索，对每个公司的搜索只计数一次；
    \item ⑤ 限制每个用户每天下载的独特公司数量不超过 50 家，以减少批量下载者的影响；
    \item ⑥ 仅考虑特定类型文件的页面浏览量，如 10-K、10-Q、8-K、14-A、S-1 等报告。
\end{itemize}

\subsection{说明}
年度搜索比例反映了投资者在搜索引擎中 \( i \) 之后，进一步搜索公司 \( j \) 的频率，借助投资者搜索行为捕捉投资者认为这两家公司之间存在经济关联的关联程度；投资者搜索行为有助于识别可能未被传统方法发现但具有密切经济联系的企业，而这些联系相比行业分类等传统关联对于解释股票回报、估值倍数等财务指标的变化更加有效。

\subsection{参考文献}
\begin{enumerate}
    \item Lee, Charles M. C., Ma, Paul, \& Wang, Charles C. Y. (2015). \textit{Search-Based Peer Firms: Aggregating Investor Perceptions Through Internet Co-Searches}. Journal of Financial Economics, 116(2), 410-431.
\end{enumerate}

\section{修正的日内反转(高频因子-动量反转类)}

\subsection{因子定义}

1. 对股票 \( i \)，计算 \( t \) 日的日内收益率 = \( t \) 日收盘价 / \( t \) 日开盘价 - 1；

2. 每月末计算当月的日内收益率和日内收益率的波动率，即月末最近 20 天的日内收益率序列的均值和标准差；

3. 对股票 \( i \)，若上述日内收益率的波动率小于横截面均值，则将其当月日内收益率取相反数（未来更可能发生动量效应），进而实现了反转的修正。

\subsection{变量定义}
\begin{itemize}
    \item ① 上述因子是借助波动率对反转做修正，也可借助换手率来做修正，即换手率变化量（\( t \) 日换手率与 \( t-1 \) 日换手率的差值）低于横截面均值，未来更可能发生动量；
    \item ② 可将日内收益率替换成日间收益率，对日间反转做修正。
\end{itemize}

\subsection{说明}
提高投资者对股价未来趋势判断的“可知性”，有助于更清晰的判断股票未来是否发生反转效应还是动量效应。波动率和换手率指标可代表股票的“可知性”，低波动（股价走势稳定）或低换手（投资分歧降低）的股票可知性更高，未来更可能发生动量效应；反之，可知性更低，未来更可能发生反转。通过将预判的动量做“翻转”，即可提高反转因子的表现。

\subsection{参考文献}
\begin{enumerate}
    \item 曹春晓. 2022. 个股动量效应的识别及“球状硬币”因子构建. 方正证券.
    \item Moskowitz, T. J. (2021). \textit{Asset pricing and sports betting}. Journal of Finance, Forthcoming.
\end{enumerate}

\section{价量相关性平均数因子(高频因子-量价相关性类)}

\subsection{因子定义}

① 每月月底，回溯每只股票过去 20 个交易日的价量信息，每日计算该股票分钟收盘价 \( p_t \) 与分钟成交量 \( v_t \) 的相关系数：
\begin{equation}
    \mathrm{PV\_corr} = \mathrm{corr}(v_t, p_t)
\end{equation}

② 每只股票取 20 日相关系数的平均值，做横截面市值中性化处理，得到的结果即为价量相关性平均数因子 \( \mathrm{PV\_corr\_avg} \)。

\subsection{说明}
\( \mathrm{PV\_corr\_avg} \) 利用成交量的信息，修正了传统反转因子对股票价格涨跌的判断，即价格涨跌的反转不完全由价格自己决定，还需看成交量的信息，若有量的确认，则月度行情的反转效应更强；若没有量的确认，则月度行情更容易呈现动量效应。

\subsection{参考文献}
\begin{enumerate}
    \item 高子剑. 2020. 高频价量相关性：意想不到的选股因子. 东吴证券.
\end{enumerate}

\section{活动筹码占比(行为金融因子-筹码分布)}

\subsection{因子定义}
\begin{equation}
    \mathrm{ASR}_t = \frac{\sum \mathrm{vol}_{\mathrm{price}_i \in [\text{低位价格}, \text{高位价格}],t}}{\sum vol_{i,t}}
\end{equation}
其中，
\begin{equation}
    \text{高位价格} = \mathrm{close}_t \times (1 + \mathrm{volatility})
\end{equation}
\begin{equation}
    \text{低位价格} = \mathrm{close}_t \times (1 - \mathrm{volatility})
\end{equation}

\subsection{变量定义}
\begin{itemize}
    \item ① 假设某只股票在 \( t \) 日的筹码分布共有 \( n \) 个价位（\( i = 1, 2, \ldots, n \)），\( \mathrm{vol}_{i,t} \) 为第 \( t \) 日 \( i \) 价位上的筹码峰高度，\( \mathrm{price}_{i,t} \) 为该筹码峰对应的价格；
    \item ② \( \mathrm{volatility} \) 为股价最大涨跌幅，主板股票的 \( \mathrm{volatility} = 10\% \)，创业板和科创板的 \( \mathrm{volatility} = 20\% \)；
    \item ③ \( \mathrm{close}_t \) 为 \( t \) 日收盘价。
\end{itemize}

\subsection{说明}
活动筹码指的是当天涨跌停范围内的筹码，这部分筹码更可能被主动交易，活动筹码占比即为位于涨跌停价格区间内的筹码占比。可用于衡量交易活跃程度；占比越高，筹码越集中在近期交易区间，流动性高，与未来收益负相关。

\subsection{参考文献}
\begin{enumerate}
    \item 陈升锐, 姚紫薇. 2025. 筹码分布因子系统构建. 中信出版社.
\end{enumerate}

\section{股价最低时刻笔均成交金额(高频因子-成交分布类)}

\subsection{因子定义}

1. 将日内 240 个 1 分钟数据集合标记为 \( T = \{t_1, t_2, \dots, t_{240}\} \)，将分钟收盘价最低的 10\% 时刻标记为集合 \( P = \{t_{i1}, t_{i2}, \dots, t_{ik}\} \)；

2. 对收盘价最低时刻进行汇总，根据其成交额 \( \mathrm{Amt} \) 之和除以成交笔数 \( \mathrm{DealNum} \) 之和，构造股价最低时刻的笔均成交金额 \( \mathrm{ATD}_{\mathrm{PriceLowest10\%}} \)：
\begin{equation}
    \mathrm{ATD}_{\mathrm{PriceLowest10\%}} = \frac{\sum_{t \in P} \mathrm{Amt}_t}{\sum_{t \in P} \mathrm{DealNum}_t}
\end{equation}

3. 同理，将全天的成交金额除以全天的成交笔数，得到全天的笔均成交金额 \( \mathrm{ATD}_T \)：
\begin{equation}
    \mathrm{ATD}_T = \frac{\sum_{t \in T} \mathrm{Amt}_t}{\sum_{t \in T} \mathrm{DealNum}_t}
\end{equation}

4. 将股价最低时刻的笔均成交金额除以全天的笔均成交金额，即得到去量纲化后的标准化因子 \( \mathrm{SATD}_{\mathrm{PriceLowest10\%}} \)：
\begin{equation}
    \mathrm{SATD}_{\mathrm{PriceLowest10\%}} = \frac{\mathrm{ATD}_{\mathrm{PriceLowest10\%}}}{\mathrm{ATD}_T}
\end{equation}

\subsection{说明}
笔均成交额类因子通过汇总“特殊的、具有较多信息含量”时刻的成交金额情况来刻画主力资金的行为动向；笔均成交金额越大，说明该特殊时间内资金优势越明显，其属于主力资金的可能性越大。“股价最低时刻”从股价高低角度追踪主力资金抄底意愿，股价越低，抄底意愿越强。

\subsection{参考文献}
\begin{enumerate}
    \item 张欣鹏, 张宇. 2025. 日内特殊时刻蕴含的主力资金 Alpha 信息. 国信证券.
\end{enumerate}

\section{大成交量已实现偏度(高频因子-收益分布类)}

\subsection{因子定义}
\begin{equation}
    \mathrm{real\_skew}_{D,i} = \frac{1}{T-1}\sum_{t=2}^{T}\frac{  \left( r_{t,D,i} - \bar{r}_{D,i} \right)^3}{\mathrm{real\_var}_{D,i}^{3/2}}
\end{equation}
\begin{equation}
    \mathrm{real\_var}_{D,i} = \frac{1}{T-2} \sum_{t=2}^{T} \left( r_{t,D,i} - \bar{r}_{D,i} \right)^2
\end{equation}

\subsection{变量定义}
\begin{itemize}
    \item ① 将个股在每个交易日的分钟成交量时间序列按照成交量大小排序，将分钟成交量排名前 \(1/3\) 的成交量定为“大成交量”；
    \item ② \(t=1,2,3,\ldots,T\) 为交易日内相应分钟频率的第 \(t\) 个时间区间，此处 \(T\) 为大成交量对应的那些区间；
    \item ③ \( r_{t,D,i} \) 为股票 \(i\) 在交易日 \(D\) 内大成交区间下的收益率，\(\bar{r}_{D,i}\) 为收益率均值。
\end{itemize}

\subsection{说明}
在传统已实现偏度计算逻辑基础上加入了成交量信息，大单成交更能代表主力行为，蕴含的信息更多。已实现偏度刻画了股价日内快速拉升或下跌的特征，与股票未来收益负相关；短期大幅拉升的股票，未来更有可能出现反转；而日内快速下跌的股票，往往具有更高的下行风险溢价。

\subsection{参考文献}
\begin{enumerate}
    \item 文巧钧, 安宁宁, 罗军. 2021. 高频量化数据因子化方法. 广发证券.
\end{enumerate}

\section{客户行业集中度(图谱网络-网络结构)}

\subsection{因子定义}


客户行业集中度 \( \mathrm{CH} \) 的计算公式如下：
\begin{equation}
    \mathrm{CH}_i = \sum_{i=1}^{I} w_{ij} \times H_j
\end{equation}

其中 \( H_j \) 为客户公司所属行业的赫芬达尔指数：
\begin{equation}
    H_j = \sum_{i=1}^{I} s_{ij}^2
\end{equation}

\subsection{变量定义}
\begin{itemize}
    \item ① \( H_j \) 为客户公司所属行业的赫芬达尔指数；
    \item ② \( s_{ij} \) 为属于行业 \( j \) 的公司在该行业的市场份额，市场份额可以通过营业收入或市值来计算；
    \item ③ \( w_{ij} \) 为供应商 \( i \) 对客户 \( j \) 的销售额占自身总销售额的占比。
\end{itemize}

\subsection{说明}
客户行业集中度为公司所有客户所属行业集中度的加权平均，反映了公司在其客户行业中的竞争地位；若客户行业集中度过高，意味着公司客户处于相对垄断地位，对客户的议价能力小，竞争激烈；若客户行业集中度过低，意味着公司可选客户较为分散，有更多的客户选择空间和议价能力。

\subsection{参考文献}
\begin{enumerate}
    \item Jing Wu, John R. Birge. (2014). \textit{Supply Chain Network Structure and Firm Returns}. SSRN Electronic Journal.
\end{enumerate}

\section{日内跳跌度(高频因子-波动跳跃类)}

\subsection{因子定义}
对股票 \( i \)，计算 \( n \) 日内 \( t \) 分钟的单利收益率、\( t \) 分钟的连续复利收益率：
\begin{equation}
    t \text{ 分钟单利收益率} = \frac{t \text{ 分钟收盘价}}{t-1 \text{ 分钟收盘价}} - 1
\end{equation}
\begin{equation}
    t \text{ 分钟连续复利收益率} = \ln\left(\frac{t \text{ 分钟收盘价}}{t-1 \text{ 分钟收盘价}}\right)
\end{equation}

计算 \( t \) 分钟的“单复利差”，进而计算“跳跌度”因子：
\begin{equation}
    t \text{ 分钟单复利差} = t \text{ 分钟单利收益率} - t \text{ 分钟连续复利收益率}
\end{equation}
\begin{equation}
    \text{跳跌度} = (t \text{ 分钟单复利差}) \times 2 - (t \text{ 分钟连续复利收益率})^2
\end{equation}

\subsection{变量定义}
\begin{itemize}
    \item ① 计算因子前需剔除开盘和收盘数据，仅使用日内数据；
    \item ② 日跳跌度为日内所有分钟跳跌度序列的均值；
    \item ③ 日跳跌度因子本质上为单利收益率 \( x \) 的泰勒残差项：
        \begin{equation*}
            e^x = 1 + \frac{1}{1!}x + \frac{1}{2!}x^2 + \frac{1}{3!}x^3 + o(x^3)
        \end{equation*}
        \begin{equation*}
            ((e^x - 1) - \frac{1}{1!}x) \times 2 - x^2 = \left(\frac{1}{2!}x^2 + \frac{1}{3!}x^3 + o(x^3)\right) \times 2 - x^2 = \left(\frac{1}{3!}x^3 + o(x^3)\right) \times 2
        \end{equation*}
    \item ④ 每月末可对最近 20 天的上述因子求均值和标准差并将两者等权合并，即得到月度跳跌度。
\end{itemize}

\subsection{说明}
跳跌度因子衡量了股价在利好或利空消息冲击下的跳跌程度，与未来收益负相关。正负向好的股票通常是非跳跌的平稳上涨；而那些当前突然跳跌上涨或跳下跌的股票，通常伴随投资者的过度反应，未来会出现补跌或补涨。

\subsection{参考文献}
\begin{enumerate}
    \item 曹春晓. 个股股价跳跃及其对振幅因子的改进, 2022, 方正证券.
    \item Jiang, G. J., \& Oomen, R. C. A. (2008). \textit{Testing for jumps when asset prices are observed with noise - a "swap variance" approach}. Journal of Econometrics, 144(2), 352-370.
\end{enumerate}

\section{漫长订单交易占比（高频因子-资金流类）}

\subsection{因子定义}
\begin{equation}
    \mathrm{VolumeLong} = \frac{\sum_i \mathrm{VolumeLongBuy}_i + \sum_i \mathrm{VolumeLongSell}_i}{\mathrm{Volume}}
\end{equation}

\subsection{变量定义}
\begin{itemize}
    \item ① 订单成交时长为订单第一次成交到最后一次成交间隔的连续竞价时长：
        \begin{equation*}
            \text{订单成交时长} = \text{最后第 } n \text{ 次成交时间 } T_n - \text{首次成交时间 } T_1
        \end{equation*}
    \item ② 漫长订单判断标准为：先计算不同委买 ID（或委卖 ID）订单的成交时长，再将全部实际成交的委买 ID（或委卖 ID）成交时长进行降序排列，最后将成交时长大于前 10\% 分位点的委买单（或委卖单）记为“漫长买单”（或“漫长卖单”），要剔除开盘集合竞价的成交记录；
    \item ③ \( \mathrm{VolumeLongBuy}_i \)、\( \mathrm{VolumeLongSell}_i \)、\( \mathrm{Volume} \) 分别为漫长买单实际成交量、漫长卖单实际成交量、单日总成交量。
\end{itemize}

\subsection{说明}
漫长订单交易占比因子从订单成交时长维度对逐笔订单进行划分和归整，与未来收益正相关。订单成交时长通常可以反映投资者分歧程度，分歧较大时，买卖单成交较为迅速，成交时长偏短；分歧较小时，买卖单成交较为缓慢，成交时长偏长。

\subsection{参考文献}
\begin{enumerate}
    \item 张欣慧, 张宇. 2024. 高频订单成交数据蕴含的 Alpha 信息. 国信证券.
\end{enumerate}

\section{买入浮亏占比小单因子（行为金融因子-遗憾规避）}

\subsection{因子定义}
\begin{equation}
    \mathrm{HCVOLS} = \frac{\sum_{i}^{N} \mathrm{volume}_{\mathrm{buy},i} \cdot \mathbf{I}_{P_{\mathrm{sell},i} > \mathrm{close}} \cdot \mathbf{I}_{\mathrm{vol} < \mathrm{vol}_{\mathrm{mean}}}}{\mathrm{total\_volume}}
\end{equation}

\subsection{变量定义}
\begin{itemize}
    \item ① \( \mathrm{volume}_{\mathrm{buy},i} \) 为主买成交量；
    \item ② \( P_{\mathrm{sell},i} \) 为主卖成交价；
    \item ③ \( \mathrm{vol}_{\mathrm{mean}} \) 为当日所有实际成交量均值；
    \item ④ 逐笔中每笔成交买卖方向确定方式：若该笔成交为买单被拆分与多个卖单成交，则记为买方交易，反之为卖方交易；
    \item ⑤ 大小单划分方式：以当日所有实际成交量均值为区分大小单的阈值，即
        \begin{equation*}
            \mathrm{vol} < \mathrm{vol}_{\mathrm{mean}}
        \end{equation*}
        为小单，反之为大单；
    \item ⑥ 此外，还可只统计尾盘期间 \([14:30,14:57)\) 的成交量，计算得到买入浮亏占比尾盘因子。
\end{itemize}

\subsection{说明}
遗憾规避理论认为非理性的投资者在做决策时，会倾向于避免产生后悔情绪并追求自满感，避免承认之前的决策失误；对于买入后浮亏（当天买入的股票收盘价低于买入的成本价）的投资者会避免承认决策失误而存在惜售心理，不愿卖出下跌浮亏的股票；买入浮亏占比因子值越高，高于收盘价买入的成交量占比越高，浮亏占比较高，未来抛压较低，会有更高的预期收益。

\subsection{参考文献}
\begin{enumerate}
    \item 高智威. 2023. 基于逐笔成交数据的遗憾规避因子. 国金证券.
\end{enumerate}

\section{有效价差(高频因子-流动性类)}

\subsection{因子定义}
\begin{equation}
    \mathrm{ES}_{i,t} = 2 D_{i,k} \frac{P_{i,t} - \left( \frac{\mathrm{Ask}_{i,t} + \mathrm{Bid}_{i,t}}{2} \right)}{\left( \mathrm{Ask}_{i,t} + \mathrm{Bid}_{i,t} \right)/2}
\end{equation}

\subsection{变量定义}
\begin{itemize}
    \item ① \( \mathrm{Ask}_{i,t}, \mathrm{Bid}_{i,t}, P_{i,t} \) 分别为股票 \( i \) 在 \( t \) 时刻的卖 1 价、买 1 价、收盘价；
    \item ② \( D_{i,k} \) 为买卖标识，买单记为 1，卖单记为 -1；
    \item ③ \( D_{i,k} \) 买卖单判断标准：若当前收盘价高于（低于）最优买卖报价平均值，则认为是买（卖）单；当前收盘价等于最优买卖报价平均值时，若收盘价高于（低于）上一笔价格，则认为是买（卖）单。
\end{itemize}

\subsection{说明}
有效价差表示成交价相对于市场中间价的相对偏差，反映执行交易所需要的交易成本；股票的交易成本越高，流动性越差，未来会有正向的流动性风险溢价。

\subsection{参考文献}
\begin{enumerate}
    \item 陈升悦. 2024. 流动性高频因子构建与投资者注意力因子分析报告. 中信建投.
\end{enumerate}

\section{喷发成交额跟随比例(高频因子-成交分布类)}

\subsection{因子定义}

① 对过去20日同时点成交量按照1倍标准差的标准划分，高于1倍标准差划为“喷发成交量”，低于1倍标准差划为“温和成交量”；

② 计算如下喷发成交额跟随比例因子：
\begin{equation}
    \text{喷发成交额跟随比例} = \frac{\text{过去20日喷发成交时点的下一分钟总成交额}}{\text{过去20日喷发成交时点的总成交额}}
\end{equation}

\subsection{说明}
喷发成交额跟随比例因子衡量了大额成交后的资金跟随程度，资金跟随程度越高，则个股交易中的个人投资者交易占比越高，个股未来表现越弱，是一个负向因子。

\subsection{参考文献}
\begin{enumerate}
    \item 魏建榕, 王志豪. 2025. 高频成交量的峰、岭、谷信息. 开源证券.
\end{enumerate}

\section{超大单买入占比(高频因子-资金流类)}

\subsection{因子定义}

\begin{equation}
    \text{超大单买入占比}_{i,t} = \frac{\mathrm{SuperBigBuy}_{i,t}}{\mathrm{SuperBigBuy}_{i,t} + \mathrm{SuperBigSell}_{i,t}}
\end{equation}

\subsection{变量定义}
\begin{itemize}
    \item ① 按照委买单号和委卖单号对 level2 逐笔成交数据中的成交金额（已成交部分）进行汇总求和得到当天股票 \( i \) 所有买单和卖单的成交金额；
    \item ② 对股票 \( i \)，将订单成交金额占当天总成交额比例超过 1\% 的订单定义为超大单（阈值在 0.5\%~2\% 区间较好）；
    \item ③ \( \mathrm{SuperBigBuy}_{i,t} \) 和 \( \mathrm{SuperBigSell}_{i,t} \) 分别为股票 \( i \) 在交易日 \( t \) 的所有超大买单、超大卖单的成交金额；
    \item ④ 因为超大单买卖金额在大票中不稳定，所以在计算全天、早盘超大单买入占比因子时，分子、分母均用 20 天超大单金额的均值替代单日的超大单金额。
\end{itemize}

\subsection{说明}
超大单刻画了潜在的过度投机和流动性冲击的负向 Alpha 影响，通常与未来收益负相关，超大买单占比越大，未来反转几率越大，小市值的大单反转更明显。

\subsection{参考文献}
\begin{enumerate}
    \item 朱剑涛, 王星星. 超大单冲击对大单因子的影响, 2022, 东方证券.
\end{enumerate}

\section{基于涨跌停的注意力捕捉(行为金融因子-投资者注意力)}

\subsection{因子定义}

将涨跌停股票的数量除以当日交易的股票数得到每日全市场涨跌停的比率作为市场关注代理指标：
\begin{equation}
    \mathrm{LIMIT\_PROP}_d = \frac{n_{\mathrm{upper},d} + n_{\mathrm{lower},d}}{N_d}
\end{equation}

将市场关注代理指标与股票日收益进行月度时序回归，代理指标新的系数绝对值 \(|\beta_{i,t}|\) 即为股票对市场关注的敏感程度：
\begin{equation}
    r_{i,d} = \mu_{i,t} + \beta_{i,t} \mathrm{LIMIT\_PROP}_d + \rho_{1,i,t} \mathrm{Mkt}_d + \rho_{2,i,t} \mathrm{SMB}_d + \rho_{3,i,t} \mathrm{HML}_d + \rho_{4,i,t} \mathrm{UMD}_d + \varepsilon_{i,t}
\end{equation}

\subsection{变量定义}
\begin{itemize}
    \item ① \(n_{\mathrm{upper},d}\)、\(n_{\mathrm{lower},d}\)、\(N_d\) 分别为 \(d\) 日涨停股票数量、跌停股票数量、当日交易股票总数；
    \item ② \(r_{i,d}\) 为个股日度收益；
    \item ③ \(\mathrm{Mkt}_d\)、\(\mathrm{SMB}_d\)、\(\mathrm{HML}_d\) 为经典的 Fama-French 三因子；
    \item ④ \(\mathrm{UMD}_d\) 为动量因子收益，即全市场过去 11 个月累积收益前后 30\% 的组合在 \(d\) 日的收益差。
\end{itemize}

\subsection{说明}
不同股票对同一市场关注度事件的反映程度往往是不同的，可借助个股收益对市场关注度事件的敏感程度来衡量股票对该事件的注意力捕捉效应；捕捉能力越强（敏感程度越高），表明个股更能吸引市场注意，从而导致买盘压力增大，股价高估；上述因子以涨跌停作为市场关注度事件，占比越高说明市场关注度越高。

\subsection{参考文献}
\begin{enumerate}
    \item Lin, F., Qiu, Z., \& Zheng, W. (2023). \textit{Cranes among chickens: The general-attention-grabbing effect of daily price limits in China's stock market}. Journal of Banking \& Finance, 150, 106818.
    \item 陈丹妮. 2024. 投资者有限关注及注意力捕捉与溢出. 中信出版社.
\end{enumerate}

\section{尾盘成交占比(高频因子-成交分布类)}

\subsection{因子定义}
\begin{equation}
    \frac{1}{T} \sum_{n=t}^{t-T+1} \frac{\mathrm{Vol}_{i,14:30-15:00,n}}{\mathrm{Vol}_{i,n}}
\end{equation}

\subsection{变量定义}
\begin{itemize}
    \item ① \( \mathrm{Vol}_{i,14:30-15:00,n} \) 为第 \( i \) 只股票在第 \( n \) 个交易日内 14:30-15:00 尾盘的成交量；
    \item ② \( \mathrm{Vol}_{i,n} \) 为第 \( i \) 只股票在第 \( n \) 个交易日的总成交量；
    \item ③ 月度选股下，\( T=20 \) 个交易日；周度选股下，\( T=5 \) 个交易日。
\end{itemize}

\subsection{说明}
尾盘成交量占比因子和股票未来收益负相关，其较好的选股效果可能源于：1. 尾盘投机度高，容易出现价格操纵行为；2. 非知情交易者（散户）不愿承担日内波动，更倾向于尾盘交易，而知情交易者（机构）则倾向于在早盘交易。

\subsection{参考文献}
\begin{enumerate}
    \item 冯佳睿等. 2020. 高频因子的现实与幻想. 海通证券.
\end{enumerate}

\section{订单失衡(高频因子-流动性类)}

\subsection{因子定义}
\begin{equation}
    \mathrm{VOI}_t = \Delta V_t^{\mathrm{WB}} - \Delta V_t^{\mathrm{WA}}
\end{equation}
其中，
\begin{equation}
    \Delta V_t^{\mathrm{WB}} = 
    \begin{cases}
        V_t^{\mathrm{WB}} , & P_t^{\mathrm{B}} > P_{t-1}^{\mathrm{B}} \\
        V_t^{\mathrm{WB}} - V_{t-1}^{\mathrm{WB}}, & P_t^{\mathrm{B}} = P_{t-1}^{\mathrm{B}} \\
        0, & P_t^{\mathrm{B}} < P_{t-1}^{\mathrm{B}}
    \end{cases}
\end{equation}
\begin{equation}
    \Delta V_t^{\mathrm{WA}} = 
    \begin{cases}
        V_t^{\mathrm{WA}}, & P_t^{\mathrm{A}} < P_{t-1}^{\mathrm{A}} \\
        V_t^{\mathrm{WA}} - V_{t-1}^{\mathrm{WA}}, & P_t^{\mathrm{A}} = P_{t-1}^{\mathrm{A}} \\
        0, & P_t^{\mathrm{A}} > P_{t-1}^{\mathrm{A}}
    \end{cases}
\end{equation}

\subsection{变量定义}
\begin{itemize}
    \item ① \( V^{\mathrm{WB}(\mathrm{WA})} = \sum w_i \times V_{i,t}^{\mathrm{B}(\mathrm{A})} \)，\( w_i = 1 - (i-1)/5, i=1,2,3,4,5 \) 为衰减加权委托量；
    \item ② \( P_t^{\mathrm{B}} \) 和 \( P_t^{\mathrm{A}} \) 分别为 \( t \) 时刻的买一价和卖一价；
    \item ③ \( V_t^{\mathrm{B}} \) 和 \( V_t^{\mathrm{A}} \) 分别为 \( t \) 时刻的买一和卖一的委托量。
\end{itemize}

\subsection{说明}
订单失衡因子为委托买入量与委托卖出量之差，衡量了盘口订单的不平衡程度，刻画了交易者的交易意图（知情交易者的信息差异使得订单不平衡）；订单失衡因子短期与收益率呈正相关，长期会出现反转（存在散户的追涨杀跌和主力的短期市场操纵）。

\subsection{参考文献}
\begin{enumerate}
    \item 丁鲁明, 陈升锐. 2020. 高频量价选股因子初探. 中信建投证券.
\end{enumerate}

\section{CSAD 模型（基于“小型”板块）(行为金融因子-羊群效应)}

\subsection{因子定义}

1. 每个月末提取当前可交易股票过去 120 个交易日的涨跌幅数据，计算股票 \( S_t \) 和其他股票（不活跃交易日不超过 20 天）之间日涨跌幅相关系数；

2. 选取相关系数最大的 9 只股票，连同 \( S_t \) 一起构成一个板块：
   \begin{equation}
       S_i \ (i = 1, \dots, 10)
   \end{equation}

3. 计算板块过去 120 个交易日的横截面绝对偏差 CSAD：
   \begin{equation}
       \mathrm{CSAD}_t = \frac{1}{10} \sum_{i=1}^{10} |r_{it} - \bar{r}_t|
   \end{equation}

4. 对 CSAD 进行时序标准化，并计算过去 20 个交易日的均值作为羊群效应因子：
   \begin{equation}
       Z_{jt} = \frac{\mathrm{CSAD}_{jt} - \mu_{\mathrm{CSAD}_j}}{\sigma_{\mathrm{CSAD}_j}}
   \end{equation}
   \begin{equation}
       f_{jt} = -\frac{1}{20} \sum_{k=1}^{20} Z_{j,t-k+1}
   \end{equation}

\subsection{说明}
CSAD 利用股票收益在横截面的分散度来衡量市场羊群效应的强弱；为了让因子对个股有选股能力，将与个股走势相近的股票组成“小型板块”，借助小板块上的羊群效应的强弱程度来近似反映该股票近期关注度的变化。因子值越大，板块近期的 CSAD 越小（即板块羊群效应越强），所在板块受关注程度提升越明显，未来收益越高。

\subsection{参考文献}
\begin{enumerate}
    \item Chang, E. C., Cheng, J., \& Khorana, A. (2000). \textit{An Examination of Herd Behavior in Equity Markets: An International Perspective}. Journal of Banking \& Finance, 24(10), 1651-1679.
    \item 丁鲁明, 陈元鹏. 2019. 基于市场羊群效应的股票 alpha 研究. 中信出版社.
\end{enumerate}

\section{耀眼波动率(高频因子-波动跳跃类)}

\subsection{因子定义}

1. 对股票 \( i \)，计算 t 日内 1 分钟成交量序列的环比差值，将成交量差值大于“该日差值序列均值+1倍标准差”的时刻定义为“成交量激增时刻 n”；

2. 利用分钟收盘价，计算股票 \( i \) 在 t 日内 1 分钟收益率序列；计算成交量激增时刻及其随后 4 分钟 \([n, n+4]\) 区间内 1 分钟收益率的标准差，即为“耀眼波动率”；

3. 将 t 日内所有耀眼波动率求均值即为“日耀眼波动率”；

4. 将各股票的“日耀眼波动率”减去该指标的横截面均值再取绝对值，即为“适度日耀眼波动率”。

\subsection{变量定义}
\begin{itemize}
    \item ① 计算因子需剔除开盘和收盘数据，仅考虑日内分钟频数据；
    \item ② 可对最近 20 个交易日的日度“适度日耀眼波动率”求均值和标准差，并等权合成，即可得“月耀眼波动率”因子。
\end{itemize}

\subsection{说明}
“耀眼波动率”衡量了日内成交量激增所引起的价格波动程度；“耀眼波动率”与未来收益负相关，若成交量激增后的价格波动较小，说明利好消息未被市场广泛认可，适度承担该风险能获得风险补偿；若引起的价格波动很大，可能代表投资者已反应过度，未来会补跌。

\subsection{参考文献}
\begin{enumerate}
    \item 曹春晓. 2022. 成交额激增时刻蕴含的 alpha 信息. 方正证券.
\end{enumerate}

\section{每笔成交量收益率相关性(高频因子-量价相关性类)}

\subsection{因子定义}
\begin{equation}
     \frac{\sum w_i \times ret_i}{\mathrm{std}(\mathrm{pvol}_i) \times \mathrm{std}(\mathrm{ret}_i)} = \mathrm{corr}(\mathrm{pvol}, \mathrm{ret})
\end{equation}

\begin{equation}
     w_i = \mathrm{pvol}_i - mean(\mathrm{pvol_i})
\end{equation}

\subsection{变量定义}
\begin{itemize}
    \item ① \(\mathrm{ret}_i\) 和 \(\mathrm{pvol}_i\) 分别为每个时间段的收益率和每笔成交量=成交量/成交笔数，时间段为分钟频率。
\end{itemize}

\subsection{说明}
每笔成交量收益率相关性本质为以每笔成交量与自身均值之差为权重对收益率序列进行加权，是利用每笔成交量对反转因子做了“提纯”，相关系数相比局部筛选保留了全时间段的所有信息。每笔成交量较大、呈现反转效应的时间段收益率的权重大于0，且成交量越大，权重越大，进而加强反转效应；反之，反向加强动量效应。

\subsection{参考文献}
\begin{enumerate}
    \item 覃川桃, 郑起. 2021. 高频因子（十）：量价关系中的反转微观结果. 长江证券.
\end{enumerate}

\section{海外运营信息(图谱网络-动量溢出)}

\subsection{因子定义}
\begin{equation}
    \mathrm{InfoProxy}_{i,j,t-1}{\mathrm{(Foreign)}} = \sum_{c\nep U.S.} f_{i,t-1}^c R_{j,t-1}^c
\end{equation}

\subsection{变量定义}
\begin{itemize}
    \item ① \( f_{i,t-1}^c \) 为跨国公司在 \( c \) 国家的销售额占比；
    \item ② \( R_{j,t-1}^c \) 为 \( c \) 国家的 \( j \) 行业回报（从 Datastream 全球股权部门指数中获取跨国公司运营所在国家的行业回报数据）。
\end{itemize}

\subsection{说明}
海外运营信息因子反映了跨国公司海外业务活动情况，跨国公司海外业务所对应的各国行业回报信息可以用来预测跨国公司自身股价的未来收益（因子方向为正）；投资者对海外运营信息的关注度较低，存在反应不足，进而导致股价未能及时反映跨国公司的实际经营状况变化，投资者可通过捕捉这些股价尚未完全反映的新信息而获得超额收益。

\subsection{参考文献}
\begin{enumerate}
    \item Huang, Xing. (2015). \textit{Thinking outside the borders: investors' underreaction to foreign operations information}. Review of Financial Studies, 28(11), 3109-3152.
\end{enumerate}

\section{开盘后买入意愿占比(高频因子-资金流类)}

\subsection{因子定义}


\begin{equation}
    \text{开盘后买入意愿占比} = \frac{1}{T} \sum_{n=t}^{t-T+1} \frac{\sum_{j=1}^{N} \text{买入意愿}_{i,j,n}}{\sum_{j=1}^{N} \text{成交额}_{i,j,n}}
\end{equation}
\begin{equation}
    \text{买入意愿}_{i,j,n} = \text{净主买成交额}_{i,j,n} - \text{净委买增额}_{i,j,n}
\end{equation}
\begin{equation}
    \text{净主买成交额}_{i,j,n} = \text{主动买入成交额}_{i,j,n} - \text{主动卖出成交额}_{i,j,n}
\end{equation}
\begin{equation}
    \text{净委买增额}_{i,j,n} = \text{委托买单增加量}_{i,j,n} - \text{委托卖单增加量}_{i,j,n}
\end{equation}

\subsection{变量定义}
\begin{itemize}
    \item ① 净主买成交额由逐笔成交数据计算得到，具体可参考开盘后净主买占比/强度等因子；
    \item ② 净委买增额由盘口委托快照数据计算得到，具体可参考开盘后净委买增额占比因子；
    \item ③ \( i, j, n \) 分别表示第 \( i \) 只股票在第 \( n \) 个交易日内的第 \( j \) 分钟的数据；
    \item ④ 开盘后的买入意愿占比用的是 9:30-10:00 范围内的数据；
    \item ⑤ 月度选股下，\( T=20 \) 个交易日；周度选股下，\( T=5 \) 个交易日。
\end{itemize}

\subsection{说明}
净委买变化额刻画了投资者还未释放的买入意愿，而净主买成交额则刻画了投资者已经释放的买入意愿，两者结合则得到广义的投资者主动买入意愿；开盘后 30 分钟内买入意愿占比越高，投资者的买入意愿越强。

\subsection{参考文献}
\begin{enumerate}
    \item 冯佳睿, 林青. 2020. 基于直观逻辑和机器学习的高频数据低频化应用. 海通证券.
\end{enumerate}

\section{P型成交量分布}

\subsection{因子定义}
高频因子-成交分布类

1. 计算同价成交量：将日内相同分钟收盘价的成交量累加至一起，得到在当前价格上的成交量总和；

2. 将成交量累计最大的价格定义为该股当日的成交量支撑点 \( \mathrm{VolumeSupportPrice} \)，包含该支撑点的附近区域定义为成交量支撑区域 \( \mathrm{VolumeSupportArea} \)；

3. 逐步计算 \( \mathrm{VSP} \) 周围的成交量累计值（价格由近及远，成交量由大及小），该累计值与全天成交量总和的比值超 50\% 的最小区域，即为成交量支撑区域 \( \mathrm{VSA} \)，该区域上限价格为 \( \mathrm{VSA\_High} \)，下限价格为 \( \mathrm{VSA\_Low} \)；

4. 计算下限价格 \( \mathrm{VSA\_Low} \) 与当日最高价之间的差异即得到 \( \mathrm{vsa\_low2max} \)：


\subsection{说明}
若下限价格与当日最高价越接近，说明成交量支撑区域越接近日内的高价区域，则该股日内的成交量分布越接近于 P 型成交量分布，预期上涨；当个股价格出现急剧上涨然后盘整，或短期急剧下跌但日内整体处于高位时，通常会出现 P 型成交量曲线。

\subsection{参考文献}
\begin{enumerate}
    \item 郑兆磊. 2022. 成交量分布中的 Alpha. 兴业证券.
\end{enumerate}

\section{客户节点度(图谱网络-网络结构)}

\subsection{因子定义}

基于供应链图谱网络计算公司节点的入度即为客户节点度因子：
\begin{equation}
    d^{\mathrm{(in)}}_i = \sum_{j (\neq i)} w_{ij}^{in}
\end{equation}

其中，
\begin{equation}
    w_{ij} = 
    \begin{cases}
        \mathbb{I}_{[i \in N(j)]} \text{边无权重}\\
        \frac{\mathrm{sales}_{ij}}{\mathrm{Total\_Procurement}_i}, & \text{边有权重}
    \end{cases}
\end{equation}

\subsection{变量定义}
\begin{itemize}
    \item ① 供应链图谱网络边的方向为供应商指向客户：\(\mathrm{supplier} \rightarrow \mathrm{customer}\)；
    \item ② \( w_{ij} \) 为公司 \( i \) 与其供应商 \( j \) 之间的边；若边无权重，取值为 1；若边有权重，取值为公司 \( i \) 对供应商 \( j \) 的采购额在公司 \( i \) 总采购额中的占比；
    \item ③ 因子可能存在板块和市值偏好，建议对其做市值行业中心化处理。
\end{itemize}

\subsection{说明}
若公司的供应商节点度较高（即公司有较多供应商），说明公司的供应商较为冗余，有助于减少单一供应商带来的风险，增强公司抗风险能力；但对投资者来说，供应商节点度较高的公司，因其风险较低，能从中得到的风险补偿也就相对较低。

\subsection{参考文献}
\begin{enumerate}
    \item Jussa et al. (2015). \textit{The Logistics of Supply Chain Alpha}. Deutsche Bank Markets Research.
\end{enumerate}

\section{小的下行跳跃波动率(高频因子-波动跳跃类q)}

\subsection{因子定义}
\begin{equation}
    \mathrm{RVSJN}_t = \mathrm{RVJN}_t - \mathrm{RVLJN}_{\gamma,t}
\end{equation}
\begin{equation}
    \mathrm{RVLJN}_{\gamma,t} = \min\left(\mathrm{RVJN}_t, \sum_{i=1}^n r_{i,t}^2 \mathrm{I}_{\{r_{i,t} < -\gamma\}}\right)
\end{equation}
\begin{equation}
    \mathrm{RVJN}_t = \max\left(\mathrm{RS}_t^- - \frac{1}{2}\hat{\mathrm{IV}}_t, 0\right)
\end{equation}
\begin{equation}
    \gamma = \alpha \Delta_n^{0.49} \sqrt{\hat{\mathrm{IV}}_t}
\end{equation}

\subsection{变量定义}
\begin{itemize}
    \item ① \( \mathrm{RVJN}_t \) 和 \( \mathrm{RVLJN}_{\gamma,t} \) 分别为交易日 \( t \) 的下行跳跃波动率和大的下行跳跃波动率，计算细节可参考相关日历页；
    \item ② \( \gamma \) 为区分大、小跳跃的阈值，\( \Delta_n \) 为日内分钟数据个数的倒数，如 5 分钟频率的 \( \Delta_n = 1/48 \)，经验参数 \( \alpha = 4 \)；
    \item ③ \( r_{i,t} \) 为某股票交易日 \( t \) 内的分钟对数收益率序列，如 5 分钟对数收益率序列。
\end{itemize}

\subsection{说明}
跳跃波动率可以根据阈值进一步分解为大、小两个部分，即用于区分信息冲击导致的大跳跃波动和无法用信息冲击解释的小跳跃波动。下行波动率也可以做上述划分；由于小跳跃波动通常由短期流动性风险导致，投资者预期并未改变，对未来收益的预测能力要弱于大跳跃波动。

\subsection{参考文献}
\begin{enumerate}
    \item Duong, D., \& Swanson, N. R. (2015). \textit{Empirical evidence on the importance of aggregation, asymmetry, and jumps for volatility prediction}. Journal of Econometrics, 187, 606–621.
    \item Bo Yu, Bruce Mizrach, Norman R. Swanson. (2020). \textit{New Evidence of the Marginal Predictive Content of Small and Large Jumps in the Cross-Section}. Econometrics, MDPI, 8(2), 1–52.
\end{enumerate}

\section{修正的隔夜反转(高频因子-动量反转类)}

\subsection{因子定义}

1. 对股票 \( i \)，计算 t 日的隔夜距离 = \(\mathrm{abs}(t\text{日隔夜收益率} - \text{横截面均值})\)；

2. 每月末计算当月的隔夜距离和隔夜距离波动率，即月末最近 20 天的隔夜距离序列的均值和标准差；

3. 对股票 \( i \)，若上述隔夜距离波动率小于横截面均值，则将其当月隔夜距离取相反数（未来更可能发生动量效应），进而实现了隔夜反转的修正。

\subsection{变量定义}
\begin{itemize}
    \item ① 上述因子是借助波动率对隔夜反转做修正，也可借助换手率距离来做修正，即换手距离 = \(\mathrm{abs}(t\text{日换手率变化量} - \text{其横截面均值})\) 低于其横截面均值，未来更可能发生动量，其中，换手率变化量 = \(t\text{日换手率} - (t-1)\text{日换手率}\)。
\end{itemize}

\subsection{说明}
提高投资者对股价未来趋势判断的“可知性”，有助于更清晰的判断股票未来是否发生反转效应还是动量效应。波动率和换手率指标可代表股票的“可知性”，低波动（股价走势稳定）或低换手（投资分歧降低）的股票可知性更高，未来更可能发生动量效应；反之，可知性更低，未来更可能发生反转。通过将预判的动量做“翻转”，即可提高反转因子的表现。由于隔夜涨跌为“两边差，中间好”的因子，还需要通过“均值距离化”来修正单调性。

\subsection{参考文献}
\begin{enumerate}
    \item 曹春晓. 2022. 个股动量效应的识别及“球队硬币”因子构建. 方正证券.
    \item Moskowitz, T. J. (2021). \textit{Asset pricing and sports betting}. Journal of Finance, Forthcoming.
\end{enumerate}

\section{上下影线/收盘价的标准差(量价因子改进)}

\subsection{因子定义}

每月末利用日度高开低收行情序列计算如下因子：
\begin{equation}
    \mathrm{std}\left( \frac{\text{每日K线的上影线长度} + \text{每日K线的下影线长度}}{\text{每日收盘价}} \right)
\end{equation}

\subsection{说明}
上下影线表现了多空博弈状态，通过选择上下影线总长度较短的标的（因子方向为负向），规避博弈剧烈、稳定性较差的股票。

\subsection{参考文献}
\begin{enumerate}
    \item 王琦. 2023. JASON's alpha：基本面+量价复合策略. 东北证券.
\end{enumerate}

\section{“午蔽古木”因子(高频因子-量价相关性类)}

\subsection{因子定义}

① 对股票 \( i \)，用每日 1 分钟收盘价计算得到第 \( t \) 分钟收益率序列 \( \mathrm{ret}_t \)；用每日 1 分钟成交量计算得到 1 分钟增量成交量 \( \mathrm{voldiff}_t \)：
\begin{equation}
    \mathrm{ret}_t = \frac{\mathrm{close}_t}{\mathrm{close}_{t-1}} - 1
\end{equation}
\begin{equation}
    \mathrm{voldiff}_t = \mathrm{volume}_t - \mathrm{volume}_{t-1}
\end{equation}

② 对每天第 6 分钟至第 240 分钟的上述数据进行带截距项的最小二乘回归：
\begin{equation}
    \mathrm{ret}_t = \alpha_0 + \alpha_1 \mathrm{voldiff}_t + \alpha_2 \mathrm{voldiff}_{t-1} + \cdots + \alpha_6 \mathrm{voldiff}_{t-5} + \varepsilon_t
\end{equation}

③ 记上述回归得到的截距项的 \( t \) 值为 \( t-{\mathrm{intercept}} \)，回归方程的 \( F \) 值为 \( F-{\mathrm{all}} \)；

④ 对股票 \( i \) 在 \( T \) 日的 \( t-{\mathrm{intercept}} \) 取绝对值，记为 \( \mathrm{abst}-{\mathrm{intercept}} \)；利用 \( F-{\mathrm{all}} \) 对 \( \mathrm{abst}-{\mathrm{intercept}} \) 做修正，即对 \( F-{\mathrm{all}} \) 小于其横截面均值的股票的 \( \mathrm{abst}-{\mathrm{intercept}} \) 乘以 -1，其余股票的 \( \mathrm{abst}-{\mathrm{intercept}} \) 保持不变，记得到“日午蔽古木”因子；

\subsection{说明}
“午蔽古木”因子衡量了市场信息和个股中长期基本面信息对价格的影响程度，且利用突然到来信息对因子做了修正，与未来收益负相关；因子值越小，表明近期突然到来的信息未对股价产生影响（\( F-{\mathrm{all}} \) 越小），且噪声对价格的推动力量也越小（\( \mathrm{abst}-{\mathrm{intercept}} \) 值越大），个股未来越容易产生高收益。

\subsection{参考文献}
\begin{enumerate}
    \item 曹春晓. 推动个股价格变化的因素分解与“花隐林间”因子. 2023, 方正证券.
\end{enumerate}

\section{盘口平均深度(高频因子-流动性类)}

\subsection{因子定义}
\begin{equation}
    avg_depth = \frac{\mathrm{av}_1 + \mathrm{bv}_1}{2}
\end{equation}

\subsection{变量定义}
\begin{itemize}
    \item ① \( \mathrm{av}_1 \) 和 \( \mathrm{bv}_1 \) 分别为盘口委托快照数据的卖一量和买一量，若挂单量为0，则令盘口深度为0；
    \item ② 可对上述因子日内序列求均值得到日度因子值，再进行20个交易日指数移动平均得到月度因子值，也可计算标准差。
\end{itemize}

\subsection{说明}
盘口平均深度反映了整体市场的深度，因子取值越大，市场整体挂单量越大，市场总体深度越大，市场流动性越高，与未来收益负相关。也可对日内该因子序列求标准差，衡量其在当日均匀程度，与流动性负相关。

\subsection{参考文献}
\begin{enumerate}
    \item 胡骏聪等. 2021. 高频视角下的微观流动性与波动性. 中金公司.
\end{enumerate}

\section{上方/下方活动筹码占比（行为金融因子-筹码分布）}

\subsection{因子定义}

上方活动筹码占比：
\begin{equation}
    \mathrm{ASR\_upper}_t = \frac{\sum \mathrm{vol}_{\mathrm{price}_{i} \in [\mathrm{close}_t, \text{高位价格}],t}}{\sum \mathrm{vol}_{i,t}}
\end{equation}

下方活动筹码占比：
\begin{equation}
    \mathrm{ASR\_lower}_t = \frac{\sum \mathrm{vol}_{\mathrm{price}_{i} \in [\text{低位价格},\mathrm{close}_t],t}}{\sum \mathrm{vol}_{i,t}}
\end{equation}

其中，
\begin{equation}
    \text{高位价格} = \mathrm{close}_t \times (1 + \mathrm{volatility})
\end{equation}
\begin{equation}
    \text{低位价格} = \mathrm{close}_t \times (1 - \mathrm{volatility})
\end{equation}

\subsection{变量定义}
\begin{itemize}
    \item ① 假设某只股票在 t 日的筹码分布共有 n 个价位（\( i = 1, 2, \ldots, n \)），\( \mathrm{vol}_{i,t} \) 为第 t 日 i 价位上的筹码峰高度，\( \mathrm{price}_{i,t} \) 为该筹码峰对应的价格；
    \item ② \( \mathrm{volatility} \) 为股价最大涨跌幅，主板股票的 \( \mathrm{volatility} = 10\% \)，创业板和科创板的 \( \mathrm{volatility} = 20\% \)；
    \item ③ \( \mathrm{close}_t \) 为 t 日收盘价。
\end{itemize}

\subsection{说明}
上方活动筹码占比指的是收盘价与最高价之间的筹码占比，对应上方买盘强度；下方活动筹码占比指的是收盘价与最低价之间的筹码占比，对应下方卖压强度。

\subsection{参考文献}
\begin{enumerate}
    \item 陈升锐, 姚紫薇. 2025. 筹码分布因子系统构建. 中信建投.
\end{enumerate}

\section{平均单笔流出金额占比(高频因子-资金流类)}

\subsection{因子定义}
\begin{equation}
    \frac{1}{T} \sum_{n=t}^{t-T+1} \frac{\sum_{j=1}^{N} \mathrm{Amt}_{i,j,n} \cdot \mathrm{I}_{r_{i,j,n} < 0} \Big/ \sum_{j=1}^{N} \mathrm{TrdNum}_{i,j,n} \cdot \mathrm{I}_{r_{i,j,n} < 0}}{\sum_{j=1}^{N} \mathrm{Amt}_{i,j,n} \Big/ \sum_{j=1}^{N} \mathrm{TrdNum}_{i,j,n}}
\end{equation}

\subsection{变量定义}
\begin{itemize}
    \item ① \( \mathrm{Amt}_{i,j,n} \) 为股票 \( i \) 在第 \( n \) 个交易日内第 \( j \) 分钟的成交额序列；
    \item ② \( r_{i,j,n} \) 为股票 \( i \) 在第 \( n \) 个交易日内第 \( j \) 分钟的收益率序列；
    \item ③ \( \mathrm{TrdNum}_{i,j,n} \) 为股票 \( i \) 在第 \( n \) 个交易日内第 \( j \) 分钟的成交笔数序列；
    \item ④ 月度选股下，\( T = 20 \) 个交易日；周度选股下，\( T = 5 \) 个交易日；
    \item ⑤ 还可计算平均单笔流入金额占比，只需将公式中的 \( \mathrm{I}_{\{r_{i,j,n} < 0\}} \) 改为 \( \mathrm{I}_{\{r_{i,j,n} > 0\}} \) 即可。
\end{itemize}

\subsection{说明}
平均单笔成交金额占比类因子可以捕捉日内大单的交易行为，大单交易会使得股票在当日的平均单笔成交金额占比在全市场中处于相对较高水平。平均单笔流出金额占比因子具有较强的选股能力，与股票未来收益正相关；在股票下跌时，如果单笔成交金额大，说明委买有大单，是一种抄底行为。

\subsection{参考文献}
\begin{enumerate}
    \item 冯佳睿, 姚石. 2019. 日内分时成交中的玄机. 海通证券.
\end{enumerate}

\section{喷发成交额相关性(高频因子-成交分布类)}

\subsection{因子定义}

① 对过去20日同时点成交量按照1倍标准差的标准划分，高于1倍标准差划为“喷发成交量”，低于1倍标准差划为“温和成交量”；

② 计算过去20日喷发成交时点分钟成交额与下一分钟成交额的相关系数：
\begin{equation}
    \mathrm{Corr}(\text{过去20日喷发成交时点分钟成交额序列}, \text{下一分钟成交额序列})
\end{equation}

\subsection{变量定义}
\begin{itemize}
    \item ① 还可以过去20日喷发成交时点分钟成交额为 \( x \)，以下一分钟成交额为 \( y \)，进行回归，回归系数记为喷发成交额跟随比例因子，用于衡量分钟成交额对大额成交的敏感度。
\end{itemize}

\subsection{说明}
喷发成交额相关性因子衡量了大额成交后的资金的协同程度，资金协同程度越高，则个股交易中的个人投资者交易跟随的越多，个股未来表现越弱，是一个负向因子。

\subsection{参考文献}
\begin{enumerate}
    \item 魏建榕, 王志豪. 2025. 高频成交量的峰、岭、谷信息. 开源证券.
\end{enumerate}

\section{喷发成交额相关性(高频因子-成交分布类)}

\subsection{因子定义}

① 对过去20日同时点成交量按照1倍标准差的标准划分，高于1倍标准差划为“喷发成交量”，低于1倍标准差划为“温和成交量”；

② 计算过去20日喷发成交时点分钟成交额与下一分钟成交额的相关系数：
\begin{equation}
    \mathrm{Corr}(\text{过去20日喷发成交时点分钟成交额序列}, \text{下一分钟成交额序列})
\end{equation}

\subsection{变量定义}
\begin{itemize}
    \item ① 还可以过去20日喷发成交时点分钟成交额为 \( x \)，以下一分钟成交额为 \( y \)，进行回归，回归系数记为喷发成交额跟随比例因子，用于衡量分钟成交额对大额成交的敏感度。
\end{itemize}

\subsection{说明}
喷发成交额相关性因子衡量了大额成交后的资金的协同程度，资金协同程度越高，则个股交易中的个人投资者交易跟随的越多，个股未来表现越弱，是一个负向因子。

\subsection{参考文献}
\begin{enumerate}
    \item 魏建榕, 王志豪. 2025. 高频成交量的峰、岭、谷信息. 开源证券.
\end{enumerate}

\section{年度搜索比例 (symmetric)(图谱网络-动量溢出)}

\subsection{因子定义}
\begin{equation}
    f_{ij}^t = \frac{\sum_{d=1}^{365} (\text{Daily Searches for } j \text{ after or before } i)_d}{\sum_{d=1}^{365} (\text{Daily Searches for } i \text{ and then any firm } j \neq i)_d}
\end{equation}
\begin{equation}
    R_{i,t} = \frac{\sum_{j=1}^N f_{ij}^t \times \mathrm{Ret}_j}{\sum_{j=1}^N f_{ij}^t}
\end{equation}

\subsection{变量定义}
\begin{itemize}
    \item ① \( f_{ij}^t \) 为年度搜索比、\( \mathrm{Ret}_j \) 为月度涨跌幅；
    \item ② 原始数据为从美国证券交易委员会（SEC）的 EDGAR 网站上获取的用户对公司信息的搜索流量；
    \item ③ 只保留同一天内至少搜索了两个不同公司（基于 CIK 编号识别的不同公司）的用户记录；
    \item ④ 剔除同一天内，用户对同一公司的连续搜索，对每个公司的搜索只计数一次；
    \item ⑤ 限制每个用户每天下载的独特公司数量不超过 50 家，以减少批量下载者的影响；
    \item ⑥ 仅考虑特定类型文件的页面浏览量，如 10-K、10-Q、8-K、14-A、S-1 等报告。
\end{itemize}

\subsection{说明}
对称的年度搜索比例反映了投资时间顺序在搜索公司 \( i \) 之后或之前紧接着搜索公司 \( j \) 的频率，借助投资者搜索行为捕捉投资者认为这两家公司之间存在经济关联的关联程度；投资者搜索行为有助于识别出可能未被传统方法发现但具有密切经济联系的企业，而这些联系相比行业分类等传统关联对于解释股票回报、估值倍数等财务指标的变化更加有效。

\subsection{参考文献}
\begin{enumerate}
    \item Lee, Charles M. C., Ma, Paul, \& Wang, Charles C. Y. (2015). \textit{Search-Based Peer Firms: Aggregating Investor Perceptions Through Internet Co-Searches}. Journal of Financial Economics, 116(2), 410-431.
\end{enumerate}

\section{高频下行波动占比(高频因子-波动跳跃类)}

\subsection{因子定义}
\begin{equation}
    \text{高频下行波动占比} = \frac{\sum_{t} (r_{i}^tI_{\{r_i^t<0\}})^2  }{\sum_{t} (r_{i}^t)^2}
\end{equation}

\subsection{变量定义}
\begin{itemize}
    \item ① \( r_{i}^t \) 为分钟频率下的股票收益序列，如1分钟、5分钟、10分钟等；
    \item ② 在任意选股时刻，因子值为前 \( N \) 日指标的均值，如月度选股，计算因子值时使用的是股票过去20个交易日的均值。
\end{itemize}

\subsection{说明}
高频下行波动占比与高频上行波动占比相对，衡量了股票价格下跌的特征，通常与未来收益正相关；日内快速下跌的股票，往往具有更高的下行风险溢价。

\subsection{参考文献}
\begin{enumerate}
    \item 冯佳睿, 袁林青. 2017. 高频因子之已实现波动分解. 海通证券.
\end{enumerate}

\section{日内信噪比(高频因子-收益分布类)}

\subsection{因子定义}
从日内分钟价格序列 \( P \) 中，利用经验模态分解 EMD 算法分解出价格序列中的噪音序列 \( \mathrm{IMF}_i(t) \) 和信号序列 \( r_n(t) \)：
\begin{equation}
    P(t) = \sum_{i=1}^n \mathrm{IMF}_i(t) + r(t)
\end{equation}

计算信号序列的标准差与噪声序列标准差的比值的对数，即为日内信噪比：
\begin{equation}
    \mathrm{SNR} = \log \left( \frac{\mathrm{std}(r_n(t))}{\mathrm{std}(P(t) - r_n(t))} \right)
\end{equation}

\subsection{变量定义}
\begin{itemize}
    \item ① EMD 算法分解详细步骤见参考文献；
    \item ② 通常分解至二或三层时，信号序列的平滑特性较为明显，能够接近地描述原始序列，通常选二层更优。
\end{itemize}

\subsection{说明}
日内信噪比因子与股票未来收益正相关，信噪比越大，股票未来收益越高；从行为金融学角度解释，股票存在虚假信息或噪声交易者越多，价格偏离内在价值越频繁或偏离度越高，未来收益越差。

\subsection{参考文献}
\begin{enumerate}
    \item 严佳炜, 朱定豪. 信息提取、寻找高质量反转因子. 华安证券.
\end{enumerate}

\section{筹码收益（量价因子改进）}

\subsection{因子定义}
利用过去 250 日的数据，对股票 \( i \) 在 \( T \) 日计算得到 \([T-250, T]\) 之间每日的留存筹码 \( \mathrm{RsdAmt}_{T-k}^T \) 和成交均价 \( \mathrm{Price}_{T-k}^{\mathrm{avg}} \)，即为该股票在 \( T \) 日的筹码分布估计：
\begin{equation}
    \mathrm{RsdAmt}_{T-k}^T = \mathrm{Amt}_{T-k} \cdot \prod_{T-k+1}^{T} (1 - \mathrm{Turnover}_t), \quad k \in (0, 250]
\end{equation}

计算最新日成交均价相对历史筹码的加权成本的相对收益，即为最新日的筹码收益因子：
\begin{equation}
    \mathrm{holding\_ret} = \frac{\sum_t \mathrm{Price}_t^{\mathrm{avg}} \times \mathrm{RsdAmt}_t^T}{\sum_t \mathrm{RsdAmt}_t^T}
\end{equation}

\subsection{变量定义}
\begin{itemize}
    \item ① \( \mathrm{Amt}_{T-k} \) 为股票在 \( T-k \) 日的成交额；
    \item ② \( \mathrm{Turnover}_t \) 为股票在 \( t \) 日的换手率。
\end{itemize}

\subsection{说明}
筹码收益因子衡量了当前价格相对于全部筹码历史成本的收益，指标越高代表当前赚钱效应越好，未来表现越弱，是一个反转类因子。

\subsection{参考文献}
\begin{enumerate}
    \item 魏建榕, 张翔. 2025. 筹码结构视角下的动量反转融合. 开源证券.
\end{enumerate}

\section{信噪比增强反转（高频因子-动量反转类）}

\subsection{因子定义}
\begin{equation}
    \text{加强反转}_i = \mathrm{Weight}_i \times \text{反转}_i
\end{equation}
\begin{equation}
    \mathrm{Weight}_i = \frac{\mathrm{SNR}_i - \min(\mathrm{SNR})}{\max(\mathrm{SNR}) - \min(\mathrm{SNR})}
\end{equation}

\subsection{变量定义}
\begin{itemize}
    \item ① \( \mathrm{SNR}_i \) 为股票 \( i \) 的日度信噪比因子，具体计算逻辑见相关日历页；
    \item ② \( \text{反转}_i \) 为股票 \( i \) 的 20 日反转因子。
\end{itemize}

\subsection{说明}
对于高信号趋势、低噪声扰动的股票的反转因子值可以给予更高的信噪比乘数，因为日内信号含量高的股票更有可能反转，而日内买卖力量博弈程度高的股票，未来方向相对更不确定；信噪比带来的增量 Alpha 信息可以修正传统反转因子多头失效的部分。

\subsection{参考文献}
\begin{enumerate}
    \item 严佳炜, 朱定豪. 信息提纯，寻找高质量反转因子. 华安证券.
\end{enumerate}

\section{价量相关性波动性因子（高频因子-量价相关性类）}

\subsection{因子定义}

1. 每月月底，回溯每只股票过去 20 个交易日的价量信息，每日计算该股票分钟收盘价 \( p_t \) 与分钟成交量 \( v_t \) 的相关系数：
\begin{equation}
    \mathrm{PV\_corr} = \mathrm{corr}(p_t, v_t)
\end{equation}

2. 每只股票取 20 日相关系数的标准差，做横截面市值中性化处理，得到的结果即为价量相关性波动性因子 \( \mathrm{PV\_corr\_std} \)。

\subsection{说明}
\( \mathrm{PV\_corr\_std} \) 从价量形态稳定性的角度，对反转因子进行了改进，即无论股票价格过去的涨跌，只要每日价量关系维持某种稳定形态，下个月就更有可能上涨；而价量关系在多种形态间频繁切换的股票，下个月更有可能下跌。

\subsection{参考文献}
\begin{enumerate}
    \item 高子剑. 2020. 高频价量相关性：意想不到的选股因子. 东吴证券.
\end{enumerate}

\section{百分比实现价差（高频因子-流动性类）}

\subsection{因子定义}
\begin{equation}
    \mathrm{PRS}_{i,t} = 2D_{i,k} \left( \ln(P_{i,t}) - \ln\left(\frac{\mathrm{Ask}_{i,t+5} + \mathrm{Bid}_{i,t+5}}{2}\right) \right)
\end{equation}

\subsection{变量定义}
\begin{itemize}
    \item ① \( \mathrm{Ask}_{i,t+5}, \mathrm{Bid}_{i,t+5} \) 分别为股票 \( i \) 在 \( t+5 \) 时刻的卖 \( 1 \) 价、买 \( 1 \) 价；
    \item ② \( P_{i,t} \) 为股票 \( i \) 在 \( t+5 \) 时刻的收盘价；
    \item ③ \( D_{i,k} \) 为买卖标识，买单记为 \( 1 \)，卖单记为 \( -1 \)；
    \item ④ \( D_{i,k} \) 买卖单判断标准：若当前收盘价高于（低于）最优买卖报价平均值，则认为是买（卖）单；当前收盘价等于最优买卖报价平均值时，若收盘价高于（低于）上一笔价格，则认为是买（卖）单。
\end{itemize}

\subsection{说明}
百分比实现价差属于“百分比 percent-cost”形式的流动性，代表执行交易所需要的交易成本；股票的交易成本越高，流动性越差，未来会有正向的流动性风险溢价。

\subsection{参考文献}
\begin{enumerate}
    \item 陈升悦. 2024. 流动性高频因子构建与投资者注意力因子分析报告. 中信建投.
\end{enumerate}

\section{卖出反弹占比小单因子（行为金融因子-遗憾规避）}

\subsection{因子定义}
\begin{equation}
    \mathrm{LCVOLS} = \frac{\sum_{i}^{N} \mathrm{volume}_{\mathrm{sell},i} \cdot \mathbf{I}_{\{P_{\mathrm{sell},i} < \mathrm{close}\}} \cdot \mathbf{I}_{\{\mathrm{vol} < \mathrm{vol}_{\mathrm{mean}}\}}}{\mathrm{total\_volume}}
\end{equation}

\subsection{变量定义}
\begin{itemize}
    \item ① \( \mathrm{volume}_{\mathrm{sell},i} \) 为主卖成交量；
    \item ② \( P_{\mathrm{sell},i} \) 为主卖成交价；
    \item ③ \( \mathrm{vol}_{\mathrm{mean}} \) 为当日所有实际成交量均值；
    \item ④ 逐笔中每笔成交买卖方向确定方式：若该笔成交为买单被拆分为多个卖单成交，则记为买方交易，反之为卖方交易；
    \item ⑤ 大小单划分方式：以当日所有实际成交量均值作为区分大小单的阈值，即
        \begin{equation*}
            \mathrm{vol} < \mathrm{vol}_{\mathrm{mean}}
        \end{equation*}
        为小单，反之为大单；
    \item ⑥ 此外，还可只统计尾盘期间 \([14:30,14:57)\) 的成交量，计算得到卖出反弹占比尾盘因子。
\end{itemize}

\subsection{说明}
遗憾规避理论认为非理性的投资者在做决策时，会倾向于避免产生后悔情绪并追求自满感，避免承认之前的决策失误；对于卖出后价格反弹（当天卖出的股票收盘价高于卖出的成本价）的投资者会倾向于继续坚持卖出判断，即使后续价格回升也不会考虑再次买回；卖出反弹占比因子值越高，低于收盘价卖出的成交量占比越高，股票未来买入动力越小，预期收益越低。

\subsection{参考文献}
\begin{enumerate}
    \item 高智威. 2023. 基于逐笔成交数据的遗憾规避因子. 国金证券.
\end{enumerate}

\section{早尾盘复合交易占比(高频因子-资金流类)}

\subsection{因子定义}

\begin{equation}
    \text{早尾盘复合交易占比} = \frac{\text{非早尾盘买单\&非早尾盘卖单的成交量} - \text{早盘买单\&早盘卖单的成交量}}{\text{总成交量}}
\end{equation}

\subsection{变量定义}
\begin{itemize}
    \item ① 早尾盘订单划分方式为：第一笔成交时间在 10:00 之前的订单为早盘订单，第一笔成交时间在 14:30 之后的订单为尾盘订单，其余时刻成交的订单为非早尾盘订单。
\end{itemize}

\subsection{说明}
早尾盘复合交易占比因子从订单早尾盘维度对逐笔订单进行划分和归整，意在追踪投资者早尾盘操纵行为。测试发现，“早盘买+早盘卖”成交量占比和“尾盘买+尾盘卖”成交量占比越高，其未来反转效应越明显，且前者比后者更有效；“非早尾盘买+非早尾盘卖”成交量占比越高，其未来的动量效应越强；早尾盘复合交易占比因子为两者的复合，与未来收益正相关。

\subsection{参考文献}
\begin{enumerate}
    \item 张欣慰, 张宇. 2024. 高频订单成交数据蕴含的 Alpha 信息. 国信证券.
\end{enumerate}

\section{股价最低时刻-主卖笔均成交金额(高频因子-成交分布类)}

\subsection{因子定义}

1. 将日内 240 个 1 分钟数据集集合标记为 \( T = \{t_1, t_2, \dots, t_{240}\} \)，将分钟收盘价最低的 10\% 时刻标记为集合 \( P = \{t_{i1}, t_{i2}, \dots, t_{ik}\} \)；

2. 利用逐笔成交数据，进一步对 \( P \) 中时刻的所有成交记录划分为主动买入成交 \( P\_Buy \) 和主动卖出成交 \( P\_Sell \) 两类；

3. 对股价最低时刻下的主动卖出成交额和成交笔数进行汇总，根据其成交额 \( \mathrm{Amt}_t \) 之和除以成交笔数 \( \mathrm{DealNum}_t \) 之和，构造股价最低时刻且主卖的笔均成交金额 \( \mathrm{ATD}_{\mathrm{PriceLowest10\%\_Sell}} \)：
\begin{equation}
    \mathrm{ATD}_{\mathrm{PriceLowest10\%\_Sell}} = \frac{\sum_{t \in P\_\mathrm{Sell}} \mathrm{Amt}_t}{\sum_{t \in P\_\mathrm{Sell}} \mathrm{DealNum}_t}
\end{equation}

4. 同理，将全天的成交金额除以全天的成交笔数，得到全天的笔均成交金额 \( \mathrm{ATD}_T \)：
\begin{equation}
    \mathrm{ATD}_T = \frac{\sum_{t \in T} \mathrm{Amt}_t}{\sum_{t \in T} \mathrm{DealNum}_t}
\end{equation}

5. 将股价最低时刻且主卖的笔均成交金额除以全天的笔均成交金额，即得到去量纲化后的标准化因子 \( \mathrm{SATD}_{\mathrm{PriceLowest10\%\_Sell}} \)：
\begin{equation}
    \mathrm{SATD}_{\mathrm{PriceLowest10\%\_Sell}} = \frac{\mathrm{ATD}_{\mathrm{PriceLowest10\%\_Sell}}}{\mathrm{ATD}_T}
\end{equation}

\subsection{说明}
笔均成交额类因子通过汇总“特殊的、具有较多信息含量”时刻的成交金额情况来刻画主力资金的行为动向；笔均成交金额越大，说明该特殊时间内资金优势越明显，其属于主力资金的可能性越大。“股价最低时刻”从股价高低角度追踪主力资金抄底意愿，叠加逐笔订单主卖意愿做更细切分能进一步提升因子表现（卖方主导下，股价越低，主力抄底意愿越强，未来收益越高）。

\subsection{参考文献}
\begin{enumerate}
    \item 张欣慰, 张宇. 2024. 高频订单成交数据蕴含的 Alpha 信息. 国信证券.
\end{enumerate}

\section{基于 PCA 的特质波动率(高频因子-波动跳跃类)}

\subsection{因子定义}

截面上对所有股票过去 20 天的日度收益率序列进行主成分分析 PCA，提取第一主成分 \( F_1 \) 为隐式因子溢价；

对每只股票 \( i \)，以过去 20 天的日度收益率序列为 \( y \)，以隐式因子 \( F_1 \) 序列为 \( x \)，进行时序回归：
\begin{equation}
    r_{i,t} = \alpha_{i,t} + \beta_1^{F_1} \times F_1 + \cdots + \varepsilon_i
\end{equation}

时序回归返回的残差序列的标准差即为隐式因子框架下的特质波动因子：
\begin{equation}
    \mathrm{IV} = \mathrm{std}(\varepsilon_i)
\end{equation}

\subsection{变量定义}
\begin{itemize}
    \item ① 在进行截面主成分分析时，若第一主成分的解释力度低于 20\%，则选取解释力最大的前两个主成分进行回归；
    \item ② 还可进一步计算特质因子：
        \begin{equation*}
            \mathrm{IVR} = 1 - R^2
        \end{equation*}
        上述回归模型拟合优度；
    \item ③ 还可进一步计算特质偏度因子：
        \begin{equation*}
            \mathrm{IS} = \frac{\sum \varepsilon_{it}^3}{\mathrm{IV}^3}
        \end{equation*}
\end{itemize}

\subsection{说明}
计算隐式因子框架下的特质波动率时，直接通过主成分分析对资产收益序列降维来提取共同波动，而无需显式的借助多因子模型估计共同风险，系统性的解决传统多因子模型遗漏变量、估计误差的问题，计算简单。

\subsection{参考文献}
\begin{enumerate}
    \item 张欣旭, 杨怡玲, 刘璐. 2022. 隐式框架下的特质类因子改进. 国信证券.
\end{enumerate}

\section{委托成交相关性(高频因子-量价相关性类)}

\subsection{因子定义}

使用盘口前1档委托挂单数据计算股票 \( i \) 在 \( T \) 日的净委买变化率：
\begin{equation}
    \text{净委买变化率}_{k,t}^T = \frac{\text{净委买变化量}_{k,t}^T}{\text{流通股本}_T}
\end{equation}
\begin{equation}
    \text{净委买变化量}_{k,t}^T = \sum_{j=1}^{k} \text{委买变化量}_{j,t}^T - \sum_{j=1}^{k} \text{委卖变化量}_{j,t}^T
\end{equation}

进一步计算 \( T \) 日高频收益序列与净委买变化率序列的相关系数：
\begin{equation}
    \text{委托成交相关性}^t_T = \mathrm{corr}(r_{T,t}^i, \mathrm{netBid}_{T,t}^i)
\end{equation}

\subsection{变量定义}
\begin{itemize}
    \item ① \( \text{委买变化量}_{j,t}^T \) 和 \( \text{委卖变化量}_{j,t}^T \) 分别为 \( T \) 日内 \( t \) 至 \( t+1 \) 时刻的第 \( j \) 档委买的变化量和第 \( j \) 档委卖的变化量；
    \item ② \( \text{流通股本}_T \) 为 \( T \) 日股票的流通股本；
    \item ③ 可过去5日、20日上述因子的均值作为周度和月度的因子值。
\end{itemize}

\subsection{说明}
委托成交相关性刻画了投资者买入意愿与股价走势间的关系，与未来收益负相关，前期委托成交相关性越低的股票，未来收益表现越好。其中“股价下跌，净委买上升”的股票在未来的超额收益最强（投资者托底意愿较强），而“股价下跌，净委买下降”的股票在未来的超额收益最弱（投资者托底意愿较弱）。

\subsection{参考文献}
\begin{enumerate}
    \item 冯佳奇, 袁林青. 2019. 当下跌遇到托底. 海通证券.
\end{enumerate}

\section{剔除超大单后的普通大单买入占比(高频因子-资金流类)}

\subsection{因子定义}

\begin{equation}
    \text{剔除超大单后的普通大单买入占比}_{i,t} = \frac{\mathrm{NormalBigBuy}_{i,t}}{\mathrm{NormalBigBuy}_{i,t} + \mathrm{NormalBigSell}_{i,t}}
\end{equation}

\subsection{变量定义}
\begin{itemize}
    \item ① 按照委买单号和委卖单号对 level2 逐笔成交数据中的成交金额（已成交部分）进行汇总求和得到当天股票 \( i \) 所有买单和卖单的成交金额；
    \item ② 对股票 \( i \)，将订单成交金额占当天总成交额比例超过 1\% 的订单定义为超大单（阈值在 0.5\%~2\% 区间较好），将所有订单成交金额的前 30\% 分位数的订单定义为大单；
    \item ③ \( \mathrm{NormalBigBuy}_{i,t} \) 和 \( \mathrm{NormalBigSell}_{i,t} \) 分别为股票 \( i \) 在交易日 \( t \) 的剔除超大单后大买单和大卖单成交金额。
\end{itemize}

\subsection{说明}
超大单买入因子具有负向 Alpha，剔除超大单后的普通大单买入占比的正向 Alpha 会在原始大单占比基础上得到加强。

\subsection{参考文献}
\begin{enumerate}
    \item 朱剑涛, 王星星. 超大单冲击对大单因子的影响, 2022, 东方证券.
\end{enumerate}

\section{基于沪深总成交额的注意力捕捉(行为金融因子-投资者注意力)}

\subsection{因子定义}

统计每日沪深两市的总成交额作为市场关注代理指标：
\begin{equation}
    \mathrm{mkt\_vol}_d = \text{沪市总成交额}_d + \text{深市总成交额}_d
\end{equation}

将市场关注代理指标与股票日收益进行月度时序回归，代理指标前的系数绝对值 \( |\beta_{i,t}| \) 即为股票对市场关注的敏感程度：
\begin{equation}
    r_{i,d} = \mu_{i,t} + \beta_{i,t} \mathrm{mkt\_vol}_d + \rho_{1,i,t} \mathrm{Mkt}_d + \rho_{2,i,t} \mathrm{SMB}_d + \rho_{3,i,t} \mathrm{HML}_d + \rho_{4,i,t} \mathrm{UMD}_d + \varepsilon_{i,t}
\end{equation}

\subsection{变量定义}
\begin{itemize}
    \item  \( n_{upper,d}、n_{lower,d}、N_d\)分别为d日涨停股票数量、跌停股票数量、当日交易股票总数
    \item  \( r_{i,d} \) 为个股日度收益；
    \item  \( \mathrm{Mkt}_d, \mathrm{SMB}_d, \mathrm{HML}_d \) 为经典的 Fama-French 三因子；
    \item  \( \mathrm{UMD}_d \) 为动量因子收益，即全市场过去 11 个月累积收益前后 30\% 的组合在 \( d \) 日的收益差。
\end{itemize}

\subsection{说明}
不同股票对同一市场关注度事件的反映程度往往是不同的，可借助个股收益对市场关注度事件的敏感程度来衡量股票对该事件的注意力捕捉效应；捕捉能力越强（敏感程度越高），表明个股更能吸引市场注意，从而导致买盘压力增大，股价高估；上述因子以两市成交额为市场关注度事件，成交额越高表明关注度越高。

\subsection{参考文献}
\begin{enumerate}
    \item Lin, F., Qiu, Z., \& Zheng, W. (2023). \textit{Cranes among chickens: The general-attention-grabbing effect of daily price limits in China's stock market}. Journal of Banking \& Finance, 150, 106818.
    \item 陈升锐. 2024. 投资者有限关注及注意力捕捉与溢出. 中信建投.
\end{enumerate}

\section{b 型成交量分布（高频因子-成交分布类）}

\subsection{因子定义}

1. 计算同价成交量：将日内相同分钟收盘价的成交量累加至一起，得到在当前价格上的成交量总和；

2. 将成交量累计最大的价格定义为该股当日的成交量支撑点 \( \mathrm{VolumeSopportPrice} \)，包含该支撑点的附近区域定义为成交量支撑区域 \( \mathrm{VolumeSopportArea} \)；

3. 逐步计算 \( \mathrm{VSP} \) 周围的成交量累计值（价格由近及远，成交量由大及小），该累计值与全天成交量总和的比值超 50\% 的最小区域，即为成交量支撑区域 \( \mathrm{VSA} \)，该区域上限价格为 \( \mathrm{VSA\_High} \)，下限价格为 \( \mathrm{VSA\_Low} \)；

4. 计算上限价格 \( \mathrm{VSA\_High} \) 与当日最低价之间的差异即得到 \( \mathrm{vsa\_high2min} \)：

\subsection{说明}
若上限价格与当日最低价越接近，说明成交量支撑区域越接近日内的低价区域，则该股日内的成交量分布越接近于 b 型成交量分布，预期下跌；当个股价格出现急剧下跌然后盘整，或者短期急剧上涨但整体处于低位时，成交量分布将出现 b 型曲线。

\subsection{参考文献}
\begin{enumerate}
    \item 郑兆森. 2022. 成交量分布中的 Alpha. 兴业证券.
\end{enumerate}

\section{凸显性收益(高频因子-动量反转类)}

\subsection{因子定义}
计算股票日收益率的凸显性系数 \( \sigma \)，\( r_{i,s} \) 为股票日收益率，\( \overline{r_s} \) 可选当日至全市场收益率中位数：
\begin{equation}
    \sigma(r_{i,s}, \overline{r_s}) = \frac{|r_{i,s} - \overline{r_s}|}{|r_{i,s}| + |\overline{r_s}| + \theta}
\end{equation}

根据计算的凸显性系数，在时序上对个股当月的日收益率进行排序，排名记作 \( k_{i,s} \)，令凸显性最强的排名等于 1，计算凸显性权重 \( \omega_{i,s} \)：
\begin{equation}
    \omega_{i,s} = \frac{\delta^{k_{i,s}}}{\sum_{s'} \delta^{k_{i,s'}} \times \pi_{s'}}
\end{equation}

计算股票收益率序列 \( r_{i,s} \) 和权重序列 \( \omega_{i,s} \) 的协方差即得到凸显性收益 \( \mathrm{STR} \)：
\begin{equation}
    \mathrm{STR}_i = \mathrm{cov}[\omega_{i,s}, r_{i,s}]
\end{equation}

\subsection{变量定义}
\begin{itemize}
    \item ① 参数 \( \theta \) 用于控制零收益的影响，可取 0.1；
    \item ② 参数 \( \delta \) 用于衡量投资者认知能力，取值范围在 (0,1] 之间，取值越小，投资者越关注凸显收益，可取 0.7；
    \item ③ \( \pi_{s'} \) 为日收益率发生的客观概率，可取窗口长度的倒数。
\end{itemize}

\subsection{说明}
凸显理论认为，投资者的注意力更容易被具有“凸显性”的收益所吸引；当股票 \( \mathrm{STR} \) 值较高时，其历史最高收益具有较高的凸显性，此时投资者过度关注其上升潜力，严重高估股价；\( \mathrm{STR} \) 值较低时，历史最低收益具有明显凸显性时，此时投资者过度强调其下行风险，严重低估股价。

\subsection{参考文献}
\begin{enumerate}
    \item Bordalo, P., Gennaioli, N., \& Shleifer, A. (2012). \textit{Salience Theory of Choice under Risk}. Quarterly Journal of Economics, 127(3), 1243-1285.
    \item 任健, 李元勋, 李世杰. 2022. 行为金融新视角——“凸显性收益”因子 STR. 招商证券.
\end{enumerate}

\section{修正后的年度搜索比例(图谱网络-动量溢出)}

\subsection{因子定义}

\begin{equation}
    f_{ij}^t = \frac{\sum_{d=1}^{365} (\text{Daily Searches for } j \text{ after } i)_d}{\sum_{d=1}^{365} (\text{Daily Searches for } i \text{ and then any firm } j \neq i)_d}
\end{equation}
\begin{equation}
    \hat{f}_{ij}^t = \frac{f_{ij}^t \equiv{\frac{ithenj}{i}}}{\frac{!ithenj}{!i}}
\end{equation}
\begin{equation}
    R_{i,t} = \frac{\sum_{j=1}^{N} \hat{f}_{ij}^t \times \mathrm{Ret}_j}{\sum_{j=1}^{N} \hat{f}_{ij}^t}
\end{equation}

\subsection{变量定义}
\begin{itemize}
    \item ① \( f_{ij}^t \) 为年度搜索比，\( \mathrm{Ret}_j \) 为月度涨跌幅；
    \item ② 原始数据为从美国证券交易委员会（SEC）的 EDGAR 网站上获取的用户对公司信息的搜索流量；
    \item ③ 只保留同一天内至少搜索了两个不同公司（基于 CIK 编号识别的不同公司）的用户记录；
    \item ④ 剔除同一天内，用户对同一公司的连续搜索，对每个公司的搜索只计数一次；
    \item ⑤ 限制每个用户每天下载的独特公司数量不超过 50 家，以减少批量下载者的影响；
    \item ⑥ 仅考虑特定类型文件的页面浏览量，如 10-K、10-Q、8-K、14-A、S-1 等报告。
\end{itemize}

\subsection{说明}
修正后的年度搜索比例中的分母表示：在没有搜索基础公司的情况下，紧接着搜索公司 \( j \) 的用户总数与没有搜索基础公司 \( i \) 的用户总数之比（在没有搜索公司 \( i \) 的情况下，公司 \( j \) 出现在其他任何公司的搜索序列中的平均概率）；目的是为了消除由于公司 \( j \) 自身受欢迎程度导致的偏差，从而更准确地反映公司 \( i \) 和 \( j \) 之间的经济相关性。

\subsection{参考文献}
\begin{enumerate}
    \item Lee, Charles M. C., Ma, Paul, \& Wang, Charles C. Y. (2015). \textit{Search-Based Peer Firms: Aggregating Investor Perceptions Through Internet Co-Searches}. Journal of Financial Economics, 116(2), 410-431.
\end{enumerate}

\section{大的上下行跳跃波动不对称性(高频因子-波动跳跃类)}

\subsection{因子定义}
\begin{equation}
    \mathrm{SRVLJ}_t = \mathrm{RVLJP}_{\gamma,t} - \mathrm{RVLJN}_{\gamma,t}
\end{equation}
\begin{equation}
    \mathrm{RVLJN}_{\gamma,t} = \min\left(\mathrm{RVJN}_t, \sum_{i=1}^n r_{i,t}^2 \mathrm{I}_{\{r_{i,t} < -\gamma\}}\right)
\end{equation}
\begin{equation}
    \mathrm{RVLJP}_{\gamma,t} = \min\left(\mathrm{RVJP}_t, \sum_{i=1}^n r_{i,t}^2 \mathrm{I}_{\{r_{i,t} > \gamma\}}\right)
\end{equation}

\subsection{变量定义}
\begin{itemize}
    \item ① \( \mathrm{RVLJP}_{\gamma,t} \) 和 \( \mathrm{RVLJN}_{\gamma,t} \) 分别为交易日 \( t \) 的大的上行跳跃波动率和大的下行跳跃波动率，计算细节可参考相关日历页。
\end{itemize}

\subsection{说明}
跳跃波动率可以根据阈值进一步分解为大、小两个部分，即用于区分信息冲击导致的大跳跃波动和无法用信息冲击解释的小跳跃波动；大的上下行跳跃波动不对称性主要对应大跳跃波动的不对称性。

\subsection{参考文献}
\begin{enumerate}
    \item Duong, D., \& Swanson, N. R. (2015). \textit{Empirical evidence on the importance of aggregation, asymmetry, and jumps for volatility prediction}. Journal of Econometrics, 187, 606-621.
    \item Bo Yu, Bruce Mizrach, Norman R. Swanson. (2020). \textit{New Evidence of the Marginal Predictive Content of Small and Large Jumps in the Cross-Section}. Econometrics, MDPI, 8(2), 1-52.
\end{enumerate}

\section{有效深度(高频因子-流动性类)}

\subsection{因子定义}
\begin{equation}
    \mathrm{effective\_depth} = \min(\mathrm{av}_1, \mathrm{bv}_1)
\end{equation}

\subsection{变量定义}
\begin{itemize}
    \item ① \( \mathrm{av}_1 \) 和 \( \mathrm{bv}_1 \) 分别为盘口委托快照数据的卖一量和买一量，若挂单量为0，则令盘口深度为0；
    \item ② 可对上述因子日内序列求均值得到日度因子值，再进行20个交易日指数移动平均得到月度因子值，也可计算标准差。
\end{itemize}

\subsection{说明}
通过取最优买卖单挂单量的最小者可以衡量市场实际的有效深度，因子取值越大，市场有效深度越大，市场流动性越高，与未来收益负相关。也可对日内该因子序列求标准差，衡量其在当日均匀程度，与流动性负相关。

\subsection{参考文献}
\begin{enumerate}
    \item 胡骏聪等. 2021. 高频视角下的微观流动性与波动性. 中金公司.
\end{enumerate}

\section{CSAD 模型（基于板块中地位）(行为金融因子-羊群效应)}

\subsection{因子定义}

1. 每个月末提取当前可交易股票过去 120 个交易日的涨跌幅数据，计算股票 \( S_* \) 和其他股票（不活跃交易日不超过 20 天）之间日涨跌幅相关系数；

2. 选取相关系数最大的 9 只股票，连同 \( S_* \) 一起构成一个板块：
   \begin{equation}
       S_i \ (i = 1, \dots, 10)
   \end{equation}

3. 以 \( S_* \) 为中心，计算板块过去 120 个交易日的横截面绝对偏差 \( \mathrm{CSAD} \)：
   \begin{equation}
       \mathrm{CSAD}_t = \frac{1}{10} \sum_{i=1}^{10} |r_{it} - r_{*t}|
   \end{equation}

4. 对 CSAD 进行时序标准化，并计算过去 20 个交易日的均值作为羊群效应因子：
   \begin{equation}
       Z_{jt} = \frac{\mathrm{CSAD}_{jt} - \mu_{\mathrm{CSAD}_j}}{\sigma_{\mathrm{CSAD}_j}}
   \end{equation}
   \begin{equation}
       f_{jt} = -\frac{1}{20} \sum_{k=1}^{20} Z_{j,t-k+1}
   \end{equation}

\subsection{说明}
CSAD 利用股票收益在横截面的分散度来衡量市场羊群效应的强弱；为了让因子对个股有选股能力，将与个股走势相近的股票组成“小型板块”，借助小板块上的羊群效应的强弱程度来近似反映该股票近期关注度的变化。因子值越大，中心股票收益越接近中位数，板块近期的 CSAD 越小（即板块羊群效应越强），所在板块受关注程度提升越明显，未来收益越高。

\subsection{参考文献}
\begin{enumerate}
    \item Chang, E. C., Cheng, J., \& Khorana, A. (2000). \textit{An Examination of Herd Behavior in Equity Markets: An International Perspective}. Journal of Banking \& Finance, 24(10), 1651-1679.
    \item 丁鲁明, 陈元鹏. 2019. 基于市场羊群效应的股票 alpha 研究. 中信出版社.
\end{enumerate}

\section{单一成交额占比熵(高频因子-量价相关性类)}

\subsection{因子定义}
单位一成交额占比熵 = \( H\left(\frac{\mathrm{vol}_1close_1 }{\mathrm{VOL} \times \mathrm{CLOSE}}, \frac{\mathrm{vol}_2close_2}{\mathrm{VOL} \times \mathrm{CLOSE}}, \dots, \frac{\mathrm{vol}_Nclose_N}{\mathrm{VOL} \times \mathrm{CLOSE}}\right) \)
\begin{equation}
    H(p_1, p_2, \dots, p_n) = -\sum_{i=1}^N p_i \ln(p_i)
\end{equation}

\subsection{变量定义}
\begin{itemize}
    \item ① \(\mathrm{close}_t\) 和 \(\mathrm{vol}_t\) 分别为 \(t\) 时间段的收盘价和成交量；
    \item ② \(\displaystyle \mathrm{VOL} = \sum_{t=1}^T \mathrm{vol}_t\)，为整个时间段总成交量；
    \item ③ \(\displaystyle \mathrm{CLOSE} = \sum_{t=1}^T \mathrm{close}_t\)，为整个时间段收盘价之和；
    \item ④ \(T\) 为划分时间段的个数，如日内 5 分钟频率的 \(T=48\)。
\end{itemize}

\subsection{说明}
以单位一成交额占比作为每一个时间段状态的发生概率，以熵的定义刻画成交的混乱程度，即得到单一成交额占比熵因子；当因子值越小时，体系越混乱，成交量占比和收盘价占比排序相对一致（单位一成交额占比彼此之间存在较大分歧），成交在价格高位较为密集，易造成股价的高估。

\subsection{参考文献}
\begin{enumerate}
    \item 周小川, 郑超. 2020. 高频因子(八): 高位成交因子——从量价匹配说起. 长江证券.
\end{enumerate}

\section{开盘后买入意愿强度(高频因子-资金流类)}

\subsection{因子定义}

\begin{equation}
    \frac{1}{T} \sum_{n=t}^{t-T+1} \frac{\mathrm{mean}(\text{买入意愿}_{i,j,n})}{\mathrm{std}(\text{买入意愿}_{i,j,n})}
\end{equation}

\begin{equation}
    \text{买入意愿}_{i,j,n} = \text{净主买成交额}_{i,j,n} - \text{净委买增额}_{i,j,n}
\end{equation}
\begin{equation}
    \text{净主买成交额}_{i,j,n} = \text{主动买入成交额}_{i,j,n} - \text{主动卖出成交额}_{i,j,n}
\end{equation}
\begin{equation}
    \text{净委买增额}_{i,j,n} = \text{委托买单增加量}_{i,j,n} - \text{委托卖单增加量}_{i,j,n}
\end{equation}


\subsection{变量定义}
\begin{itemize}
    \item ① 净主买成交额由逐笔成交数据计算得到，具体可参考开盘后净主买占比/强度等因子；
    \item ② 净委买增额由盘口委托快照数据计算得到，具体可参考开盘后净委买增额占比因子；
    \item ③ \( i, j, n \) 分别表示第 \( i \) 只股票在第 \( n \) 个交易日内的第 \( j \) 分钟的数据；
    \item ④ 开盘后的买入意愿占比用的是 9:30-10:00 范围内的数据；
    \item ⑤ 月度选股下，\( T=20 \) 个交易日；周度选股下，\( T=5 \) 个交易日。
\end{itemize}

\subsection{说明}
净委买变化额刻画了投资者还未释放的买入意愿，而净主买成交额则刻画了投资者已经释放的买入意愿，两者结合则得到广义的投资者主动买入意愿；开盘后 30 分钟内的买入意愿强度越高，投资者的买入意愿越稳健。

\subsection{参考文献}
\begin{enumerate}
    \item 冯佳睿, 袁林需. 2020. 基于直观逻辑和机器学习的高频数据低频化应用. 海通证券.
\end{enumerate}

\section{特异质换手波动率(量价因子改进)}

\subsection{因子定义}
每只股票在每个换仓日（如每个月初）取近1个月的日度数据进行如下时序ols回归：
\begin{equation}
    \mathrm{turnover}_{i,t} = \alpha_{i,t} + \beta_{1t} \mathrm{TMKT}_t + \beta_{2t} \mathrm{TSMB}_t + \beta_{3t} \mathrm{THML}_t + \beta_{4t} \mathrm{TMOM1}_{m,t} + 
 
     \varepsilon_{i,t}\mathrm{Ivol\_turnover}_{i,t} = \mathrm{std}(\varepsilon_{i,t})
\end{equation}

取最近一期回归的日度残差序列 \( \varepsilon_{i,t} \) 计算标准差，可得到 \( t \) 时刻近1个月股票的特异质换手波动率因子：
\begin{equation}
    \mathrm{Ivol\_turnover}_{i,t} = \mathrm{std}(\varepsilon_{i,t})
\end{equation}

\subsection{变量定义}
\begin{itemize}
    \item ① \( \mathrm{turnover}_{i,t} \)：\( t \) 时刻股票 \( i \) 基于自由流通股换手率的平均值；
    \item ② \( \mathrm{TMKT}_t \)：\( t \) 时刻全市场截面股票日度自由流通股换手率的平均值；
    \item ③ \( \mathrm{TSMB}_t \)：基于 \( t-1 \) 日总市值将全市场截面股票排序，分别计算处于前后 \( 1/3 \) 的股票池在 \( t \) 日的自由流通股日度换手率的均值，计算两者均值之差（后 \( 1/3 \) - 前 \( 1/3 \)）；
    \item ④ \( \mathrm{THML}_t \)：基于 \( t-1 \) 日 PB 值将全市场截面股票排序，分别计算处于前后 \( 1/3 \) 的股票池在 \( t \) 日的自由流通股日度换手率的均值，计算两者均值之差（前 \( 1/3 \) - 后 \( 1/3 \)）；
    \item ⑤ \( \mathrm{TMOM1}_{m,t} \)：基于 \( t-1 \) 日近一个月动量效应（21个交易日区间收益率）将全市场截面股票排序，分别计算处于前后 \( 1/3 \) 的股票池在 \( t \) 日的自由流通股日度换手率的均值，计算两者均值之差（前 \( 1/3 \) - 后 \( 1/3 \)）。
\end{itemize}

\subsection{说明}
借鉴特质波动率的构建理念，特异质换手波动率剥离了市场主流风格（风险因子）对于个股换手率的影响，可以更好的测算市场对错误定价的修正效率。

\subsection{参考文献}
\begin{enumerate}
    \item 郭策. 2024. 价值驱动因子. 中银证券.
\end{enumerate}

\section{日内异常交易量(高频因子-成交分布类)}

\subsection{因子定义}
\begin{equation}
    \mathrm{ATV}_{t,i} = \frac{\text{实际交易量}}{\text{预期交易量}} = \frac{t \text{日第} i \text{分钟交易量}}{\mathrm{mean}(\text{过去一段时间分钟交易量})}
\end{equation}


\subsection{变量定义}
\begin{itemize}
    \item ① 计算预期交易量的“过去一段时间”可以为过去1个交易日；
    \item ② 分钟交易量的频率可以为5分钟频率的成交数据。
\end{itemize}

\subsection{说明}
异常交易量通常用于捕捉股票的异常交易活动，与股票未来收益负相关；异常交易量较高时，股票价格会被高涨的非理性情绪化交易不断推高造成错误定价，未来股价会为此出现补跌。

\subsection{参考文献}
\begin{enumerate}
    \item 任瞳. 2023. “持续异常交易量”选股因子 PATV. 招商证券.
\end{enumerate}

\section{惊恐因子(高频因子-动量反转类)}

\subsection{因子定义}

1. 计算各股票的日度“惊恐度”：
\begin{equation}
    \text{惊恐度} = \frac{\text{偏离度}}{\text{基准项}} = \frac{\mathrm{abs}(\text{个股收益率} - \text{市场收益率})}{\mathrm{abs}(\text{个股收益率}) + \mathrm{abs}(\text{市场收益率}) + 0.1}
\end{equation}

2. 对各股票，将每日惊恐度与每日收益率相乘，得到加权调整收益率；

3. 对各股票，每月末计算过去20日加权调整收益率序列的均值和标准差，分别称为“惊恐收益”因子和“惊恐波动”因子，再将二者等权合成为“原始惊恐”因子。

\subsection{变量定义}
\begin{itemize}
    \item ① 选取中证全指 000985 代表市场水平；
    \item ② “偏离度”衡量了个股收益率对市场收益率的偏离水平；“基准项”衡量了市场总体收益率水平。
\end{itemize}

\subsection{说明}
原始惊恐因子以“惊恐度”为权重改进了传统反转因子。“惊恐度”刻画了极端收益对投资决策权重的扭曲程度；即使某股票相邻两天收益率相同，但相对市场水平的偏离程度却是不同的，即投资者的惊恐度（过度反应程度、收益显著性）是不同的，对此也应给出不同的投资决策权重。

\subsection{参考文献}
\begin{enumerate}
    \item 曹春晓. 显著效应、极端收益扭曲决策权重和“草木皆兵”因子, 2022, 方正证券.
    \item Cosemans, M., \& Frehen, R. (2021). \textit{Salience theory and stock prices: Empirical evidence}. Journal of Financial Economics, 140(2), 460-483.
\end{enumerate}

\section{供应商节点度(图谱网络-网络结构)}

\subsection{因子定义}
基于供应链图谱网络计算公司节点的出度即为供应商节点度因子：
\begin{equation}
    d_i^{(\mathrm{out})} = \sum_{j （\neq i）} w_{ij}^{\mathrm{out}}
\end{equation}

边无权重：
\begin{equation}
    w_{ij}^{\mathrm{out}} = \mathbb{1}_{[j \in N(i)]}
\end{equation}

边有权重：
\begin{equation}
    w_{ij}^{\mathrm{out}} = \frac{\mathrm{sales}_{ij}}{\mathrm{Total\_Sales}_i}
\end{equation}

\subsection{变量定义}
\begin{itemize}
    \item ① 供应链图谱网络边的方向为供应商指向客户： \(\mathrm{supplier} \rightarrow \mathrm{customer}\)
    \item ② \( w_{ij}^{\mathrm{out}} \) 为公司 \( i \) 与其客户 \( j \) 之间的边；若边无权重，取值为 1；若边有权重，取值为公司 \( i \) 对客户 \( j \) 的销售额在公司 \( i \) 总销售额中的占比；
    \item ③ 因子可能存在板块和市值偏好，建议对其做市值行业中心化处理。
\end{itemize}

\subsection{说明}
若公司的客户节点度较高（即公司有较多客户），说明公司的客户较为多样化，有助于减少单一客户带来的风险，增强公司抗风险能力；但对投资者来说，客户节点度较高的公司，因其风险较低，能从中得到的风险补偿也就相对较低。

\subsection{参考文献}
\begin{enumerate}
    \item Jussa et al. (2015). \textit{The Logistics of Supply Chain Alpha}. Deutsche Bank Markets Research.
\end{enumerate}

\section{订单失衡率(高频因子-流动性类)}

\subsection{因子定义}
\begin{equation}
    \mathrm{OIR}_t = \frac{V_t^{\mathrm{WB}} - V_t^{\mathrm{WA}}}{V_t^{\mathrm{WB}} + V_t^{\mathrm{WA}}}
\end{equation}
\begin{equation}
    V_t^{\mathrm{WB(WA)}} = \frac{\sum w_i \times V_{i,t}^{\mathrm{B(A)}}}{\sum w_i}
\end{equation}

\subsection{变量定义}
\begin{itemize}
    \item ① \( V_t^{\mathrm{WB(WA)}} \) 为衰减加权委托量；
    \item ② \( V_{t}^{\mathrm{B}} \) 和 \( V_{t}^{\mathrm{A}} \) 分别是 \( t \) 时刻第 \( i \) 档的买量和卖量；
    \item ③ \( w_i = 1 - (i-1)/5, \, i = 1, 2, 3, 4, 5 \)。
\end{itemize}

\subsection{说明}
订单失衡率因子为买卖委托量差与其和的比值，衡量了潜在买家和卖家在市场上的相对实力；订单失衡率为正，说明市场买压大于卖压，未来价格趋向上涨；订单失衡因子短期与收益率呈正相关，长期会出现反转（存在散户的追涨杀跌和主力的短期市场操纵）。

\subsection{参考文献}
\begin{enumerate}
    \item 丁鲁明, 陈开锐. 2020. 高频溢价选股因子初探. 中信建投证券.
\end{enumerate}

\section{已实现波动率(高频因子-波动跳跃类)}

\subsection{因子定义}
\begin{equation}
    \mathrm{RV}_t = \sum_{i=1}^N r_{t,i}^2
\end{equation}

\subsection{变量定义}
\begin{itemize}
    \item ① \( r_{t,i} \) 为某股票交易日 \( t \) 内的分钟对数收益率序列，如 5 分钟对数收益率序列；
    \item ② \( N \) 表示该股票在交易日 \( t \) 中特定频率的分钟级别数据个数，如 5 分钟频率通常有 48 个数据点。
\end{itemize}

\subsection{说明}
作者提出的上述方法是利用高频数据估计波动率的重要研究成果；非参数估计，无需构建模型，计算简单，且在一般条件下几乎无测量误差。

\subsection{参考文献}
\begin{enumerate}
    \item Andersen, T. G., Bollerslev, T., Diebold, F. X., \& Labys, P. (2001). \textit{The distribution of exchange rate volatility}. Journal of the American Statistical Association, 96, 42–55.
\end{enumerate}

\section{改进的客户动量因子（基于锚定偏差）(图谱网络-动量溢出)}

\subsection{因子定义}
利用供应链的相关数据，计算如下客户动量因子：
\begin{equation}
    \mathrm{cmmom}_i^{1M} = \sum_{j=1}^{N_i} w_{ij}^{\mathrm{sales}} \cdot \mathrm{mom}_j^{1M}, \quad i = 1, 2, \ldots, N
\end{equation}

计算最近 52 周最高价因子作为衡量“锚定偏差”的代理变量：
\begin{equation}
    \mathrm{nearness\ to\ 52\_week\ high} = \frac{\text{月末收盘价}}{\text{过去52周内最高价}}
\end{equation}

先基于客户动量因子（正向）进行分组，在各组内基于 nearness to 52-week high 因子（正向）进行分组，以此构建投资组合。

\subsection{变量定义}
\begin{itemize}
    \item ① \( \mathrm{mom}_j^{1M} \) 为公司 \( i \) 的客户 \( j \) 过去一个月收益率；
    \item ② \( w_{ij}^{\mathrm{sales}} \) 为销售占比。
\end{itemize}

\subsection{说明}
客户动量因子可以从“投资者有限注意力或信息反应不足”这个维度来解释，投资者对客户信息反应越不足，客户动量效应可能越强；“锚定偏差”会延缓投资者对新信息（包括客户等各种关联公司信息）的反应，将其与客户动量结合使用会提升客户动量因子的表现。

\subsection{参考文献}
\begin{enumerate}
    \item Cohen, L., \& Frazzini, A. (2008). \textit{Economic Links and Predictable Returns}. Journal of Finance, 63(4), 1977-2011.
    \item Huang, S., Lin, T.-C., \& Xiang, H. (2021). \textit{Psychological barrier and cross-firm return predictability}. Journal of Financial Economics, 142(1), 338-356.
\end{enumerate}

\section{“夜眠霜路”因子(高频因子-量价相关性类)}

\subsection{因子定义}

1. 对股票 \( i \)，用每日 1 分钟收盘价计算得到第 \( t \) 分钟收益率序列 \( \mathrm{ret}_t \)；用每日 1 分钟成交量计算得到 1 分钟增量成交量 \( \mathrm{voldiff}_t \)：
\begin{equation}
    \mathrm{ret}_t = \frac{\mathrm{close}_t}{\mathrm{close}_{t-1}} - 1
\end{equation}
\begin{equation}
    \mathrm{voldiff}_t = \mathrm{volume}_t - \mathrm{volume}_{t-1}
\end{equation}

2. 对每天第 6 分钟至第 240 分钟的上述数据进行带截距项的最小二乘回归：
\begin{equation}
    \mathrm{ret}_t = \alpha_0 + \alpha_1 \mathrm{voldiff}_t + \alpha_2 \mathrm{voldiff}_{t-1} + \cdots + \alpha_6 \mathrm{voldiff}_{t-5} + \varepsilon_t
\end{equation}

3. 记上述回归得到的截距项的 \( t \) 值为 \( t-{\mathrm{intercept}} \)，回归方程的 \( F \) 值为 \( F-{\mathrm{all}} \)；

4. 每月末，分别计算每只股票过去 20 天 \( t-{\mathrm{intercept}} \) 序列与当期横截面所有股票过去 20 天的 \( t-{\mathrm{intercept}} \) 序列之间的相关系数的绝对值，并取均值，即为“夜眠霜路”因子。

\subsection{说明}
“夜眠霜路”因子剥离了 \( t-{\mathrm{intercept}} \) 中的市场层面的信息与个股中长期的基本面信息，与未来收益正相关；因子值越大，表明该股票的其他信息中，与其余所有股票的其他信息共同的部分越多（共同部分即市场信息），即该股票的市场层面的信息占其他信息的比重越大，个股中长期的基本面信息占比就越小（投资者分歧较小），个股未来越容易产生高收益。

\subsection{参考文献}
\begin{enumerate}
    \item 曹春晓. 推动个股价格变化的因素分解与“花旗株间”因子, 2023, 方正证券.
\end{enumerate}

\section{大单资金净流入率(高频因子-资金流类)}

\subsection{因子定义}


\begin{equation}
    \frac{1}{T} \sum_{n=t}^{t-T+1} \frac{\sum_{j=1}^{N} \mathrm{Amt}_{i,j,n} \cdot \mathrm{I}_{\{r_{i,j,n}>0, j \in \mathrm{IdxSet}\}} - \sum_{j=1}^{N} \mathrm{Amt}_{i,j,n} \cdot \mathrm{I}_{\{r_{i,j,n}<0, j \in \mathrm{IdxSet}\}}}{\sum_{j=1}^{N} \mathrm{Amt}_{i,j,n}}
\end{equation}

其中，平均单笔成交金额为：
\begin{equation}
    \text{平均单笔成交金额} = \frac{\sum_{j=1}^{N} \mathrm{Amt}_{i,j,n}}{\sum_{j=1}^{N} \mathrm{TrdNum}_{i,j,n}}
\end{equation}

\subsection{变量定义}
\begin{itemize}
    \item ① \( \mathrm{Amt}_{i,j,n} \) 为股票 \( i \) 在第 \( n \) 个交易日内第 \( j \) 分钟的成交额序列；
    \item ② \( r_{i,j,n} \) 为股票 \( i \) 在第 \( n \) 个交易日内第 \( j \) 分钟的收益率序列；
    \item ③ \( \mathrm{IdxSet} \) 为第 \( n \) 个交易日内平均单笔成交金额最大的 30\% 的分钟 \( K \) 线的序号；
    \item ④ \( \mathrm{TrdNum}_{i,j,n} \) 为股票 \( i \) 在第 \( n \) 个交易日内第 \( j \) 分钟的成交笔数序列；
    \item ⑤ 月度选股下，\( T=20 \) 个交易日；周度选股下，\( T=5 \) 个交易日。
\end{itemize}

\subsection{说明}
平均单笔成交金额较大的 \( K \) 线多空博弈激烈，未来的反转效应更强；该因子与股票未来收益负相关。

\subsection{参考文献}
\begin{enumerate}
    \item 冯佳睿, 姚石. 2019. 日内分时成交中的玄机. 海通证券.
\end{enumerate}

\section{已实现盈利比率(行为金融因子-处置效应)}

\subsection{因子定义}

1. 提取股票分钟级前复权收盘价、成交量，以及日频流通换手率数据，计算股票每日分钟级量价分布情况，其中 \( t \) 日股票 \( i \) 在每个价位的筹码峰计算公式如下：
\begin{equation}
    \mathrm{Chip}_{i,t,p} = \mathrm{Vol}_{i,t,p} \times T_{i,t} + \mathrm{Chip}_{i,t-1,p} \times (1 - T_{i,t})
\end{equation}

2. 基于 \( t \) 日换手和 \( t-1 \) 日筹码峰，计算当天换手的筹码的盈利总金额作为已实现盈利：
\begin{equation}
    \text{已实现盈利}_{i,t} = \sum_p (\mathrm{Close}_{i,t} - \mathrm{Price}_{i,p}) \times (\mathrm{Chip}_{i,t-1,p} \times T_{i,t}) \times \mathrm{I}_{\mathrm{Close}_{i,t} \geq \mathrm{Price}_{i,p}}
\end{equation}

3. 基于 \( t \) 日筹码峰，计算当天收盘时所有盈利筹码的账面盈利金额作为账面盈利总金额：
\begin{equation}
    \text{账面盈利}_{i,t} = \sum_p (\mathrm{Close}_{i,t} - \mathrm{Price}_{i,p}) \times \mathrm{Chip}_{i,t,p} \times \mathrm{I}_{\mathrm{Close}_{i,t} \geq \mathrm{Price}_{i,p}}
\end{equation}

4. 计算当日兑现的盈利占筹码总盈利的比例，即为已实现盈利比率 \( \mathrm{CPGR} \)：
\begin{equation}
    \mathrm{CPGR}_{i,t} = \frac{\text{已实现盈利}_{i,t}}{\text{已实现盈利}_{i,t} + \text{账面盈利}_{i,t}}
\end{equation}

\subsection{说明}
已实现盈利比率衡量了投资者盈利时的卖出倾向，因子值越大，投资者越倾向于卖出赚钱的股票，抛压大，股价上涨受阻，与未来收益负相关。

\subsection{参考文献}
\begin{enumerate}
    \item 陈升斌, 姚家敏. 2025. 处置效应与 V 型处置效应在量化选股中的应用. 中信建投.
    \item Odean, T. (1998). \textit{Are investors reluctant to realize their losses?}. Journal of Finance, 53(5), 1775-1790.
\end{enumerate}

\section{多层加权日内信噪比(高频因子-收益分布类)}

\subsection{因子定义}
\begin{equation}
    \text{新}\mathrm{SNR} = \mathrm{Weight}_{\mathrm{vol}} \times \mathrm{SNR}_{\mathrm{layer2}} + (1 - \mathrm{Weight}_{\mathrm{vol}}) \times \mathrm{SNR}_{\mathrm{layer3}}
\end{equation}
\begin{equation}
    \mathrm{Weight}_{\mathrm{vol}} = 0.5 + \frac{\mathrm{Vol}_{\mathrm{std}} - 0.5}{\delta}
\end{equation}
\begin{equation}
    \mathrm{Vol}_{\mathrm{std}} = \frac{\mathrm{Vol} - \min(\mathrm{Vol})}{\max(\mathrm{Vol}) - \min(\mathrm{Vol})}
\end{equation}

\subsection{变量定义}
\begin{itemize}
    \item ① \( \mathrm{SNR}_{\mathrm{layer}n} \) 为分解至 \( n \) 层后计算得到的信噪比，计算逻辑见相关日历页；
    \item ② \( \mathrm{Vol} \) 为日内价格波动率，先在横截面上对其做极值处理和归一化处理，然后将标准化后的波动率缩至合适的权重区间。
\end{itemize}

\subsection{说明}
新的日内信噪比因子由日内价格波动为权重对第二层和第三层信噪比因子加权得到；日内价格波动越大，剥离更多噪声，赋予第三层更大的权重；加权后的复合信噪比的表现有所提升。

\subsection{参考文献}
\begin{enumerate}
    \item 严佳炜, 朱定豪. 信息提取: 信息提取与信号分析. 华安证券.
\end{enumerate}

\section{同时点峰岭数相关性（高频因子-成交分布类）}

\subsection{因子定义}

① 对过去20日同时点成交量按照1倍标准差的标准划分，高于1倍标准差划为“喷发成交量”，低于1倍标准差划为“温和成交量”；

② 对“喷发成交量”进一步按照连续与否划分：若当前分钟为“喷发成交量”，且前后1分钟均为“温和成交量”，则划为“孤立喷发成交量”，称为“量峰”；若当前分钟为“喷发成交量”，且前后1分钟也存在“喷发成交量”，则划为“连续喷发成交量”，称为“量岭”；将“温和成交量”称为“量谷”；

③ 计算过去20日同一时点的量峰与量岭样本数（对应240个分钟点）的相关系数：
\begin{equation}
    \mathrm{Corr}(\text{过去20日同一时点的量峰数序列}, \text{过去20日同一时点的量岭数序列})
\end{equation}

\subsection{说明}
同时点峰岭数相关性因子衡量了知情交易者与个人投资者在同一交易时点的参与度一致性。一致性越高，个人投资者对知情交易者跟随越紧密，知情交易者的投资动向可能已被个人投资者识别，个股未来表现会减弱，是一个负向因子。

\subsection{参考文献}
\begin{enumerate}
    \item 魏建榕, 王志豪. 2025. 高频成交量的峰、岭、谷信息. 开源证券.
\end{enumerate}

\section{买入浮亏偏离小单因子（行为金融因子-遗憾规避）}

\subsection{因子定义}
\begin{equation}
    \mathrm{HCPS} = \frac{\sum_{i}^{N} \overline{\mathrm{price}}_{\mathrm{buy},i} \cdot \mathbf{I}_{\{P_{\mathrm{buy},i} > \mathrm{close}\}} \cdot \mathbf{I}_{\{\mathrm{vol} < \mathrm{volume}_{\mathrm{mean}}\}}}{close} - 1
\end{equation}

\subsection{变量定义}
\begin{itemize}
    \item ① \( \mathrm{price}_{\mathrm{buy},i} 、p_{sell,i}\) 为主买成交价；
    \item ② \( \mathrm{volume}_{\mathrm{mean}} \) 为当日所有实际成交量均值；
    \item ③ 逐笔中每笔成交买卖方向确定方式：若该笔成交为买单被拆分与多个卖单成交，则记为买方交易，反之为卖方交易；
    \item ④ 大小单划分方式：以当日所有实际成交量均值作为区分大小单的阈值，即
        \begin{equation*}
            \mathrm{vol} < \mathrm{volume}_{\mathrm{mean}}
        \end{equation*}
        为小单，反之为大单；
    \item ⑤ 此外，还可只统计尾盘期间 \([14:30,14:57)\) 的成交量，计算得到买入浮亏偏离尾盘因子。
\end{itemize}

\subsection{说明}
遗憾规避理论认为非理性的投资者在做决策时，会倾向于避免产生后悔情绪并追求自我安慰，避免承认之前的决策失误；对于买入后浮亏（当天买入的股票收盘价低于买入的成本价）的投资者会避免承认决策失误而存在惜售心理，不愿卖出下跌浮亏的股票；买入浮亏偏离因子从成交价格偏离度来量化上述心理，因子值越高，买入成本越高，成交价格偏离程度越大，未来抛压较低，会有更高的预期收益。

\subsection{参考文献}
\begin{enumerate}
    \item 高智威. 2023. 基于逐笔成交数据的遗憾规避因子. 国金证券.
\end{enumerate}

\section{个股收益最高时的价变共振(高频因子-动量反转类)}

\subsection{因子定义}

1. 计算个股在日内中最高的 5 个分钟收益率期间，对应全市场等权收益率的均值：
\begin{equation}
    \mathrm{RM\_MAX}(N)_i = \frac{\sum_{t=n1}^{n_N} \mathrm{MEAN\_MKT}_t}{N}
\end{equation}
\begin{equation}
    \mathrm{MEAN\_MKT}_t = \frac{\sum_{i=1}^I r_{i,t}}{I}
\end{equation}

2. 计算个股在日内中最高的 5 个分钟收益率期间，全市场收益率标准差的均值：
\begin{equation}
    \mathrm{STD\_MAX}_i = \frac{\sum_{t=n11}^{nN} \mathrm{std}(r_t)}{N}
\end{equation}

3. 对 \( \mathrm{RM\_MAX} \) 做截面上的正序排名名次归一化，并除以 \( \mathrm{STD\_MAX} \) 即为个股收益最高时的市场共振因子：
\begin{equation}
    \mathrm{CO\_MAX}_i = \frac{\mathrm{RM\_MAX}(N)_i}{\mathrm{STD\_MAX}_i}
\end{equation}

\subsection{变量定义}
\begin{itemize}
    \item ① \( r_{i,t} \) 为个股 \( i \) 在 \( t \) 分钟的收益率，\( r_t \) 为 \( t \) 时刻的全部个股收益率序列；
    \item ② \( n1, n2, \ldots, nN \) 为日内分钟收益率最高的 \( N \) 个分钟，\( N=5 \)；
    \item ③ \( \mathrm{MEAN\_MKT}_t \) 为在 \( t \) 分钟的全市场等权收益率均值。
\end{itemize}

\subsection{说明}
\( \mathrm{RM\_MAX} \) 刻画了个股与市场趋势的一致性，因子值越大，说明个股是在市场表现较好时同步出现价格上涨，个股上涨与全市场上涨趋势一致，个股与市场形成有效共振，股价上涨被普遍认可，更容易延续，未来收益表现较好；\( \mathrm{STD\_MAX} \) 刻画了市场整体收益的分歧度，分歧度越小，市场趋势更具市场代表性，个股与市场共振越有效，将其叠加后有助于提升 \( \mathrm{RM\_MAX} \) 的表现。

\subsection{参考文献}
\begin{enumerate}
    \item 任继敏, 袁元勋, 许继宏. 2024. 基于分钟数据的价变共振因子. 招商证券.
\end{enumerate}

\section{单位交易额的有效价差（高频因子-流动性类）}

\subsection{因子定义}
\begin{equation}
    \mathrm{ESI}_{i,t} = \frac{\mathrm{ES}_{i,t}}{\$ \mathrm{vol}_{i,t}}
\end{equation}
\begin{equation}
    \mathrm{ES}_{i,t} = 2 D_{i,k} \cdot \frac{P_{i,t} - \left( \frac{\mathrm{Ask}_{i,t} + \mathrm{Bid}_{i,t}}{2} \right)}{\left( \mathrm{Ask}_{i,t} + \mathrm{Bid}_{i,t} \right)/2}
\end{equation}

\subsection{变量定义}
\begin{itemize}
    \item ① \( \mathrm{Ask}_{i,t}, \mathrm{Bid}_{i,t}, P_{i,t} \) 分别为股票 \( i \) 在 \( t \) 时刻的卖 1 价、买 1 价、收盘价；
    \item ② \( D_{i,k} \) 为买卖标识，买单记为 1，卖单记为 -1；
    \item ③ \( D_{i,k} \) 买卖单判断标准：若当前收盘价高于（低于）最优买卖报价平均值，则认为是买（卖）单；当前收盘价等于最优买卖报价平均值时，若收盘价高于（低于）上一笔价格，则认为是买（卖）单；
    \item ④ \( \$ \mathrm{vol}_{i,t} \) 为 \( t \) 日当天成交额。
\end{itemize}

\subsection{说明}
单位交易额的有效价差衡量了交易中每单位交易额的边际交易成本，反映流动性效率；股票的交易成本越高，流动性越差，未来会有正向的流动性风险溢价。

\subsection{参考文献}
\begin{enumerate}
    \item 陈升悦. 2024. 流动性高频因子构建与投资者注意力因子分析报告. 中信建投.
\end{enumerate}

\section{特质日内收益波动率（高频因子-波动跳跃类）}

\subsection{因子定义}

截面上对所有股票过去 20 天的日内收益率序列进行主成分分析 PCA，提取第一主成分 \( F_1 \) 为隐式因子溢价：
\begin{equation}
    \mathrm{IntradayRet}_{i,t} = \frac{\mathrm{Close}_{i,t}}{\mathrm{Open}_{i,t}} - 1
\end{equation}

对每只股票 \( i \)，以过去 20 天的日内收益率序列为 \( y \)，以隐式因子 \( F_1 \) 序列为 \( x \)，进行时序回归：
\begin{equation}
    \mathrm{IntradayRet}_{i,t} = \alpha_{i,t} + \beta_i^{F_1} \times F_1 + \cdots + \varepsilon_i
\end{equation}

时序回归返回的残差序列的标准差即为隐式因子框架下的特质波动因子：
\begin{equation}
    \mathrm{IntradayRetIV} = \mathrm{std}(\varepsilon_i)
\end{equation}

\subsection{变量定义}
\begin{itemize}
    \item ① 在进行截面主成分分析时，若第一主成分的解释力度低于 20\%，则选取解释力最大的前两个主成分进行回归；
    \item ② 还可进一步计算特质度因子：
        \begin{equation*}
            \mathrm{IVR} = 1 - R^2
        \end{equation*}
        上述回归模型拟合优度；
    \item ③ 还可进一步计算特质偏度因子：
        \begin{equation*}
            \mathrm{IS} = \frac{\sum e_{lt}^3}{\mathrm{IV}^3}
        \end{equation*}
\end{itemize}

\subsection{说明}
计算隐式因子框架下的特质波动率时，直接通过主成分分析对资产收益率序列降维来提取共同波动，而无需显式的借助多因子模型估计共同风险。系统性的解决传统多因子模型遗漏变量、估计误差的问题，计算简单；此处提取的是日内收益率序列的共同波动。

\subsection{参考文献}
\begin{enumerate}
    \item 张欣慧, 杨怡玲, 刘璐. 2022. 隐式框架下的特质因子改进. 国信证券.
\end{enumerate}

\section{筹码收益调整(量价因子改进)}

\subsection{因子定义}

利用过去 250 日的数据，对股票 \( i \) 在 \( T \) 日计算得到 \([T-250, T]\) 之间每日的留存筹码 \( \mathrm{RsdAmt}_{T-k}^T \) 和成交均价 \( \mathrm{Price}_{T-k}^{\mathrm{avg}} \)，即为该股票在 \( T \) 日的筹码分布估计：
\begin{equation}
    \mathrm{RsdAmt}_{T-k}^T = \mathrm{Amt}_{T-k} \cdot \prod_{t=T-k+1}^{T} (1 - \mathrm{Turnover}_t), \quad k \in (0, 250]
\end{equation}

计算最新日成交均价相对历史筹码的加权成本的相对收益，即为最新日的筹码收益因子：
\begin{equation}
    \mathrm{holding\_ret} = \frac{\sum_t \mathrm{Price}_t^{\mathrm{avg}} \times \mathrm{RsdAmt}_t^T}{\sum_t \mathrm{RsdAmt}_t^T} 
\end{equation}

对全市场的筹码收益因子加权计算均值作为 A 股市场的筹码收益 \( \mathrm{mkt\_holding\_ret} \)，用以调节个股筹码收益因子 \( \mathrm{holding\_ret} \) 的动量效应和反转效应，即得到筹码收益调整因子：
\begin{equation}
    \mathrm{holding\_ret\_adj} = \mathrm{holding\_ret} \times \mathrm{sign}(\mathrm{mkt\_holding\_ret})
\end{equation}

\subsection{变量定义}
\begin{itemize}
    \item ① \( \mathrm{Amt}_{T-k} \) 为股票在 \( T-k \) 日的成交额；
    \item ② \( \mathrm{Turnover}_t \) 为股票在 \( t \) 日的换手率；
    \item ③ 市场筹码收益切分阈值也可取 -2\%，即：\( \mathrm{sign}(\mathrm{mkt\_holding\_ret} - (-2\%)) \)。
\end{itemize}

\subsection{说明}
原始筹码收益因子在牛市中通常表现为动量特征，在熊市中通常表现为反转特征，可借助市场筹码收益指标刻画市场赚钱效应，保留市场赚钱效应较好时（市场筹码收益为正）个股筹码收益的动量效应，扭转市场赚钱效应较差（市场筹码收益为负）时个股筹码收益的反转效应，调整后的筹码收益因子为正向因子。

\subsection{参考文献}
\begin{enumerate}
    \item 魏建福, 张翔. 2025. 筹码结构视角下的动量反转融合. 开源证券.
\end{enumerate}

\section{价量相关性综合因子(高频因子-量价相关性类)}

\subsection{因子定义}

价量相关性综合因子 \( \mathrm{PV\_corr} \) 为横截面标准化后的价量相关性平均数因子和价量相关性波动性因子的等权线性相加：
\begin{equation}
    \mathrm{PV_corr} = \frac{\mathrm{PV_{corravg}} - \mathrm{mean}(\mathrm{PV_{corravg}})}{\mathrm{std}(\mathrm{PV_{corravg}})} + \frac{\mathrm{PV_\section{高低位放量选股因子（五分钟波动率版本）}

\subsection{因子定义}
\begin{equation}
\begin{aligned}
\mathrm{Factor\_Low}_{i} &= \frac{\mathrm{Avg\_Std\_Low}}{\mathrm{Avg\_Std\_Total}} \\[12pt]
\mathrm{Factor\_High}_{i} &= \frac{\mathrm{Avg\_Std\_High}}{\mathrm{Avg\_Std\_Total}} \\[12pt]
\mathrm{Factor\_Diff}_{i} &= \mathrm{Factor\_Low}_{i} - \mathrm{Factor\_High}_{i}
\end{aligned}
\end{equation}

\subsection{变量定义}
\begin{itemize}
    \item 本因子基于**分钟频收益率**与**波动率**构建。若从原始价格数据开始，需计算以下前置变量：
\end{itemize}
\subsubsection{前置计算}
\paragraph{分钟收益率}
\begin{equation}
r_{t} = \frac{P^{close}_{t}}{P^{open}_{t}} - 1
\end{equation}
其中，$P^{close}_{t}$ 为第 $t$ 分钟的收盘价，$P^{open}_{t}$ 为第 $t$ 分钟的开盘价，$r_{t}$ 为第 $t$ 分钟的收益率。

\paragraph{五分钟滚动波动率}
\begin{equation}
\sigma_{t} = \mathrm{Std}(r_{t-4}, r_{t-3}, r_{t-2}, r_{t-1}, r_{t})
\end{equation}
其中，$\sigma_{t}$ 为截至第 $t$ 分钟的、过去5分钟收益率的标准差。

\subsubsection{核心变量}
\begin{itemize}
    \item $T$： 回看期，默认为20个交易日。
    \item $M$： 回看期 $T$ 内所有股票 $i$ 的分钟数据点集合。
    \item $P_{m}$： 分钟 $m$ 的收盘价，其中 $m \in M$。
    \item $\mathrm{Thd}_{low}$： 集合 $M$ 中分钟收盘价的20\%分位数，即 $\mathrm{Thd}_{low} = \mathrm{Quantile}(P_{m}, 0.2)$。
    \item $\mathrm{Thd}_{high}$： 集合 $M$ 中分钟收盘价的80\%分位数，即 $\mathrm{Thd}_{high} = \mathrm{Quantile}(P_{m}, 0.8)$。
    \item $S_{low}$： 低位分钟集合，$S_{low} = \{m | P_{m} \le \mathrm{Thd}_{low}, m \in M\}$。
    \item $S_{high}$： 高位分钟集合，$S_{high} = \{m | P_{m} \ge \mathrm{Thd}_{high}, m \in M\}$。
    \item $\mathrm{Avg\_Std\_Low}$： 低位分钟集合 $S_{low}$ 中所有分钟对应的五分钟波动率 $\sigma_{m}$ 的均值。
    \item $\mathrm{Avg\_Std\_High}$： 高位分钟集合 $S_{high}$ 中所有分钟对应的五分钟波动率 $\sigma_{m}$ 的均值。
    \item $\mathrm{Avg\_Std\_Total}$： 回看期 $T$ 内所有分钟对应的五分钟波动率 $\sigma_{m}$ 的均值。
\end{itemize}

\subsection{计算步骤（改为1分钟频率）}
将上述所有计算中的“5分钟滚动波动率”计算基准改为“1分钟滚动波动率”，即计算**每1分钟的波动率**。具体修改如下：

\paragraph{一分钟滚动波动率（前置计算修改）}
\begin{equation}
\sigma^{1min}_{t} = \mathrm{Std}(r_{t})
\end{equation}
此处 $\sigma^{1min}_{t}$ 是第 $t$ 分钟的收益率 $r_{t}$ 自身的标准差。在单点情况下，其值为0，这在实际计算中无意义。因此，**文档未详述此点，但基于我所掌握的知识**，对于1分钟频率，通常需要定义一个滚动窗口来计算波动率，例如使用过去N分钟的收益率序列（如N=5或N=20）。为保持与原逻辑一致，建议将波动率计算窗口定义为过去 $W$ 分钟（例如 $W=5$）：
\begin{equation}
\sigma^{1min}_{t} = \mathrm{Std}(r_{t-W+1}, r_{t-W+2}, ..., r_{t})
\end{equation}
其中 $W$ 为滚动窗口长度。后续所有计算中的 $\sigma_{m}$ 均替换为此 $\sigma^{1min}_{m}$。变量 $S_{low}$、$S_{high}$ 的划分依据（分钟收盘价分位数）以及 $\mathrm{Avg\_Std\_Low}$、$\mathrm{Avg\_Std\_High}$、$\mathrm{Avg\_Std\_Total}$ 的计算方式保持不变。

\subsection{说明}
该因子通过比较股票在历史价格高位区间与低位区间的日内波动率相对强弱来选股。计算时，首先识别出过去20个交易日内，所有分钟收盘价处于最高20\%（高位）和最低20\%（低位）的时段，然后分别计算这两个时段内每分钟波动率的平均值，最后将它们分别除以整个回看期内所有分钟波动率的平均值，得到“高位波动率占比”和“低位波动率占比”因子。两者相减可得到“差分”因子。文档测试表明，使用分钟频（5分钟）数据计算的该因子，在分层回测上表现优于日频版本。因子值可采用过去一段时间的均值或标准差进行合成，其中“低位”和“差分”因子使用均值合成效果较好。

\subsection{参考文献}
\begin{enumerate}
    \item 本文档内容整理自《高低位放量选股因子之五分钟波动率版本，效果炸裂！只能说分钟频>>日频！》。
    \item 原文档未提供外部学术参考文献。
\end{enumerate}{corrstd}} - \mathrm{mean}(\mathrm{PV_{corrstd}})}{\mathrm{std}(\mathrm{PV_{corrstd}})}
\end{equation}

\subsection{说明}
\( \mathrm{PV\_corr} \) 综合反映了价量关系对传统反转因子的修正。

\subsection{参考文献}
\begin{enumerate}
    \item 高子剑. 2020. 高频价量相关性，意想不到的选股因子. 东吴证券.
\end{enumerate}

\section{高低价复合交易占比(高频因子-资金流类)}

\subsection{因子定义}

\begin{equation}
    \text{高低价复合交易占比} = \frac{\text{非高低价买单\&非高低价卖单的成交量} - \text{高价买单\&高价卖单的成交量}}{\text{总成交量}}
\end{equation}

\subsection{变量定义}
\begin{itemize}
    \item ① 高低价订单划分方式为：委托价格处于同类型订单委托价格前 10\% 分位点的订单为高价单，处于后 10\% 分位点的订单为低价单。
\end{itemize}

\subsection{说明}
高低价复合交易占比因子从订单委托价格高低维度对逐笔订单进行划分和归整。测试发现，“高价买\&高价卖”成交量占比和“低价买\&低价卖”成交量占比越高，其未来反转效应越明显，且前者比后者更有效；“非高低价买+非高低价卖”成交量占比越高，其未来的动量效应越强；高低价复合交易占比因子为两者的复合，与未来收益正相关。

\subsection{参考文献}
\begin{enumerate}
    \item 张欣慰, 张宇. 2024. 高频订单成交数据蕴含的 Alpha 信息. 国信证券.
\end{enumerate}

\section{横盘时刻-主买笔均成交金额(高频因子-成交分布类)}

\subsection{因子定义}

1. 将日内 240 个 1 分钟数据集合标记为 \( T = \{t_1, t_2, \dots, t_{240}\} \)，判断股票 i 的每分钟涨跌幅，将涨跌幅为 0 时刻标记为集合 \( P = \{t_{i1}, t_{i2}, \dots, t_{ik}\} \)；

2. 利用逐笔成交数据，进一步对 P 中时刻的所有成交记录划分为主动买入成交 \( P\_Buy \) 和主动卖出成交 \( P\_Sell \) 两类；

3. 对横盘时刻下的主动买入成交额和成交笔数进行汇总，根据其成交额 \( \mathrm{Amt} \) 之和除以成交笔数 \( \mathrm{DealNum}\) 之和，构造横盘时刻且主买的笔均成交金额 \( \mathrm{ATD}_{\mathrm{Zero\_Buy}} \)：
\begin{equation}
    \mathrm{ATD}_{\mathrm{Zero\_Buy}} = \frac{\sum_{t \in P\_\mathrm{Buy}} \mathrm{Amt}_t}{\sum_{t \in P\_\mathrm{Buy}} \mathrm{DealNum}_t}
\end{equation}

4. 同理，将全天的成交金额除以全天的成交笔数，得到全天的笔均成交金额 \( \mathrm{ATD}_T \)：
\begin{equation}
    \mathrm{ATD}_T = \frac{\sum_{t \in T} \mathrm{Amt}_t}{\sum_{t \in T} \mathrm{DealNum}_t}
\end{equation}

5. 将横盘时刻且主买的笔均成交金额除以全天的笔均成交金额，即得到去量纲化后的标准化因子 \( \mathrm{SATD}_{\mathrm{Zero\_Buy}} \)：
\begin{equation}
    \mathrm{SATD}_{\mathrm{Zero\_Buy}} = \frac{\mathrm{ATD}_{\mathrm{Zero\_Buy}}}{\mathrm{ATD}_T}
\end{equation}

\subsection{说明}
笔均成交额类因子通过汇总“特殊的、具有较多信息含量”时刻的成交金额情况来刻画主力资金的行为动向；笔均成交金额越大，说明该特殊时间内资金优势越明显，属于主力资金的可能性越大。买方主导的横盘时刻，成交越多，未来收益越低。

\subsection{参考文献}
\begin{enumerate}
    \item 张欣鹏, 张宇. 2025. 日内特殊时刻蕴含的主力资金 Alpha 信息. 国信证券.
\end{enumerate}

\section{结构化反转(高频因子-动量反转类)}

\subsection{因子定义}

① 将一段时间内每个时间段的成交量，从小到大进行排序，将成交量小于等于 10\% 的时间段记为动量时间段 \(\mathrm{period}_{\mathrm{mom}}\)，大于 10\% 的时间段记为反转时间段 \(\mathrm{period}_{\mathrm{rev}}\)；

② 以成交量倒数为权重构建动量时间段反转因子：
\begin{equation}
    \mathrm{Rev}_{\mathrm{mom}} = \sum_{i=1}^{\mathrm{period}_{\mathrm{mom}}} w_i \log \left( \frac{\mathrm{Close}_{t-i+1}}{\mathrm{Close}_{t-i}} \right), \quad w_i \propto \frac{1}{\mathrm{volume}_i}
\end{equation}

③ 以成交量为权重构建反转时间段反转因子：
\begin{equation}
    \mathrm{Rev}_{\mathrm{rev}} = \sum_{i=1}^{\mathrm{period}_{\mathrm{mom}}} w_i \log \left( \frac{\mathrm{Close}_{t-i+1}}{\mathrm{Close}_{t-i}} \right), \quad w_i \propto \mathrm{volume}_i
\end{equation}

④ 以动量时间段反转因子和反转时间段反转因子合成结构化反转因子：
\begin{equation}
    \mathrm{Rev}_{\mathrm{struct}} = \mathrm{Rev}_{\mathrm{rev}} - \mathrm{Rev}_{\mathrm{mom}}
\end{equation}

\subsection{变量定义}
\begin{itemize}
    \item ① \(\mathrm{period}\) 为过去一段时间按一定规则划分时间段的总时间段个数，如过去 21 个交易日的日内 10 分钟频率数据的 \(\mathrm{period}=24 \times 21\)，\(\mathrm{period} = \mathrm{period}_{\mathrm{mom}} \cup \mathrm{period}_{\mathrm{rev}}\)；大小成交量划分阈值也可选 10\%、15\%、20\% 等。
\end{itemize}

\subsection{说明}
结构化反转因子改进了高频反转因子头部分组收益的线性单调性。多空双方的博弈状态会影响反转效应的强弱；在好坏信息发生初期，多空双方可能几乎不存在博弈，价格呈现很强的确定性，不可逆性强，此时成交量低，动量效应强；随着时间推移，逐渐进入新的博弈阶段，成交量开始恢复，反转效应加强；即存在一个成交量的阈值，小于该阈值时的价格变动呈现动量效应，大于该阈值时的价格变动呈现反转效应。

\subsection{参考文献}
\begin{enumerate}
    \item 骤川桃, 郑起. 2019. 高频因子（二）：结构化反转因子. 长江证券.
\end{enumerate}

\section{跳跃强度(高顿因子-波动跳跃类)}

\subsection{因子定义}

计算检验是否存在跳跃的检验统计量，原假设为不存在跳跃：
\begin{equation}
    \mathrm{T}_{SwV} = \frac{\mathrm{BV}_t}{N^{-1}\sqrt{\hat{\Omega}_{\mathrm{SwV}}}} \left(1 - \frac{\mathrm{RV}_t}{\mathrm{SwV}_t}\right) \xrightarrow{d} N(0,1)
\end{equation}
\begin{equation}
    \hat{\Omega}_{\mathrm{SwV}} = \frac{\mu_6}{9} \frac{N^3}{N - 5} \int_{t-1}^{t} \sigma_s^6 ds = \frac{\mu_6}{9} \frac{N^3}{N - 3}\mu_{3/2}^{-4}\sum_{i=0}^{N-3} (\prod_{k=1}^4|r_{t_{i+k}}|^{3/2})
\end{equation}
\begin{equation}
    \mathrm{SwV}_t = 2 \sum_{i=1}^N (R_{t_i} - r_{t_i})
\end{equation}

构建用于表示日内价格是否存在跳跃的示性函数：
\begin{equation}
    I_{\mathrm{jump\_SwV}} = I(\mathrm{T}_{SwV} > \Phi_{1-\alpha}^{-1})
\end{equation}

计算观察期内跳跃强度，如月度的 \( D = 20 \)：
\begin{equation}
    \mathrm{JArr} = \frac{1}{D} \sum_{t=1}^D I_{\mathrm{jump\_SwV},t}
\end{equation}

\subsection{变量定义}
\begin{itemize}
    \item ① \( R_{t_i} \) 和 \( r_{t_i} \) 分别为交易日 \( t \) 内第 \( i \) 个时间区间内的简单收益率和对数收益率；
    \item ② \( \Phi \) 为标准正态分布累计分布函数，\( \alpha \) 为显著性水平；
    \item ③ 可根据日内收益的正负方向，分开统计正向跳跃强度和负向跳跃强度。
\end{itemize}

\subsection{说明}
跳跃强度为一段时间内出现价格跳跃的天数占比，衡量了股价异常波动的持续性，可以用于刻画信息冲击的强度。

\subsection{参考文献}
\begin{enumerate}
    \item Tauchen, G., \& Zhou, H. (2011). \textit{Realized jumps on financial markets and predicting credit spreads}. Journal of Econometrics, 160(1), 102-118.
    \item 周飞鹏, 罗军, 安宁宁. 2022. 再谈股价跳跃因子. 广发证券.
\end{enumerate}

\section{基于市场波动率的注意力捕捉(行为金融因子-投资者注意力)}

\subsection{因子定义}

对每个行业内分别统计每日行业成分股截面收益标准差作为市场关注代理指标：
\begin{equation}
    \mathrm{indus\_vol}_d = \mathrm{std}(\text{行业内成分股的日度收益序列 } r_{i,d})
\end{equation}

将市场关注代理指标与股票日收益进行月度时序回归，代理指标前的系数绝对值 \(|\beta_{i,t}|\) 即为股票对市场关注的敏感程度：
\begin{equation}
    r_{i,d} = \mu_{i,t} + \beta_{i,t} \mathrm{indus\_vol}_d + \rho_{1,i,t} \mathrm{Mkt}_d + \rho_{2,i,t} \mathrm{SMB}_d + \rho_{3,i,t} \mathrm{HML}_d + \rho_{4,i,t} \mathrm{UMD}_d + \varepsilon_{i,d}
\end{equation}

\subsection{变量定义}
\begin{itemize}
     \item ① \(n_{\mathrm{upper},d}\)、\(n_{\mathrm{lower},d}\)、\(N_d\) 分别为 \(d\) 日涨停股票数量、跌停股票数量、当日交易股票总数；
    \item ① \( r_{i,d} \) 为个股日度收益；
    \item ② \( \mathrm{Mkt}_d, \mathrm{SMB}_d, \mathrm{HML}_d \) 为经典的 Fama-French 三因子；
    \item ③ \( \mathrm{UMD}_d \) 为动量因子收益，即全市场过去 11 个月累积收益前后 30\% 的组合在 \( d \) 日的收益差。
\end{itemize}

\subsection{说明}
不同股票对同一市场关注度事件的反映程度往往是不同的，可借助个股收益对市场关注度事件的敏感程度来衡量股票对该事件的注意力捕捉效应；捕捉能力越强（敏感程度越高），表明个股更能吸引市场注意，从而导致买盘压力增大，股价高估；上述因子以市场波动率作为市场关注度事件，收益的分化也表明市场交易活跃，关注度高。

\subsection{参考文献}
\begin{enumerate}
    \item Lin, F., Qiu, Z., \& Zheng, W. (2023). \textit{Cranes among chickens: The general-attention-grabbing effect of daily price limits in China's stock market}. Journal of Banking \& Finance, 150, 106818.
    \item 陈升锐. 2024. 投资者有限关注及注意力捕捉与溢出. 中信出版社.
\end{enumerate}

\section{模糊金额比(高频因子-成交分布类)}

\subsection{因子定义}

1. 对股票 \( i \) 计算 \( n \) 日内的 1 分钟收益率；计算第 \( t \) 分钟的波动率，即 \([t-4, t]\) 区间内分钟收益率的标准差；计算第 \( t \) 分钟的模糊性，即 \([t-4, t]\) 区间内上述分钟波动率的标准差；

2. 将日内模糊性大于日内模糊性均值的时间段记为“起雾时刻”；

3. 计算“起雾时刻”的分钟成交额均值为“雾中金额”，计算日内所有时刻的分钟成交额均值为“总体金额”；

4. 计算“日模糊金额比”：
\begin{equation}
    \text{日模糊金额比} = \frac{\text{雾中金额}}{\text{总体金额}}
\end{equation}

\subsection{变量定义}
\begin{itemize}
    \item ① 计算上述因子时，剔除开盘和收盘数据，仅保留日内交易数据，所以开盘前 9 分钟的模糊性为空值；
    \item ② 每月末可对最近 20 天的上述因子求均值和标准差并将两者等权合并，即得到月度“模糊金额比”因子。
\end{itemize}

\subsection{说明}
模糊金额比因子从金额维度统计了投资者在模糊性较大时的成交程度，同样衡量了投资者对模糊性的厌恶程度，与未来收益负相关；投资者对波动率模糊性的厌恶心理会使得投资者急于卖出而产生过度反应，未来很可能会补涨。

\subsection{参考文献}
\begin{enumerate}
    \item 曹春晓. 2022. 波动率的波动率与投资者模糊性厌恶. 方正证券.
    \item Kostopoulos, D., Meyer, S., \& Uhr, C. (2021). \textit{Ambiguity about volatility and investor behavior}. Journal of Financial Economics.
\end{enumerate}

\section{知情交易概率(高频因子-资金流类)}

\subsection{因子定义}
\begin{equation}
    \mathrm{PIN} = \frac{\alpha \mu}{\alpha \mu + 2\varepsilon}
\end{equation}

\subsection{变量定义}
\begin{itemize}
    \item ① 假设信息事件发生的概率为 \( \alpha \)，且坏消息发生的概率为 \( \delta \)；
    \item ② 假设知情交易者的委托单到达速率为 \( \mu \)，非知情交易者的委托单到达速率为 \( \varepsilon \)，上述过程满足泊松分布；
    \item ③ 上述参数需通过求解买卖委托单混合泊松分布的极大似然函数估计得到。
\end{itemize}

\subsection{说明}
PIN 是最早对知情交易概率进行直接测度的方法；知情交易概率指在一段时间内，拥有信息优势的知情交易者提交的订单数量占总委托单数量的比例，用以刻画该段时间内的信息不对称程度，通常与未来收益正相关，取值越大，信息不对称程度越大，投资者要求的风险回报越高。

\subsection{参考文献}
\begin{enumerate}
    \item Easley, D., Kiefer, N. M., O'Hara, M., \& Paperman, J. B. (1996). \textit{Liquidity, information and infrequently traded stocks}. Journal of Finance, 51(4), 1405-1436.
\end{enumerate}

\section{CSAD模型（基于时序回归变形）(行为金融因子-投资者注意力)}

\subsection{因子定义}

1. 每个月末提取当前可交易股票过去 120 个交易日的涨跌幅数据，计算股票 \( S_* \) 和其他股票（不活跃交易日不超过 20 天）之间日涨跌幅相关系数；

2. 选取相关系数最大的 9 只股票，连同 \( S_* \) 一起构成一个板块：
   \begin{equation}
       S_i \ (i = 1, \dots, 10)
   \end{equation}

3. 以 \( S_* \) 为中心，计算板块过去 120 个交易日的横截面绝对偏差 \( \mathrm{CSAD} \)：
   \begin{equation}
       \mathrm{CSAD}_t = \frac{1}{10} \sum_{i=1}^{10} |r_{it} - r_{*t}|
   \end{equation}

4. 对 CSAD 波动率相对个股波动率比值的变化作为羊群效应因子：
   \begin{equation}
       \mathrm{ratio}[t_1, t_2] = \frac{\sigma_{[t_1, t_2]}(\mathrm{CSAD})}{\sigma_{[t_1, t_2]}(r_*)}
   \end{equation}
   \begin{equation}
       f_{jt} = -\frac{\mathrm{ratio}[t-19, t]}{\mathrm{ratio}[t-119, t]}
   \end{equation}

\subsection{说明}
CSAD 模型对羊群效应的检验主要基于时序回归上市场收益平方项前回归系数正负及显著性（系数显著为负，说明股票收益分散度的增速随市场极端收益的增加而放缓，股价走势越一致，羊群效应越强）；考虑到时序长度不易确定和回归易受异常值影响等问题，可以对时序回归进行变形，通过 CSAD 方差相对收益方差的比值变动来衡量羊群效应的强弱，比值越小，板块羊群效应越强，所在板块受关注程度提升越明显，未来收益越高。

\subsection{参考文献}
\begin{enumerate}
    \item Chang, E. C., Cheng, J., \& Khorana, A. (2000). \textit{An Examination of Herd Behavior in Equity Markets: An International Perspective}. Journal of Banking \& Finance, 24(10), 1651-1679.
    \item 丁鲁明, 陈元勋. 2019. 基于市场羊群效应的股票 alpha 研究. 中信出版社.
\end{enumerate}

\section{高频弹性因子(高频因子-流动性类)}

\subsection{因子定义}
使用 Hodrick-Prescott 算法将自然对数后的 1 分钟或 5 分钟股票价格分解为基础价格 \( q_t \) 和暂时价格 \( z_t \)：
\begin{equation}
    p_t = q_t + z_t
\end{equation}

计算 \( z_t \) 离散傅里叶变换后的频谱函数 \( Z_k \)，其中 \( k \) 为频域单位，\( D \) 为交易日总天数，\( i \) 为虚数单位：
\begin{equation}
    Z_k = \sum_{t=1}^{D} z_t e^{-\frac{i2\pi kt}{D}}, \quad (k = 1, 2, \cdots, D)
\end{equation}

计算归一化后的频谱函数的幅度 \( |\bar{Z}_k| \)：
\begin{equation}
    |\bar{Z}_k| = \left| \frac{1}{D} Z_k \right|
\end{equation}

计算暂时价格恢复的平均速度即为弹性因子，其中 \( D_{i,t} \) 为股票 \( i \) 在滚动窗口中的可用样本天数，符号 \([\cdot]\) 表示取最近整数：
\begin{equation}
    \mathrm{Resiliency}_{i,t} = \frac{1}{\left[ \frac{D_{i,t}}{2} \right]} \sum_{k=1}^{\left[ D_{i,t}/2 \right]} \frac{2|\bar{Z}_{k,i,t}|}{T_{k,i,t}} = \frac{1}{\left[ \frac{D_{i,t}}{2} \right]} \sum_{k=1}^{\left[ D_{i,t}/2 \right]} 2|\bar{Z}_{k,i,t}| \cdot f_{k,i,t}
\end{equation}

\subsection{说明}
弹性可以描述为价格从信息优势交易者驱动的暂时价格影响恢复到其基本价格的速率，而该弹性因子直接衡量了暂时价格的恢复速度，即通过频域的频谱分析得到暂时价格的距离和恢复时间，然后用距离除以恢复时间来计算速度；与未来收益负相关，因子值越高，说明股价从先前暂时价格影响中恢复越快，流动性越高。

\subsection{参考文献}
\begin{enumerate}
    \item Kim, J., \& Kim, Y. (2019). \textit{Transitory prices, resiliency, and the cross-section of stock returns}. International Review of Financial Analysis, 63, 243-256.
    \item 陈原文, 罗军, 安宁宁. 2023. 弹性因子研究：从高频数据说起. 广发证券.
\end{enumerate}

\section{“待著而救”因子（高频因子-量价相关性类）}

\subsection{因子定义}

1. 对股票 \( i \) 在交易日 \( T \) 内，将当日成交量最大的 10 个分钟时刻定义为“海量时刻”；

2. 从时间上最靠前的“海量时刻”开始，若相邻两个“海量时刻”间隔小于 5 分钟，认为第二个“海量时刻”大概率仍是上个“海量时刻”导致的跟随交易，需剔除这些“海量时刻”，剩余时刻定义为“优势时刻”，“优势时刻”之后的 5 分钟定义为“跟随时刻”；

3. 计算每个“跟随时刻”的成交量总和，并除以对应的“优势时刻”的成交量，得到“跟随系数”；对每只股票，将其日内所有“跟随系数”求均值，即为“日跟随系数”；

4. 每月末，分别计算每只股票过去 20 天的“日跟随系数”的均值和标准差，并将两者等权合成得到“待著而救”因子。

\subsection{变量定义}
\begin{itemize}
    \item ① 计算该因子时，仅使用开盘后第 16 分钟开始的 1 分钟交易数据。
\end{itemize}

\subsection{说明}
“待著而救”因子衡量了股票日内普通投资者跟随优势投资者的反应程度，与未来收益负相关；因子值越大，表明优势投资者短时间内大量买入造成的成交量激增将导致普通投资者的明显跟随买入，使得股价短期内反应过度，股价未来更可能出现较大回落。

\subsection{参考文献}
\begin{enumerate}
    \item 曹春晓. 大单成交后的跟随效应与“待著而救”因子, 2023, 方正证券.
\end{enumerate}

\section{改进的客户动量因子（基于信息离散度）（图谱网络-动量溢出）}

\subsection{因子定义}
利用供应链的相关数据，计算如下客户动量因子：
\begin{equation}
    \mathrm{comm}_i^{1M} = \sum_{j=1}^{N_i} w_{ij}^{\mathrm{sales}} \mathrm{mom}_j^{1M}, \quad i=1,2,\ldots,N
\end{equation}

对每个客户 \( c \) 计算信息离散度 ID：
\begin{equation}
    \mathrm{ID}_{c,t} = \mathrm{sign}(\mathrm{CR}_{c,t}) \times [\%\mathrm{neg}_{c,t} - \%\mathrm{pos}_{c,t}]
\end{equation}

先基于信息离散度 ID（负向）进行分组，在各组内基于客户动量因子（正向）进行分组，以此构建投资组合。

\subsection{变量定义}
\begin{itemize}
    \item ① \( \mathrm{mom}_j^{1M} \) 为公司 \( j \) 的客户过去一个月收益率；
    \item ② \( w_{ij}^{\mathrm{sales}} \) 为销售占比；
    \item ③ \( \mathrm{CR}_{c,t} \) 为客户 \( c \) 在过去 3 个月的累计收益率；
    \item ④ \( \%\mathrm{neg}_{c,t} \)、\( \%\mathrm{pos}_{c,t} \) 分别为客户 \( c \) 在过去 3 个月内负收益天数占比和正收益天数占比。
\end{itemize}

\subsection{说明}
客户动量因子可以从“投资者有限注意力或信息反应不足”这个维度来解释，投资者对客户信息反应越不足，客户动量效应可能越强；信息离散度统计了收益形成期间日度股价涨跌天数占比的差异，间接反映了投资者对信息的反应程度，当信息连续时（即收益由小幅但大量的上涨天数来推动），投资者可能由于有限注意力而未能及时反应，股票动量效应更强；当信息离散时（即收益由大幅但少量的上涨天数来推动），投资者能够更快反应，股票动量效应更弱。同样的，投资者对经济相关公司（如供应链、竞争对手等）提供的连续信息反应不足，而对离散信息则反应迅速，所以可以用信息离散度来改进客户动量等关联性因子。

\subsection{参考文献}
\begin{enumerate}
    \item Huang, S., Lee, C. M. C., Song, Y., \& Xiang, H. (2022). \textit{A Frog in Every Pan: Information Discreteness and the Lead-Lag Returns Puzzle}. Journal of Financial Economics, 145(2), 83-102.
    \item Da, Z., Gurun, U. G., \& Warachka, M. (2014). \textit{Frog in the Pan: Continuous Information and Momentum}. Review of Financial Studies, 27(7), 528-554.
\end{enumerate}

\section{波动率加剧-惊恐因子（高频因子-动量反转类）}

\subsection{因子定义}

1. 计算各股票的日度“惊恐度”：
\begin{equation}
    \text{惊恐度} = \frac{\text{偏离度}}{\text{基准项}} = \frac{\mathrm{abs}(\text{个股收益率} - \text{市场收益率})}{\mathrm{abs}(\text{个股收益率}) + \mathrm{abs}(\text{市场收益率}) + 0.1}
\end{equation}

2. 计算各股票日内波动率，即对各股票日内的 1 分钟收益率序列求标准差；

3. 对各股票，将每日惊恐度、每日收益率、日内波动率相乘，得到加权调整收益率；

4. 对各股票，每月末计算过去 20 日加权调整收益率序列的均值和标准差，分别称为“波动率加剧-惊恐收益”因子和“波动率加剧-惊恐波动”因子，再将二者等权合成为“波动率加剧-惊恐”因子。

\subsection{变量定义}
\begin{itemize}
    \item ① 选取中证全指 000985 代表市场水平；
    \item ② “偏离度”衡量了个股收益率对市场收益率的偏离水平；“基准项”衡量了市场总体收益率水平。
\end{itemize}

\subsection{说明}
波动率加剧-惊恐因子在原始惊恐因子上新增了日内波动率权重，因为“显著效应（投资者会被具有显著收益的股票所吸引，导致它们被错误定价）在波动率加剧时更加强烈”。

\subsection{参考文献}
\begin{enumerate}
    \item 曹春晓. 显著效应, 极端收益扭曲决策权重和“草木皆兵”因子, 2022, 方正证券.
    \item Cakici, N., \& Zaremba, A. (2022). \textit{Salience theory and the cross-section of stock returns: International and further evidence}. Journal of Financial Economics, 146(2), 689-725.
\end{enumerate}

\section{下行已实现波动率（高频因子-波动跳跃类）}

\subsection{因子定义}
\begin{equation}
    \mathrm{RS}_t^- = \sum_{i=1}^N r_{t,i}^2 \cdot \mathrm{I}_{r_{t,i} \leq 0}
\end{equation}

\subsection{变量定义}
\begin{itemize}
    \item ① \( r_{t,i} \) 为某股票交易日 \( t \) 内的分钟对数收益率序列，如 5 分钟对数收益率序列；
    \item ② \( \mathrm{I}_{r_{t,i} \leq 0} \) 为示性函数，对数收益率小于等于 0 时取 1，否则取 0；
    \item ③ \( N \) 表示该股票在交易日 \( t \) 中特定频率的分钟级别数据个数，如 5 分钟频率通常有 48 个数据点。
\end{itemize}

\subsection{说明}
下行已实现波动率对应资产价格波动中的下行波动；下行已实现波动率较高的股票伴随着较高的下行风险，可以期望未来能获得更高的风险补偿。

\subsection{参考文献}
\begin{enumerate}
    \item Barndorff-Nielsen, O. E., Kinnebrock, S., \& Shephard, N. (2010). \textit{Measuring downside risk: realised semivariance}. In T. Bollerslev, J. Russell, \& M. Watson (Eds.), Volatility and Time Series Econometrics: Essays in Honor of Robert F. Engle (pp. 117-136). Oxford University Press.
\end{enumerate}

\section{供应链节点度（图谱网络-网路结构）}

\subsection{因子定义}
基于供应链图谱网络计算公司的节点度：
\begin{equation}
    d_i = d_i^{(\mathrm{in})} + d_i^{(\mathrm{out})}
\end{equation}
\begin{equation}
    d_i^{(\mathrm{out})} = \sum_{j \neq i} w_{ij}^{\mathrm{out}}; \quad d_i^{(\mathrm{in})} = \sum_{j \neq i} w_{ij}^{\mathrm{in}}
\end{equation}

\subsection{变量定义}
\begin{itemize}
    \item ① 供应链图谱网络边的方向为供应商指向客户：\(\mathrm{supplier} \rightarrow \mathrm{customer}\)；
    \item ② \( d_i^{(\mathrm{in})}, d_i^{(\mathrm{out})} \) 分别为公司 \( i \) 的入度（供应商数量）和出度（客户数量），计算详情见相关日历页；
    \item ③ 因子可能存在板块和市值偏好，建议对其做市值行业中心化处理。
\end{itemize}

\subsection{说明}
供应链节点度较高的公司，说明公司的供应链关系更复杂，公司可以通过多样化供应链或建立冗余系统来减少潜在的供应链中断风险（物流对冲）；供应链简单的公司，由于缺乏物流对冲会面临更高的风险，但投资者因承担了更高的风险可获得更高的风险溢价。

\subsection{参考文献}
\begin{enumerate}
    \item Jussa et al. (2015). \textit{The Logistics of Supply Chain Alpha}. Deutsche Bank Markets Research.
\end{enumerate}

\section{持续异常交易量(高频因子-成交分布类)}

\subsection{因子定义}

计算日内异常交易量：
\begin{equation}
    \mathrm{ATV}_{t,i} = \frac{\text{实际交易量}}{\text{预期交易量}} = \frac{t\text{日第}i\text{分钟交易量}}{\mathrm{mean}(\text{过去一段时间分钟交易量})}
\end{equation}

在各分钟全市场横截面上，计算每只股票该分钟异常交易量的排名百分位 \( \mathrm{rank\_ATV}_{t,i} \)；

计算日内持续异常交易量：
\begin{equation}
    \mathrm{PATV}_t = \frac{\mathrm{mean}(\mathrm{rank\_ATV}_{t,i})}{\mathrm{std}(\mathrm{rank\_ATV}_{t,i})} + \mathrm{kurt}(\mathrm{rank\_ATV}_{t,i})
\end{equation}

\subsection{变量定义}
\begin{itemize}
    \item ① 计算预期交易量的“过去一段时间”可以为过去1个交易日；分钟交易量的频率可以为5分钟频率的成交数据；
    \item ② \( \mathrm{rank\_ATV}_{t,i} \) 的均值越高，日内的异常交易量相对更高；标准差越小，日内异常交易量相对更稳定；高峰点的数据分布往往更集中。
\end{itemize}

\subsection{说明}
持续异常交易量衡量了股票日内异常交易量的持续性，取值越高，持续性越强，与未来收益负相关；持续异常的高交易量才真正代表了非理性的情绪化交易，当极端情绪持续越长，累计的股价偏差越大，随后的错误定价修正幅度也会越大，反转趋势的确定性才更大。

\subsection{参考文献}
\begin{enumerate}
    \item 任继. 2023. “持续异常交易量”选股因子 PATV. 招商证券.
\end{enumerate}

\section{羊群行为(高频因子-资金流类)}

\subsection{因子定义}

\begin{equation}
    H(i, T) = \left| \frac{B(i, T)}{B(i, T) + S(i, T)} - p_T \right| - \mathrm{AF}(i, T)
\end{equation}

\begin{equation}
    \mathrm{AF}(i, T) = \sum_{k=0}^{N_{i,T}} \binom{N_{i,T}}{k} p_T^k (1 - p_T)^{N_{i,T}-k} \left| \frac{k}{N_{i,T}} - p_T \right|
\end{equation}

\begin{equation}
    \mathrm{HB}(i, T) = H(i, T) \quad \text{if} \quad \frac{B(i, T)}{B(i, T) + S(i, T)} > p_T
\end{equation}
\begin{equation}
    \mathrm{HS}(i, T) = H(i, T) \quad \text{if} \quad \frac{B(i, T)}{B(i, T) + S(i, T)} < p_T
\end{equation}

\subsection{变量定义}
\begin{itemize}
    \item ① \( B(i, T) \) 和 \( S(i, T) \) 分别为股票 \( i \) 在交易日 \( T \) 的所有买方驱动单数量和卖方驱动单数量；
    \item ② \( p_T \) 为交易日 \( T \) 每只股票买单占其交易单位比例在该日横截面上的均值；
    \item ③ \( \mathrm{AF}(i, T) \) 为调整项，在投资者各自独立进行交易的假设下，其为个股买单比例占比相对所有股票平均买单比例占比的偏离程度的期望值；买方驱动单服从二项分布 \( B(i, T) \sim B(N_{i,T}, p_T) \)，其中，\( N_{i,T} = B(i, T) + S(i, T) \)；
    \item ④ 可根据买单占比相对横截面均值的高低，将羊群行为划分为买入羊群行为 \( \mathrm{HB}(i, T) \) 和卖出羊群行为 \( \mathrm{HS}(i, T) \)。
\end{itemize}

\subsection{说明}
日内羊群行为衡量了股票的日内买卖压力；通常认为短期羊群行为会伴随着价格的反转，即出现买入（卖出）羊群行为时，未来股价可能会下跌（上涨），而且羊群行为越剧烈，未来价格变动的幅度越大。

\subsection{参考文献}
\begin{enumerate}
    \item Bikhchandani, S., Hirshleifer, D., \& Welch, I. (1992). \textit{A Theory of Fads, Fashion, Custom, and Cultural Change as Informational Cascades}. Journal of Political Economy, 100(5), 992-1026.
    \item Christofersen, S. E., \& Tang, Y. (2010). \textit{Institutional Herding and Information Cascades: Evidence from Daily Trades}. SSRN Working Paper, No. 1572726.
    \item 朱菲菲等. 2019. 短期羊群行为的影响因素与价格效应——基于高频数据的实证验证. 金融研究, 469, 191-206.
\end{enumerate}

\section{上方/下方锁定筹码占比(行为金融因子-筹码分布)}

\subsection{因子定义}

上方锁定筹码占比：
\begin{equation}
    \mathrm{ASR\_above}_t = \frac{\sum \mathrm{vol}_{\mathrm{price}_{i,t} \in [\text{高位价格}, \mathrm{Chip\_Highest}],t}}{\sum \mathrm{vol}_{i,t}}
\end{equation}

下方锁定筹码占比：
\begin{equation}
    \mathrm{ASR\_below}_t = \frac{\sum \mathrm{vol}_{\mathrm{price}_{i,t} \in [ \mathrm{Chip\_Lowest}，\text{低位价格}],t}}{\sum \mathrm{vol}_{i,t}}
\end{equation}

其中，
\begin{equation}
    \text{高位价格} = \mathrm{close}_t \times (1 + \mathrm{volatility})
\end{equation}
\begin{equation}
    \text{低位价格} = \mathrm{close}_t \times (1 - \mathrm{volatility})
\end{equation}

\subsection{变量定义}
\begin{itemize}
    \item ① 假设某只股票在 \( t \) 日的筹码分布共有 \( n \) 个价位（\( i = 1, 2, \ldots, n \)），\( \mathrm{vol}_{i,t} \) 为第 \( t \) 日 \( i \) 价位上的筹码峰高度，\( \mathrm{price}_{i,t} \) 为该筹码峰对应的价格，\( \mathrm{Chip\_Highest} \)、\( \mathrm{Chip\_Lowest} \) 分别对应筹码最高价和筹码最低价；
    \item ② \( \mathrm{volatility} \) 为股价最大涨跌幅，主板股票的 \( \mathrm{volatility} = 10\% \)，创业板和科创板的 \( \mathrm{volatility} = 20\% \)；
    \item ③ \( \mathrm{close}_t \) 为 \( t \) 日收盘价。
\end{itemize}

\subsection{说明}
锁定筹码指的是涨跌停价格以外的筹码，该部分筹码不易交易，通常较为稳定；上方/下方锁定筹码占比是对锁定筹码占比的进一步细分，分别可用于衡量高位套牢程度和低价筹码集中程度。

\subsection{参考文献}
\begin{enumerate}
    \item 陈升锐, 姚紫薇. 2025. 筹码分布因子系统构建. 中信建投.
\end{enumerate}

\section{逐档订单失衡率（高频因子-流动性类）}

\subsection{因子定义}
\begin{equation}
    \mathrm{SOIR}_t = \frac{\sum_{i=1}^{5} w_i \mathrm{SOIR}_{i,t}}{\sum_{i=1}^{5} w_i}
\end{equation}
\begin{equation}
    \mathrm{SOIR}_{i,t} = \frac{V_{i,t}^B - V_{i,t}^A}{V_{i,t}^B + V_{i,t}^A}
\end{equation}
\begin{equation}
    w_i = 1 - (i - 1)/5, \quad i = 1, 2, 3, 4, 5
\end{equation}

\subsection{变量定义}
\begin{itemize}
    \item ① \( V_{i,t}^B \) 和 \( V_{i,t}^A \) 分别为 \( t \) 时刻第 \( i \) 档的委买量和委卖量。
\end{itemize}

\subsection{说明}
SOIR 是对各档订单失衡率的加权，可以避免某一档订单量过大对总体比率的影响；其综合衡量了盘口各档买卖委托量的不均衡程度，SOIR 为正说明市场买压大于卖压，未来价格趋向上涨，且 SOIR 的值越大，上涨的概率越高，反之亦然。

\subsection{参考文献}
\begin{enumerate}
    \item 丁鲁明, 陈升锐. 2021. 高频订单失衡及价差因子. 中信建投证券.
\end{enumerate}

\section{价量相关性趋势因子（高频因子-量价相关性类）}

\subsection{因子定义}

1. 每月月底，回溯每只股票过去 20 个交易日的价量信息，每日计算该股票分钟收盘价 \( p_t \) 与分钟成交量 \( v_t \) 的相关系数：
\begin{equation}
    \rho_t = \mathrm{corr}(v_t, p_t)
\end{equation}

2. 将每只股票的 20 个相关系数 \( \rho_t \) 对时间 \( t \) 回归，取回归系数 \( \beta \)，即
\begin{equation}
    \rho_t = \beta t + \varepsilon_t
\end{equation}
其中，\( t = 1, 2, 3, \ldots, 20 \)；

3. 将所有股票的回归系数 \( \beta \) 在横截面上剔除市值、传统价量类因子（20 日反转、20 日换手率、20 日波动率因子），得到的结果即为价量相关性趋势因子 \( \mathrm{PV\_corr\_trend} \)。

\subsection{说明}
\( \mathrm{PV\_corr\_trend} \) 越小，即价量相关系数随时间推移变小的股票，未来收益倾向于越高。

\subsection{参考文献}
\begin{enumerate}
    \item 高子剑. 2020. 高频价量相关性：意想不到的选股因子. 东吴证券.
\end{enumerate}

\section{基于收益规模的信息离散度（量价因子改进）}

\subsection{因子定义}
\begin{equation}
    \mathrm{ID_{MAG}} = -\frac{1}{N} \mathrm{sgn}(\mathrm{PRET}) \times \sum_{i=1}^{N} \mathrm{sgn}(\mathrm{Return}_i) \times w_i
\end{equation}

\subsection{变量定义}
\begin{itemize}
    \item ① \( \mathrm{PRET} \) 为过去 12 个月的累计收益率，剔除最近 1 个月的收益数据；
    \item ② \( \mathrm{Return}_i \) 为区间内第 \( i \) 个交易日的日度收益率；
    \item ③ \( w_i \) 为权重：每日对个股日度收益绝对值进行横截面上的排序分为 5 组（第 1 组取值最小），将为第 1 组到第 5 组中赋予单调递减的权重（日度收益越小权重越高），分别为 \( 5/15 \)、\( 4/15 \)、\( 3/15 \)、\( 2/15 \) 和 \( 1/15 \)。
\end{itemize}

\subsection{说明}
作者认为“一系列频繁但微小的变化对于人的吸引力远不如少数却显著的变化，投资者对于连续信息造成的股价变化是反应不足的”，信息离散度越低（信息连续性越强）越好；上述计算方式考虑了日度收益规模对信息离散度的影响。

\subsection{参考文献}
\begin{enumerate}
    \item Da, Z., Gurun, U. G., \& Warachka, M. (2014). \textit{Frog in the pan: Continuous information and momentum}. Review of Financial Studies, 27(7), 2171-2218.
\end{enumerate}

\section{日内条件在险价值（高频因子-收益分布类）}

\subsection{因子定义}

1. 给定置信区间 \( \alpha \in (0, 1) \)，使用 \( T \) 交易日的分钟 VWAR 序列计算 VCVaR：
\begin{equation}
    \mathrm{VCVaR}_T = \mathrm{CVaR}_\alpha (\mathrm{VWAR}_t)
\end{equation}
其中，
\begin{equation}
    \mathrm{VWAR} = \frac{\sum \mathrm{Return}_t \times \mathrm{Volume}_t}{\sum \mathrm{Volume}_t}
\end{equation}

2. 每日计算左侧 5\% CvaR 值，每月末对当月所有交易日的 CvaR 值求算术平均，得到月频的 CvaR 因子。

\subsection{变量定义}
\begin{itemize}
    \item ① \( \mathrm{Return}_t \) 和 \( \mathrm{Volume}_t \) 分别为股票日内的分钟收益率序列和分钟成交量序列。
\end{itemize}

\subsection{说明}
在分钟高频数据环境下，激烈成交的时段往往比稀疏成交的时段更具有价格发现功能。利用成交量加权平均收益率可以淡化成交稀疏的时间段，强化真正由于买卖力量形成的收益率，由此计算的条件在险价值表现更优。

\subsection{参考文献}
\begin{enumerate}
    \item 严佳炜. 2020. 分钟线的尾部特征. 方正证券.
\end{enumerate}

\section{个股收益最低时的价变共振（高频因子-动量反转类）}

\subsection{因子定义}

1. 计算个股在日内中最低的 5 个分钟收益率期间，对应全市场等权收益率的均值：
\begin{equation}
    \mathrm{RM\_MIN}(N)_i = \frac{\sum_{t=m_1}^{mM} \mathrm{MEAN\_MKT}_t}{M}
\end{equation}
\begin{equation}
    \mathrm{MEAN\_MKT}_t = \frac{\sum_{i=1}^I r_{i,t}}{I}
\end{equation}

2. 计算个股在日内中最低的 5 个分钟收益率期间，全市场收益率标准差的均值：
\begin{equation}
    \mathrm{STD\_MIN}_i = \frac{\sum_{t=m_1}^{mM} \mathrm{std}(r_t)}{M}
\end{equation}

3. 对 \( \mathrm{RM\_MIN} \) 做截面上的逆序排名名次归一化，并除以 \( \mathrm{STD\_MIN} \) 即为个股收益最低时的市场共振因子：
\begin{equation}
    \mathrm{CO\_MIN}_i = \frac{\mathrm{Rank}^{-1}(\mathrm{RM\_MIN}(M)_i)}{\mathrm{STD\_MIN}_i}
\end{equation}

\subsection{变量定义}
\begin{itemize}
    \item ① \( r_{i,t} \) 为个股 \( i \) 在 \( t \) 分钟的收益率，\( r_t \) 为 \( t \) 时刻的全部个股收益率序列；
    \item ② \( m_1, m_2, \ldots, m_M \) 为日内分钟收益率最低的 \( M \) 个分钟，\( M=5 \)；
    \item ③ \( \mathrm{MEAN\_MKT}_t \) 为在 \( t \) 分钟的全市场等权收益均值。
\end{itemize}

\subsection{说明}
\( \mathrm{RM\_MIN} \) 刻画了个股与市场趋势的一致性，因子值越小，说明个股是在市场表现较差时同步出现价格下跌，与全市场下跌趋势一致。个股与市场形成有效共振，股价下跌可能是市场恐慌性错杀带来的，未来股价更容易被修正，未来收益表现较好；\( \mathrm{STD\_MIN} \) 刻画了市场整体收益的分歧度，分歧度越小，市场趋势更具市场代表性，个股与市场共振越有效，将其叠加后有助于提升 \( \mathrm{RM\_MIN} \) 的表现。

\subsection{参考文献}
\begin{enumerate}
    \item 任耀, 袁元勋, 许继宏. 2024. 基于分钟数据的价变共振因子. 招商证券.
\end{enumerate}

\section{卖出反弹偏离小单因子(行为金融因子-遗憾规避)}

\subsection{因子定义}
\begin{equation}
    \mathrm{LCPS} = \frac{\sum_{i}^{N} \overline{\mathrm{price}}_{\mathrm{buy},i} \cdot \mathbf{I}_{\{P_{\mathrm{buy},i} < \mathrm{close}\}} \cdot \mathbf{I}_{\{\mathrm{vol} < \mathrm{vol}_{\mathrm{mean}}\}}}{\mathrm{close} } - 1
\end{equation}

\subsection{变量定义}
\begin{itemize}
    \item ① \( \mathrm{price}_{\mathrm{buy},i} \) 为主卖成交价；
    \item ② \( \mathrm{vol}_{\mathrm{mean}} \) 为当日所有实际成交量均值；
    \item ③ 逐笔中每笔成交买卖方向确定方式：若该笔成交为买单被拆分与多个卖单成交，则记为买方交易，反之为卖方交易；
    \item ④ 大小单划分方式：以当日所有实际成交量均值作为区分大小单的阈值，即
        \begin{equation*}
            \mathrm{vol} < \mathrm{vol}_{\mathrm{mean}}
        \end{equation*}
        为小单，反之为大单；
    \item ⑤ 此外，还可只统计尾盘期间 \([14:30,14:57)\) 的成交量，计算得到卖出反弹偏离尾盘因子。
\end{itemize}

\subsection{说明}
遗憾规避理论认为非理性的投资者在做决策时，会倾向于避免产生后悔情绪并追求自满感，避免承认之前的决策失误；对于卖出后价格反弹（当天卖出的股票收盘价高于卖出的成本价）的投资者会避免承认决策失误而存在惜售心理，不愿卖出下跌浮亏的股票；卖出反弹偏离因子从成交价格偏离维度来量化上述心理，因子值越高，卖出价格越低，成交价格偏离程度越大，未来买回动力较弱，预期收益较低。

\subsection{参考文献}
\begin{enumerate}
    \item 高智威. 2023. 基于逐笔成交数据的遗憾规避因子. 国金证券.
\end{enumerate}

\section{反转残差成交不平衡因子(高频因子-成交分布类)}

\subsection{因子定义}

1. 每个交易日，统计每只股票当日所有逐笔成交中主动买入和主动卖出的成交单数，并计算如下每日成交单不平衡指标：
\begin{equation}
    \text{成交单不平衡} = \frac{\text{主买成交单数} - \text{主卖成交单数}}{\text{主买成交单数} + \text{主卖成交单数}}
\end{equation}

2. 每月底回看过去 20 个交易日，计算 20 日成交单不平衡指标的平均值，做横截面市值中性化处理，得到成交不平衡因子；

3. 将成交不平衡因子，对过去 20 个交易日的累计涨跌幅进行正交化处理，得到反转残差成交不平衡因子。

\subsection{说明}
成交不平衡因子衡量了净买单的强弱程度，若股票在过去一段时间是买方占据主动成交的优势，其未来收益表现会较好；当月的净买单数量较多通常会造成当月涨跌幅较高，而当月涨跌幅与未来涨跌幅呈现显著的负相关性，因此会削弱成交不平衡因子对未来收益的正向预测能力，可将成交不平衡因子对同期涨跌幅进行正交化处理，得到反转残差成交不平衡因子，提升因子稳定性。

\subsection{参考文献}
\begin{enumerate}
    \item 沈芷琦, 刘富兵, 阮俊烨. 2024. 逐笔买卖差异中的选股信息: 条件成交不平衡因子. 国盛证券.
\end{enumerate}

\section{特质隔夜收益波动率(高频因子-波动跳跃类)}

\subsection{因子定义}

截面上对所有股票过去 20 天的隔夜收益率序列进行主成分分析 PCA，提取第一主成分 \( F_1 \) 为隐式因子溢价：
\begin{equation}
    \mathrm{OvernightRet}_{i,t} = \frac{\mathrm{Open}_{i,t}}{\mathrm{Close}_{i,t-1}} - 1
\end{equation}

对每只股票 \( i \)，以过去 20 天的隔夜收益率序列为 \( y \)，以隐式因子 \( F_1 \) 序列为 \( x \)，进行时序回归：
\begin{equation}
    \mathrm{OvernightRet}_{i,t} = \alpha_{i,t} + \beta_i^{F_1} \times F_1 + \cdots + \varepsilon_i
\end{equation}

时序回归返回的残差序列的标准差即为隐式因子框架下的特质波动因子：
\begin{equation}
    \mathrm{OvernightRetIV} = \mathrm{std}(\varepsilon_i)
\end{equation}

\subsection{变量定义}
\begin{itemize}
    \item ① 在进行截面主成分分析时，若第一主成分的解释力度低于 20\%，则选取解释力最大的前两个主成分进行回归；
    \item ② 还可进一步计算特异度因子：\( \mathrm{IVR} = 1 - R^2 \) 上述回归模型拟合优度；
    \item ③ 还可进一步计算特质偏度因子：\( \mathrm{IS} = \frac{\sum \varepsilon_{it}^3}{\mathrm{IV}^3} \)。
\end{itemize}

\subsection{说明}
计算隐式因子框架下的特质波动率时，直接通过主成分分析对资产收益序列降维来提取共同波动，而无需显式的借助多因子模型估计共同风险，系统性的解决传统多因子模型遗漏变量、估计误差的问题，计算简单；此处提取的是隔夜收益率序列的共同波动。

\subsection{参考文献}
\begin{enumerate}
    \item 张欣鹏, 杨怡玲, 刘璐. 2022. 隐式框架下的特质因子改进. 国信证券.
\end{enumerate}

\section{基于投资者交流的关联动量(图谱网络-动量溢出)}

\subsection{因子定义}
\begin{equation}
    \mathrm{MRR}_{m,t} = \frac{1}{5} \sum_{j=1}^{5} \mathrm{Ret}_{j,m,t}
\end{equation}

\subsection{变量定义}
\begin{itemize}
    \item ① \( \mathrm{Ret}_{j,m,t} \) 为目标公司 \( m \) 相关的关联公司 \( j \) 在 \( t \) 日的日度收益率；
    \item ② 投资者交流信息来自于东方财富论坛 \url{http://guba.eastmoney.com/} 中的帖子，当前论坛对应的股票为目标股票 \( m \)，投资者在该论坛上发帖交流时提及的其他股票为关联股票；每月统计目标股票 \( m \) 论坛上关联股票被提及的次数，取被提及次数最多的 5 只关联股票进行收益等权复合即得到上述关联动量。
\end{itemize}

\subsection{说明}
投资者交流过程中被提及的关联公司与目标公司间存在收益联动，而且关联股票被投资者提及的越频繁，股票间收益联动越强。

\subsection{参考文献}
\begin{enumerate}
    \item Jiang, L., Liu, J., \& Yang, B. (2019). \textit{Communication and comovement: Evidence from online stock forums}. Financial Management, 48(3), 805-847.
\end{enumerate}

\section{大单买入占比(高频因子-资金流类)}

\subsection{因子定义}
\begin{equation}
    \frac{1}{T} \sum_{n=t}^{t-T+1} \frac{\sum_{j=1}^{N} \text{大买单成交额}_{i,j,n}}{\sum_{j=1}^{N} \text{成交额}_{i,j,n}}
\end{equation}

\subsection{变量定义}
\begin{itemize}
    \item ① 基于逐笔成交数据中的叫买与叫卖单号将逐笔成交数据合成为买卖单数据；
    \item ② 根据买卖单分布界定大小单，如可将多日买卖单成交单数据对数调整后的“均值+1倍标准差”作为大单筛选阈值；
    \item ③ \( i, j, n \) 分别表示第 \( i \) 只股票在第 \( n \) 个交易日内的第 \( j \) 分钟的成交额；
    \item ④ 月度选股下，\( T=20 \) 个交易日；周度选股下，\( T=5 \) 个交易日。
\end{itemize}

\subsection{说明}
大单类因子刻画了大资金的交易行为，大资金往往具有信息优势，一般受到大资金关注的股票，未来通常具有更好的表现。

\subsection{参考文献}
\begin{enumerate}
    \item 冯佳睿, 袁林青. 2021. 大单的精细化处理与大单因子重构. 海通证券.
    \item 冯佳睿, 袁林青. 2019. 买卖单数据中的 Alpha. 海通证券.
\end{enumerate}

\section{单位成交额的百分比实现价差(高频因子-流动性类)}

\subsection{因子定义}
\begin{equation}
    \mathrm{PRSI}_{i,t} = \frac{\mathrm{PRS}_{i,t}}{\$ \mathrm{vol}_{i,t}}
\end{equation}
\begin{equation}
    \mathrm{PRS}_{i,t} = 2 D_{i,k} \left( \ln(P_{i,t}) - \ln\left(\frac{\mathrm{Ask}_{i,t+5} + \mathrm{Bid}_{i,t+5}}{2}\right) \right)
\end{equation}

\subsection{变量定义}
\begin{itemize}
    \item ① \( \mathrm{Ask}_{i,t+5}, \mathrm{Bid}_{i,t+5} \) 分别为股票 \( i \) 在 \( t+5 \) 时刻的卖 1 价、买 1 价；
    \item ② \( P_{i,t} \) 为股票 \( i \) 在 \( t+5 \) 时刻的收盘价；
    \item ③ \( D_{i,k} \) 为买卖标识，买单记为 1，卖单记为 -1；
    \item ④ \( D_{i,k} \) 买卖单判断标准：若当前收盘价高于（低于）最优买卖报价平均值，则认为是买（卖）单；当前收盘价等于最优买卖报价平均值时，若收盘价高于（低于）上一笔价格，则认为是买（卖）单；
    \item ⑤ \( \$ \mathrm{vol}_{i,t} \) 为 \( t \) 日当天成交额。
\end{itemize}

\subsection{说明}
单位成交额的百分比实现价差属于“每美元成本 cost-per-dollar-volume”形式的流动性，代表交易中每单位交易额的边际交易成本，反映流动性效率；股票的交易成本越高，流动性越差，未来会有正向的流动性风险溢价。

\subsection{参考文献}
\begin{enumerate}
    \item 陈升悦. 2024. 流动性高频因子构建与投资者注意力因子分析报告. 中信建投.
\end{enumerate}

\section{筹码穿透率(行为金融因子-筹码分布)}

\subsection{因子定义}
\begin{equation}
    \mathrm{PTR} = \frac{\mathrm{winner}_t - \mathrm{winner}_{t-1}}{\mathrm{turnover}_t}
\end{equation}
\begin{equation}
    \mathrm{winner}_t = t \text{ 日盈利筹码占总筹码的比例}
\end{equation}

\subsection{变量定义}
\begin{itemize}
    \item ① \( \mathrm{turnover}_t \) 为股票 \( t \) 日换手率。
\end{itemize}

\subsection{说明}
该筹码穿透率的计算方式为当天新增解套筹码比例除以当天换手率，用于反映穿透筹码的能力；因子值越大，说明相同换手率使套牢筹码解套的效率越高（少量换手就能让大量历史筹码解套），股票筹码的通透性越好，趋势延续性强。

\subsection{参考文献}
\begin{enumerate}
    \item 陈升锐, 姚紫薇. 2025. 筹码分布因子系统构建. 中信建投.
\end{enumerate}

\section{成交量最高时刻笔均成交金额(高频因子-成交分布类)}

\subsection{因子定义}

① 将日内 240 个 1 分钟数据集集合标记为 \( T = \{t_1, t_2, \dots, t_{240}\} \)，将分时成交量最高的 10\% 时刻标记为集合 \( P = \{t_{i1}, t_{i2}, \dots, t_{ik}\} \)；

② 对成交量最高时刻进行汇总，根据其成交额 \( \mathrm{Amt}_t \) 之和除以成交笔数 \( \mathrm{DealNum}_t \) 之和，构造成交量最高时刻的笔均成交金额 \( \mathrm{ATD}_{\mathrm{VolumeHigh10\%}} \)：
\begin{equation}
    \mathrm{ATD}_{\mathrm{VolumeHigh10\%}} = \frac{\sum_{t \in P} \mathrm{Amt}_t}{\sum_{t \in P} \mathrm{DealNum}_t}
\end{equation}

③ 同理，将全天的成交金额除以全天的成交笔数，得到全天的笔均成交金额 \( \mathrm{ATD}_T \)：
\begin{equation}
    \mathrm{ATD}_T = \frac{\sum_{t \in T} \mathrm{Amt}_t}{\sum_{t \in T} \mathrm{DealNum}_t}
\end{equation}

④ 将成交量最高时刻的笔均成交金额除以全天的笔均成交金额，即得到去量纲化后的标准化因子 \( \mathrm{SATD}_{\mathrm{VolumeHigh10\%}} \)：
\begin{equation}
    \mathrm{SATD}_{\mathrm{VolumeHigh10\%}} = \frac{\mathrm{ATD}_{\mathrm{VolumeHigh10\%}}}{\mathrm{ATD}_T}
\end{equation}

\subsection{说明}
笔均成交额类因子通过汇总“特殊的、具有较多信息含量”时刻的成交金额情况来刻画主力资金的行为动向；笔均成交金额越大，说明该特殊时间内资金优势越明显，属于主力资金的可能性越大。“成交量最高时刻”从成交热度角度刻画主力资金交易行为，也可按同样逻辑计算“成交量最低时刻”相关因子，前者对未来收益有正向预测能力，后者有负向预测能力。

\subsection{参考文献}
\begin{enumerate}
    \item 张欣鹏, 张宇. 2025. 日内特殊时刻蕴含的主力资金 Alpha 信息. 国信证券.
\end{enumerate}

\section{筹码收益增强(量价因子改进)}

\subsection{因子定义}
利用过去 250 日的成交额 \( \mathrm{Amt} \)、换手率 \( \mathrm{Turnover} \) 等数据，对股票 \( i \) 在 \( T \) 日计算得到 \([T-250, T]\) 之间每日的留存筹码 \( \mathrm{RsdAmt}_{T-k}^T \) 和成交均价 \( \mathrm{Price}_{T-k}^{\mathrm{avg}} \)，即为该股票在 \( T \) 日的筹码分布估计：
\begin{equation}
    \mathrm{RsdAmt}_{T-k}^T = \mathrm{Amt}_{T-k} \cdot \prod_{t=T-k+1}^{T} (1 - \mathrm{Turnover}_t), \quad k \in (0, 250]
\end{equation}

计算最新日成交均价相对历史筹码的加权成本的相对收益，即为最新日的筹码收益因子：
\begin{equation}
    \mathrm{holding\_ret} = \frac{\sum_t \mathrm{Price}_t^{\mathrm{avg}} \times \mathrm{RsdAmt}_t^T}{\sum_t \mathrm{RsdAmt}_t^T}
\end{equation}

对全市场的筹码收益因子加权计算均值作为 A 股市场的筹码收益 \( \mathrm{mkt\_holding\_ret} \)，用以调节个股筹码收益因子 \( \mathrm{holding\_ret} \) 的动量效应和反转效应，即得到筹码收益调整因子：
\begin{equation}
    \mathrm{holding\_ret\_adj} = \mathrm{holding\_ret} \times \mathrm{sign}(\mathrm{mkt\_holding\_ret})
\end{equation}

将筹码收益调整因子的多头端与反转因子 \( \mathrm{ret20} \) 的空头端进行合成，构建筹码收益增强因子：
\begin{equation}
     (1 - \mathrm{ret20}) \times \mathrm{ret20} + (1 - \mathrm{ret20}) \times \mathrm{holding\_ret\_adj}
\end{equation}

\subsection{说明}
筹码收益调整因子优点在于能够对多头组有较好区分，反转因子优点在于对于空头组具有显著的区分效果，筹码收益增强因子是两者的合成。

\subsection{参考文献}
\begin{enumerate}
    \item 魏建榕, 张翔. 2025. 筹码结构视角下的动量反转融合. 开源证券.
\end{enumerate}

\section{基于大小单和成交时长的精选复合交易占比(高频因子-资金流类)}

\subsection{因子定义}
\begin{equation}
    \mathrm{VolumeLongBigSelect} = \frac{\text{大买\&漫长的大卖} + \text{漫长的大买\&非漫长的大卖}}{\text{总成交量}} 
\end{equation}   
\begin{equation}
    - \frac{\text{非漫长的非大买\&非漫长的大卖} + \text{非漫长的大买\&非漫长的非大卖}}{\text{总成交量}}
\end{equation}

\subsection{变量定义}
\begin{itemize}
    \item ① 大单和非大单的判断标准为：先将同一委买 ID（或委卖 ID）对应的实际成交量进行加总，再将全部实际成交的委买 ID（或委卖 ID）成交量进行降序排列，最后将成交量大于前 10\% 分位点的委买单（或委卖单）记为“大买单”（或“大卖单”），要剔除开盘集合竞价成交记录；
    \item ② 漫长订单判断标准为：先计算不同委买 ID（或委卖 ID）订单的成交时长，再将全部实际成交的委买 ID（或委卖 ID）成交时长进行降序排列，最后将成交时长大于前 10\% 分位点的委买单（或委卖单）记为“漫长买单”（或“漫长卖单”），要剔除开盘集合竞价的成交记录。
\end{itemize}

\subsection{说明}
精选复合交易占比因子同时考虑了逐笔委托单的“大小单”和“成交时长”属性，是对成交单更细致的划分；基于细致划分下的因子表现，选取有效性较高的成交单类型进行复合，进而得到表现更优的复合因子。

\subsection{参考文献}
\begin{enumerate}
    \item 张欣鹏, 张宇. 2024. 高频订单成交数据蕴含的 Alpha 信息. 国信证券.
\end{enumerate}

\section{模糊数量比(高频因子-成交分布类)}

\subsection{因子定义}

1. 对股票 \( i \) 计算 \( n \) 日内的 1 分钟收益率；计算第 \( t \) 分钟的波动率，即 \([t-4, t]\) 区间内分钟收益率的标准差；计算第 \( t \) 分钟的模糊性，即 \([t-4, t]\) 区间内上述分钟波动率的标准差；

2. 将日内模糊性大于日内模糊性均值的时间段记为“起雾时刻”；

3. 计算“起雾时刻”的分钟成交量均值为“雾中数量”，计算日内所有时刻的分钟成交量均值为“总体数量”；

4. 计算“日模糊数量比”：
\begin{equation}
    \text{日模糊数量比} = \frac{\text{雾中数量}}{\text{总体数量}}
\end{equation}

\subsection{变量定义}
\begin{itemize}
    \item ① 计算上述因子时，剔除开盘和收盘数据，仅保留日内交易数据，所以开盘前 9 分钟的模糊性为空值；
    \item ② 每月末可对最近 20 天的上述因子求均值和标准差并将两者等权合并，即得到月度“模糊数量比”因子。
\end{itemize}

\subsection{说明}
模糊数量比因子从量维度统计了投资者在模糊性较大时的成交程度，同样衡量了投资者对模糊性的厌恶程度，与未来收益负相关；投资者对波动率模糊性的厌恶心理会使得投资者急于卖出而产生过度反应，未来很可能会补涨。

\subsection{参考文献}
\begin{enumerate}
    \item 曹春晓. 2022. 波动率的波动率与投资者模糊性厌恶. 方正证券.
    \item Kostopoulos, D., Meyer, S., \& Uhr, C. (2021). \textit{Ambiguity about volatility and investor behavior}. Journal of Financial Economics.
\end{enumerate}

\section{收盘委托价差(高频因子-流动性类)}

\subsection{因子定义}
\begin{equation}
    \mathrm{CPQSI}_{i,t} = \frac{\mathrm{CPQS}_{i,t}}{\$ \mathrm{vol}_{i,t}}
\end{equation}
\begin{equation}
    \mathrm{CPQS}_{i,t} = \frac{\mathrm{Closing\_Ask}_{i,t} - \mathrm{Closing\_Bid}_{i,t}}{(\mathrm{Closing\_Ask}_{i,t} + \mathrm{Closing\_Bid}_{i,t})/2}
\end{equation}

\subsection{变量定义}
\begin{itemize}
    \item ① \( \mathrm{Closing\_Ask}_{i,t} \)、\( \mathrm{Closing\_Bid}_{i,t} \) 分别为股票 \( i \) 在 \( t \) 日收盘时的卖 1 价、买 1 价；
    \item ② \( \$ \mathrm{vol}_{i,t} \) 为 \( t \) 日当天成交额。
\end{itemize}

\subsection{说明}
收盘委托价差衡量了收盘前的交易成本，因子值越大，买卖价差越大，流动性越差，未来可获得流动性风险溢价。

\subsection{参考文献}
\begin{enumerate}
    \item 陈升悦. 2024. 流动性高频因子构建与投资者注意力因子分析报告. 中信建投.
\end{enumerate}

\section{最大异常日收益率(行为金融因子-投资者注意力)}

\subsection{因子定义}
\begin{equation}
    \mathrm{ABNRETD} = \max_t |r_{i,t} - \bar{r}_t|
\end{equation}

\subsection{变量定义}
\begin{itemize}
    \item ① \( r_{i,t} \) 为股票 \( i \) 在 \( t \) 交易日的收益率；
    \item ② \( \bar{r}_t \) 为交易日 \( t \) 的全市场收益率；
    \item ③ \( m \) 为交易日窗口，一般取为月度。
\end{itemize}

\subsection{说明}
最大异常日收益率 \( \mathrm{ABNRETD} \) 为每月日频收益率与市场收益率绝对差值的最大值，属于极端日收益类因子；极端收益（过高收益或过低收益）的股票相较其他股票更能吸引投资者的注意，因子值越高，投资者对该股票的关注度越高；而有限注意力导致投资者更倾向于交易引起他们关注的股票，投资者关注度越高的股票通常意味着越多的投资购买；但 A 股市场的做多与做空不对称，使得投资者非理性买入高于非理性卖出，进而导致关注度高的股票由于投资者净买入而存在溢价，未来也更可能出现反转。

\subsection{参考文献}
\begin{enumerate}
    \item Cosemans, M., \& Frehen, R. (2021). \textit{Salience theory and stock prices: Empirical evidence}. Journal of Financial Economics, 140(2), 460-483.
    \item 陈升锐. 2024. 投资者有限关注及注意力捕捉与溢出. 中信建投.
\end{enumerate}

\section{交易量同步的知情交易概率(高频因子-资金流类)}

\subsection{因子定义}

\begin{equation}
    \mathrm{VPIN} = \frac{\alpha\mu}{\alpha\mu + 2\varepsilon} \approx \frac{E[|V_\tau^S - V_\tau^B|]}{V} = \frac{\sum_{\tau=1}^n |V_\tau^S - V_\tau^B|}{nV}
\end{equation}

\begin{equation}
    V_\tau^B = \sum_{i=t(\tau-1)+1}^{t(\tau)} V_i \cdot Z \left( \frac{P_i - P_{i-1}}{\sigma_{\Delta P}} \right)
\end{equation}

\begin{equation}
    V_\tau^S = \sum_{i=t(\tau-1)+1}^{t(\tau)} V_i \cdot \left[ 1 - Z \left( \frac{P_i - P_{i-1}}{\sigma_{\Delta P}} \right) \right] = V - V_\tau^B
\end{equation}

\subsection{变量定义}
\begin{itemize}
    \item ① 作者将每日的交易划分为 \( n \) 个交易量相同的时间段或交易量桶 volume bucket，每个交易桶的交易量为 \( V \)；
    \item ② \( V_\tau^S \) 和 \( V_\tau^B \) 分别为第 \( \tau \) 个交易量桶中的卖方和买方交易量；
    \item ③ \( V_i \) 和 \( P_i \) 为某个交易量桶内第 \( i \) 笔交易的交易量和价格，\( \Delta P = P_i - P_{i-1} \) 为价格序列的环比增量，\( \sigma_{\Delta P} \) 为当日价格环比增量序列的标准差；
    \item ④ \( Z \left( \frac{P_i - P_{i-1}}{\sigma_{\Delta P}} \right) \) 为每个价格变动的正态分布分位数。
\end{itemize}

\subsection{说明}
VPIN 是知情交易概率的一种非参数估计方法，衡量了等交易量时间段内 volume time（而非物理时间 clock time）的交易量不平衡性；知情交易概率指在一段时间内，拥有信息优势的知情交易者提交的订单数量占总委托单数量的比例。用以刻画该段时间内的信息不对称程度，通常与未来收益正相关，取值越大，信息不对称程度越大，投资者要求的风险回报越高。

\subsection{参考文献}
\begin{enumerate}
    \item Easley, D., Prado, M., \& O'Hara, M. (2012). \textit{Flow toxicity and liquidity in a high frequency world}. Review of Financial Studies, 25(5), 1457-1493.
\end{enumerate}

\section{跳跃显著程度加权的上下行跳跃波动不对称性(高频因子-波动跳跃类)}

\subsection{因子定义}
\begin{equation}
    \mathrm{TSRJV} = \frac{\sum_{t=1}^{D} \left( \frac{T_{t,i}}{\Phi^{-1}_{1-\alpha}} \cdot \mathrm{SRJV}_{t,i} \right)}{\sum_{t=1}^{D} \frac{T_{t,i}}{\Phi^{-1}_{1-\alpha}}}
\end{equation}
\begin{equation}
    \mathrm{SRJV}_t = \mathrm{RVLP}_t - \mathrm{RVJN}_t
\end{equation}
\begin{equation}
    T_{t,i} = \frac{\mathrm{BV}_t}{N^{-1}\sqrt{\hat{\Omega}_{\mathrm{SwV}}}} \left(1 - \frac{\mathrm{RV}_t}{\mathrm{SwV}_t}\right) \xrightarrow{d} N(0, 1)
\end{equation}

\subsection{变量定义}
\begin{itemize}
    \item ① \( \mathrm{SRJV}_{t,i} \) 为股票 \( i \) 在交易日 \( t \) 内的上下行跳跃波动不对称性，计算细节可参考相关日历页；
    \item ② \( T_{t,i} \) 为检验股票 \( i \) 在交易日 \( t \) 的跳跃是否显著的检验统计量，计算细节可参考相关日历页；
    \item ③ \( D = 20 \) 个交易日。
\end{itemize}

\subsection{说明}
通过跳跃显著程度加权能提升上下行跳跃波动不对称性因子的表现。

\subsection{参考文献}
\begin{enumerate}
    \item 周飞鹏, 罗军, 安宁宁. 2022. 再谈股价跳跃因子. 广发证券.
\end{enumerate}

\section{相似业务收益联动因子(图谱网络-动量溢出)}

\subsection{因子定义}
\begin{equation}
    \mathrm{Linkage}_{i,t} = \frac{\sum_{j=1}^N \mathrm{SIM}_{i,j,t} * \mathrm{Ret}_{j,t}}{\sum_{j=1}^N \mathrm{SIM}_{i,j,t}} - \mathrm{Ret}_{i,t}
\end{equation}
\begin{equation}
    \mathrm{SIM}_{i,j,t} = \frac{P_{i,t} \cdot P_{j,t}}{\|P_{i,t}\| \cdot \|P_{j,t}\|}
\end{equation}

\subsection{变量定义}
\begin{itemize}
    \item ① \( P_{i,t} \) 为取值只有 0 和 1 的业务关键词向量，通过自然语言处理技术处理 A 股上市公司年度财务报告附注中关于公司基本情况的描述部分内容得到；
    \item ② \( \mathrm{SIM}_{i,j,t} \) 为公司 \( i \) 和公司 \( j \) 在 \( t \) 时期的业务相似度；
    \item ③ \( \mathrm{Ret}_{j,t} \) 为公司 \( j \) 在 \( t \) 月的收益率。
\end{itemize}

\subsection{说明}
相似业务收益联动因子度量了 \( t \) 期下与 \( i \) 公司相似的公司在当月超过 \( i \) 公司的收益率，该值越大，说明 \( i \) 公司越有可能在下一期出现一定的补涨行情。

\subsection{参考文献}
\begin{enumerate}
    \item 叶尔乐. 2023. 财报文本中公司竞争信息刻画与 Alpha 构建. 民生证券.
\end{enumerate}

\section{极端收益反转(高频因子-动量反转类)}

\subsection{因子定义}

基于股票 \( i \) 的分钟行情数据，每日寻找出日内收益最极端的那根 bar，即 \( S \) 最大的那根 bar：
\begin{equation}
    S = |r_{i,t} - \mathrm{median}(\{r_{i,t}\})|
\end{equation}

每月底，将过去 20 天的每天 \( S \) 最高那根 bar 收益率求均值得到最极端收益；同理，每月底，将过去 20 天的每天 \( S \) 最高那根 bar 前一分钟的收益率求均值得到前一分钟收益；将两者截面排序值等权求和，即得极端收益率反转因子：
\begin{equation}
    \mathrm{rank}\left(\frac{1}{20} \sum r_{i,t_{\max(S)}}\right) + \mathrm{rank}\left(\frac{1}{20} \sum r_{i,t_{\max(S)}-1}\right)
\end{equation}

\subsection{说明}
日内，最极端收益前呈反转特性，最极端收益后呈动量特性；切割出的极端收益反转因子具有强反转效应。

\subsection{参考文献}
\begin{enumerate}
    \item 魏建榕, 盛少成, 苏良. 2022. 日内极端收益前后的反转特性与因子构建. 开源证券.
\end{enumerate}

\section{分钟成交额自相关性(高频因子-量价相关性类)}

\subsection{因子定义}
\begin{equation}
    \mathrm{ACMA}_i = \mathrm{corr}(\mathrm{Amount}_{t-i}, \mathrm{Amount}_t)
\end{equation}

\subsection{变量定义}
\begin{itemize}
    \item ① \( \mathrm{Amount}_t \) 为交易日内的分钟成交额；
    \item ② \( i \) 为滞后阶数，如 \( i=1 \)，对应滞后 1 阶的自相关系数。
\end{itemize}

\subsection{说明}
日内分钟成交额序列有显著的自相关性；自相关性越强的股票，微观上可理解为股票受到事件刺激频率更高，对信息反应时容易羊群交易，羊群效应显著，未来收益表现差。

\subsection{参考文献}
\begin{enumerate}
    \item 尹治技. 2020. 高频视角下成交额蕴藏的 Alpha. 华安证券.
\end{enumerate}

\section{LSV 模型(行为金融因子-羊群效应)}

\subsection{因子定义}
\begin{equation}
    \mathrm{LSV}_{i,t} = |p_{i,t} - p_t| - E|p_{i,t} - p_t|
\end{equation}
\begin{equation}
    E|p_{i,t} - p_t| = \mathrm{AF}(i,t) = \sum_{k=1}^{N_{i,t}} \binom{N_{i,t}}{k} p_t^k (1 - p_t)^{N_{i,t} - k} \left| \frac{k}{N_{i,t}} - p_t \right|
\end{equation}
\begin{equation}
    p_{i,t} = \frac{B(i,t)}{B(i,t) + S(i,t)}, \quad p_t = \frac{\sum_{i=1}^N p_{i,t}}{N}, \quad N_{i,t} = B(i,t) + S(i,t)
\end{equation}

\subsection{变量定义}
\begin{itemize}
    \item ① \( B(i,t), S(i,t) \) 分别为区间 \( t \) 内买入股票和卖出股票的总数量；
    \item ② \( N_{i,t} \) 为持仓发生变动的基金数量；
    \item ③ \( p_{i,t} \) 为股票 \( i \) 在区间 \( t \) 内买入股票基金占持仓变动基金的比例；
    \item ④ 此处使用的是基金年报和半年报持仓数据，假设有 300 只基金加仓股票 \( i \)，则 \( B(i,t) = 300 \)；
    \item ⑤ \( N \) 为全市场股票总数。
\end{itemize}

\subsection{说明}
LSV 模型从基金持股变动数据出发判断一只股票是否被集中买入或集中卖出，捕捉基金“抱团”现象。基金“抱团”现象是股票出现羊群效应的潜在驱动因子之一。

\subsection{参考文献}
\begin{enumerate}
    \item Lakonishok, J., Shleifer, A., \& Vishny, R. W. (1992). \textit{The Impact of Institutional Trading on Stock Prices}. Journal of Financial Economics, 32(1), 23-43.
\end{enumerate}

\section{尾盘成交额占比（高频因子-成交分布类）}

\subsection{因子定义}
\begin{equation}
    \mathrm{APL}_{t,d} = \frac{\sum_{i=240-t}^{240} \mathrm{Amount}_{i,d}}{\mathrm{TotalAmount}_d}
\end{equation}
\begin{equation}
    \mathrm{APL}_t = \frac{\sum_{d=1}^{N} \mathrm{APL}_{t,d}}{N}
\end{equation}

\subsection{变量定义}
\begin{itemize}
    \item ① \( \mathrm{Amount}_{i,d} \) 为交易日 \( d \) 内第 \( i \) 分钟的成交额；
    \item ② \( t \) 为尾盘最后的 \( t \) 分钟，如 \( t=30 \) 对应尾盘 30 分钟；
    \item ③ \( N \) 为日频以上的时间区间内的交易日天数，如 \( N=20 \) 对应的月度因子。
\end{itemize}

\subsection{说明}
尾盘成交额占比比早盘成交额占比有更强的收益预测能力。与未来收益负相关，尾盘成交额占比高的股票存在知情交易者比例更多，资金获得利好消息后选择在噪声交易者更少的尾盘买入；取值高可识别为交易驱动型股票，反之为基本面驱动型股票。

\subsection{参考文献}
\begin{enumerate}
    \item 尹治柱. 2020. 高频视角下成交额蕴藏的 Alpha. 华安证券.
\end{enumerate}

\section{多层次订单失衡(高频因子-流动性类)}

\subsection{因子定义}

\begin{equation}
    \mathrm{MOFI}_{\mathrm{weight}} = \frac{\sum_{i=1}^{5} w_i \times \mathrm{OFI}_t^{(i)}}{\sum_{i=1}^{5} w_i}, \quad w_i = \frac{i}{5}, \quad i = 1, 2, 3, 4, 5
\end{equation}

\begin{equation}
    \Delta V_t^B = 
    \begin{cases} 
        -V_{t-1}^B, & P_t^B < P_{t-1}^B \\ 
        V_t^B - V_{t-1}^B, & P_t^B = P_{t-1}^B \\ 
        V_t^B, & P_t^B > P_{t-1}^B 
    \end{cases}
\end{equation}

\begin{equation}
    \Delta V_t^A = 
    \begin{cases} 
        V_t^A, & P_t^A < P_{t-1}^A \\ 
        V_t^A - V_{t-1}^A, & P_t^A = P_{t-1}^A \\ 
        -V_{t-1}^A, & P_t^A > P_{t-1}^A 
    \end{cases}
\end{equation}

\begin{equation}
    \mathrm{OFI}_t = \Delta V_t^B - \Delta V_t^A
\end{equation}

\subsection{变量定义}
\begin{itemize}
    \item ① \( \mathrm{OFI}_t^{(i)} \) 为第 \( i \) 档下的订单失衡因子；
    \item ② \( P_t^B \) 和 \( P_t^A \) 分别为 \( t \) 时刻的买一价和卖一价；
    \item ③ \( V_t^B \) 和 \( V_t^A \) 分别为 \( t \) 时刻的买一量和卖一量。
\end{itemize}

\subsection{说明}
MOFI 衡量了不同档位订单失衡的加权累积影响，能更准确量化订单失衡对股价的短期和长期影响。短期内，订单失衡通常与未来收益正相关；中长期内，随着买卖压力失衡的消失，股票的超额收益出现均值回复，与收益相关性会减弱。

\subsection{参考文献}
\begin{enumerate}
    \item 丁鲁明, 陈升锐. 2021. 多层次订单失衡及订单斜率因子. 中信建投证券.
\end{enumerate}

\section{趋势资金相对均价因子(量价因子改进)}

\subsection{因子定义}

1. 每个交易日 \( t \)，回看过去 \( k \) 个交易日 \( t-k, t-(k-1), \ldots, t-1 \)，计算过去 \( k \) 日分钟成交量的 \( m\% \) 分位数，定义为成交量阈值，\( k=5, m=90 \)；

2. 将交易日 \( t \) 的每分钟成交量与上述阈值进行对比，若某一分钟的成交量大于上述阈值，则该分钟的交易被定义为趋势资金的交易；

3. 每个交易日，计算如下因子：
\begin{equation}
    \frac{\text{趋势资金的成交量加权平均价格}}{\text{所有分钟的成交量加权平均价格}} - 1
\end{equation}

4. 每月月底回看过去 20 个交易日，计算上述因子的平均值，并做横截面市值、中信一级行业中性化处理。

\subsection{说明}
趋势资金相对均价因子值较大，则表明趋势资金的交易价格相对较高，其出货的可能性更大，后市看空；若因子值较小，则表明趋势资金的交易价格较低，其交易行为更有可能是逢低买入，后市看涨。

\subsection{参考文献}
\begin{enumerate}
    \item 沈亚琦, 刘富兵. 2024. 基于趋势资金日内交易行为的选股因子. 国盛证券.
\end{enumerate}

\section{修正的净流入率(高频因子-资金流类)}

\subsection{因子定义}

1. 分别计算主力资金的买入金额 \( B_P \) 和卖出金额 \( S_P \)，即 Wind 的超大单、大单和中单的合并定义为具有市场定价能力的“广义主力资金”：
\begin{equation}
    B_P = B_{\mathrm{Extra}} + B_{\mathrm{Large}} + B_{\mathrm{Med}}
\end{equation}
\begin{equation}
    S_P = S_{\mathrm{Extra}} + S_{\mathrm{Large}} + S_{\mathrm{Med}}
\end{equation}

2. 利用日涨跌幅来修正主力资金的买入金额和卖出金额：
\begin{equation}
    \ln(B_P / S_P) = \alpha + \beta \times \mathrm{Ret} + \varepsilon
\end{equation}
\begin{equation}
    \hat{B}_P = \frac{e^{\varepsilon}}{1 + e^{\varepsilon}} \times (B_P + S_P)
\end{equation}
\begin{equation}
    \hat{S}_P = \frac{1}{1 + e^{\varepsilon}} \times (B_P + S_P)
\end{equation}

3. 计算修正后的主力资金的净流入占比 \( \mathrm{CNIR} \)：
\begin{equation}
    \mathrm{CNIR} = \frac{\sum_{i=1}^N (\hat{B}_{P,i} - \hat{S}_{P,i})}{\sum_{i=1}^N (\hat{B}_{P,i} + \hat{S}_{P,i})}
\end{equation}

\subsection{说明}
从资金流动力学的角度理解，大单资金流与涨跌幅呈正相关缘自主力资金买卖的不平衡，通过选取买入卖出金额比作为代理变量（IMB），并做截面回归的方法，可以剥离反转因素的影响以提纯资金流的 Alpha 信息；具有定价权的主力资金并非仅是 20 万元以上的大单，基于小金额标准（如 2 万元）识别的主力资金因子表现更好。

\subsection{参考文献}
\begin{enumerate}
    \item 魏建榕, 苏良. 2022. 大小单重定标与资金流因子改进. 开源证券.
\end{enumerate}

\section{上行已实现波动率(高频因子-波动跳跃类)}

\subsection{因子定义}
\begin{equation}
    \mathrm{RS}_t^+ = \sum_{i=1}^N r_{t,i}^2 \cdot \mathrm{I}_{r_{t,i} \geq 0}
\end{equation}

\subsection{变量定义}
\begin{itemize}
    \item ① \( r_{t,i} \) 为某股票交易日 \( t \) 内的分钟对数收益率序列，如 5 分钟对数收益率序列；
    \item ② \( \mathrm{I}_{r_{t,i} \geq 0} \) 为示性函数，对数收益率大于等于 0 时取 1，否则取 0；
    \item ③ \( N \) 表示该股票在交易日 \( t \) 中特定频率的分钟级别数据个数，如 5 分钟频率通常有 48 个数据点。
\end{itemize}

\subsection{说明}
上行已实现波动率对应资产价格波动中的上行波动；上行波动率较高的股票短期股价出现大幅拉升，未来更有可能补跌，收益出现反转。

\subsection{参考文献}
\begin{enumerate}
    \item Barndorff-Nielsen, O. E., Kinnebrock, S., \& Shephard, N. (2010). \textit{Measuring downside risk: realised semivariance}. In T. Bollerslev, J. Russell, \& M. Watson (Eds.), Volatility and Time Series Econometrics: Essays in Honor of Robert F. Engle (pp. 117-136). Oxford University Press.
\end{enumerate}

\section{客户接近中心性(图谱网络-网络结构)}

\subsection{因子定义}
\begin{equation}
    c_i = \frac{1}{\frac{1}{N-1} \sum_{j(\neq i)} D_{ij}}
\end{equation}

\subsection{变量定义}
\begin{itemize}
    \item ① \( D_{ij} \) 为供应商公司 \( j \) 到客户公司 \( i \) 的最短路径；
    \item ② \( N \) 为与客户公司 \( i \) 相连的所有供应商公司的总数。
\end{itemize}

\subsection{说明}
接近中心性衡量了供应链网络中能达到某个客户节点的所有供应商节点的平均最短距离的倒数，接近中心性越高，说明该企业与网络中其他企业建立连接的路径长度越短；供应链网络中占据中心位置的企业更能吸引投资者的注意力，市场对其信息反应更快，有较低的股票回报可预测性；中心性较低的企业由于信息传播较慢，有较高的股票回报可预测性。

\subsection{参考文献}
\begin{enumerate}
    \item Bozok, İ., \& Özyıldırım, S. (2022). \textit{Firm centrality and limited attention}. International Review of Economics \& Finance, 78, 483-500.
\end{enumerate}

\section{个人投资者交易比-惊恐因子(高频因子-动量反转类)}

\subsection{因子定义}

1. 计算各股票的日度“惊恐度”：
\begin{equation}
    \text{惊恐度} = \frac{\text{偏离度}}{\text{基准项}} = \frac{\mathrm{abs}(\text{个股收益率} - \text{市场收益率})}{\mathrm{abs}(\text{个股收益率}) + \mathrm{abs}(\text{市场收益率}) + 0.1}
\end{equation}

2. 计算各股票的个人投资者交易比：
\begin{equation}
    \text{个人投资者交易比} = \frac{\text{成交额小于4万元的个人投资者买入和卖出金额均值}}{\text{个股总成交额}}
\end{equation}

3. 对各股票，将每日惊恐度、每日收益率、个人投资者交易比相乘，得到加权调整收益率；

4. 对各股票，每月末计算过去20日加权调整收益率序列的均值和标准差，分别称为“个人投资者交易比-惊恐收益”因子和“个人投资者交易比-惊恐波动”因子，再将二者等权合成为“个人投资者交易比-惊恐”因子。

\subsection{变量定义}
\begin{itemize}
    \item ① 选取中证全指 000985 代表市场水平；
    \item ② “偏离度”衡量了个股收益率对市场收益率的偏离水平；“基准项”衡量了市场总体收益率水平。
\end{itemize}

\subsection{说明}
个人投资者交易比-惊恐因子在原始惊恐因子上新增了个人投资者交易比重，因为“显著效应（投资者会被具有显著收益的股票所吸引，导致它们被错误定价）在小市值股票上更为强烈”，而小市值股票的个人投资者交易占比相对更高，更易受情绪影响而出现反应过度。

\subsection{参考文献}
\begin{enumerate}
    \item 曹春晓, 温晋欣. 极端收益扭曲决策权重和“草木皆兵”因子. 2022, 方正证券.
    \item Cakici, N., \& Zaremba, A. (2022). \textit{Salience theory and the cross-section of stock returns: International and further evidence}. Journal of Financial Economics, 146(2), 689-725.
\end{enumerate}

\section{价量相关性因子(高频因子-量价相关性类)}

\subsection{因子定义}

① 将 \( \mathrm{PV\_corr\_avg} \) 和 \( \mathrm{PV\_corr\_std} \) 分别剔除反转因子（以过去 20 日收益率），再各自横截面标准化，等权线性相加构建综合因子，将得到新因子 \( \mathrm{PV\_corr\_deRet20} \)：
\begin{equation}
    \mathrm{PV\_corr\_deRet20} = \frac{\mathrm{PV\_corr\_avg}_{\perp} - \mathrm{mean}(\mathrm{PV\_corr\_avg}_{\perp})}{\mathrm{std}(\mathrm{PV\_corr\_avg}_{\perp})} 
\end{equation}
\begin{equation}
    +\frac{\mathrm{PV\_corr\_std}_{\perp} - \mathrm{mean}(\mathrm{PV\_corr\_std}_{\perp})}{\mathrm{std}(\mathrm{PV\_corr\_std}_{\perp})}
\end{equation}

② 将 \( \mathrm{PV\_corr\_trend} \) 带来的增量信息叠加到 \( \mathrm{PV\_corr\_deRet20} \) 之上，得到 \( \mathrm{cpv} \)：
\begin{equation}
    \mathrm{cpv} = \frac{\mathrm{PV\_corr\_deRet20} - \mathrm{mean}(\mathrm{PV\_corr\_deRet20})}{\mathrm{std}(\mathrm{PV\_corr\_deRet20})} + \frac{\mathrm{PV\_corr\_trend} - \mathrm{mean}(\mathrm{PV\_corr\_trend})}{\mathrm{std}(\mathrm{PV\_corr\_trend})}
\end{equation}

\subsection{说明}
价量相关性因子综合了价量相关性平均强度、波动性、趋势性 3 方面的信息，对传统反转因子进行了修正，具备良好的选股能力。

\subsection{参考文献}
\begin{enumerate}
    \item 高子剑. 2020. 高频价量相关性，意想不到的选股因子. 东吴证券.
\end{enumerate}

\section{公司特征相似度动量因子(图谱网络-动量溢出)}

\subsection{因子定义}

对于任意两只股票 \(i\) 和 \(j\)，计算它们在特征空间中的欧几里得距离：
\begin{equation}
    D_{i,j} = \sqrt{(\mathrm{Prc}_i - \mathrm{Prc}_j)^2 + (\mathrm{SIZE}_i - \mathrm{SIZE}_j)^2 + (\mathrm{BM}_i - \mathrm{BM}_j)^2 + (\mathrm{OP}_i - \mathrm{OP}_j)^2 + (\mathrm{INV}_i - \mathrm{INV}_j)^2}
\end{equation}

对每只股票，将其距离最近的50只股票（最相似的50只股票）的过去一个月的超额回报进行加权平均即为公司特征相似度动量因子：
\begin{equation}
    \mathrm{SIM} = \frac{\sum_{k \in N(i)} w_k \cdot R_{\mathrm{excess},k}}{\sum_{k \in N(i)} w_k}
\end{equation}

\subsection{变量定义}
\begin{itemize}
    \item ① 计算距离的特征分别为：月末收盘价（\(\mathrm{Prc}\)）、对数市值（\(\mathrm{SIZE}\)）、账面市值比（\(\mathrm{BM}\)）、经营利润/净资产（\(\mathrm{OP}\)）和总资产同比增速（\(\mathrm{INV}\)）；
    \item ② \(w_k\) 为股票 \(k\) 的市值权重；
    \item ③ \(R_{\mathrm{excess},k}\) 为股票 \(k\) 在过去一个月的超额收益率。
\end{itemize}

\subsection{说明}
具有相似基本面特征或是相似因子暴露的公司之间有着相似的预期收益，可以利用相似公司的历史收益来预测自身的未来收益；与其相似的那些股票的收益越高，该股票未来收益也越高。

\subsection{参考文献}
\begin{enumerate}
    \item He, W., Wang, Y., \& Yu, J. (2021). \textit{Similar Stocks}.
\end{enumerate}

\section{反向日内逆转的频率(高频因子-收益分布类)}

\subsection{因子定义}
\begin{equation}
    \mathrm{NR}_{i,t} = \frac{\sum_{d=1}^{T} \mathrm{I}\{\mathrm{RET\_CO}_{i,d}>0\} \times \mathrm{I}\{\mathrm{RET\_OC}_{i,d}<0\}}{T}
\end{equation}

\subsection{变量定义}
\begin{itemize}
    \item ① \( \mathrm{RET\_CO}_{i,d} \) 和 \( \mathrm{RET\_OC}_{i,d} \) 分别为股票 \( i \) 在交易日 \( d \) 的隔夜收益率和日内收益率；
    \item ② \( T \) 为股票 \( i \) 在 \( t \) 月内的实际交易天数。
\end{itemize}

\subsection{说明}
当正的隔夜收益后伴随负的日内收益，即出现了反向日内逆转，通过统计月内出现反向日内逆转的频率，可以衡量隔夜交易者和盘中交易者间“拉锯战”的激烈程度；与未来收益正相关，反向日内逆转出现频率越高，导致日内价格被过度修正拉低，未来越有可能补涨。

\subsection{参考文献}
\begin{enumerate}
    \item Cheema, M. A., Chiah, M., \& Man, Y. (2022). \textit{Overnight returns, daytime reversals, and future stock returns: Is China different?}. Pacific-Basin Finance Journal, 74, 101809.
\end{enumerate}

\section{反转残差非孤立成交不平衡因子(高频因子-成交分布类)}

\subsection{因子定义}

1. 每个交易日，将当日所有成交单分为孤立成交、非孤立成交，其中是否孤立的判断标准为：对于某笔成交单，若在其成交时间点 \( t \) 的邻域 \((t - \sigma, t + \sigma)\) 内不存在任何一笔其他的成交单，则该笔订单即为孤立成交单，反之为非孤立成交单，\(\sigma = 10\mathrm{ms}\)；

2. 对当日所有的非孤立成交单计算如下成交单不平衡指标：
\begin{equation}
    \text{成交单不平衡} = \frac{\text{主买成交单数} - \text{主卖成交单数}}{\text{主买成交单数} + \text{主卖成交单数}}
\end{equation}

3. 每月底回看过去 20 个交易日，计算 20 日非孤立成交单不平衡指标的平均值，做横截面市值中性化处理，得到非孤立成交不平衡因子；

4. 将非孤立成交不平衡因子，对过去 20 个交易日的累计涨跌幅进行正交化处理，得到反转残差非孤立成交不平衡因子。

\subsection{说明}
成交不平衡因子衡量了净买单的强弱程度，若股票在过去一段时间是买方占据主动成交的优势，其未来收益表现会较好；将逐笔成交单划分为“孤立成交”和“非孤立成交”，其中非孤立成交对应的时刻，可能交易火热，或者存在一些较为激进的交易行为，包含更多市场信息，经过划分后计算的因子表现有所提升。

\subsection{参考文献}
\begin{enumerate}
    \item 沈芷琦, 刘富兵, 阮俊烨. 2024. 逐笔买卖差异中的选股信息: 条件成交不平衡因子. 国盛证券.
\end{enumerate}

\section{V 型处置效应（行为金融因子-处置效应）}

\subsection{因子定义}
\begin{equation}
    \mathrm{VNSP}_t = \mathrm{Gain}_t + \sigma |\mathrm{Loss}_t|
\end{equation}
\begin{equation}
    \mathrm{Gain}_t = \sum_{n=0}^{N-1} \omega_{t-n} \cdot \mathrm{gain}_{t-n}
\end{equation}
\begin{equation}
    \mathrm{Loss}_t = \sum_{n=0}^{N-1} \omega_{t-n} \cdot \mathrm{loss}_{t-n}
\end{equation}

\subsection{变量定义}
\begin{itemize}
    \item ① \( \mathrm{Gain}_t \) 为资本盈利量，具体计算公式见“盈利卖出倾向因子”相关因子页；
    \item ② \( \mathrm{Loss}_t \) 为资本亏损量，具体计算公式见“亏损卖出倾向因子”相关因子页；
    \item ③ \( \sigma \) 用来衡量非对称的卖出意愿，可取 \( \sigma = 0.23 \)。
\end{itemize}

\subsection{说明}
V 型处置效应指的是随着持股盈利或亏损幅度的增加，投资者出售股票的意愿也随之增加，且盈利股票卖出意愿比亏损股票卖出意愿更强；V 型处置效应因子由盈利卖出倾向因子和亏损卖出倾向因子组合而成，用于衡量股票的卖出压力。

\subsection{参考文献}
\begin{enumerate}
    \item 陈升锐, 姚紫薇. 2025. 处置效应与 V 型处置效应在量化选股中的应用. 中信建投.
    \item An, L. (2016). \textit{Asset Pricing When Traders Sell Extreme Winners and Losers}. The Review of Financial Studies, 29(3), 823-861.
\end{enumerate}

\section{开盘后大单净买入强度(高频因子-资金流类)}

\subsection{因子定义}

\begin{equation}
    \frac{1}{T} \sum_{n=t}^{t-T+1} \frac{\mathrm{mean}(\text{大买单成交额}_{i,j,n} - \text{大卖单成交额}_{i,j,n})}{\mathrm{std}(\text{大买单成交额}_{i,j,n} - \text{大卖单成交额}_{i,j,n})}
\end{equation}

\subsection{变量定义}
\begin{itemize}
    \item ① 基于逐笔成交数据中的叫买与叫卖单号将逐笔成交数据合成为买卖单数据；
    \item ② 根据买卖单分布特征界定大小单，比如可将多日买卖单成交单数据对数调整后的“均值+1倍标准差”作为大单筛选阈值；
    \item ③ \( i, j, n \) 分别表示第 \( i \) 只股票在第 \( n \) 个交易日内的第 \( j \) 分钟的成交额；
    \item ④ 开盘后的大单净买入强度用的是 9:30-10:00 范围内的数据；
    \item ⑤ 月度选股下，\( T=20 \) 个交易日；周度选股下，\( T=5 \) 个交易日。
\end{itemize}

\subsection{说明}
大单类因子刻画了大资金的交易行为，大资金往往具有信息优势，一般受到大资金关注的股票，未来通常具有更好的表现；大单净买入强度因子同时衡量了大单买入强度和大单买入稳健性两方面的信息。

\subsection{参考文献}
\begin{enumerate}
    \item 冯佳睿, 袁林青. 2021. 大单的精细化处理与大单因子重构. 海通证券.
\end{enumerate}

\section{特质换手波动率(高频因子-波动跳跃类)}

\subsection{因子定义}

截面上对所有股票过去 20 天的日度换手率序列进行主成分分析 PCA，提取第一主成分 \( F_1 \) 为隐式因子溢价；

对每只股票 \( i \)，以过去 20 天的日度换手率序列为 \( y \)，以隐式因子 \( F_1 \) 序列为 \( x \)，进行时序回归：
\begin{equation}
    \mathrm{Turnover}_{i,t} = \alpha_{i,t} + \beta_i^{F_1} \times F_1 + \cdots + \varepsilon_i
\end{equation}

时序回归返回的残差序列的标准差即为隐式因子框架下的特质波动因子：
\begin{equation}
    \mathrm{TurnoverIV} = \mathrm{std}(\varepsilon_i)
\end{equation}

\subsection{变量定义}
\begin{itemize}
    \item ① 在进行截面主成分分析时，若第一主成分的解释力度低于 20\%，则选取解释力最大的前两个主成分进行回归；
    \item ② 还可进一步计算特异度因子：
        \begin{equation*}
            \mathrm{IVR} = 1 - R^2
        \end{equation*}
        上述回归模型拟合优度；
    \item ③ 还可进一步计算特质偏度因子：
        \begin{equation*}
            \mathrm{IS} = \frac{\sum \varepsilon_{it}^3}{\mathrm{IV}^3}
        \end{equation*}
\end{itemize}

\subsection{说明}
计算隐式因子框架下的特质波动率时，直接通过主成分分析对资产收益序列降维来提取共同波动，而无需显式的借助多因子模型估计共同风险。系统性的解决传统多因子模型遗漏变量、估计误差的问题，计算简单；此处提取的是日度换手率序列的共同波动。

\subsection{参考文献}
\begin{enumerate}
    \item 张欣鹏, 杨怡玲, 刘璐. 2022. 隐式框架下的特质因子改进. 国信证券.
\end{enumerate}

\section{筹码换手率(行为金融因子-筹码分布)}

\subsection{因子定义}
\begin{equation}
    \mathrm{PTR} = \frac{\text{当日股票最高价与最低价之间筹码占比}}{\mathrm{turnover}_t}
\end{equation}

\subsection{变量定义}
\begin{itemize}
    \item ① \( \mathrm{turnover}_t \) 为股票 \( t \) 日换手率。
\end{itemize}

\subsection{说明}
该筹码换手率（即筹码穿透率）的计算方式为当天最高价和最低价范围内的筹码占比除以当天换手率，同样用于反映穿透筹码的能力；因子值越大，说明相同换手率穿透筹码占总筹码的比例越高；经测试，该因子与未来收益正相关。

\subsection{参考文献}
\begin{enumerate}
    \item 陈升锐, 姚紫薇. 2025. 筹码分布因子系统构建. 中信建投.
\end{enumerate}

\section{个股收益最高时的领先成交量异动(高频因子-动量反转类)}

\subsection{因子定义}

\begin{equation}
    \mathrm{AMT\_MAX}(N)_i = \frac{\sum_{t=n1}^{nN} \mathrm{AMT\_Percent}_{i,t-1}}{\frac{1}{T}\sum_{t=1}^{T} \mathrm{AMT\_Percent}_{i,t}}
\end{equation}
\begin{equation}
    \mathrm{AMT\_Percent}_{i,t} = \frac{\mathrm{AMT}_{i,t}}{\sum_{i=1}^{I} \mathrm{AMT}_{i,t}}
\end{equation}

\subsection{变量定义}
\begin{itemize}
    \item ① \( \mathrm{AMT}_{i,t} \) 为个股 \( i \) 在 \( t \) 分钟的成交额；
    \item ② \( \mathrm{AMT\_Percent}_{i,t} \) 为个股 \( i \) 在 \( t \) 时刻的成交额占全市场总成交额的比例；
    \item ③ \( n1, n2, \ldots, nN \) 为日内分钟收益率最高的 \( N \) 个分钟，\( N=30 \)；
    \item ④ \( T \) 为日内分钟数 = 240 个 1 分钟。
\end{itemize}

\subsection{说明}
\( \mathrm{AMT\_MAX} \) 衡量了股价大幅上涨时的成交量在全市场成交量中的占比情况；若个股的股价上涨没有伴随市场成交量异动，而是由自身成交量异动造成的（因子值越大，占比越大），未在成交量上与市场形成有效共振关系，则容易受到自身交易过热的反噬，未来收益表现较弱。

\subsection{参考文献}
\begin{enumerate}
    \item 任晓, 麦元勋, 许继宏. 2024. 基于分钟数据的价变共振因子. 招商证券.
\end{enumerate}

\section{流动性溢价因子-等权(高频因子-流动性类)}

\subsection{因子定义}

1. 每月底，以 \( T=21 \) 个交易日为周期，以 \( n=48,120 \) 次为频率（对应 5 分钟线和 2 分钟线），按日结算，进行评估；

2. 分别以每只股票当日成交额的 1\% 或当日成交额的开方，配置每日个股资金 \( M_s \)，并将个股资金按频率等分得到每个交易周期下的资金量 \( M_s/n \)；

3. 获得每个交易周期下收盘时的买单信息（买一到买五），以撮合机制模拟交易股票的操作。当资金不能以五档的挂单量完全成交时，采用线性插值的方式设置虚拟档进行模拟撮合交易，具体算法见参考文献；

4. 按上述算法得到每个周期下，需求方定价下交易的股票数量 \( \mathrm{Vol}_{i,\mathrm{need}} \)，并以当时交易周期的均价得到实际平均交易的股票数量 \( \mathrm{Vol}_{i,\mathrm{actual}} \)，再按当前交易周期收盘价计算股票市值，并每日对上述 2 个市值序列求和得到当日总市值：
\begin{equation}
    \mathrm{Cap}_{\mathrm{need}} = \sum_{i=1}^{n} (\mathrm{Vol}_{i,\mathrm{need}} \times \mathrm{Close}_{i},n)
\end{equation}
\begin{equation}
    \mathrm{Cap}_{\mathrm{actual}} = \sum_{i=1}^{n} (\mathrm{Vol}_{i,\mathrm{actual}} \times \mathrm{Close}_i,n)
\end{equation}

5. 每月底汇总过去 \( T \) 天的数据，计算平均相对差距，即为流动性溢价因子：
\begin{equation}
    f = \frac{\sum_{d=1}^{T} \mathrm{Cap}_{\mathrm{need},d}}{\sum_{d=1}^{T} \mathrm{Cap}_{\mathrm{actual},d}} - 1
\end{equation}

\subsection{说明}
高频流动性溢价因子借助高频挂单数据和模拟撮合机制衡量了需求方定价下交易股票数量与实际平均交易股票数量的相对差距进而刻画流动性风险溢价；理论上，市场的平均成交价高于买单价格（即 \( \mathrm{Vol}_{\mathrm{actual}} < \mathrm{Vol}_{\mathrm{need}} \)），高出部分即为风险溢价，作为持有资产者收益的补偿。

\subsection{参考文献}
\begin{enumerate}
    \item 覃川桃, 郑超. 2019. 高频因子（一）：流动性溢价因子. 长江证券.
\end{enumerate}

\section{趋势清晰度(量价因子改进)}

\subsection{因子定义}
\begin{equation}
    \mathrm{TC}_{i,T} = R^2
\end{equation}
\begin{equation}
    P_{i,t} = \beta_{i,0} + \beta_{i,1} \times T + \varepsilon_t
\end{equation}

\subsection{变量定义}
\begin{itemize}
    \item ① \( P_{i,t} \) 为股票 \( i \) 在第 \( t \) 日的收盘价；
    \item ② \( t \in [T-240, T-20] \) 用 \( T-240 \) 交易日到 \( T-20 \) 交易日的收盘价时序序列进行回归（即不包含最近一个月），时序回归至少需要 200 个交易日的可用数据。
\end{itemize}

\subsection{说明}
“股价时间趋势回归的 \( R^2 \) 值”与“投资者对股价形成的趋势清晰感知度”之间存在正相关关系，可借助回归的 \( R^2 \) 对趋势清晰度进行量化。相对下跌股票中，趋势清晰度越高，股票未来表现越差；相对上涨股票中，趋势清晰度越高，股票未来表现越好。

\subsection{参考文献}
\begin{enumerate}
    \item 郑琳琳, 盛宝丹. 2024. 价格形成路径与趋势清晰度因子. 西南证券.
\end{enumerate}

\section{成交量最高时刻-主卖笔均成交金额(高频因子-成交分布类)}

\subsection{因子定义}

1. 将日内 240 个 1 分钟数据集合标记为 \( T = \{t_1, t_2, \dots, t_{240}\} \)，将分钟成交量最高的 10\% 时刻标记为集合 \( P = \{t_{i1}, t_{i2}, \dots, t_{ik}\} \)；

2. 利用逐笔成交数据，进一步对 \( P \) 中时刻的所有成交记录划分为主动买入成交 \( P\_\mathrm{Buy} \) 和主动卖出成交 \( P\_\mathrm{Sell} \) 两类；

3. 对成交量最高时刻下的主动卖出成交额和成交笔数进行汇总，根据其成交额 \( \mathrm{Amt}_t \) 之和除以成交笔数 \( \mathrm{DealNum}_t \) 之和，构成成交量最高时刻且主卖的笔均成交金额 \( \mathrm{ATD}_{\mathrm{VolumeHigh10\%,Sell}} \)：
\begin{equation}
    \mathrm{ATD}_{\mathrm{VolumeHigh10\%,Sell}} = \frac{\sum_{t \in P\_\mathrm{Sell}} \mathrm{Amt}_t}{\sum_{t \in P\_\mathrm{Sell}} \mathrm{DealNum}_t}
\end{equation}

4. 同理，将全天的成交金额除以全天的成交笔数，得到全天的笔均成交金额 \( \mathrm{ATD}_T \)：
\begin{equation}
    \mathrm{ATD}_T = \frac{\sum_{t \in T} \mathrm{Amt}_t}{\sum_{t \in T} \mathrm{DealNum}_t}
\end{equation}

5. 将成交量最高时刻且主卖的笔均成交金额除以全天的笔均成交金额，即得到去量纲化后的标准化因子 \( \mathrm{SATD}_{\mathrm{VolumeHigh10\%,Sell}} \)：
\begin{equation}
    \mathrm{SATD}_{\mathrm{VolumeHigh10\%,Sell}} = \frac{\mathrm{ATD}_{\mathrm{VolumeHigh10\%,Sell}}}{\mathrm{ATD}_T}
\end{equation}

\subsection{说明}
笔均成交额类因子通过汇总“特殊的、具有较多信息含量”时刻的成交金额情况来刻画主力资金的行为动向；笔均成交金额越大，说明该特殊时间内资金优势越明显，属于主力资金的可能性越大。“成交量最高时刻”从成交热度角度刻画主力资金交易行为，叠加逐笔订单主卖意愿做更细切分能进一步提升因子表现。

\subsection{参考文献}
\begin{enumerate}
    \item 张欣鹏, 张宇. 2025. 日内特殊时刻蕴含的主力资金 Alpha 信息. 国信证券.
\end{enumerate}

\section{基于异常收益的注意力捕捉(行为金融因子-投资者注意力)}

\subsection{因子定义}

对每个行业内分别统计每日行业成分股收益率绝对值高于 7\% 的股票数量占比作为市场关注代理指标：
\begin{equation}
    \mathrm{indus\_abn}_d = \frac{n_{|\mathrm{ret}| \geq 7\%, d}}{N_d}
\end{equation}

将市场关注代理指标与股票日收益进行月度时序回归，代理指标前的系数绝对值 \(|\beta_{i,t}|\) 即为股票对市场关注的敏感程度：
\begin{equation}
    r_{i,d} = \mu_{i,t} + \beta_{i,t} \mathrm{indus\_abn}_d + \rho_{1,i,t} \mathrm{Mkt}_d + \rho_{2,i,t} \mathrm{SMB}_d + \rho_{3,i,t} \mathrm{HML}_d + \rho_{4,i,t} \mathrm{UMD}_d + \varepsilon_{i,d}
\end{equation}

\subsection{变量定义}
\begin{itemize}
    \item ① \( n_{|\mathrm{ret}| \geq 7\%, d} \) 为 \( d \) 日收益率绝对值大于等于 7\% 的行业内股票数量，\( N_d \) 为行业内当日交易股票总数；
    \item ① \(n_{\mathrm{upper},d}\)、\(n_{\mathrm{lower},d}\)、\(N_d\) 分别为 \(d\) 日涨停股票数量、跌停股票数量、当日交易股票总数；
    \item ② \( r_{i,d} \) 为个股日度收益；
    \item ③ \( \mathrm{Mkt}_d, \mathrm{SMB}_d, \mathrm{HML}_d \) 为经典的 Fama-French 三因子；
    \item ④ \( \mathrm{UMD}_d \) 为动量因子收益，即全市场过去 11 个月累积收益前后 30\% 的组合在 \( d \) 日的收益差。
\end{itemize}

\subsection{说明}
不同股票对同一市场关注度事件的反映程度往往是不同的，可借助个股收益对市场关注度事件的敏感程度来衡量股票对该事件的注意力捕捉效应；捕捉能力越强（敏感程度越高），表明个股更能吸引市场注意，从而导致买盘压力增大，股价高估；上述因子以异常收益占比作为市场关注度事件，极端日收益是衡量关注度的一个重要指标，出现较高或较低收益都代表关注度高。

\subsection{参考文献}
\begin{enumerate}
    \item Lin, F., Qiu, Z., \& Zheng, W. (2023). \textit{Cranes among chickens: The general-attention-grabbing effect of daily price limits in China's stock market}. Journal of Banking \& Finance, 150, 106818.
    \item 陈升锐. 2024. 投资者有限关注及注意力捕捉与溢出. 中信建投.
\end{enumerate}

\section{引入主动买卖特征的改进大单交易占比(高频因子-资金流类)}

\subsection{因子定义}

\begin{equation}
    \text{改进大单交易占比} = \frac{\text{大买\&大卖\&主卖} - \text{大买\&非大卖\&主买} - \text{非大买\&大卖\&主卖}}{\text{总成交量}}
\end{equation}

\subsection{变量定义}
\begin{itemize}
    \item ① 大单和非大单的判断标准为：先将同一委买ID（或委卖ID）对应的实际成交量进行加总，再将全部实际成交的委买ID（或委卖ID）成交量进行降序排列，最后将成交量大于前10\%分位点的委买单（或委卖单）记为“大买单”（或“大卖单”），要剔除开盘集合竞价成交记录；
    \item ② 主动买卖单的判断标准为：采用的是一种结合成交价格及委托时间的主动买卖特征划分法，结合当前成交价及委托订单第一次出现时的主动买卖特征进行综合判断，具体细节见报告原文。
\end{itemize}

\subsection{说明}
在订单大小特征基础上再叠加主动买卖的特征，能更准确的追踪知情交易者的交易意图，提升因子表现。

\subsection{参考文献}
\begin{enumerate}
    \item 张欣鹏, 张宇. 2024. 基于主动买卖特征的高频订单因子改进. 国信证券.
\end{enumerate}

\section{基于筹码分布的处置效应(行为金融因子-处置效应)}

\subsection{因子定义}

1. 基于筹码分布等相关数据，计算当日兑现的亏损占筹码总亏损的比例 \( \mathrm{CPLR} \) 和当日兑现的盈利占筹码总盈利的比例 \( \mathrm{CPGR} \)：
\begin{equation}
    \mathrm{CPGR}_{i,t} = \frac{\text{已实现盈利}_{i,t}}{\text{已实现盈利}_{i,t} + \text{账面盈利}_{i,t}}
\end{equation}
\begin{equation}
    \mathrm{CPLR}_{i,t} = \frac{\text{已实现亏损}_{i,t}}{\text{已实现亏损}_{i,t} + \text{账面亏损}_{i,t}}
\end{equation}

2. 计算如下筹码分布下的处置效应因子：
\begin{equation}
    \mathrm{CDE}_{i,t} = \mathrm{CPGR}_{i,t} - \mathrm{CPLR}_{i,t}
\end{equation}

\subsection{变量定义}
\begin{itemize}
    \item ① 已实现盈利、账面盈利、已实现亏损、账面亏损等指标的计算逻辑，见“已实现盈利比率”和“已实现亏损比率”等相关因子日历页。
\end{itemize}

\subsection{说明}
处置效应因子由已实现盈利比率和已实现亏损比率组合而成，用于衡量投资者“卖盈持亏”行为意愿；因子值越大，则兑现盈利并卖出股票的处置效应越明显，因子值越小，则继续持有亏损股票的处置效应越明显。

\subsection{参考文献}
\begin{enumerate}
    \item 陈升锐, 姚紫薇. 2025. 处置效应与 V 型处置效应在量化选股中的应用. 中信建投.
    \item Odean, T. (1998). \textit{Are investors reluctant to realize their losses?}. Journal of Finance, 53(5), 1775–1798.
\end{enumerate}

\section{修正模糊价差(高频因子-成交分布类)}

\subsection{因子定义}

1. 计算各股票每日的“日模糊价差”因子：
   \begin{equation}
       \text{日模糊价差} = \text{日模糊金额比} - \text{日模糊数量比}
   \end{equation}

2. 每日计算各股票过去 10 日的“日模糊价差”的标准差；

3. 每日将横截面上“日模糊价差”为负的“日模糊价差”求和得 \( s_1 \)：

4. 将“日模糊价差”为负的“日模糊价差”除以其过去 10 日标准差，得到“修正日模糊价差”，“日模糊价差”为正部分不变：

5. 每日将横截面上“修正日模糊价差”为负的“修正日模糊价差”求和得 \( s_2 \)；

6. 将“修正日模糊价差”为负的部分加以调整，除以 \( s_2 \)，乘以 \( s_1 \)，使其负的部分数量级与原来保持一致，避免市值行业中心化时受影响。

\subsection{变量定义}
\begin{itemize}
    \item ① “日模糊金额比”和“日模糊数量比”因子计算逻辑见相应日历页；
    \item ② 每月末可对最近 20 天的上述因子求均值和标准差并将两者等权合并，即得到月度“修正模糊价差”因子。
\end{itemize}

\subsection{说明}
模糊价差因子衡量了投资者因厌恶波动模糊而急于出售股票时所付出的流动性成本，与未来收益负相关；上述计算逻辑包含了对该因子的修正细节，识别出了多头组中基本面恶化而不会发生补涨的股票（利用标准差对负的日模糊价差进行调整）。

\subsection{参考文献}
\begin{enumerate}
    \item 曹春晓. 2022. 波动率的波动率与投资者模糊性厌恶. 方正证券.
    \item Kostopoulos, D., Meyer, S., \& Uhr, C. (2021). \textit{Ambiguity about volatility and investor behavior}. Journal of Financial Economics.
\end{enumerate}

\section{流动性溢价因子-加权(高频因子-流动性类)}

\subsection{因子定义}

1. 每月底，以 \( T=21 \) 个交易日为周期，以 \( n=48,120 \) 次为频率（对应 5 分钟线和 2 分钟线），按日结算，进行评估；

2. 分别以每只股票当日成交额的 1\% 或当日成交额的开方，配置每日个股资金 \( M_s \)，并将个股资金按频率等分得到每个交易周期下的资金量 \( M_s/n \)；

3. 获得每个交易周期下收盘时的买单信息（买一到买五），以撮合机制模拟交易股票的操作。当资金不能以五档的挂单量完全成交时，采用线性插值的方式设置虚拟档进行模拟撮合交易，具体算法见参考文献；

4. 按上述算法得到每个周期下，需求方定价下交易的股票数量 \( \mathrm{Vol}_{i,\mathrm{need}} \)，并以当时交易周期的均价得到实际平均交易的股票数量 \( \mathrm{Vol}_{i,\mathrm{actual}} \)，再按当前交易周期收盘价计算股票市值，并每日对上述 2 个市值序列求和得到当日总市值：
\begin{equation}
    \mathrm{Cap}_{\mathrm{need}} = \sum_{i=1}^{n} (\mathrm{Vol}_{i,\mathrm{need}} \times \mathrm{Close}_i,n)
\end{equation}
\begin{equation}
    \mathrm{Cap}_{\mathrm{actual}} = \sum_{i=1}^{n} (\mathrm{Vol}_{i,\mathrm{actual}} \times \mathrm{Close}_i,n)
\end{equation}

5. 每月底汇总过去 \( T \) 天的数据，计算指数加权或波动率加权 \( \mathrm{weight}^t \propto |\mathrm{return}^t| \) 下的平均相对差距，半衰期为 21 天，即为流动性溢价因子：
\begin{equation}
    f = \frac{\sum_{t=1}^{T} \mathrm{Cap}_{\mathrm{need}}^t \times \mathrm{weight}^t}{\sum_{t=1}^{T} \mathrm{Cap}_{\mathrm{actual}}^t \times \mathrm{weight}^t} - 1
\end{equation}

\subsection{说明}
高频流动性溢价因子借助高频挂单数据和模拟撮合机制衡量了需求方定价下交易股票数量与实际平均交易股票数量的相对差距进而刻画流动性风险溢价；理论上，市场的平均成交价高于买单价格（即 \( \mathrm{Vol}_{\mathrm{actual}} < \mathrm{Vol}_{\mathrm{need}} \)），高出部分即为风险溢价，作为持有资产者收益的补偿；时间加权（时间越近权重越大）后的因子有更强的空头区分能力，波动加权（收益变动越大权重越大）后的因子表现更优。

\subsection{参考文献}
\begin{enumerate}
    \item 黎川桃, 郑超. 2019. 高频因子（一）：流动性溢价因子. 长江证券.
\end{enumerate}

\section{平均异常日收益率(行为金融因子-投资者注意力)}

\subsection{因子定义}
\begin{equation}
    \mathrm{ABNRETAVG} = \frac{1}{m} \sum_{t}^{t-m+1} (r_{i,t} - \bar{r}_t)^2
\end{equation}

\subsection{变量定义}
\begin{itemize}
    \item ① \( r_{i,t} \) 为股票 \( i \) 在 \( t \) 交易日的收益率；
    \item ② \( \bar{r}_t \) 为交易日 \( t \) 的全市场收益率；
    \item ③ \( m \) 为交易日窗口，一般取为月度。
\end{itemize}

\subsection{说明}
平均异常日收益率为日频收益与市场中位数差值的均方，属于极端日收益类因子；极端收益（过高收益或过低收益）的股票相较其他股票更能吸引投资者的注意，因子值越高，投资者对该股票的关注度越高；而有限注意力导致投资者更倾向于交易引起他们关注的股票，投资者关注度越高的股票通常意味着越多的投资购买；但 A 股市场的做多与做空不对称，使得投资者非理性买入高于非理性卖出，进而导致关注度高的股票由于投资者净买入而存在溢价，未来也更可能出现反转。

\subsection{参考文献}
\begin{enumerate}
    \item Cosemans, M., \& Frehen, R. (2021). \textit{Salience theory and stock prices: Empirical evidence}. Journal of Financial Economics, 140(2), 460-483.
    \item 陈升锐. 2024. 投资者有限关注及注意力捕捉与溢出. 中信建投.
\end{enumerate}

\section{基于物理时间交易量加权的知情交易概率(高频因子-资金流类)}

\subsection{因子定义}
\begin{equation}
    \mathrm{VWPIN} = \frac{\alpha\mu}{\alpha\mu + 2\varepsilon} \approx \sum_{i=1}^{n} w_i \frac{|S_i - B_i|}{S_i + B_i}
\end{equation}
\begin{equation}
    w_i = \frac{\mathrm{Volume}_i}{\sum_{k=1}^{n} \mathrm{Volume}_k}
\end{equation}

\subsection{变量定义}
\begin{itemize}
    \item ① \( S_i \) 和 \( B_i \) 分别为第 \( i \) 个交易时段内的卖单数和买单数；
    \item ② \( w_i \) 为第 \( i \) 个交易时段内成交量占全天成交量的比例，\( \mathrm{Volume}_i \) 为第 \( i \) 个交易时段内的成交量；
    \item ③ \( n \) 为交易日 \( t \) 内的交易时段的数量，比如 5 分钟的频率，则共有 \( n=48 \) 个交易时段；
    \item ④ 若想构建日度以上频率的因子，只需在周期内取均值即可。
\end{itemize}

\subsection{说明}
VWPIN 是一种同时包含了订单数量不平衡程度（如拆单操作）和交易量不平衡程度的知情交易概率的计算方式，可以动态考察因信息不对称导致的逆向选择风险在日内的变化情况；知情交易概率指在一段时间内，拥有信息优势的知情交易者提交的订单数量占总委托单数量的比例，用以刻画该段时间内的信息不对称程度，通常与未来收益正相关，取值越大，信息不对称程度越大，投资者要求的风险回报越高。

\subsection{参考文献}
\begin{enumerate}
    \item 李平等. 2020. 知情交易概率与风险定价：基于不同 PIN 测度方法的比较研究. 管理科学学报, 23(1), 33-46.
\end{enumerate}

\section{日内收益波动比(高频因子-波动跳跃类)}

\subsection{因子定义}

① 对于股票 \( i \)，计算 \( n \) 日内第 \( t \) 分钟的“更优波动率”，即对第 \( t-4, t-3, t-2, t-1, t \) 分钟的开盘价、最高价、最低价和收盘价，共 20 个价格数据求标准差，再除以这 20 个价格数据的均值，最后对该比值取平方：

② 计算 \( t \) 分钟的收益波动比，即 \( t \) 分钟收益率 / \( t \) 分钟“更优波动率”：


③ 计算每日收益波动比序列与“更优波动率”序列间的协方差，即为日度“灾后重建”因子：

\subsection{变量定义}
\begin{itemize}
    \item ① 计算因子需剔除开盘和收盘数据，仅考虑日内分钟频数据；
    \item ② 每月末可对最近 20 天的上述因子求均值和标准差并将两者等权合并，即得到月度“灾后重建”因子。
\end{itemize}

\subsection{说明}
“灾后重建”因子衡量了投资者对波动变化反应不足的程度，与未来收益呈反比关系；当股价波动加剧时，人们期初往往对此反应不足，导致收益波动比下降，引起那些偏爱收益波动比高的投资者抛售股票，造成股价下跌，但未来很有可能会补涨。

\subsection{参考文献}
\begin{enumerate}
    \item 曹春晓. 2022. 个股波动率的变动及“勇攀高峰”因子构建. 方正证券.
    \item Moreira, A., \& Muir, T. (2017). \textit{Volatility-managed portfolios}. The Journal of Finance, 72(4), 1611-1644.
\end{enumerate}

\section{修正的相似业务收益联动因子(图谱网络-动量溢出)}

\subsection{因子定义}

\begin{equation}
    \mathrm{Linkage}_{i,t} = \frac{\sum_{j=1}^{N} \mathrm{SIM}_{i,j,t} \times \mathrm{VOL}_{j,t} \times \mathrm{EV}_{j,t} \times \mathrm{Ret}_{j,t}}{\sum_{j=1}^{N} (\mathrm{SIM}_{i,j,t} \times \mathrm{VOL}_{j,t} \times \mathrm{EV}_{j,t})} - \mathrm{Ret}_{i,t}
\end{equation}
\begin{equation}
    \mathrm{SIM}_{i,j,t} = \frac{P_{i,t} \cdot P_{j,t}}{\|P_{i,t}\| \cdot \|P_{j,t}\|}
\end{equation}

\subsection{变量定义}
\begin{itemize}
    \item ① \( P_{i,t} \) 为取值只有 0 和 1 的业务关键词向量，通过自然语言技术处理 A 股上市公司年度财务报告附注中关于公司基本情况的描述部分内容得到；
    \item ② \( \mathrm{SIM}_{i,j,t} \) 为公司 \( i \) 和公司 \( j \) 在 \( t \) 时期的业务相似度；
    \item ③ \( \mathrm{Ret}_{j,t} \) 为公司 \( j \) 在 \( t \) 月的收益率；
    \item ④ \( \mathrm{EV}_{j,t} \) 为公司 \( j \) 在 \( t \) 月的总市值；
    \item ⑤ \( \mathrm{VOL}_{j,t} \) 为公司 \( j \) 在 \( t \) 月的成交额。
\end{itemize}

\subsection{说明}
相似业务收益联动因子度量了 \( t \) 期与 \( i \) 公司相似的公司在当月超过 \( i \) 公司的收益率，该值越大，说明 \( i \) 公司越有可能在下一期出现一定的补涨行情；修正后的因子剔除了一些小市值的股票短期受游资炒作拉升的影响。

\subsection{参考文献}
\begin{enumerate}
    \item 叶尔乐. 2023. 财报文本中公司竞争信息刻画与 Alpha 构建. 民生证券.
\end{enumerate}

\section{超大单涨跌幅(高频因子-动量反转类)}

\subsection{因子定义}

计算每笔成交数据对应的对数价格变动：
\begin{equation}
    \Delta \ln P_i = \text{当前成交价格的对数}\ln P_i - \text{上一笔成交价格的对数}\ln P_{i-1}
\end{equation}

对于每笔成交数据中涉及到的买卖双方，将时间上靠后的订单标记为主动订单；计算所有主动订单为超大单的对数价格变动之和即为超大单涨跌幅：
\begin{equation}
    \mathrm{SuperBigRet} = \sum_{i \in S} \Delta \ln P_i
\end{equation}

\subsection{变量定义}
\begin{itemize}
    \item ① 按照委买单号和委卖单号对 level2 逐笔成交数据中的成交金额（已成交部分）进行汇总求和得到当天股票 \( i \) 所有买单和卖单的成交金额；
    \item ② 对股票 \( i \)，将订单成交金额占当天总成交额比例超过 1\% 的订单定义为超大单（阈值在 0.5\%~2\% 区间较好）。
\end{itemize}

\subsection{说明}
超大单涨跌幅因子刻画了是由超大单主导的已成交订单的涨跌幅，与未来收益负相关；大单是为了找出有信息优势的订单，超大单是为了刻画潜在的过度投机和流动性冲击的负向冲击效应。

\subsection{参考文献}
\begin{enumerate}
    \item 朱剑涛, 王星星. 超大单冲击对大单因子的影响, 2022, 东方证券.
\end{enumerate}

\section{绝对收益与成交量相关性(高频因子-量价相关性类)}

\subsection{因子定义}
\begin{equation}
    \mathrm{CORA} = \mathrm{corr}(| \mathrm{Ret}_t |, \mathrm{Amount}_t), \quad t \in [1, 240]
\end{equation}

\subsection{变量定义}
\begin{itemize}
    \item ① \( | \mathrm{Ret}_t | \) 为个股日内第 \( t \) 分钟的对数收益率的绝对值；
    \item ② \( \mathrm{Amount}_t \) 为个股日内第 \( t \) 分钟的成交额；
    \item ③ 月度因子可使用月末 5-15 个交易日因子的平均值来合成。
\end{itemize}

\subsection{说明}
CORA 代表单位成交额对股价推动的关系（双方通常正相关），与未来收益负相关；因子值过高，表明股票成交额放量时伴随着价格剧烈波动，该类股票投机属性强，资金对价格的扰动多，短期交易过热，该类因子月度上往往具有负向 Alpha。

\subsection{参考文献}
\begin{enumerate}
    \item 朱定豪, 严佳炜. 2020. 量价关系的高频乐章. 方正证券.
\end{enumerate}

\section{SLSV 模型(行为金融因子-羊群效应)}

\subsection{因子定义}

\begin{equation}
    \mathrm{SLSV}_{i,t} = 
    \begin{cases} 
        \mathrm{LSV}_{i,t}, & p_{i,t} > p_t \\ 
        -\mathrm{LSV}_{i,t}, & p_{i,t} < p_t 
    \end{cases}
\end{equation}

\begin{equation}
    \mathrm{LSV}_{i,t} = |p_{i,t} - p_t| - E|p_{i,t} - p_t|
\end{equation}

\begin{equation}
    E|p_{i,t} - p_t| = \mathrm{AF}(i,t) = \sum_{k=1}^{N_{i,t}} \binom{N_{i,t}}{k} p_t^k (1 - p_t)^{N_{i,t} - k} \left| \frac{k}{N_{i,t}} - p_t \right|
\end{equation}

\begin{equation}
    p_{i,t} = \frac{B(i,t)}{B(i,t) + S(i,t)}, \quad p_t = \frac{\sum_{i=1}^{N} p_{i,t}}{N}, \quad N_{i,t} = B(i,t) + S(i,t)
\end{equation}

\subsection{变量定义}
\begin{itemize}
    \item ① \( B(i,t), S(i,t) \) 分别为区间 \( t \) 内买入股票和卖出股票的基金数量；
    \item ② \( N_{i,t} \) 为持仓发生变动的基金数量；
    \item ③ \( p_{i,t} \) 为股票 \( i \) 在区间 \( t \) 内买入股票基金占持仓变动基金的比例；
    \item ④ 此处使用的是基金年报和半年报持仓数据，假设有 300 只基金加仓股票 \( i \)，则 \( B(i,t) = 300 \)；
    \item ⑤ \( N \) 为全市场股票总数。
\end{itemize}

\subsection{说明}
LSV 模型从基金持股变动数据出发判断一只股票是否被集中买入或集中卖出，进而捕捉基金“抱团”现象；而 SLSV 在 LSV 基础上，考虑了投资者的买卖方向，当 SLSV 显著为正时，对应买入从众效应，当 SLSV 显著为负时，对应卖出从众效应，SLSV 越趋于 0 时，从众效应越弱。

\subsection{参考文献}
\begin{enumerate}
    \item Wermers, R. (1999). \textit{Mutual Fund Herding and the Impact on Stock Prices}. The Journal of Finance, 54, 581-622.
\end{enumerate}

\section{分钟成交额方差(高频因子-成交分布类)}

\subsection{因子定义}
\begin{equation}
    \mathrm{VMA} = \frac{1}{N} \sum_{t=1}^{N} (\mathrm{Amount}_t - \mu)^2
\end{equation}
\begin{equation}
    \mu = \frac{1}{N} \sum_{t=1}^{N} \mathrm{Amount}_t
\end{equation}

\subsection{变量定义}
\begin{itemize}
    \item ① \( \mathrm{Amount}_t \) 为交易日内的分钟成交额；
    \item ② \( N \) 为日内不同分钟频率下的时间段数量。
\end{itemize}

\subsection{说明}
分钟成交额方差刻画了日内分钟成交额的高阶矩特征；此外，还可对日内分钟序列计算偏度和峰度；分钟成交额高阶矩因子在中心化后的表现一般。

\subsection{参考文献}
\begin{enumerate}
    \item 尹治技. 2020. 高频视角下成交额蕴藏的 Alpha. 华安证券.
\end{enumerate}

\section{买卖委托价格比率(高频因子-流动性类)}

\subsection{因子定义}
\begin{equation}
    \mathrm{PIR}_t = \frac{P_t^{\mathrm{WB}} - P_t^{\mathrm{WA}}}{P_t^{\mathrm{WB}} + P_t^{\mathrm{WA}}}
\end{equation}
\begin{equation}
    P_t^{\mathrm{WB(A)}} = \frac{\sum_{i=1}^{5} w_i \times P_{i,t}^{\mathrm{B(A)}}}{\sum_{i=1}^{5} w_i}
\end{equation}
\begin{equation}
    w_i = 1 - (i - 1)/5, \quad i = 1, 2, 3, 4, 5
\end{equation}

\subsection{变量定义}
\begin{itemize}
    \item ① \( P_t^{\mathrm{WB(A)}} \) 为衰减加权委托价格；
    \item ② \( P_{i,t}^{\mathrm{B}} \) 和 \( P_{i,t}^{\mathrm{A}} \) 分别是 \( t \) 时刻第 \( i \) 档的买一和卖一的委托价格；
    \item ③ \( w_i = 1 - (i - 1)/5, \ i = 1, 2, 3, 4, 5 \)。
\end{itemize}

\subsection{说明}
PIR 为买卖委托价差与其和的比值，从价格的角度描述了买盘的压力。

\subsection{参考文献}
\begin{enumerate}
    \item 陈升锐. 2023. 高频选股因子分类体系. 中信建投证券.
\end{enumerate}


\section{净支撑成交量因子(量价因子改进)}

\subsection{因子定义}

1. 每个交易日，计算所有分钟收盘价的均值，筛选收盘价小于该平均值的分钟，将它们的成交量加总，得到支撑成交量；

2. 将收盘价大于上述平均值的分钟成交量加总，得到阻力成交量；

3. 每个交易日，计算如下因子：
\begin{equation}
    \text{净支撑成交量} = \frac{\text{支撑成交量} - \text{阻力成交量}}{\text{当日流通股本}}
\end{equation}

4. 每月月底回看过去 20 个交易日，计算上述因子的平均值，并做横截面市值、中信一级行业中性化处理。

\subsection{说明}
净支撑成交量因子值较大，则表明支撑价格不破位的力量较强，后市上涨概率较高；若因子值较小，则表明阻碍价格上涨的力量较强，后市相对偏空。

\subsection{参考文献}
\begin{enumerate}
    \item 沈芷琦, 刘富兵. 2024. 基于趋势资金日内交易行为的选股因子. 国盛证券.
\end{enumerate}

\section{博弈因子(高频因子-资金流类)}

\subsection{因子定义}
\begin{equation}
    \mathrm{Stren} = \frac{\sum \mathrm{weight}_t \times \mathrm{Vol\_Buy}_t}{\sum \mathrm{weight}_t \times \mathrm{Vol\_Sell}_t}
\end{equation}
\begin{equation}
    \mathrm{Vol\_Buy}_t = \sum \mathrm{vol\_buy}_i
\end{equation}
\begin{equation}
    \mathrm{Vol\_Sell}_t = \sum \mathrm{vol\_sell}_i
\end{equation}

\subsection{变量定义}
\begin{itemize}
    \item ① \( \mathrm{vol\_buy}_i \) 为每日逐笔成交数据中，当前成交价大于上一笔买一价，交易以买方主导的成交量，每日加总得到主买量，用以衡量多空博弈中的多头力量；
    \item ② \( \mathrm{vol\_sell}_i \) 为每日逐笔成交数据中，当前成交价小于上一笔卖一价，交易以卖方主导的成交量，每日加总得到主卖量，用以衡量多空博弈中的空头力量；
    \item ③ \( \mathrm{weight}_t \) 为加权权重，可以等权或 21 天、42 天半衰期等指数加权。
\end{itemize}

\subsection{说明}
博弈因子衡量了多空双方博弈的相对力量，与未来收益负相关；通常买卖双方的短期博弈均存在过度反应的现象，当个股过去一段时间主买量大于主卖量时，买方力量强于卖方力量，多头筹码量增价，个股价格容易被高估；反之当个股过去一段时间主卖量大于主买量时，卖方力量强于买方力量，个股价格容易被低估。

\subsection{参考文献}
\begin{enumerate}
    \item 覃川桃, 郑超. 2019. 高频因子(五)：高频因子与交易行为. 长江证券.
\end{enumerate}

\section{上下行波动率不对称性(高频因子-波动跳跃类)}

\subsection{因子定义}
\begin{equation}
    \mathrm{RSJ}_d = \frac{\mathrm{RV}_{\mathrm{up}_d} - \mathrm{RV}_{\mathrm{down}_d}}{\mathrm{RV}_d}
\end{equation}
其中，
\begin{equation}
    \mathrm{RV}_d = \sum_{i=1}^{N_d} r_{d,i}^2
\end{equation}
\begin{equation}
    \mathrm{RV}_{\mathrm{up},d} = \sum_{i=1}^{N_d} r_{d,i}^2 \cdot \mathrm{I}_{r_{d,i}>0}
\end{equation}
\begin{equation}
    \mathrm{RV}_{\mathrm{down},d} = \sum_{i=1}^{N_d} r_{d,i}^2 \cdot \mathrm{I}_{r_{d,i}<0}
\end{equation}

\subsection{变量定义}
\begin{itemize}
    \item ① \( r_{d,i} \) 为某股票交易日 \( d \) 内的分钟对数收益率序列，如 5 分钟对数收益率序列；
    \item ② \( N_d \) 表示该股票在交易日 \( d \) 中特定频率的分钟级别数据个数，如 5 分钟频率通常有 48 个数据点。
\end{itemize}

\subsection{说明}
RSJ 衡量了上下行波动的不对称性；RSJ 通常有负风险溢价，短期日内的情绪不稳定的大幅上涨往往跟着未来的补跌（RSJ 越大），短期日内情绪不稳定的大幅下跌也往往跟着未来的补涨。

\subsection{参考文献}
\begin{enumerate}
    \item Barndorff-Nielsen, O. E., Kinnebrock, S., \& Shephard, N. (2010). \textit{Measuring downside risk: realised semivariance}. In T. Bollerslev, J. Russell, \& M. Watson (Eds.), Volatility and Time Series Econometrics: Essays in Honor of Robert F. Engle (pp. 117–136). Oxford University Press.
\end{enumerate}

\section{客户 Katz-Bonacich 中心性(图谱网络-网络结构)}

\subsection{因子定义}
\begin{equation}
    k_i = \lambda \sum_{j \neq i} x_{ij} k_j + \beta
\end{equation}

\subsection{变量定义}
\begin{itemize}
    \item ① \(\lambda\) 为介于 0-1 之间的衰减因子，通常小于邻接矩阵最大特征值的倒数，用于减少随着路径长度增加的贡献；
    \item ② \(x_{ij}\) 为供应链邻接矩阵的元素，若公司 \(i\) 和公司 \(j\) 相连，取值可以为 1，或是销售占比等；
    \item ③ \(\beta\) 为基础中心度，通常为一个小的正数；
    \item ④ \(k_i\) 为节点公司 \(i\) 的 Katz 中心度，初始取值可以为 \(\beta\)，然后不断迭代，直到中心度得分收敛。
\end{itemize}

\subsection{说明}
Katz 中心性在衡量供应链节点重要性时考虑了节点间的所有路径长度，可以捕捉节点间较长路径的影响；供应链网络中占据中心位置的企业更能吸引投资者的注意力，市场对其信息反应更快，有较低的股票回报可预测性；中心性较低的企业由于信息传播较慢，有较高的股票回报可预测性。

\subsection{参考文献}
\begin{enumerate}
    \item Bozok, İ., \& Özyıldırım, S. (2022). \textit{Firm centrality and limited attention}. International Review of Economics \& Finance, 78, 483-500.
\end{enumerate}

\section{注意力衰减-惊恐因子(高频因子-动量反转类)}

\subsection{因子定义}

1. 计算各股票的日度“惊恐度”：
   \begin{equation}
       \text{惊恐度} = \frac{\text{偏离度}}{\text{基准项}} = \frac{\mathrm{abs}(\text{个股收益率} - \text{市场收益率})}{\mathrm{abs}(\text{个股收益率}) + \mathrm{abs}(\text{市场收益率}) + 0.1}
   \end{equation}

2. 计算 \( t \) 日衰减后的“惊恐度”，并将负值替换成空值：
   \begin{equation}
       \text{衰减惊恐度}_t = \text{惊恐度}_t - \frac{\text{惊恐度}_{t-1} + \text{惊恐度}_{t-2}}{2}
   \end{equation}

3. 对各股票，将每日衰减惊恐度与每日收益率相乘，得到加权调整收益率；

4. 对各股票，每月末计算过去 20 日加权调整收益率序列的均值和标准差，分别称为“注意力衰减-惊恐收益”因子和“注意力衰减-惊恐波动”因子，再将二者等权合成为“注意力衰减-惊恐因子”。

\subsection{变量定义}
\begin{itemize}
    \item ① 选取中证全指 000985 代表市场水平；
    \item ② “偏离度”衡量了个股收益率对市场收益率的偏离水平；“基准项”衡量了市场总体收益水平。
\end{itemize}

\subsection{说明}
注意力衰减-惊恐因子对原始惊恐度做了时间衰减，考虑了投资者对短暂连续显著收益的注意力会随时间减弱，过度反应也会逐渐缓解。

\subsection{参考文献}
\begin{enumerate}
    \item 曹春晓. 显著效应、极端收益扭曲决策权重和“草木皆兵”因子, 2022, 方正证券.
    \item Cosemans, M., \& Frehen, R. (2021). \textit{Salience theory and stock prices: Empirical evidence}. Journal of Financial Economics, 140(2), 460-483.
\end{enumerate}

\section{大成交量价量相关性(高频因子-量价相关性类)}

\subsection{因子定义}
\begin{equation}
    \mathrm{corr\_VP}_{D,i} = \mathrm{corr} \left( v_{t,D,i}, \; p_{t,D,i} \right), \quad t \in \{ \text{大成交量对应的分钟区间} \}
\end{equation}

\subsection{变量定义}
\begin{itemize}
    \item ① 将个股在每个交易日的分钟成交量时间序列按照成交量大小排序，将分钟成交量排名前 \( \frac{1}{3} \) 的成交量定为“大成交量”；
    \item ② \( t=1,2,3,\ldots,T \) 为交易日内的相应分钟频率的第 \( t \) 个时间区间，此处 \( T \) 为大成交量对应的那些区间；
    \item ③ \( v_{t,D,i} \) 和 \( p_{t,D,i} \) 分别为股票 \( i \) 在交易日 \( D \) 内分钟频率下的成交量和价格。
\end{itemize}

\subsection{说明}
在传统价量相关性因子计算逻辑的基础上加入了成交量大小的信息，大单成交更能代表主力行为，蕴含的信息更多。

\subsection{参考文献}
\begin{enumerate}
    \item 文巧钧, 安宁宁, 罗军. 2021. 高频量化数据因子化方法. 广发证券.
\end{enumerate}

\section{新闻网络领先收益(图谱网络-动量溢出)}

\subsection{因子定义}
以标题中股票为 lead，以同时出现在新闻正文中的其他股票为 follower，确定领边 \( l_{ij,T} \)，并构建股票共现图的邻接矩阵 \( W_T \)：
\begin{equation}
    l_{ij,T} \stackrel{\mathrm{def}}{=} \bigcup_{m_d \in D_T} \{(i,j)_{m_d} \mid j \text{ in } m_d \text{ title, } i \text{ in } m_d \text{ headline, } i \neq j\}
\end{equation}
\begin{equation}
    W_T \stackrel{\mathrm{def}}{=} [\omega_{ij,T}]_{n \times n} = \left[ \frac{l_{ij,T}}{\sum_{j=1}^n l_{ij,T}} \right]_{n \times n}
\end{equation}

将邻接矩阵进行拆解 \( \omega_{ij,T}^{\omega} \)（同属一个行业的股票的邻接矩阵）和 \( \omega_{ij,T}^{\varepsilon} \)（不属于同一行业股票的邻接矩阵）：
\begin{equation}
    \omega_{ij,T}^{\omega} \stackrel{\mathrm{def}}{=} \mathbb{I}\{i,j \text{ from same sector}\} \cdot \omega_{ij,T}
\end{equation}
\begin{equation}
    \omega_{ij,T}^{\varepsilon} \stackrel{\mathrm{def}}{=} \mathbb{I}\{i,j \text{ from diff sector}\} \cdot \omega_{ij,T}
\end{equation}

将 lead 股票的收益进行加权即得到属于 follower 股票的领先收益 lead return，可以根据 lead 股票的收益的正负，单独计算正（负）lead return：
\begin{equation}
    \mathcal{L}R_{i,t}^{+(-)} (\omega_{ij,t-l:t}) \stackrel{\mathrm{def}}{=} \sum_{j=1}^n \omega_{ij,t-l:t} \cdot |r_{j,t}^{+(-)}|
\end{equation}

由于非同一行业的 lead 股票收益有明显的反转效应，正 lead 股票收益的有明显的动量效应，也可以构建如下复合因子：
\begin{equation}
    \mathcal{L}R_{\mathrm{agg}} (\omega_{ij,t-365:t}) = \mathcal{L}R^{+} (\omega_{ij,t-365:t}^{\omega}) + \mathcal{L}R^{-} (\omega_{ij,t-365:t}^{\omega}) + \mathcal{L}R^{-} (\omega_{ij,t-365:t}^{\varepsilon}) - \mathcal{L}R^{+} (\omega_{ij,t-365:t}^{\varepsilon})
\end{equation}

\subsection{说明}
对于新闻中的 lead 股票和 follower 股票，无论业务关系如何，他们之间的收益都具有较强的共振效应，lead 股票的收益率对 follower 股票的收益率有预测作用。

\subsection{参考文献}
\begin{enumerate}
    \item Hu, J., \& Härdle, W. (2021). \textit{Networks of news and cross-sectional returns}. IRTG 1792 Discussion Papers 2021-023, Humboldt University of Berlin, International Research Training Group 1792 "High Dimensional Nonstationary Time Series".
\end{enumerate}

\section{正向日内逆转的频率(高频因子-收益分布类)}

\subsection{因子定义}
\begin{equation}
    \mathrm{PR}_{i,t} = \frac{\sum_{d=1}^{T} \mathrm{I}_{\{\mathrm{RET\_CO}_{i,d}<0\}} \times \mathrm{I}_{\{\mathrm{RET\_OC}_{i,d}>0\}}}{T}
\end{equation}

\subsection{变量定义}
\begin{itemize}
    \item ① \( \mathrm{RET\_CO}_{i,d} \) 和 \( \mathrm{RET\_OC}_{i,d} \) 分别为股票 \( i \) 在交易日 \( d \) 的隔夜收益率和日内收益率；
    \item ② \( T \) 为股票 \( i \) 在 \( t \) 月内的实际交易天数。
\end{itemize}

\subsection{说明}
当负的隔夜收益后伴随正的日内收益，即出现了正向日内逆转，通过统计月内出现正向日内逆转的频率，也可以衡量隔夜交易者和盘中交易者间“拉锯战”的强度；与未来收益负相关，正向日内逆转出现频率越高，导致日内价格被过度修正拉高，未来越有可能补跌。

\subsection{参考文献}
\begin{enumerate}
    \item Cheema, M. A., Chiah, M., \& Man, Y. (2022). \textit{Overnight returns, daytime reversals, and future stock returns: Is China different?}. Pacific-Basin Finance Journal, 74, 101809.
\end{enumerate}

\section{成交量占比(高频因子-成交分布类)}

\subsection{因子定义}
开盘集合竞价成交量占比因子 \( \mathrm{OCVP} \)：
\begin{equation}
    \mathrm{OCVP}_t = \frac{1}{d} \sum_{i=1}^{d} \omega_{t-i}\left( \frac{\text{开盘集合竞价阶段成交量 } \mathrm{VOL}_{\mathrm{call}}}{\text{日内个股总成交量 } \mathrm{VOL}_{\mathrm{total}}} \right)_{t-i}
\end{equation}

收盘集合竞价成交量占比因子 \( \mathrm{BCVP} \)：
\begin{equation}
    \mathrm{BCVP}_t = \frac{1}{d} \sum_{i=1}^{d} \left( \frac{\text{收盘前5分钟内个股成交量 } \mathrm{VOL}_{\mathrm{call}}}{\text{日内个股总成交量 } \mathrm{VOL}_{\mathrm{total}}} \right)_{t-i}
\end{equation}

成交量占比复合因子：
\begin{equation}
    \mathrm{OBCVP}_t = \alpha \cdot \mathrm{OCVP}_t + (1 - \alpha) \cdot \mathrm{BCVP}_t
\end{equation}

\subsection{变量定义}
\begin{itemize}
    \item ① \( w_i \) 为开盘集合竞价期间不同时间点上的成交量加权权重，可以是等权，也可以是时间衰减权重；
    \item ② \( \alpha \) 是两大成交量占比因子的复合权重。
\end{itemize}

\subsection{说明}
集合竞价阶段是反映投资者行为信息的重要时点，集合竞价成交量因子能反映多空双方之间的博弈，集合竞价成交量占比越低，股票次月的收益率越高。

\subsection{参考文献}
\begin{enumerate}
    \item刘均伟.2017.见微知著：成交量占比高频因子解析.光大证券
\end{enumerate}

\section{盈利卖出倾向(行为金融因子-处置效应)}

\subsection{因子定义}

1. 计算 \( t \) 时刻收盘价 \( \mathrm{Close}_t \) 与前 \( n \) 天 VWAP 均价 \( P_{t-n} \) 之间的涨跌幅，取涨跌幅为正的部分作为盈利收益率：
\begin{equation}
    \mathrm{gain}_t = \frac{\mathrm{Close}_t - P_{t-n}}{\mathrm{Close}_t} \cdot \mathbf{1}_{\{\mathrm{Close}_t \geq P_{t-n}\}}
\end{equation}

2. 以换手率 \( V_t \) 为权重加权盈利收益率得到资本盈利量即为盈利卖出倾向因子 \( \mathrm{Gain} \)：
\begin{equation}
    \mathrm{Gain}_t = \sum_{n=0}^{N-1} w_{t-n} \cdot \mathrm{gain}_{t-n}
\end{equation}
\begin{equation}
    w_{t-n} = \frac{1}{k} V_{t-n} \prod_{s=1}^{n-1} (1 - V_{t-n+s})
\end{equation}
\begin{equation}
    k = \sum_{n=1}^{T} V_{t-n} \prod_{s=1}^{n-1} (1 - V_{t-n+s})
\end{equation}

\subsection{说明}
盈利卖出倾向因子衡量了大部分投资者“未实现”的盈利，因子值越大，投资者未实现收益越大，卖出意愿越强，抛压大，股价上涨受阻，与未来收益负相关。

\subsection{参考文献}
\begin{enumerate}
    \item 陈升锐, 姚紫薇. 2025. 处置效应与 V 型处置效应在量化选股中的应用. 中信建投.
    \item An, L. (2016). \textit{Asset Pricing When Traders Sell Extreme Winners and Losers}. The Review of Financial Studies, 29(3), 823-861.
\end{enumerate}

\section{基于高低价的买卖价差(高频因子-流动性类)}

\subsection{因子定义}
\begin{equation}
    \mathrm{HLI}_{i,t} = \frac{\mathrm{HL}_{i,t}}{\$ \mathrm{vol}_{i,t}}
\end{equation}
\begin{equation}
    \mathrm{HL}_{i,t} = \frac{2 \times (e^{\alpha_t} - 1)}{1 + e^{\alpha_t}}
\end{equation}
\begin{equation}
    \alpha_t = \frac{\sqrt{2 \times \beta_t} - \sqrt{\beta_t}}{3 - 2\sqrt{2}} - \sqrt{\frac{\gamma_t}{3 - 2\sqrt{2}}}
\end{equation}
\begin{equation}
    \beta_t = \frac{\sum_{j=0}^{1} \left[ \ln \left( \frac{H_{t+j}^0}{L_{t+j}^0} \right) \right]^2}{2}, \quad \gamma_t = \left[ \ln \left( \frac{H_{t,t+1}^0}{L_{t,t+1}^0} \right) \right]^2
\end{equation}

\subsection{变量定义}
\begin{itemize}
    \item ① \( H_{t+j}^0, L_{t+j}^0 \) 分别为 \( t+j \) 日股票的最高价和最低价；
    \item ② \( H_{t,t+1}^0, L_{t,t+1}^0 \) 分别为股票 \( t \) 日至 \( t+1 \) 日这两天的最高价和最低价；
    \item ③ \( \$ \mathrm{vol}_{i,t} \) 为 \( t \) 日当天成交额。
\end{itemize}

\subsection{说明}
每日的最高价几乎总是由买方发起的交易，而每日的最低价几乎总是由卖方发起的交易，高低价格比反映了股票价格的真实方差和买卖价差，所以可以此估计买卖价差；买卖价差越大，流动性越弱，未来可获得流动性风险溢价。

\subsection{参考文献}
\begin{enumerate}
    \item 陈升悦. 2024. 流动性高频因子构建与投资者注意力因子分析报告. 中信建投.
\end{enumerate}

\section{买单集中度(高频因子-资金流类)}

\subsection{因子定义}
\begin{equation}
    \frac{1}{T} \sum_{n=t}^{t-T+1} \frac{\sum_{j=1}^{N} \text{买单成交额}_{i,j,n}^2}{(\text{总成交额}_{i,n})^2}
\end{equation}

\subsection{变量定义}
\begin{itemize}
    \item ① 基于逐笔成交数据中的叫买与叫卖单号将逐笔成交数据合成为买卖单数据；
    \item ② \( i, j, n \) 分别表示第 \( i \) 只股票在第 \( n \) 个交易日的第 \( j \) 分钟的成交额；
    \item ③ 月度选股下，\( T=20 \) 个交易日；
    \item ④ 将买单替换成卖单即可得卖单集中度。
\end{itemize}

\subsection{说明}
成交集中度类因子刻画了日内买卖单成交金额分布的均匀程度，集中度越高，未来超额收益表现越好；买单集中度因子与卖单集中度因子的 IC 方向相同，并未呈现出“买单集中度看多，卖单集中度看空”的现象。

\subsection{参考文献}
\begin{enumerate}
    \item 冯佳睿, 袁林青. 2019. 买卖单数据中的 Alpha. 海通证券.
\end{enumerate}

\section{筹码乖离率(行为金融因子-筹码分布)}

\subsection{因子定义}
\begin{equation}
    \mathrm{BIAS}_t = \mathrm{winner}_t \times \mathrm{turnover}_t + \mathrm{BIAS}_{t-1} \times (1 - \mathrm{turnover}_t)
\end{equation}

\subsection{变量定义}
\begin{itemize}
    \item ① \( \mathrm{winner}_t \) 为 \( t \) 日盈利筹码占总筹码的比例，即筹码分布图中红色面积占比；
    \item ② \( \mathrm{BIAS}_{t-1} \) 为 \( t-1 \) 日乖离率；
    \item ③ \( \mathrm{turnover}_t \) 为股票 \( t \) 日换手率。
\end{itemize}

\subsection{说明}
筹码乖离率采用动态加权平均的方式，用于刻画盈利筹码的持续性变化，即股价与筹码持仓成本间的偏离与回归情况；一般取值越高，盈利程度越高，趋势稳定，投资者情绪相对乐观。

\subsection{参考文献}
\begin{enumerate}
    \item 陈升锐, 姚紫薇. 2025. 筹码分布因子系统构建. 中信建投.
\end{enumerate}

\section{个股收益最低时的领先成交量异动(高频因子-动量反转类)}

\subsection{因子定义}
\begin{equation}
    \mathrm{AMT\_MIN}(M)_i = \frac{\sum_{t=m1}^{mM} \mathrm{AMT\_Percent}_{i,t-1}}{\frac{1}{T}\sum_{t=1}^{T} \mathrm{AMT\_Percent}_{i,t}}
\end{equation}
\begin{equation}
    \mathrm{AMT\_Percent}_{i,t} = \frac{\mathrm{AMT}_{i,t}}{\sum_{k=1}^{I} \mathrm{AMT}_{k,t}}
\end{equation}

\subsection{变量定义}
\begin{itemize}
    \item ① \( \mathrm{AMT}_{i,t} \) 为个股 \( i \) 在 \( t \) 分钟的成交额；
    \item ② \( \mathrm{AMT\_Percent}_{i,t} \) 为个股 \( i \) 在 \( t \) 时刻的成交额占全市场总成交额的比例；
    \item ③ \( m1, m2, \ldots, mM \) 为日内分钟收益率最低的 \( M \) 个分钟，\( M=30 \)；
    \item ④ \( T \) 为日内分钟数 = 240 个 1 分钟。
\end{itemize}

\subsection{说明}
\( \mathrm{AMT\_MIN} \) 衡量了股价大幅下跌时的成交量在全市场成交量中的占比情况；若个股的股价下跌没有伴随市场成交量异动，而是由自身成交量异动造成的（因子值越大，占比越大），未在成交量上与市场形成有效共振关系，则容易受到自身交易过热的反噬，未来收益表现较弱。

\subsection{参考文献}
\begin{enumerate}
    \item 任峻, 麦元勋, 许继宏. 2024. 基于分钟数据的价差共振因子. 招商证券.
\end{enumerate}

\section{分钟理想振幅(高频因子-波动跳跃类)}

\subsection{因子定义}

① 对单只股票，回溯取其最近 \( N \) 个交易日的 1 分钟数据（如：\(10 \times 240\) 个分钟 K 线），计算每分钟的分钟振幅（最高价 / 最低价 - 1）；

② 选择所有分钟收盘价较高的 \( \lambda \)（如 25\%）有效交易分钟，计算振幅均值得到高价振幅因子 \( V_{\text{high}}(\lambda) \)；选择所有分钟收盘价较低的 \( \lambda \)（如 25\%）有效分钟，计算振幅均值得到低价振幅因子 \( V_{\text{low}}(\lambda) \)；

③ 计算理想振幅因子：
\begin{equation}
    V(\lambda) = V_{\text{high}}(\lambda) - V_{\text{low}}(\lambda)
\end{equation}

\subsection{说明}
基于股价将振幅因子进行切割得到的理想振幅因子，考虑了不同价格区间的振幅分布信息差异（高价振幅因子具有更强的负向选股能力），可用于刻画不同价格位置的资金多空博弈情况；将日频数据提升至分钟级，对交易信息提取的更充分，选股效果更优。

\subsection{参考文献}
\begin{enumerate}
    \item 魏建榕, 高鹏. 2025. 高频振幅因子的内部切割. 开源证券.
\end{enumerate}

\section{趋势清晰度动量(趋势清晰度动量)}

\subsection{因子定义}
\begin{equation}
    \mathrm{TM}_{i,T} = -|\mathrm{MOM}'_{i,T} - \mathrm{TC}'_{i,T}|
\end{equation}
\begin{equation}
    \mathrm{TC}'_{i,T} = R^2
\end{equation}
\begin{equation}
    P_{i,t} = \beta_{i,0} + \beta_{i,1} \times t + \varepsilon_t
\end{equation}
\begin{equation}
    \mathrm{Mom}_{i,T} = \sum_{t=T-20}^{T-240} r_{i,t}
\end{equation}
其中，\(\mathrm{MOM}'_{i,T}\) 和 \(\mathrm{TC}'_{i,T}\) 为经过横截面标准化（减均值除以标准差）后的动量和趋势清晰度。

\subsection{变量定义}
\begin{itemize}
    \item ① \( P_{i,t} \) 为股票 \( i \) 在第 \( t \) 日的收盘价；
    \item ② \( t \in [T-240, T-20] \) 用 \( T-240 \) 交易日到 \( T-20 \) 交易日的收盘价时序序列进行回归（即不包含最近一个月），时序回归至少需要 200 个交易日的可用数据；
    \item ③ \( r_{i,t} \) 为股票 \( i \) 在第 \( t \) 日的收益率。
\end{itemize}

\subsection{说明}
“股价时间趋势回归的 \( R^2 \) 值”与“投资者对股价形成的趋势清晰感知度”之间存在正相关关系，可借助回归的 \( R^2 \) 对趋势清晰度进行量化。过去下跌股票中，下跌路径越明确、趋势越清晰的股票，未来更有可能维持下跌；下跌路径不明确、趋势不清晰的股票，未来更有可能反转上涨；过去上涨的股票中，上涨路径越明确、趋势越清晰的股票，未来更有可能维持上涨趋势；上涨路径不明确、趋势不清晰的股票，未来更有可能下跌；可借助趋势清晰度因子来加强动量。

\subsection{参考文献}
\begin{enumerate}
    \item 郑琳琳, 盛宝丹. 2024. 价格形成路径与趋势清晰度因子. 西南证券.
\end{enumerate}

\section{量价背离时刻笔均成交金额(高频因子-成交分布类)}

\subsection{因子定义}

1. 将日内 240 个 1 分钟数据集合标记为 \( T = \{t_1, t_2, \dots, t_{240}\} \)，计算每分钟所有成交记录的价格与成交量的相关系数（若该分钟内成交笔数不足 10 笔，则将相关系数置为空），将相关系数最小的 50\% 时刻作为量价背离时刻，标记为集合 \( P = \{t_{i1}, t_{i2}, \dots, t_{ik}\} \)；

2. 对量价背离时刻进行汇总，根据其成交额 \( \mathrm{Amt}_t \) 之和除以成交笔数 \( \mathrm{DealNum}_t \) 之和，构造量价背离时刻的笔均成交金额：
\begin{equation}
    \mathrm{ATD}_{\mathrm{PriceVolCorrLower50\%}} = \frac{\sum_{t \in P} \mathrm{Amt}_t}{\sum_{t \in P} \mathrm{DealNum}_t}
\end{equation}

3. 同理，将全天的成交金额除以全天的成交笔数，得到全天的笔均成交金额 \( \mathrm{ATD}_T \)：
\begin{equation}
    \mathrm{ATD}_T = \frac{\sum_{t \in T} \mathrm{Amt}_t}{\sum_{t \in T} \mathrm{DealNum}_t}
\end{equation}

4. 将量价背离时刻的笔均成交金额除以全天的笔均成交金额，即得到去量纲化后的标准化因子 \( \mathrm{SATD}_{\mathrm{PriceVolCorrLower50\%}} \)：
\begin{equation}
    \mathrm{SATD}_{\mathrm{PriceVolCorrLower50\%}} = \frac{\mathrm{ATD}_{\mathrm{PriceVolCorrLower50\%}}}{\mathrm{ATD}_T}
\end{equation}

\subsection{说明}
笔均成交额类因子通过汇总“特殊的、具有较多信息含量”时刻的成交金额情况来刻画主力资金的行为动向；笔均成交金额越大，说明该特殊时间内资金优势越明显，属于主力资金的可能性越大。基于“量价背离时刻”计算的笔均成交金额因子对未来收益有较为显著的正向预测能力。

\subsection{参考文献}
\begin{enumerate}
    \item 张欣鹏, 张宇. 2025. 日内特殊时刻蕴含的主力资金 Alpha 信息. 国信证券.
\end{enumerate}

\section{基于换手率的注意力捕捉(行为金融因子-投资者注意力)}

\subsection{因子定义}
对每个行业内分别统计每日行业成分股的平均换手率作为市场关注代理指标：
\begin{equation}
    \mathrm{indus\_turn}_d = \mathrm{mean}(\text{行业内成分股的日度换手率序列 } \mathrm{turnover}_{i,d})
\end{equation}

将市场关注代理指标与股票日收益进行月度时序回归，代理指标前的系数绝对值 \(|\beta_{i,t}|\) 即为股票对市场关注的敏感程度：
\begin{equation}
    r_{i,d} = \mu_{i,t} + \beta_{i,t} \mathrm{indus\_turn}_d + \rho_{1,i,t} \mathrm{Mkt}_d + \rho_{2,i,t} \mathrm{SMB}_d + \rho_{3,i,t} \mathrm{HML}_d + \rho_{4,i,t} \mathrm{UMD}_d + \varepsilon_{i,d}
\end{equation}

\subsection{变量定义}
\begin{itemize}
    \item ① \(n_{\mathrm{upper},d}\)、\(n_{\mathrm{lower},d}\)、\(N_d\) 分别为 \(d\) 日涨停股票数量、跌停股票数量、当日交易股票总数；
    \item ① \( r_{i,d} \) 为个股日度收益；
    \item ② \( \mathrm{Mkt}_d, \mathrm{SMB}_d, \mathrm{HML}_d \) 为经典的 Fama-French 三因子；
    \item ③ \( \mathrm{UMD}_d \) 为动量因子收益，即全市场过去 11 个月累积收益前后 30\% 的组合在 \( d \) 日的收益差。
\end{itemize}

\subsection{说明}
不同股票对同一市场关注度事件的反映程度往往是不同的，可借助个股收益对市场关注度事件的敏感程度来衡量股票对该事件的注意力捕捉效应；捕捉能力越强（敏感程度越高），表明个股更能吸引市场注意，从而导致买盘压力增大，股价高估；上述因子以换手率为市场关注度事件，换手率衡量投资者交易热度，反映关注水平。

\subsection{参考文献}
\begin{enumerate}
    \item Lin, F., Qiu, Z., \& Zheng, W. (2023). \textit{Cranes among chickens: The general-attention-grabbing effect of daily price limits in China's stock market}. Journal of Banking \& Finance, 150, 106818.
    \item 陈升锐. 2024. 投资者有限关注及注意力捕捉与溢出. 中信建投.
\end{enumerate}

\section{主买成交特异性(高频因子-资金流类)}

\subsection{因子定义}

1. 计算主买成交量占比：对于每一分钟，计算每个股票 \( t \) 分钟主买成交量占 \( t \) 分钟全市场主买成交总量的比例；

2. 确定主买时刻：对于每个股票，计算“主买成交量占比”大于该股票当日“主买成交占比”序列 80\% 分位数的时刻；

3. 修正主买时刻：进一步要求“主买时刻”的主买成交量大于前一分钟和后一分钟的主买成交量（对于第一分钟和最后一分钟，则分别考虑其最后一分钟和前一分钟）；

4. 统计主买次数：对于每一只股票，剔除开盘的前 15 分钟和收盘前的 15 分钟，统计剩余 210 分钟内“修正主买时刻”出现次数；

5. 高频因子低频化：每月月底，对“主买次数”进行截面标准化，然后计算过去 20 日的均值，得到最终的“主买成交特异性”因子。

\subsection{说明}
主买成交特异性因子刻画了持续且特异性的主买资金流入，代表资金看好该股票，未来股价上涨的可能性更大。“特异性”：个股相对于市场主买成交量的比例在截面上损益相消，“主买成交占比”的扩大代表主买资金“特异性”流入个股；“持续性”：修正主买时刻对连续出现的主买时刻进行修正剔除，而后内经过剔除之后剩余的“主买时刻”出现次数越多，代表其在全天分布比较分散，即股票当前需求不是一时兴起而是持续存在，未来上涨的概率就越大。

\subsection{参考文献}
\begin{enumerate}
    \item 叶尔乐. 2025. 资金流潮汐与“引力场”因子构建. 民生证券.
\end{enumerate}

\section{基于筹码分布的 V 型处置效应(行为金融因子-处置效应)}

\subsection{因子定义}

① 基于筹码分布等相关数据，计算当日兑现的亏损占筹码总亏损的比例 \( \mathrm{CPLR} \) 和当日兑现的盈利占筹码总盈利的比例 \( \mathrm{CPGR} \)：
\begin{equation}
    \mathrm{CPGR}_{i,t} = \frac{\text{已实现盈利}_{i,t}}{\text{已实现盈利}_{i,t} + \text{账面盈利}_{i,t}}
\end{equation}
\begin{equation}
    \mathrm{CPLR}_{i,t} = \frac{\text{已实现亏损}_{i,t}}{\text{已实现亏损}_{i,t} + \text{账面亏损}_{i,t}}
\end{equation}

② 计算如下几种筹码分布下的 V 型处置效应因子：
\begin{equation}
    \mathrm{VCDE1}_{i,t} = |\mathrm{CPGR}_{i,t} - \mathrm{CPLR}_{i,t}|
\end{equation}
\begin{equation}
    \mathrm{VCDE2}_{i,t} = \mathrm{CPGR}_{i,t} + 0.23 \times \mathrm{CPLR}_{i,t}
\end{equation}
\begin{equation}
    \mathrm{VCDE3}_{i,t} = \mathrm{CPGR}_{i,t} + \mathrm{CPLR}_{i,t}
\end{equation}

\subsection{变量定义}
\begin{itemize}
    \item ① 已实现盈利、账面盈利、已实现亏损、账面亏损等指标的计算逻辑，见“已实现盈利比率”和“已实现亏损比率”等相关因子日历页。
\end{itemize}

\subsection{说明}
V 型处置效应因子同时包含了盈利或亏损下投资者的卖出意愿，盈利或亏损幅度越大，V 型处置效应越明显，投资者的出售意愿越强。

\subsection{参考文献}
\begin{enumerate}
    \item 陈升锐, 姚紫薇. 2025. 处置效应与 V 型处置效应在量化选股中的应用. 中信建投.
    \item Odean, T. (1998). \textit{Are investors reluctant to realize their losses?}. Journal of Finance, 53(5), 1775-1798.
\end{enumerate}

\section{剔除超大单影响后的大单涨跌幅(高频因子-动量反转类)}

\subsection{因子定义}

计算每笔成交数据对应的对数价格变动：
\begin{equation}
    \Delta \ln P_i = \text{当前成交价格的对数}\ln P_i - \text{上一笔成交价格的对数}\ln P_{i-1}
\end{equation}

对于每笔成交数据中涉及到的买卖双方，将时间上靠后的订单标记为主动订单；计算所有主动订单为普通大单（剔除超大单）的对数价格变动之和即为普通大单涨跌幅：
\begin{equation}
    \mathrm{NormalBigRet} = \sum_{i \in S_{\mathrm{normal}}} \Delta \ln P_i
\end{equation}

\subsection{变量定义}
\begin{itemize}
    \item ① 按照委买单号和委卖单号对 level2 逐笔成交数据中的成交金额（已成交部分）进行汇总求和得到当天股票 \( i \) 所有买单和卖单的成交金额；
    \item ② 对股票 \( i \)，将订单成交金额占当天总成交额比例超过 1\% 的订单定义为超大单（阈值在 0.5\%~2\% 区间较好），将所有订单成交金额的前 30\% 分位数的订单定义为大单。
\end{itemize}

\subsection{说明}
超大单涨跌幅因子具有负向 Alpha，剔除超大单后的普通大单涨跌幅的正向 Alpha 会在原始因子基础上得到加强。

\subsection{参考文献}
\begin{enumerate}
    \item 朱剑涛, 王星星. 超大单冲击对大单因子的影响. 2022, 东方证券.
\end{enumerate}

\section{动态知情交易概率(高频因子-资金流类)}

\subsection{因子定义}

构建下列自回归模型，计算股票 \( i \) 在日内区间 \( j \) 内的非预期收益率 \( \varepsilon_{i,j} \)：
\begin{equation}
    R_{i,j} = \gamma_0 + \sum_{k=1}^{4} \gamma_{1i,k} D_k^{\mathrm{Day}} + \sum_{k=1}^{48} \gamma_{2i,k} D_k^{\mathrm{Int}} + \sum_{k=1}^{12} \gamma_{3i,k} R_{i,j-k} + \varepsilon_{i,j}
\end{equation}

计算未预期收益率为正（负）时卖出（买入）交易量占比即为动态知情交易概率：
\begin{equation}
    \mathrm{DPIN}_{\mathrm{BASE}}^{i,j} = \frac{\mathrm{NB}_{i,j}}{\mathrm{NT}_{i,j}} \cdot \mathbb{I}_{(\varepsilon_{i,j} < 0)} + \frac{\mathrm{NS}_{i,j}}{\mathrm{NT}_{i,j}} \cdot \mathbb{I}_{(\varepsilon_{i,j} > 0)}
\end{equation}

\subsection{变量定义}
\begin{itemize}
    \item ① \( R_{i,j} \) 为股票 \( i \) 在交易日 \( t \) 内区间 \( j \) 的收益率，\( R_{i,j-k} \) 为自回归的滞后收益率，滞后阶数 \( k=12 \)；
    \item ② \( D_k^{\mathrm{Day}} \) 为周内效应虚拟变量，当前交易日属于周 \( k \) 就取 1，其余取 0；同理，\( D_k^{\mathrm{Int}} \) 为日内效应虚拟变量，日内区间 \( j \) 的 \( D_j^{\mathrm{Int}} \) 取 1，其余取 0；
    \item ③ \( \mathrm{NB}_{i,j}, \mathrm{NS}_{i,j}, \mathrm{NT}_{i,j} \) 分别为股票 \( i \) 在日内区间 \( j \) 的主买成交笔数、主卖成交笔数、总成交笔数；
    \item ④ \( \mathbb{I}_{(\varepsilon_{i,j} < 0)} \) 为虚拟变量，未预期收益为负，取值为 1，否则取值为 0，反之亦然；
    \item ⑤ 自回归的时间窗口可以为从交易日 \( t \) 的区间 \( j \) 开始回溯过去 \( T \) 个交易日，如 5 分钟频率，则回溯过去 \( T \times 48 \) 个区间；
    \item ⑥ 可以进一步计算日度、周度内 DPIN 序列的均值、标准差等统计指标，刻画 DPIN 的时序特征。
\end{itemize}

\subsection{说明}
DPIN 基于日内交易活动将得知新信息后的反转交易占比作为知情交易概率的衡量指标，同时考虑了知情买入交易和知情卖出交易；知情交易概率指在一段时间内，拥有信息优势的知情交易者提交的订单数量占总委托单数量的比例，用以刻画该段时间内的信息不对称程度，通常与未来收益正相关，取值越大，信息不对称程度越大，投资者要求的风险回报越高。

\subsection{参考文献}
\begin{enumerate}
    \item Chang, S. S., Chang, L. V., \& Wang, F. A. (2014). \textit{A dynamic intraday measure of the probability of informed trading and firm-specific return variation}. Journal of Empirical Finance, 29, 80-94.
\end{enumerate}

\section{基于异常收益的注意力溢出(行为金融因子-投资者注意力)}

\subsection{因子定义}

计算每月平均异常日收益率作为注意力代理指标：
\begin{equation}
    \mathrm{ATTN}_{i,t} = \frac{1}{m} \sum_{t=1}^{t-m+1} (r_{i,t} - \bar{r}_t)^2
\end{equation}

计算当期相近股票的注意力均值与个股注意力之差即为注意力溢出程度：
\begin{equation}
    \mathrm{SPILL} = \overline{\mathrm{ATTN}}_{i} - \mathrm{ATTN}_{i}
\end{equation}
其中，\(\overline{\mathrm{ATTN}}_{i}\) 为与股票 \( i \) 相近的所有股票的注意力均值。

\subsection{变量定义}
\begin{itemize}
    \item ① 相近股票的确定标准：先以中信一级行业作为分类标准，再在每个行业内按 3:4:3 的比例将股票分为大、中、小市值；
    \item ② \( r_{i,t} \) 为股票 \( i \) 在 \( t \) 交易日的收益率；
    \item ③ \( \bar{r}_t \) 为交易日 \( t \) 的全市场收益率；
    \item ④ \( m \) 为交易日窗口，一般取为月度。
\end{itemize}

\subsection{说明}
投资者注意力存在溢出效应，当一只股票关注度较高时，投资者也会注意到与其相似的邻居股票，导致这些邻居股票的未来收益也会有不错的收益，享有注意力溢价；上述因子以异常收益为注意力代理指标，计算注意力差值，注意力差值越大，表明未来注意力溢出效应越显著，股票预期会取得正向超额。

\subsection{参考文献}
\begin{enumerate}
    \item Chen, X., An, L., Wang, Z., \& Yu, J. (2023). \textit{Attention spillover in asset pricing}. The Journal of Finance, 78(6), 3515-3559.
    \item 陈升锐. 2024. 投资者有限关注及注意力捕捉与溢出. 中信建投.
\end{enumerate}

\section{“情绪溢出”因子(高频因子-成交分布类)}

\subsection{因子定义}

1. 将股票按照代码从小到大进行排序后，对每只股票计算相邻 \( n \) 只股票在过去 \( T \) 日的日均换手率，并等权合成为焦点股票 \( i \) 的 \( \mathrm{NBR\_tov} \)；

2. 将焦点股票 \( i \) 的 \( \mathrm{NBR\_tov} \) 对其自身的日均换手率进行截面回归，回归得到的残差即为“情绪溢出”因子 \( \mathrm{RNBR\_tov} \)：
\begin{equation}
    \mathrm{turnover}_{i,T} = \alpha + \beta \cdot \mathrm{NBR\_tov}_{i,T} + \varepsilon_i
\end{equation}
\begin{equation}
    \mathrm{RNBR\_tov}_{i,T} = \varepsilon_i
\end{equation}

\subsection{变量定义}
\begin{itemize}
    \item ① 焦点股票 \( i \) 的相邻股票数对因子效果影响不大，建议使用较少数量的股票（如 \( n=10 \)）；
    \item ② 日均涨跌幅的计算窗口 \( T \) 建议取 5-20 天最为适宜。
\end{itemize}

\subsection{说明}
由于投资者的注意力有限，以及行情软件通常是根据代码顺序依次展示股票信息，投资者会较轻易地关注到与焦点股票相邻的其它股票，出现“注意力溢出”，焦点股票与相邻股票之间会存在价格与交易情绪的转移；“情绪溢出”对应股票近期引起大量关注时的强烈交易情绪溢出到相邻股票的情况，与未来收益正相关。

\subsection{参考文献}
\begin{enumerate}
    \item 任曦, 袁元勋, 李世杰. 2023. 基于股票代码有序性的“注意力溢出”. 招商证券.
\end{enumerate}

\section{绝对收益与滞后成交量相关性(高频因子-量价相关性类)}

\subsection{因子定义}
\begin{equation}
    \mathrm{CORA\_A} = \mathrm{corr}(|\mathrm{Ret}_t|, \mathrm{Amount}_{t-1}), \quad t \in [1, 240]
\end{equation}

\subsection{变量定义}
\begin{itemize}
    \item ① \( |\mathrm{Ret}_t| \) 为个股日内第 \( t \) 分钟的对数收益率的绝对值；
    \item ② \( \mathrm{Amount}_{t-1} \) 为个股日内第 \( t-1 \) 分钟的成交额，相比收益率滞后一期；
    \item ③ 月度因子可使用月末 5-15 个交易日因子的平均值来合成。
\end{itemize}

\subsection{说明}
该因子衡量了前一分钟量与后一分钟价的协同运动情况，捕捉“量在价先”的异常交易行为，与未来收益负相关；取值越大，表示信息的提前泄露程度越高，聪明钱在抢跑交易，次月负向 Alpha 越显著。

\subsection{参考文献}
\begin{enumerate}
    \item 朱定豪, 严佳炜. 2020. 量价关系的高频乐章. 方正证券.
\end{enumerate}

\section{结合业务复杂度的相似业务收益联动因子(行为金融因子-投资者注意力)}

\subsection{因子定义}
\begin{equation}
    \mathrm{Linkage}_{i,t} = \mathrm{complexity}_{i,t} \times \left( \frac{\sum_{j=1}^{N} \mathrm{SIM}_{i,j,t} \times \mathrm{Ret}_{j,t}}{\sum_{j=1}^{N} \mathrm{SIM}_{i,j,t}} - \mathrm{Ret}_{i,t} \right)
\end{equation}
\begin{equation}
    \mathrm{complexity}_{i,t} = \frac{\sum_{j=1}^{n} \mathrm{SIM}_{i,j,t}}{n} - \frac{\sum_{k=1}^{m} \mathrm{SIM}_{i,k,t}}{m}
\end{equation}
\begin{equation}
    \mathrm{SIM}_{i,j,t} = \frac{P_{i,t} \cdot P_{j,t}}{\|P_{i,t}\| \cdot \|P_{j,t}\|}
\end{equation}

\subsection{变量定义}
\begin{itemize}
    \item ① \( P_{i,t} \) 为取值只有 0 和 1 的业务关键词向量，通过自然语言技术处理 A 股上市公司年度财务报告附注中关于公司基本情况的描述部分内容得到；
    \item ② \( \mathrm{SIM}_{i,j,t} \) 为公司 \( i \) 和公司 \( j \) 在 \( t \) 时期的业务相似度；
    \item ③ \( \mathrm{complexity}_{i,t} \) 为公司 \( i \) 在 \( t \) 时期的业务复杂度，\( n \) 为不同行业公司的个数，\( m \) 为同行业公司的个数；
    \item ④ \( \mathrm{Ret}_{j,t} \) 为公司 \( j \) 在 \( t \) 月的收益率。
\end{itemize}

\subsection{说明}
相似业务收益联动因子度量了 \( t \) 期下与 \( i \) 公司相似的公司在当月超过 \( i \) 公司的收益率，该值越大，说明 \( i \) 公司越有可能在下一期出现一定的补涨行情；公司业务复杂度越高，投资者可能更难把握和判断市场信息对其的影响，带来认知资源的限制，进而对应更强的动量溢出效应。

\subsection{参考文献}
\begin{enumerate}
    \item 叶尔乐. 2023. 财报文本中公司竞争信息刻画与 Alpha 构建. 民生证券.
\end{enumerate}

\section{异常高波动下的收益波动比(高频因子-波动跳跃类)}

\subsection{因子定义}

1. 对于股票 \( i \)，计算 \( n \) 日内第 \( t \) 分钟的“更优波动率”，即对第 \( t-4, t-3, t-2, t-1, t \) 分钟的开盘价、最高价、最低价和收盘价，共 20 个价格数据求标准差，再除以这 20 个价格数据的均值，最后对该比值取平方：
\begin{equation}
    \mathrm{BetterVol}_{i,t} = \left( \frac{\mathrm{std}(P_{t-4}, P_{t-3}, P_{t-2}, P_{t-1}, P_t)}{\mathrm{mean}(P_{t-4}, P_{t-3}, P_{t-2}, P_{t-1}, P_t)} \right)^2
\end{equation}

2. 计算 \( t \) 分钟的收益波动比，即 \( t \) 分钟收益率 / \( t \) 分钟“更优波动率”：
\begin{equation}
    \mathrm{RVR}_{i,t} = \frac{r_{i,t}}{\mathrm{BetterVol}_{i,t}}
\end{equation}

3. 计算 \( n \) 日“更优波动率”的均值 \( \mu \) 和标准差 \( \sigma \)，将日内“更优波动率”大于等于 \( \mu + \sigma \) 的时刻定义为当日波动率异常高时刻；

4. 对当日波动率异常高时刻的收益波动比序列与“更优波动率”序列求协方差，即为日度“勇攀高峰”因子：
\begin{equation}
    \mathrm{PeakClimber}_{i,d} = \mathrm{cov}(\mathrm{RVR}_{i,t}, \mathrm{BetterVol}_{i,t}), \quad t \in \{ \text{波动率异常高时刻} \}
\end{equation}

\subsection{变量定义}
\begin{itemize}
    \item ① 计算因子需剔除开盘和收盘数据，仅考虑日内分钟频数据；
    \item ② 每月末可对最近 20 天的上述因子求均值和标准差并将两者等权合并，即得到月度“勇攀高峰”因子。
\end{itemize}

\subsection{说明}
“勇攀高峰”因子衡量了股价对异常高波动提供的风险补偿程度，与未来收益正相关；那些能对异常高波动及时提供风险补偿的股票，往往是持续利好的实力股，此时若能克服对异常高波动的风险厌恶，反倒能获得实力股异常高波动下的充分风险补偿。

\subsection{参考文献}
\begin{enumerate}
    \item 曹春晓. 2022. 个股波动率的变动及“勇攀高峰”因子构建. 方正证券.
    \item Lochstoer, L. A., \& Muir, T. (2022). \textit{Volatility expectations and returns}. The Journal of Finance, 77(2), 1055-1096.
\end{enumerate}

\section{边际交易成本(高频因子-流动性类)}

\subsection{因子定义}

\begin{equation}
    \mathrm{MCI}_A = \frac{\mathrm{VWAPM}_A }{\mathrm{DolVol}_A}
\end{equation}
\begin{equation}
    \mathrm{VWAPM}_A = \frac{VWAP_A-M}{M}
\end{equation}
\begin{equation}
    \mathrm{VWAP}_A = \frac{\sum_i^nP_{A,i} \times Q_{A,i}}{\sum_i^nQ_{A,i}} = \frac{DolVol_A}{\sum_i^n Q_{A,i}}
\end{equation}

\subsection{变量定义}
\begin{itemize}
    \item ① \( \mathrm{VWAPM}_A \) 为市价交易成本，即卖单五档报单的报单量加权平均价格；
    \item ② \( \mathrm{DolVol}_A \) 为卖单报单总金额，即所有卖单报单价格乘以报单量之和（共 5 档）；
    \item ③ \( M = (P_{A,1} + P_{B,1})/2 \) 为限价交易成本，即买一价和卖一价的均值；
    \item ④ \( P_{A,i} \) 为第 \( i \) 个卖单的价格，\( Q_{A,i} \) 为第 \( i \) 个卖单的报单量；
    \item ⑤ 若要计算买单的 \( \mathrm{MCI}_B \) 需将卖单数据替换成买单数据，并乘上 -1；
    \item ⑥ MCI 因子的统一单位为 bps/万元。
\end{itemize}

\subsection{说明}
MCI 衡量了市场的流动性；买单的 \( \mathrm{MCI}_B \) 衡量了发起市价交易的卖方所付出的交易费用，短期内，\( \mathrm{MCI}_B \) 较大时，卖方要付出较大的费用卖出股票，股票买压更强，较难下跌，与未来短期收益正相关；卖单的 \( \mathrm{MCI}_A \) 衡量了发起市价交易的买方所付出的交易费用，短期内，\( \mathrm{MCI}_A \) 较大时，买方要付出较大的费用买入股票，来衡量买卖报单流动性的不平衡性。

\subsection{参考文献}
\begin{enumerate}
    \item 丁鲁明, 陈升锐. 2020. 买卖报单流动性因子构建. 中信建投证券.
\end{enumerate}

\section{趋势资金净支撑量因子(量价因子改进)}

\subsection{因子定义}

1. 每个交易日 \( t \)，回看过去 \( k \) 个交易日 \( t-k, t-(k-1), \dots, t-1 \)，计算过去 \( k \) 日分钟成交量的 \( m\% \) 分位数，定义为成交量阈值，\( k=5, m=90 \)；将交易日 \( t \) 的每分钟成交量与上述阈值进行对比，若某一分钟的成交量大于上述阈值，则该分钟的交易被定义为趋势资金的交易；

2. 每个交易日，计算所有趋势资金分钟收盘价的均值，筛选趋势资金中收盘价小于该平均值的分钟，将它们的成交量加总，得到趋势资金支撑成交量；同理，将收盘价大于上述平均值的分钟成交量加总，得到趋势资金阻力成交量；

3. 每个交易日，计算如下因子：
\begin{equation}
    \text{趋势资金净支撑量} = \frac{\text{趋势资金支撑成交量} - \text{趋势资金阻力成交量}}{\text{当日流通股本}}
\end{equation}

4. 每月月底回看过去 20 个交易日，计算上述因子的平均值，并做横截面市值、中信一级行业中性化处理。

\subsection{说明}
趋势资金净支撑成交量因子只考虑趋势资金的净支撑情况，因子表现显著提升；因子值较大，则表明支撑价格不破位的力量较强，后市上涨概率较高；若因子值较小，则表明阻碍价格上涨的力量较强，后市相对偏空。

\subsection{参考文献}
\begin{enumerate}
    \item 沈亚琦, 刘富兵. 2024. 基于趋势资金日内交易行为的选股因子. 国盛证券.
\end{enumerate}

\section{趋势占比(高频因子-动量反转类)}

\subsection{因子定义}
\begin{equation}
    \mathrm{trend\_ratio}_{D,i} = \frac{P_{T,D,i} - P_{1,D,i}}{\sum_{t=2}^{T} |P_{t,D,i} - P_{t-1,D,i}|}
\end{equation}

\subsection{变量定义}
\begin{itemize}
    \item ① \( P_{t,D,i} \) 为股票 \( i \) 在交易日 \( D \) 内分钟频率下的价格；
    \item ② \( t = 1,2,3,\ldots,T \) 为交易日内相应分钟频率的第 \( t \) 个时间区间，如 5 分钟频率，\( T=48 \)。
\end{itemize}

\subsection{说明}
趋势占比可以看作是日内价格位移与路程之比，衡量了日内股价的趋势强度，取值范围为 \([-1, 1]\)，通常与未来收益负相关。

\subsection{参考文献}
\begin{enumerate}
    \item 文巧钧, 安宁宁, 罗军. 2021. 高频量化数据因子化方法. 广发证券.
\end{enumerate}

\section{朴素主动占比(高频因子-资金流类)}

\subsection{因子定义}

\begin{equation}
    \text{主动买入金额}_i = \mathrm{Amount}_i \times t \left( \frac{\mathrm{Close}_i - \mathrm{Close}_{i-1}}{\sigma_{\Delta \mathrm{close}}} ,df\right)
\end{equation}

\begin{equation}
    \text{朴素主动占比} =  \frac{\sum_{i=1}^{N}\text{主动买入金额}_i}{\sum_{i=1}^{N} \mathrm{Amount}_i} 
\end{equation}

\subsection{变量定义}
\begin{itemize}
    \item ① \( t() \) 为 t 分布的累计分布函数，取值范围在 0 到 1 之间；
    \item ② \( \sigma_{\Delta \mathrm{close}} \) 为整个时间区间内，每段时间末截面收盘价的标准差，保证每段时间价格变动相对可比；
    \item ③ df 为自由度，针对相同的股价变动，自由度越小，则根据 t 分布得到的主动买入金额占比越小；
    \item ④ 数据频率可以为 1 分钟、5 分钟等。
\end{itemize}

\subsection{说明}
绝大部分时间段的成交几乎均有买卖双方各自驱动成交的部分，在计算主动买入金额时采用连续性值作为成交额所乘系数可以提高主买卖额估计准确度；连续值在划分主动买卖成交时，价格变动方向对应买卖驱动方向（价格正向变动以买方驱动为主，价格负向变动以卖方驱动为主），价格变动幅度对应驱动力量占比情况（价格变动越大，主要驱动力量占比越大，价格变动越小，主要驱动力量占比越小）。

\subsection{参考文献}
\begin{enumerate}
    \item 川栋, 郑超. 2020. 高频因子(七): 分布估计下的主动成交占比. 长江证券.
\end{enumerate}

\section{供给端中心性变化(图谱网络-网络结构)}

\subsection{因子定义}
\begin{equation}
    C(u) = \frac{n-1}{N-1} \cdot \frac{n-1}{\sum_{v=1}^{n-1} D(u, v)}
\end{equation}
\begin{equation}
    C^{\mathrm{st}}(u) = \frac{C(u) - C_{\mathrm{min}}}{C_{\mathrm{max}} - C_{\mathrm{min}}}
\end{equation}
\begin{equation}
    D(u, v) = e^{-s(u, v)}
\end{equation}
\begin{equation}
    s(u, v) = \frac{\mathrm{sales}_{u,v}}{\mathrm{Total\_Procurement}_u}
\end{equation}

\subsection{变量定义}
\begin{itemize}
    \item ① \( s(u, v) \) 为供应关联度，即公司 \( u \) 对供应商 \( v \) 的采购额在公司 \( u \) 总采购额中的占比（边方向为供应商 $\rightarrow$ 客户）；
    \item ② \( D(u, v) \) 为基于供应关联度计算的公司间的距离；
    \item ③ \( n \) 为与客户公司 \( u \) 相连的所有供应商公司 \( v \) 的总数；
    \item ④ \( N \) 为供应链网络节点总数。
\end{itemize}

\subsection{说明}
供给端中心性变化为当期标准化后的供应商接近中心性减去上期标准化中心性，衡量了客户公司核心地位的变化程度；横截面分组选股时，供应端中心性变化对头尾组有一定的区分度（中心性变强的组为 top 组，收益更高），但组间单调性较弱；供应端中心性变化的选股效果弱于客户端中性变化。

\subsection{参考文献}
\begin{enumerate}
    \item 任健, 麦元勋, 杨帆. 2022. 供应链中心性初探. 招商证券.
\end{enumerate}

\section{买入浮亏占比因子(高频因子-成交分布类)}

\subsection{因子定义}
\begin{equation}
    \mathrm{HCVOL} = \frac{\sum_{i=1}^{N} \mathrm{volume}_{\mathrm{buy},i} \cdot \mathbf{I}_{\{P_{\mathrm{buy},i} > \mathrm{close}\}}}{\mathrm{total\_volume}}
\end{equation}

\subsection{变量定义}
\begin{itemize}
    \item ① \( \mathrm{volume}_{\mathrm{buy},i} \) 和 \( P_{\mathrm{buy},i} \) 分别为逐笔成交数据中买单的成交量和成交价；
    \item ② \( \mathrm{close} \) 和 \( \mathrm{total\_volume} \) 分别为当日收盘价和成交量；
    \item ③ 可以对该因子进行小单改进和尾盘改进，即只统计小单的占比和只统计 \([14:30,14:57)\) 尾盘区间的占比。
\end{itemize}

\subsection{说明}
买入浮亏占比因子统计了高于收盘价买入的成交量占比，与未来收益正相关；因子取值越高，则投资者当天买入情绪较高，浮亏占比较低。根据遗憾规避理论，为了避免产生遗憾和后悔情绪，投资者存在惜售心理，未来抛压较低，有着更高的预期收益；遗憾规避理论认为非理性的投资者在做决策时，会倾向于避免产生后悔情绪并追求自豪感，避免承认之前的决策失误。

\subsection{参考文献}
\begin{enumerate}
    \item 高智威. 2023. 基于逐笔成交数据的遗憾规避因子. 国金证券.
\end{enumerate}

\section{聪明钱因子(高频因子-量价相关性类)}

\subsection{因子定义}

对选定股票，回溯取其过去10个交易日的分钟行情数据，构造指标：
\begin{equation}
    S_t = \frac{|R_t|}{V_t^{0.25}}
\end{equation}
其中，\( R_t \) 为第t分钟涨跌幅，\( V_t \) 为第t分钟成交量。

将分钟数据按照指标 \( S_t \) 从大到小进行排序，取成交量累积占比20\%的分钟，视为聪明钱交易；

计算聪明钱交易的成交量加权平均价 \( \mathrm{VWAP}_{\mathrm{smart}} \)；计算所有交易的成交量加权平均价 \( \mathrm{VWAP}_{\mathrm{all}} \)；

计算如下聪明钱因子：
\begin{equation}
    \text{聪明钱因子} = \frac{\mathrm{VWAP}_{\mathrm{smart}}}{\mathrm{VWAP}_{\mathrm{all}}}
\end{equation}

\subsection{说明}
聪明钱因子是通过在分钟行情数据的价量信息中识别机构参与交易的多寡而构建的跟踪聪明钱交易的选股因子，而聪明钱的交易是基于“单笔订单数量更大、订单报价更为激进”的交易特征来识别的。

\subsection{参考文献}
\begin{enumerate}
    \item 魏建榕. 2020. 聪明钱因子模型的2.0版本. 开源证券.
\end{enumerate}

\section{筹码分布峰度（偏度）(行为金融因子-筹码分布)}

\subsection{因子定义}

\begin{equation}
    \mathrm{kurtosis}_t = \sum \left( \frac{\mathrm{price}_{i,t} - \mathrm{mean}_t}{\mathrm{std}_t} \right)^4 \times p_{i,t}
\end{equation}

\begin{equation}
    \mathrm{skewness}_t = \sum \left( \frac{\mathrm{price}_{i,t} - \mathrm{mean}_t}{\mathrm{std}_t} \right)^3 \times p_{i,t}
\end{equation}

\subsection{变量定义}
\begin{itemize}
    \item ① 假设某只股票在 \( t \) 日的筹码分布共有 \( n \) 个价位 \( (i=1,2,\ldots,n) \)，\( \mathrm{vol}_{i,t} \) 为第 \( t \) 日 \( i \) 价位上的筹码峰高度，\( \mathrm{price}_{i,t} \) 为该筹码峰对应的价格；
    \item ② 筹码加权权重为 \( p_{i,t} = \dfrac{\mathrm{vol}_{i,t}}{\sum \mathrm{vol}_{i,t}} \)；
    \item ③ \( \mathrm{mean}_t = \sum (\mathrm{price}_{i,t} \times p_{i,t}) \) 为筹码分布加权平均成本；
    \item ④ \( \mathrm{std}_t = \sqrt{\sum (\mathrm{price}_{i,t} - \mathrm{mean}_t)^2 \times p_{i,t}} \) 为筹码成本加权标准差。
\end{itemize}

\subsection{说明}
筹码分布偏度因子衡量了筹码分布偏斜程度和偏斜方向；筹码分布峰度因子衡量了筹码分布在平均值处峰值的高低程度；经测试，筹码分布峰度因子有较好的选股能力，与未来收益负相关。

\subsection{参考文献}
\begin{enumerate}
    \item 陈升锐, 姚紫薇. 2025. 筹码分布因子系统构建. 中信建投.
\end{enumerate}

\section{市价偏离度(高频因子-流动性类)}

\subsection{因子定义}
\begin{equation}
    \mathrm{MPB}_t = \overline{\mathrm{TP}}_t - \frac{M_t + M_{t-1}}{2} = \overline{\mathrm{TP}}_t - \overline{\mathrm{MP}}_t
\end{equation}

\subsection{变量定义}
\begin{itemize}
    \item ①
        \begin{equation*}
            \overline{\mathrm{TP}}_t = 
            \begin{cases} 
                \dfrac{T_t}{V_t}, & V_t \neq 0 \\[6pt]
                \overline{\mathrm{TP}}_{t-1}, & V_t = 0
            \end{cases}
        \end{equation*}
        为平均交易价格；
    \item ②
        \begin{equation*}
            M_t = \frac{P_t^B + P_t^A}{2}
        \end{equation*}
        为买一价和卖一价的均值；
    \item ③ \(P_t^B\) 和 \(P_t^A\) 为 \(t\) 时刻的买一价和卖一价；\(T_t\) 和 \(V_t\) 为 \(t\) 时刻成交额和成交量。
\end{itemize}

\subsection{说明}
市场偏离度为平均交易价格和平均委托挂单价格的差值；当市场偏离度为正，交易均价更接近卖一价，为卖方发起的交易，卖压大，未来价格更趋向于下行，即市场交易均价将向市场平均中间价回归，反之亦然。

\subsection{参考文献}
\begin{enumerate}
    \item 丁鲁明, 陈升锐. 2020. 高频量价选股因子初探. 中信建投证券.
\end{enumerate}

\section{已实现双幂次变差(高频因子-波动跳跃类)}

\subsection{因子定义}
\begin{equation}
    \mathrm{RBV}_t = \mu_1^{-2} \sum_{i=2}^n |r_{t_i}| |r_{t_{i-1}}|
\end{equation}
\begin{equation}
    \mu_q = E(|Z|^q) = 2^{q/2} \frac{\Gamma\left(\frac{1}{2}(q+1)\right)}{\Gamma\left(\frac{1}{2}\right)}
\end{equation}
\begin{equation}
    \Gamma(\alpha) = \int_0^{+\infty} x^{\alpha-1} e^{-x} dx
\end{equation}

\subsection{变量定义}
\begin{itemize}
    \item ① \( r_{t_i} \) 为某股票交易日 \( t \) 内的分钟对数收益率序列，如 5 分钟对数收益率序列；
    \item ② \( k=3 \)，\( n \) 表示该股票在交易日 \( t \) 中特定频率的分钟级别数据个数，如 5 分钟频率通常有 48 个数据点；
    \item ③ \( \mu_q \) 为标准正态分布的第 \( q \) 阶绝对矩。
\end{itemize}

\subsection{说明}
积分波动率 Integrated Variance 刻画了股价波动率中股价连续分量的波动水平（非跳跃部分），作者证明了积分波动率可以通过多幂次变差来有效估计，比如常见的已实现双幂次变差 realized bipower variation，已实现三幂次变差 realized tripower variation 等。

\subsection{参考文献}
\begin{enumerate}
    \item Barndorff-Nielsen, O. E., \& Shephard, N. (2004). \textit{Power and bipower variation with stochastic volatility and jumps}. Journal of Financial Econometrics, 2, 1-48.
\end{enumerate}

\section{每日最大异常成交量(行为金融因子-投资者注意力)}

\subsection{因子定义}
\begin{equation}
    \mathrm{ABNVOLD} = \max_t \left( \frac{\mathrm{VOL}_{i,t}}{\overline{\mathrm{VOL}}_{i,t}} \right)
\end{equation}
其中，\(\overline{\mathrm{VOL}}_{i,t}\) 为股票 \( i \) 在 \( t \) 交易日滚动 1 年的平均成交量。

\subsection{变量定义}
\begin{itemize}
    \item ① \( \mathrm{VOL}_{i,t} \) 为股票 \( i \) 在 \( t \) 交易日的成交量；
    \item ② \( \overline{\mathrm{VOL}}_{i,t} \) 为股票 \( i \) 在 \( t \) 交易日滚动 1 年的平均成交量；
    \item ③ \( m \) 为交易日窗口，一般取为月度。
\end{itemize}

\subsection{说明}
每日最大异常成交量为日频成交量除以过去一年平均成交量的最大值，属于极端交易量类因子；异常成交量也是投资者注意力的常用代理指标，因子值越高，投资者对该股票的关注度越高；而有限注意力导致投资者更倾向于交易引起他们关注的股票，投资者关注度越高的股票通常意味着越多的投资购买；但 A 股市场的做多与做空不对称，使得投资者非理性买入高于非理性卖出，进而导致关注度高的股票由于投资者净买入而存在溢价，未来也更可能出现反转。

\subsection{参考文献}
\begin{enumerate}
    \item Cosemans, M., \& Frehen, R. (2021). \textit{Salience theory and stock prices: Empirical evidence}. Journal of Financial Economics, 140(2), 460-483.
    \item 陈升锐. 2024. 投资者有限关注及注意力捕捉与溢出. 中信建投.
\end{enumerate}

\section{偏锋涨跌幅(高频因子-动量反转类)}

\subsection{因子定义}

1. 每分钟计算市场上所有上涨或者下跌股票收益率的平均值，分别作为“正向动量基准”和“负向动量基准”；

2. 对每个股票、每分钟，计算该股票1分钟收益率与当前分钟同方向的“动量基准”之差，记为1分钟偏离度，若该股票1分钟收益率为0或没有股票与之收益率方向相同，则1分钟偏离度为0；

3. 每分钟截面计算全市场股票偏离度的幅度，并将其大于0的时刻作为市场的“动量时刻”，将每只股票“动量时刻”的1分钟对数“偏离度”进行加总，得到日度因子；

4. 高频因子低频化：每月月底，分别计算步骤3中日度因子过去20天的标准差，即为“日内偏锋涨跌幅”因子。

\subsection{说明}
1分钟偏离度的偏度刻画了当前时段市场上股票处于“离群上涨”或“离群下跌”的比例，若1分钟偏离度的偏度大于0，说明这1分钟市场上股票偏离度分布呈现右偏，即当前时段市场上离群上涨（超涨）或离群下跌（少跌）的股票的数量更多，上涨的股票更有可能过度反应，下跌的股票更有可能反应不足，在动量时段内“偏离度”越高的股票，日后越有可能使得股票价格会有更低的预期收益。

\subsection{参考文献}
\begin{enumerate}
    \item 叶尔乐. 2025. 日内“动量脉冲”与股价过度反应的精细刻画. 民生证券.
\end{enumerate}

\section{主买占比(高频因子-资金流类)}

\subsection{因子定义}
\begin{equation}
    \frac{1}{T} \sum_{n=t}^{t-T+1} \frac{\sum_{j=1}^{N} \text{主动买入成交额}_{i,j,n}}{\sum_{j=1}^{N} \text{成交额}_{i,j,n}}
\end{equation}

\subsection{变量定义}
\begin{itemize}
    \item ① 基于逐笔成交数据中的 BS 标志将逐笔成交数据合成为分钟级别的主动买卖金额，B 为主动买入（卖出方先挂单，买入方主动触碰卖单并成交）；S 为主动卖出（买入方先挂单，卖出方主动触碰买单并成交）；
    \item ② 剔除了处于涨跌停分钟上的主动买入金额和主动卖出金额数据；
    \item ③ \( i, j, n \) 分别表示第 \( i \) 只股票在第 \( n \) 个交易日内的第 \( j \) 分钟的成交额；
    \item ④ 月度选股下，\( T = 20 \) 个交易日；周度选股下，\( T = 5 \) 个交易日。
\end{itemize}

\subsection{说明}
除收盘时段外，主买占比因子呈现出的是动量效应，股票前期主买占比越高，股票未来相对收益表现越好。而使用收盘前（14:27~14:56）数据计算得到的主买占比因子呈现出的是反转效应。

\subsection{参考文献}
\begin{enumerate}
    \item 冯佳睿, 袁林青. 2020. 基于主动买入行为的选股因子. 海通证券.
\end{enumerate}

\section{基于 PE 的分析师锚效益偏差(行为金融因子-锚定效应)}

\subsection{因子定义}

\begin{equation}
    \mathrm{CAF\_EP}_{i,t} = \frac{\mathrm{FEP}_{i,t} - \mathrm{IFEP}_t}{\mathrm{IFEP}_t}
\end{equation}
\begin{equation}
    \mathrm{FEP}_{i,t} = \frac{1}{\mathrm{FPE}_{i,t}}
\end{equation}

\subsection{变量定义}
\begin{itemize}
    \item ① \( \mathrm{FPE}_{i,t} \) 为分析师在时间 \( t \) 对股票 \( i \) 的 PE 指标的预测值；
    \item ② \( \mathrm{IFEP}_t \) 为股票 \( i \) 所属行业中所有成分股的 \( \mathrm{FPE}_{i,t} \) 中位数。
\end{itemize}

\subsection{说明}
锚定效应指人们在作出判断时会将看到或听到的第一印象的影响作为锚点，并根据锚点对判断做调整，使得调整不充分偏向于该锚点而造成判断偏差的现象；假设分析师预测也会有锚定效应，即分析师以公司所在行业作为定价的“锚”，当预测的财务指标超出（低于）行业中位数时，分析师会做出悲观（乐观）预测的指标，以接近行业中位数，财报公布后，被预测的公司的财务指标将高于（低于）预测值，股价表现将更优（更差）。

\subsection{参考文献}
\begin{enumerate}
    \item Ashoura, S., \& Hao, Q. (2019). \textit{Do analysts really anchor? Evidence from credit risk and suppressed negative information}. Journal of Banking \& Finance, 98, 183–197.
    \item 陈升锐. 2023. 行为金融学在量化选股中的应用. 中信出版社.
\end{enumerate}

\section{成交量比值(高频因子-成交分布类)}

\subsection{因子定义}
\begin{equation}
    \mathrm{VR}_t = \frac{1}{d} \sum_{i=1}^{d} w_{t-i} \times \left( \frac{\mathrm{Vol}_{\mathrm{morning}}}{\mathrm{Vol}_{\mathrm{afternoon}}} \right)_{t-i}
\end{equation}

\subsection{变量定义}
\begin{itemize}
    \item ① \( \dfrac{\mathrm{Vol}_{\mathrm{morning}}}{\mathrm{Vol}_{\mathrm{afternoon}}} \) 为每日上午、下午开盘的前 30 分钟成交量比值；
    \item ②
        \begin{equation*}
            w_{t-i} = 
            \begin{cases} 
                \dfrac{(1-\alpha)^{i-1}}{\sum_{j=1}^{\infty}(1-\alpha)^{j-1}}, & \text{指数加权} \\[6pt]
                1, & \text{算术平均}
            \end{cases}
        \end{equation*}
        为时间权重因子；
    \item ③ \( \alpha \) 为信息的衰减强度，即 \( \alpha = \dfrac{2}{1+d} \)；
    \item ④ \( d \) 为每月最后一个交易日向前回溯移动平均的交易日个数。
\end{itemize}

\subsection{说明}
上午开盘 30 分钟的成交量越小或者下午开盘 30 分钟的成交量越大，即成交量比值因子的值越小，对于股票次月收益预测的有效性越强。

\subsection{参考文献}
\begin{enumerate}
    \item 刘均伟. 2017. 高频因子：日内分时成交隐藏玄机. 光大证券.
\end{enumerate}

\section{量岭分钟收益(高频因子-量价相关性类)}

\subsection{因子定义}

① 对过去 20 日同时点成交量按照 1 倍标准差的标准划分，高于 1 倍标准差划分为“喷发成交量”，低于 1 倍标准差划分为“温和成交量”；

② 对“喷发成交量”进一步按照连续与否划分：若当前分钟为“喷发成交量”，且前后 1 分钟均为“温和成交量”，则划为“孤立喷发成交量”，称为“量峰”；若当前分钟为“喷发成交量”，且前后 1 分钟也存在“喷发成交量”，则划为“连续喷发成交量”，称为“量岭”；将“温和成交量”称为“量谷”。

③ 计算过去 20 日“量岭”时点的累计收益，即为量岭分钟收益因子。

\subsection{说明}
“量岭”对应的连续性大额交易，更符合个人投资者的跟随交易特征，“量岭”时刻的累计收益衡量了个人投资者交易的收益贡献，因子值越高，个人投资者交易越可能过度反应，未来收益越低。

\subsection{参考文献}
\begin{enumerate}
    \item 魏建榕, 王志豪. 2025. 高频成交量的峰、岭、谷信息. 开源证券.
\end{enumerate}

\section{最大客户动量(图谱网络-动量溢出)}

\subsection{因子定义}
\begin{equation}
    \mathrm{cmom}_i^{1M} = \mathrm{mom}_{j\_{\max\_\mathrm{sales}}}^{1M}
\end{equation}

\subsection{变量定义}
\begin{itemize}
    \item ① \( j\_{\max\_\mathrm{sales}} \) 为股票 \( i \) 的最大客户 \( j \)，最大客户 \( j \) 的动量即为最大客户动量因子；
    \item ② 动量的计算周期可以是 1 月动量、6 月动量、12 月动量等。
\end{itemize}

\subsection{说明}
最大客户动量的表现与通过完整客户加权得到的客户动量表现相近，客户动量溢出主要来源于最大客户。

\subsection{参考文献}
\begin{enumerate}
    \item Jussa et al. (2015). \textit{The Logistics of Supply Chain Alpha}. Deutsche Bank Markets Research.
\end{enumerate}

\section{反向日内逆转的异常频率(高频因子-收益分布类)}

\subsection{因子定义}

\begin{equation}
    \mathrm{AB\_NR}_{i,t} = \frac{\mathrm{NR}_{i,t}}{\frac{1}{12} \sum_{k=0}^{11} \mathrm{NR}_{i,t-k}}
\end{equation}
\begin{equation}
    \mathrm{NR}_{i,t} = \frac{\sum_{d=1}^{T} \mathrm{I}_{\{\mathrm{RET\_CO}_{i,d}>0\}} \times \mathrm{I}_{\{\mathrm{RET\_OC}_{i,d}<0\}}}{T}
\end{equation}

\subsection{变量定义}
\begin{itemize}
    \item ① \( \mathrm{RET\_CO}_{i,d} \) 和 \( \mathrm{RET\_OC}_{i,d} \) 分别为股票 \( i \) 在交易日 \( d \) 的隔夜收益率和日内收益率；
    \item ② \( T \) 为股票 \( i \) 在 \( t \) 月内的实际交易天数。
\end{itemize}

\subsection{说明}
当正的隔夜收益后伴随负的日内收益，即出现了反向日内逆转，通过统计月内出现反向日内逆转的频率，可以衡量隔夜交易者和盘中交易者间“拉锯战”的激烈程度；反向日内逆转的异常频率为当月反向日内逆转频率相对过去12个月的异常程度，考虑了“拉锯战”的持续性，取值越高，“拉锯战”强度增高，价格修正越过度。

\subsection{参考文献}
\begin{enumerate}
    \item Cheema, M. A., Chiah, M., \& Man, Y. (2022). \textit{Overnight returns, daytime reversals, and future stock returns: Is China different?}. Pacific-Basin Finance Journal, 74, 101809.
\end{enumerate}

\section{日内振幅切割因子(高频因子-波动跳跃类)}

\subsection{因子定义}


1. 对单只股票，回溯取其最近 \( N \) 个交易日（如 \( N=10 \)）的 1 分钟数据，计算每分钟的分钟振幅（最高价 / 最低价 -1）；

2. 对每个交易日的 240 个分钟，选择 1 分钟涨跌幅较高的 \( \lambda \)（如 20\%）有效交易分钟，计算振幅均值得到高价振幅因子 \( V_{\text{high}}(\lambda) \)；选择 1 分钟涨跌幅较低的 \( \lambda \)（如 20\%）有效分钟，计算振幅均值得到低价振幅因子 \( V_{\text{low}}(\lambda) \)；进一步计算日内振幅因子：
\begin{equation}
    V(\lambda) = V_{\text{high}}(\lambda) - V_{\text{low}}(\lambda)
\end{equation}

3. 回看过去 \( N \) 个交易日，计算日内振幅因子的均值 \( V_{\text{mean}} \) 和标准差 \( V_{\text{std}} \)，并进行横截面标准化，等权合成后即为最终的日内振幅切割因子。

\subsection{说明}
用于振幅切割的情绪指标除了股价外，还可以是振幅、成交额、涨跌幅等指标；日内振幅切割是对分钟理想振幅（见相关日历页）的另一种构造思路：先日内切割得到日度序列再合成，且两种计算方式得到的因子相关性并不高。

\subsection{参考文献}
\begin{enumerate}
    \item 魏建榕, 高鹏. 2025. 高频振幅因子的内部切割. 开源证券.
\end{enumerate}

\section{当日新增筹码盈利占比(行为金融因子-筹码分布)}

\subsection{因子定义}
\begin{equation}
    \mathrm{TGRatio}_{i,t} = \frac{\mathrm{TG}_{i,t}}{\mathrm{PG}_{i,t}}
\end{equation}
\begin{equation}
    \mathrm{PG}_{i,t} = \sum \left( \mathrm{IF}(p\leq close_{i,t},Chip_{i,t,p}* (close_{i,t}-p),0))
\end{equation}
\begin{equation}
    \mathrm{TG}_{i,t} = \sum (IF(p \leq close_{i,t},Vol_{i,t,p}*(close_{i,t}-p)*T_{i,t},0))
\end{equation}

\subsection{变量定义}
\begin{itemize}
    \item ① 假设某只股票 \( i \) 在 \( t \) 日的筹码分布共有 \( n \) 个价位（\( p = 1, 2, \ldots, n \)），\( \mathrm{Chip}_{i,t,p} \) 为第 \( t \) 日 \( p \) 价位上的筹码峰高度，\( p \) 为该筹码峰对应的价格；
    \item ② \( \mathrm{Vol}_{i,t} \) 为股票 \( i \) 在 \( t \) 日的成交量；
    \item ③ \( T_{i,t} \) 为股票 \( i \) 在 \( t \) 日的换手率；
    \item ④ \( \mathrm{close}_{i,t} \) 为股票 \( i \) 在 \( t \) 日的收盘价。
\end{itemize}

\subsection{说明}
当日新增筹码盈利占比为当日成交的新增筹码盈利总额占账面总盈利额的比例；  
若某只股票当日新增筹码的盈利比例较高，可能意味着该股票短期不稳定筹码比例较大，短期浮盈筹码可能快速获利了结，抛压变大，可能对股价走势产生不利影响。

\subsection{参考文献}
\begin{enumerate}
    \item 陈升锐, 姚紫薇. 2025. 筹码分布因子系统构建. 中信建投.
\end{enumerate}

\section{虹吸效应(高频因子-资金流类)}

\subsection{因子定义}

① 计算成交量放大倍数：\( \mathrm{AM} = \) 个股第 \( t \) 分钟成交量 / 个股第 \( t-1 \) 分钟成交量；

② 得到收益率修正的成交量放大倍数 \( \mathrm{RAM} \)：若个股当前分钟收益率 > 所有个股当日总成交额 \( \mathrm{TA} \) 加权的当前分钟平均收益率，则取个股成交量放大倍数 \( \mathrm{AM} \)；若个股当前分钟收益率 < 所有个股成交总额加权的当前分钟平均收益率，则为 0；

③ 计算市场炒作热度：个股当日成交额 \( \mathrm{TA} \) 加权的个股每分钟 \( \mathrm{RAM} \) 之和；

④ 标记热度时间：选取“市场炒作热度”最大的 23 分钟，记为“热度时间”，代表该分钟内成交量突然放大的股票较多，市场较为活跃（不考虑下午第 1 分钟、收盘最后 3 分钟）；

⑤ 计算虹吸效应：计算“热度时间”中，“个股每分钟主卖成交额占比”和“市场每分钟主卖成交额占比”的相关系数，其中：
    \begin{equation*}
        \text{个股每分钟主卖成交额占比} = \frac{\text{个股每分钟主卖成交额}}{\text{个股当日主卖成交额}}
    \end{equation*}
    \begin{equation*}
        \text{市场每分钟主卖成交额占比} = \frac{\text{市场每分钟主卖成交额}}{\text{市场当日主卖成交额}}
    \end{equation*}

⑥ 计算虹吸回流综合效应：将收盘前最后 5 分钟定义为“回流时刻”，将其视为 1 根 K 线，计算对应时刻的“个股每分钟主卖成交额占比”和“市场每分钟主卖成交额占比”；计算“热度时间”+“回流时刻”共 24 根 K 线中，“个股每分钟主卖成交额占比”和“市场每分钟主卖成交额占比”的相关系数即为“虹吸回流综合效应”；

⑦ 将“虹吸回流综合效应”对“虹吸效应”做差得到“净虹吸效应”：

⑧ 高频因子低频化：计算过去 20 天的“净虹吸效应”的均值和标准差，等权合成得到月度“虹吸效应”因子。

\subsection{说明}
“虹吸效应”衡量资金从没被炒作的股票被“虹吸”到被炒作的股票中，待炒作结束之后，资金“回流”到没被炒作的股票中，带动这部分股票价格产生反弹。如果资金被“虹吸”当天就产生了资金的回流，那么股票日后反弹的空间就会比较小，是一个反向因子。

\subsection{参考文献}
\begin{enumerate}
    \item 叶尔乐. 2025. 资金流潮汐与“引力场”因子构建. 民生证券.
\end{enumerate}

\section{加权盈利频率(量价因子改进)}

\subsection{因子定义}
\begin{equation}
    f_{i,t}^w = \frac{\sum_{j=1}^{M_{i,t}} w_j \cdot \mathbb{I}_{\{r_{i,j} > u\}}}{M_{i,t}}
\end{equation}
\begin{equation}
    w_{\mathrm{decay}_j} = 0.5^{\frac{t-j}{\lambda}}
\end{equation}

\subsection{变量定义}
\begin{itemize}
    \item ① \( M_{i,t} \) 为股票 \( i \) 在第 \( t \) 个回望期的天数，回望期可取 40 天；
    \item ② \( r_{i,j} \) 为股票 \( i \) 在第 \( j \) 个交易日的超额收益；
    \item ③ \( u \) 为收益阈值，可取 0.02；
    \item ④ \( w_{\mathrm{decay}_j} \) 为衰退指数加权权重，最近的收益计数获得最大的权重，而最远的收益计数获得最小的权重。
\end{itemize}

\subsection{说明}
加权盈利频率为股票过去一段时间收益率大于一定阈值的天数加权之和。因子值越大表示股票过去一段时间日超额收益大于一定阈值的天数越多并且出现的时间越近，越容易受到投资者关注，股票短期取得高收益的机会大，投机性需求较高，定价相对较高，未来预期回报较低；反之，股票并没有过多地受到投资者的关注，投机性需求小，股价处于相对低估状态，未来预期回报较高。

\subsection{参考文献}
\begin{enumerate}
    \item 郑琳琳, 盛宝丹. 2023. 计效启发法与加权盈利频率. 西南证券.
\end{enumerate}

\section{量价背离时刻-主卖笔均成交金额(高频因子-成交分布类)}

\subsection{因子定义}

1. 将日内 240 个 1 分钟数据集合标记为 \( T = \{t_1, t_2, \dots, t_{240}\} \)，计算每分钟所有成交记录的价格与成交量的相关系数（若该分钟内成交笔数不足 10 笔，则将相关系数置为空），将相关系数最小的 50\% 时刻作为量价背离时刻，标记为集合 \( P = \{t_{i1}, t_{i2}, \dots, t_{ik}\} \)；

2. 利用逐笔成交数据，进一步对 \( P \) 中时刻的所有成交记录划分为主动买入成交 \( P\_\mathrm{Buy} \) 和主动卖出成交 \( P\_\mathrm{Sell} \) 两类；

3. 对量价背离时刻下的主动卖出成交额和成交笔数进行汇总，根据其成交额 \( \mathrm{Amt}_t \) 之和除以成交笔数 \( \mathrm{DealNum}_t \) 之和，构造量价背离时刻且主卖的笔均成交金额 \( \mathrm{ATD}_{\mathrm{PriceVolCorrLower50\%\_Sell}} \)：
\begin{equation}
    \mathrm{ATD}_{\mathrm{PriceVolCorrLower50\%\_Sell}} = \frac{\sum_{t \in P\_\mathrm{Sell}} \mathrm{Amt}_t}{\sum_{t \in P\_\mathrm{Sell}} \mathrm{DealNum}_t}
\end{equation}

4. 同理，将全天的成交金额除以全天的成交笔数，得到全天的笔均成交金额 \( \mathrm{ATD}_T \)：
\begin{equation}
    \mathrm{ATD}_T = \frac{\sum_{t \in T} \mathrm{Amt}_t}{\sum_{t \in T} \mathrm{DealNum}_t}
\end{equation}

5. 将量价背离时刻且主卖的笔均成交金额除以全天的笔均成交金额，即得到去量纲化后的标准化因子 \( \mathrm{SATD}_{\mathrm{PriceVolCorrLower50\%\_Sell}} \)：
\begin{equation}
    \mathrm{SATD}_{\mathrm{PriceVolCorrLower50\%\_Sell}} = \frac{\mathrm{ATD}_{\mathrm{PriceVolCorrLower50\%\_Sell}}}{\mathrm{ATD}_T}
\end{equation}

\subsection{说明}
笔均成交额类因子通过汇总“特殊的、具有较多信息含量”时刻的成交金额情况来刻画主力资金的行为动向；笔均成交金额越大，说明该特殊时间内资金优势越明显，属于主力资金的可能性越大。“量价背离时刻”从投资者分歧角度刻画主力资金交易行为，叠加逐笔订单主卖意愿做更细切分能进一步提升因子表现。

\subsection{参考文献}
\begin{enumerate}
    \item 张欣鹏, 张宇. 2025. 日内特殊时刻蕴含的主力资金 Alpha 信息. 国信证券.
\end{enumerate}

\section{小单的动态知情交易概率(高频因子-资金流类)}

\subsection{因子定义}

构建下列自回归模型，计算股票 \( i \) 在日内区间 \( j \) 内的非预期收益率 \( \varepsilon_{i,j} \)：
\begin{equation}
    R_{i,j} = \gamma_0 + \sum_{k=1}^{4} \gamma_{1i,k} D_k^{\mathrm{Day}} + \sum_{k=1}^{48} \gamma_{2i,k} D_k^{\mathrm{Int}} + \sum_{k=1}^{12} \gamma_{3i,k} R_{i,j-k} + \varepsilon_{i,j}
\end{equation}

计算未预期收益率为正（负）时卖出（买入）交易量占比即为动态知情交易概率：
\begin{equation}
    \mathrm{DPIN}_{\mathrm{SMALL}}^{i,j} = \left[ \frac{\mathrm{NB}_{i,j}}{\mathrm{NT}_{i,j}} \cdot \mathbb{I}_{(\varepsilon_{i,j} < 0)} + \frac{\mathrm{NS}_{i,j}}{\mathrm{NT}_{i,j}} \cdot \mathbb{I}_{(\varepsilon_{i,j} > 0)} \right] \cdot \mathrm{ST}_{i,j}
\end{equation}

\subsection{变量定义}
\begin{itemize}
    \item ① \( R_{i,j} \) 为股票 \( i \) 在交易日 \( t \) 内区间 \( j \) 的收益率，\( R_{i,j-k} \) 为自回归的滞后收益率，滞后阶数 \( k=12 \)；
    \item ② \( D_k^{\mathrm{Day}} \) 为周内效应虚拟变量，当前交易日属于周 \( k \) 就取 1，其余取 0；同理，\( D_k^{\mathrm{Int}} \) 为日内效应虚拟变量，日内区间 \( j \) 的 \( D_j^{\mathrm{Int}} \) 取 1，其余取 0；
    \item ③ \( \mathrm{NB}_{i,j} \)、\( \mathrm{NS}_{i,j} \)、\( \mathrm{NT}_{i,j} \) 分别为股票 \( i \) 在日内区间 \( j \) 的主买成交笔数、主卖成交笔数、总成交笔数；
    \item ④ \( \mathbb{I}_{(\varepsilon_{i,j} < 0)} \) 为虚拟变量，未预期收益为负，取值为 1，否则取值为 0，反之亦然；
    \item ⑤ \( \mathrm{ST}_{i,j} \) 为小单的虚拟变量，若交易日 \( t \) 内，股票 \( i \) 在区间 \( j \) 内的总交易量小于当日各个区间内总交易量的中位数时，赋值为 1；
    \item ⑥ 自回归的时间窗口可以为从交易日 \( t \) 的区间 \( j \) 开始回溯过去 \( T \) 个交易日，如 5 分钟频率，则回溯过去 \( T \times 48 \) 个区间；
    \item ⑦ 可以进一步计算日度、周度内小单 DPIN 序列的均值、标准差等统计指标，刻画小单 DPIN 的时序特征。
\end{itemize}

\subsection{说明}
在市场交易不活跃的时段，知情交易者为不暴露自身的信息优势，更有可能选择将大订单拆成小订单进行交易，所以可以进一步筛选出小订单下的知情交易概率；日内交易不活跃时段的小单知情交易概率通常更高，日内因子走势整体呈倒“U”型。

\subsection{参考文献}
\begin{enumerate}
    \item Chang, S. S., Chang, L. V., \& Wang, F. A. (2014). \textit{A dynamic intraday measure of the probability of informed trading and firm-specific return variation}. Journal of Empirical Finance, 29, 80-94.
\end{enumerate}

\section{每笔成交量差分标准差(高频因子-成交分布类)}

\subsection{因子定义}

\begin{equation}
    \mathrm{diff}_i = \mathrm{volb}_i - \mathrm{volb}_{i-1}
\end{equation}
\begin{equation}
    \mathrm{diff\_mean_{day}} = \mathrm{mean} \left( \mathrm{abs} \left( \frac{\mathrm{diff}_i}{\mathrm{mean}(\{\mathrm{volb}_i\})} \right) \right)
\end{equation}
\begin{equation}
    \mathrm{diff\_std_{volb}} = \mathrm{mean}(\{ \mathrm{diff\_std_{day}} \})
\end{equation}

\subsection{变量定义}
\begin{itemize}
    \item ① \( \mathrm{volb}_i \) 为个股日内 240 根 1 分钟 K 线的每笔成交量（成交量/成交笔数），\( \mathrm{diff\_std_{day}} \) 取过去 20 个交易日数据。
\end{itemize}

\subsection{说明}
上述因子通过计算差分从成交量数据中提取出时间序列信息，是对数据局部峰值的度量，而局部峰值可能是对趋势交易者或知情交易者交易行为的刻画，其参与交易越多的个股在未来一段时间内的收益越高，因子表现出正向收益能力。

\subsection{参考文献}
\begin{enumerate}
    \item 覃川桃, 郑超. 2020. 高频因子(九): 高频波动中的时间序列信息. 长江证券.
\end{enumerate}

\section{基于换手率的注意力溢出(行为金融因子-投资者注意力)}

\subsection{因子定义}
计算每月日度换手率均值作为注意力代理指标：
\begin{equation}
    \mathrm{ATTN}_{i,t} = \frac{1}{m} \sum_{k=1}^{t-m+1} \mathrm{TURN}_{i,k}
\end{equation}

计算当期相近股票的注意力均值与个股注意力之差即为注意力溢出程度：
\begin{equation}
    \mathrm{SPILL} = \overline{\mathrm{ATTN}}_{i,t} - \mathrm{ATTN}_{i,t}
\end{equation}
其中，\(\overline{\mathrm{ATTN}}_{i,t}\) 为与股票 \( i \) 相近的所有股票的注意力均值。

\subsection{变量定义}
\begin{itemize}
    \item ① 相近股票的确定标准：先以中信一级行业作为分类标准，再在每个行业内按 3:4:3 的比例将股票分为大、中、小市值；
    \item ② \( \mathrm{TURN}_{i,k} \) 为股票 \( i \) 在交易日 \( k \) 的换手率；
    \item ③ \( m \) 为交易日窗口，一般取为月度。
\end{itemize}

\subsection{说明}
投资者注意力存在溢出效应，当一只股票关注度较高时，投资者也会注意到与其相似的邻居股票，导致这些邻居股票的未来收益也会有不错的收益，享有注意力溢价；上述因子以日均换手率为注意力代理指标，计算注意力差值，注意力差值越大，表明未来注意力溢出效应越显著，股票预期会取得正向超额。

\subsection{参考文献}
\begin{enumerate}
    \item Chen, X., An, L., Wang, Z., \& Yu, J. (2023). \textit{Attention spillover in asset pricing}. The Journal of Finance, 78(6), 3515-3559.
    \item 陈升锐. 2024. 投资者有限关注及注意力捕捉与溢出. 中信建投.
\end{enumerate}

\section{滞后绝对收益与成交量相关性(高频因子-量价相关性类)}

\subsection{因子定义}
\begin{equation}
    \mathrm{CORA\_R} = \mathrm{corr}(|\mathrm{Ret}_{t-1}|, \mathrm{Amount}_t), \quad t \in [1, 240]
\end{equation}

\subsection{变量定义}
\begin{itemize}
    \item ① \( |\mathrm{Ret}_{t-1}| \) 为个股日内第 \( t-1 \) 分钟的对数收益率的绝对值，相比成交额滞后一期；
    \item ② \( \mathrm{Amount}_t \) 为个股日内第 \( t \) 分钟的成交额；
    \item ③ 月度因子可使用月末 5-15 个交易日因子的平均值来合成。
\end{itemize}

\subsection{说明}
该因子衡量了前一分钟价与后一分钟量的协同运动情况，捕捉“价在量先”的异常交易行为，与未来收益负相关；取值越大，表示市场交易情绪受股价涨跌的影响程度越高，聪明钱在利用波动频繁交易，次月负向 Alpha 越显著。

\subsection{参考文献}
\begin{enumerate}
    \item 朱定豪, 严佳炜. 2020. 量价关系的高频乐章. 方正证券.
\end{enumerate}

\section{高频特异波动率(高频因子-波动跳跃类)}

\subsection{因子定义}
在 Fama 三因子基础上加入流动性和反转因子，对过去 21 天个股收益率做线性时序回归，得到新特异率 \( = 1 - \) 拟合优度 \( R^2 \)：
\begin{equation}
    r_t = \alpha + \beta_{\mathrm{mkt}} \mathrm{MKT}_t + \beta_{\mathrm{smb}} \mathrm{SMB}_t + \beta_{\mathrm{hml}} \mathrm{HML}_t + \beta_{\mathrm{ret}} \mathrm{RET}_t + \beta_{\mathrm{liq}} \mathrm{LIQ}_t + \varepsilon_t
\end{equation}

计算个股过去 21 天的日内分钟收益率序列（频率控制在 30 分钟内较好）的标准差，得到高频波动率；再与新特异率开方相乘得到高频特异波动率：
\begin{equation}
    \text{特异波动率} = \sqrt{\text{新特异率} \times \text{基础波动率}}
\end{equation}

\subsection{说明}
在 Fama 三因子基础上加入和波动率类因子相关性较高的流动性和反转因子，可以更纯粹的剥离价格变动相对市场的偏离情况，有助于提升因子表现；特异波动率表现出的选股能力是通过个股波动（基础波动率）和个股特异（特异率）的正向关系，杠杆性地加强了个股股价变动的异常信息，个股波动部分用高频数据进行细化，也有助于提升因子表现。

\subsection{参考文献}
\begin{enumerate}
    \item 覃川桃, 郑起. 2019. 高频因子(六): 特异视角下的波动率. 长江证券.
\end{enumerate}

\section{客户端中心性变化(图谱网络-网络结构)}

\subsection{因子定义}

\begin{equation}
    C(u) = \frac{n-1}{N-1} \times \frac{n-1}{\sum_{v=1}^{n-1} D(u, v)}
\end{equation}
\begin{equation}
    C^{\mathrm{st}}(u) = \frac{C(u) - C_{\mathrm{min}}}{C_{\mathrm{max}} - C_{\mathrm{min}}}
\end{equation}
\begin{equation}
    D(u, v) = e^{-s(u,v)} s(u, v) = \frac{\mathrm{sales}_{u,v}}{\mathrm{Total\_Sales}_u}
\end{equation}

\subsection{变量定义}
\begin{itemize}
    \item ① \( s(u, v) \) 为客户关联度，即公司 \( u \) 对客户 \( v \) 的销售额在公司 \( u \) 总销售额中的占比（边方向为客户 $\rightarrow$ 供应商）；
    \item ② \( D(u, v) \) 为基于客户关联度计算的公司间的距离；
    \item ③ \( n \) 为与供应商公司 \( u \) 相连的所有客户公司 \( v \) 的总数；
    \item ④ \( N \) 为供应链网络节点总数。
\end{itemize}

\subsection{说明}
客户端中心性变化为当期标准化后的客户接近中心性减去上期标准化中心性，衡量了供应商公司核心地位的变化程度；横截面分组选股时，客户端中心性变化对好坏股票有一定的区分度（中心性变强的组为 top 组，收益更高）；客户端中心性变化的选股效果优于供应端中心性变化。

\subsection{参考文献}
\begin{enumerate}
    \item 任晓, 麦元勋, 杨帆. 2022. 供应链中心性初探. 招商证券.
\end{enumerate}

\section{“自信溢出”因子(高频因子-动量反转类)}

\subsection{因子定义}

1. 将股票按照代码从小到大进行排序后，对每只股票计算相邻 \( n \) 只股票在过去 \( T \) 日的日均涨跌幅，并等权合成为焦点股票 \( i \) 的 \( \mathrm{NBR\_ret} \)；

2. 将焦点股票 \( i \) 的 \( \mathrm{NBR\_ret} \) 对其自身的日均涨跌幅进行截面回归，回归得到的残差即为“自信溢出”因子 \( \mathrm{NRBR\_ret} \)：
\begin{equation}
    \mathrm{ret}_{i,T} = \alpha + \beta \cdot \mathrm{NBR\_ret}_{i,T} + \varepsilon_i
\end{equation}
\begin{equation}
    \mathrm{NRBR\_ret}_{i,T} = \varepsilon_i
\end{equation}

\subsection{变量定义}
\begin{itemize}
    \item ① 焦点股票 \( i \) 的相邻股票数对因子效果影响不大，建议使用较少数量的股票（如 \( n=10 \)）；
    \item ② 日均涨跌幅的计算窗口 \( T \) 建议取 5-20 天最为适宜。
\end{itemize}

\subsection{说明}
由于投资者的注意力有限，以及行情软件通常是根据代码顺序依次展示股票信息，投资者会较容易地关注到与焦点股票相邻的其它股票，出现“注意力溢出”，焦点股票与相邻股票之间会存在价格与交易情绪的转移；“自信溢出”对应投资者因近期获利的过度自信向邻近股票外溢的情况，与未来收益正相关。

\subsection{参考文献}
\begin{enumerate}
    \item 任建, 麦元勋, 李世杰. 2023. 基于股票代码有序性的“注意力溢出”. 招商证券.
\end{enumerate}

\section{多层订单斜率(高频因子-流动性类)}

\subsection{因子定义}
\begin{equation}
    \mathrm{MLQS}_{\mathrm{Weight}} = \frac{\sum_{i=1}^{5} w_i \times \mathrm{LogquoteSlope}_t^{(i)}}{\sum_{i=1}^{5} w_i}, \quad w_i = \frac{i}{5}, \quad i = 1, 2, 3, 4, 5
\end{equation}
\begin{equation}
    \mathrm{LogquoteSlope}_t = \frac{\log\left(P_t^A\right) - \log\left(P_t^B\right)}{\log\left(V_t^A\right) + \log\left(V_t^B\right)}
\end{equation}

\subsection{变量定义}
\begin{itemize}
    \item ① \( P_t^B \) 和 \( P_t^A \) 分别为 \( t \) 时刻的买一价和卖一价；
    \item ② \( V_t^B \) 和 \( V_t^A \) 分别为 \( t \) 时刻的买一量和卖一量；
    \item ③ 对订单单价和订单量取对数是为了使序列更接近正态分布，增加因子平稳性；
    \item ④ 考虑订单方向：假设卖为正向，买为负向。
\end{itemize}

\subsection{说明}
订单斜率因子衡量的是订单价差对订单量差的敏感度，是对市场流动性的衡量。因子值越大，代表买卖价差越大，市场流动性越差，因此其长期具有流动性风险溢价，长期表现越好。

\subsection{参考文献}
\begin{enumerate}
    \item 丁鲁明, 陈升锐. 2021. 多层次订单失衡及订单斜率因子. 中信建投证券.
    \item Seppi, D. J. (2001). \textit{Common factors in prices, order flows, and liquidity}. Journal of Financial Economics.
\end{enumerate}

\section{当日新增筹码亏损占比(行为金融因子-筹码分布)}

\subsection{因子定义}

\begin{equation}
    \mathrm{TLRatio}_{i,t} = \frac{\mathrm{TL}_{i,t}}{\mathrm{PL}_{i,t}}
\end{equation}
\begin{equation}
    \mathrm{PL}_{i,t} = \sum \left( \mathrm{IF}(p> close_{i,t},Chip_{i,t,p}* (close_{i,t}-p),0))
\end{equation}
\begin{equation}
    \mathrm{TL}_{i,t} = \sum (IF(p > close_{i,t},Vol_{i,t,p}*(p-close_{i,t})*T_{i,t},0))
\end{equation}

\subsection{变量定义}
\begin{itemize}
    \item ① 假设某只股票 \( i \) 在 \( t \) 日的筹码分布共有 \( n \) 个价位（\( p = 1, 2, \ldots, n \)），\( \mathrm{Chip}_{i,t,p} \) 为第 \( t \) 日 \( p \) 价位上的筹码峰高度，\( p \) 为该筹码峰对应的价格；
    \item ② \( \mathrm{Vol}_{i,t} \) 为股票 \( i \) 在 \( t \) 日的成交量；
    \item ③ \( T_{i,t} \) 为股票 \( i \) 在 \( t \) 日的换手率；
    \item ④ \( \mathrm{close}_{i,t} \) 为股票 \( i \) 在 \( t \) 日的收盘价。
\end{itemize}

\subsection{说明}
当日新增筹码亏损占比为当日成交的新增筹码亏损总额占账面总亏损额的比例；若某只股票当日新增筹码的亏损比例较高，也可能意味着该股票短期不稳定筹码比例越大，新资金被套，筹码抛压变大，可能对股价走势产生不利影响。

\subsection{参考文献}
\begin{enumerate}
    \item 陈升锐, 姚紫薇. 2025. 筹码分布因子系统构建. 中信建投.
\end{enumerate}

\section{T分布主动占比(高频因子-资金流类)}

\subsection{因子定义}

\begin{equation}
    T \text{分布主动买入金额}_i = \mathrm{Amount}_i \times t\left(\frac{\mathrm{ret}_i}{\sigma_{\mathrm{ret}}}, \mathrm{df}\right)
\end{equation}
\begin{equation}
    T \text{分布主动占比因子} = \frac{\sum_{i=1}^{N} T \text{分布主动买入金额}_i}{\sum_{i=1}^{N} \mathrm{Amount}_i}
\end{equation}

\subsection{变量定义}
\begin{itemize}
    \item ① \( t(\cdot) \) 为 t 分布的累计分布函数，取值范围在 0 到 1 之间；
    \item ② \( \mathrm{ret}_i \) 为每个时间段收益率，\( \sigma_{\mathrm{ret}} \) 为全部时间段收益率标准差，频率可为 1 分钟、5 分钟等；
    \item ③ \( \mathrm{df} \) 为自由度，针对相同的股价变动，自由度越小，则根据 t 分布得到的主动买入金额占比越小。
\end{itemize}

\subsection{说明}
原始的批量成交划分法以价格变动作为自变量，做价格变化相对于价格标准差的标准化处理，但价格变化是指数级别上的变动，可以将价格替换成收益率，既保留了价格变动方向和变动幅度上对主动买卖占比的刻画，时间序列上又不会存在受价格高低影响的占比估计误差。

\subsection{参考文献}
\begin{enumerate}
    \item 覃川桃, 郑起. 2020. 高频因子(七): 分布估计下的主动成交占比. 长江证券.
\end{enumerate}

\section{卖出反弹占比因子(高频因子-成交分布类)}

\subsection{因子定义}
\begin{equation}
    \mathrm{LCVOL} = \frac{\sum_{i=1}^{N} \mathrm{volume}_{\mathrm{sell},i} \cdot \mathbf{I}_{\{P_{\mathrm{sell},i} < \mathrm{close}\}}}{\mathrm{total\_volume}}
\end{equation}

\subsection{变量定义}
\begin{itemize}
    \item ① \( \mathrm{volume}_{\mathrm{sell},i} \) 和 \( P_{\mathrm{sell},i} \) 分别为逐笔成交数据中卖单的成交量和成交价；
    \item ② \( \mathrm{close} \) 和 \( \mathrm{total\_volume} \) 分别为当日收盘价和成交量；
    \item ③ 可以对该因子进行小单改进和尾盘改进，即只统计小单的占比和只统计 \([14:30,14:57)\) 尾盘区间的占比。
\end{itemize}

\subsection{说明}
卖出反弹占比因子统计了低于收盘价卖出的成交量占比，与未来收益负相关；因子取值越高，根据遗憾规避理论，为了避免产生遗憾和后悔情绪，投资者卖出股票后不愿再次买回的意愿较高，未来买入动力较低，有更低的预期收益；遗憾规避理论认为非理性的投资者在做决策时，会倾向于避免产生后悔情绪并追求自豪感，避免承认之前的决策失误。

\subsection{参考文献}
\begin{enumerate}
    \item 高智威. 2023. 基于逐笔成交数据的遗憾规避因子. 国金证券.
\end{enumerate}

\section{极端跟随行为\_比值因子(量价因子改进)}

\subsection{因子定义}
① 每个交易日 \( t \)，回看过去 \( k \) 个交易日 \( t-k, t-(k-1), \ldots, t-1 \)，计算过去 \( k \) 日分钟成交量的 \( m\% \) 分位数，定义为成交量阈值，\( k=5, m=90 \)；将交易日 \( t \) 的每分钟成交量与上述阈值进行对比，若某一分钟的成交量大于上述阈值，则该分钟的交易被定义为趋势资金的交易；

② 对于趋势资金的每一分钟，取随后 \( n \) 分钟里成交额最大的一分钟，将其视为跟随趋势资金的极端行为，记录这一分钟的成交额，\( n=5 \)；

③ 将极端跟随行为的分钟成交额，除以对应趋势资金的分钟成交额，得到该次跟随行为对应的比值；

④ 月底回看过去 20 个交易日，求 20 日比值的平均值，再做横截面市值、行业中性化处理。

\subsection{说明}
A股市场中“追涨”和“杀跌”的两股力量是非对称的，若某个股票上出现了强烈的羊群效应，通常对利好消息的反应过度现象会更加严重，股价未来下跌的概率也就更高。所以羊群效应通常会对股票未来收益产生负面影响；极端跟随行为\_比值因子值越大，意味着极端跟随行为越剧烈，羊群效应越明显，从而导致股票未来收益越低。

\subsection{参考文献}
\begin{enumerate}
    \item 沈芷琦, 刘富兵, 阮俊辉. 2024. 基于趋势资金日内交易行为的选股因子. 国盛证券.
\end{enumerate}

\section{主力交易情绪(高频因子-量价相关性类)}

\subsection{因子定义}
\begin{equation}
    \mathrm{TE} = \mathrm{RankCorr}(A_{\mathrm{order}}, C_{\mathrm{minute}})
\end{equation}

\subsection{变量定义}
\begin{itemize}
    \item ① \( A_{\mathrm{order}} \) 为某只股票某个交易内的分钟单笔成交金额序列；
    \item ② \( C_{\mathrm{minute}} \) 为某只股票某个交易内的分钟收盘价序列；
    \item ③ 滚动 20 个交易日计算 \( \mathrm{TE} \) 指标的均值，即主力交易情绪因子 \( \mathrm{MTE} \)。
\end{itemize}

\subsection{说明}
MTE 因子实质上反映了主力参与交易的相对价位，因子值越大，表明主力交易更倾向于出现在高价位，这是“逢高出货”的表现，反映了主力悲观态度；而因子值越小，则表明主力交易更多出现在低价位，这是“逢低吸筹”的表现，属于乐观情绪。

\subsection{参考文献}
\begin{enumerate}
    \item 魏建榕, 苏良. 2022. 高频因子：分钟单笔金额序列中的主力行为刻画. 开源证券.
\end{enumerate}

\section{已实现三幂次变差(高频因子-波动跳跃类)}

\subsection{因子定义}

\begin{equation}
    \mathrm{RTV}_t = V_{\frac{2}{3}, \frac{2}{3}, \frac{2}{3}} \; \mu_{\frac{2}{3}}^{-3}
\end{equation}

\begin{equation}
    V_{m_1, m_2, \dots, m_k} = \sum_{i=k}^{n} |r_{t_i}|^{m_1} |r_{t_{i-1}}|^{m_2} \cdots |r_{t_{i-k+1}}|^{m_k}
\end{equation}

\begin{equation}
    \mu_q = E(|Z|^q) = 2^{q/2} \frac{\Gamma\left(\frac{1}{2}(q+1)\right)}{\Gamma\left(\frac{1}{2}\right)}
\end{equation}

\begin{equation}
    \Gamma(\alpha) = \int_0^{+\infty} x^{\alpha-1} e^{-x} dx
\end{equation}

\subsection{变量定义}
\begin{itemize}
    \item ① \( r_{t_i} \) 为某股票交易日 \( t \) 内的分钟对数收益率序列，如 5 分钟对数收益率序列；
    \item ② \( k=3 \)，\( n \) 表示该股票在交易日 \( t \) 中特定频率的分钟级别数据个数，如 5 分钟频率通常有 48 个数据点；
    \item ③ \( \mu_q \) 为标准正态分布的第 \( q \) 阶绝对矩。
\end{itemize}

\subsection{说明}
积分波动率 Integrated Variance 刻画了股价波动率中股价连续分量的波动水平（非跳跃部分），作者证明了积分波动率可以通过多幂次变差来有效估计，比如常见的已实现双幂次变差 realized bipower variation，已实现三幂次变差 realized tripower variation 等。

\subsection{参考文献}
\begin{enumerate}
    \item Barndorff-Nielsen, O. E., \& Shephard, N. (2004). \textit{Power and bipower variation with stochastic volatility and jumps}. Journal of Financial Econometrics, 2, 1–48.
\end{enumerate}

\section{最大供应商动量(图谱网络-动量溢出)}

\subsection{因子定义}
\begin{equation}
    \mathrm{smm}_i^{1M} = \mathrm{mom}_{j\_{\max\_\mathrm{sales}}}^{1M}
\end{equation}

\subsection{变量定义}
\begin{itemize}
    \item ① \( j\_{\max\_\mathrm{sales}} \) 为股票 \( i \) 的最大供应商 \( j \)，最大供应商 \( j \) 的动量即为最大供应商动量因子；
    \item ② 动量的计算周期可以是 1 月动量、6 月动量、12 月动量等。
\end{itemize}

\subsection{说明}
最大供应商动量的表现与通过完整供应商加权得到的供应商动量表现相近，供应商动量溢出主要来源于最大供应商。

\subsection{参考文献}
\begin{enumerate}
    \item Jussa et al. (2015). \textit{The Logistics of Supply Chain Alpha}. Deutsche Bank Markets Research.
\end{enumerate}

\section{日内最大回撤(高频因子-动量反转类)}

\subsection{因子定义}
\begin{equation}
    \mathrm{intraday\_maxdrawdown}_{D,i} = \min_{0 < t < T} \; \min_{0 < \tau < T - t} \left( \frac{p_{t+\tau,D,i}}{p_{t,D,i}} - 1 \right)
\end{equation}

\subsection{变量定义}
\begin{itemize}
    \item ① \( p_{t,D,i} \) 为股票 \( i \) 在交易日 \( D \) 内分钟频率下的价格；
    \item ② \( t = 1,2,3,\ldots,T \) 为交易日内相应分钟频率的第 \( t \) 个时间区间，如 5 分钟频率，\( T=48 \)。
\end{itemize}

\subsection{说明}
日内最大回撤统计了股票日内的最大回撤情况，取值范围为 \((-1,0]\)，可用于衡量一天内股价的上涨趋势程度，通常与未来收益正相关；取值越大（越接近于 0），说明日内股价上涨趋势越明显，上涨情绪越稳定。

\subsection{参考文献}
\begin{enumerate}
    \item 文巧钧, 安宁宁, 罗军. 2021. 高频量化数据因子化方法. 广发证券.
\end{enumerate}

\section{中间价变化率(高频因子-流动性类)}

\subsection{因子定义}
\begin{equation}
    \mathrm{MPC}_{t,k} = \frac{M_t - M_{t-k}}{M_{t-k}}, \quad M_t = \frac{P_t^B + P_t^A}{2}
\end{equation}

\subsection{变量定义}
\begin{itemize}
    \item ① \( \mathrm{MPC}_{t,k} \) 为从时刻 \( t-k \) 到时刻 \( t \) 中间价 \( M_t \) 的百分比变化率；
    \item ② 市场中间价 \( M_t \) 为买一价和卖一价的平均值；
    \item ③ \( M_{t-k} \) 为 \( k \) 分钟前的中间价，其中 \( k = 1, 5 \)；
    \item ④ \( P_t^B \) 和 \( P_t^A \) 分别为 \( t \) 时刻的买一价和卖一价。
\end{itemize}

\subsection{说明}
MPC 因子衡量了市场中间价的短期变动趋势，刻画了市场交易者的最新交易和挂单撤单行为；如当市场中间价上升时，可能出现了大量的主动买入交易，抬高了卖一价，或是有交易者以高于当前买一价的价格下了限价订单，抬高了买一价，这两种情况都反应了市场交易者对未来股价走势的乐观预期；MPC 与股票未来收益正相关。

\subsection{参考文献}
\begin{enumerate}
    \item 丁鲁明, 陈升锐. 2021. 高频订单失衡及价差因子. 中信建投证券.
\end{enumerate}

\section{基于跳跃频率关联的点度中心性(图谱网络-动量溢出)}

\subsection{因子定义}

1. 利用股票日内 5 分钟行情数据，计算跳跃统计量 \( \mathrm{JS} \)，识别出股价发生跳跃的分钟时刻，具体计算细节见相关日历页：
\begin{equation}
    \mathrm{JS} = N \frac{\hat{V}_{(0,1)}}{\sqrt{\hat{\Omega}_{\mathrm{SWV}}}} \left( 1 - \frac{\mathrm{RV}_N}{\mathrm{SwV}_N} \right) \sim N(0,1)
\end{equation}

2. 计算关联跳跃次数：如果股票 \( i \) 在 \( t \) 日发生跳跃，股票 \( j \) 在 \( t-1 \) 日至 \( t+1 \) 日内也发生了一次同向跳跃，则记为一次关联跳跃；

3. 计算两只股票间的跳跃关联度，即在股票 \( i \) 所有的股价跳跃 \( \mathrm{jump\_num}_i \) 中关联跳跃 \( \mathrm{cojump\_num}_{i,j} \) 所占的比例：
\begin{equation}
    \mathrm{Corr}_{i,j} = \frac{\mathrm{cojump\_num}_{i,j}}{\mathrm{jump\_num}_i}
\end{equation}

4. 对网络做稀释处理，剔除整个网络中关联度最低的 50\% 的关联关系，统计每只股票与之直接相连的股票数量即为该股票的点度中心性值。

\subsection{说明}
具有相似属性的公司往往会受到类似的信息冲击，进而导致股价的同步跳跃。公司间股价跳跃的同步程度也可用于识别公司间的潜在关联程度；跳跃关联网络的点度中心性因子衡量了股票在整个网络中的重要性或影响程度，与未来收益正相关。

\subsection{参考文献}
\begin{enumerate}
    \item 任晓, 刘凯, 杨航. 2025. 基于股价跳跃关联性的选股策略. 招商证券.
    \item 任晓, 杨航. 2024. 如何识别股价跳跃. 招商证券.
\end{enumerate}

\section{正向日内逆转的异常频率(高频因子-收益分布类)}

\subsection{因子定义}

\begin{equation}
    \mathrm{AB\_PR}_{i,t} = \frac{\mathrm{PR}_{i,t}}{\frac{1}{12} \sum_{k=0}^{11} \mathrm{PR}_{i,t-k}}
\end{equation}
\begin{equation}
    \mathrm{PR}_{i,t} = \frac{\sum_{d=1}^{T} \mathrm{I}\{\mathrm{RET\_CO}_{i,d}<0\} \times \mathrm{I}\{\mathrm{RET\_OC}_{i,d}>0\}}{T}
\end{equation}

\subsection{变量定义}
\begin{itemize}
    \item ① \( \mathrm{RET\_CO}_{i,d} \) 和 \( \mathrm{RET\_OC}_{i,d} \) 分别为股票 \( i \) 在交易日 \( d \) 的隔夜收益率和日内收益率；
    \item ② \( T \) 为股票 \( i \) 在 \( t \) 月内的实际交易天数。
\end{itemize}

\subsection{说明}
当负的隔夜收益后伴随正的日内收益，即出现了正向日内逆转，通过统计月内出现正向日内逆转的频率，也可以衡量隔夜交易者和盘中交易者间“拉锯战”的强度；正向日内逆转的异常频率为当月正向日内逆转频率相对过去 12 个月的异常程度，考虑了“拉锯战”的持续性，取值越高，“拉锯战”强度增高，价格修正越过度。

\subsection{参考文献}
\begin{enumerate}
    \item Cheema, M. A., Chiah, M., \& Man, Y. (2022). \textit{Overnight returns, daytime reversals, and future stock returns: Is China different?}. Pacific-Basin Finance Journal, 74, 101809.
\end{enumerate}

\section{一致交易(高频因子-成交分布类)}

\subsection{因子定义}

通过定义实体 K 线，将交易日中的每根 5 分钟 K 线划分为集体一致交易与非集体一致交易：
\begin{equation}
    |\mathrm{close} - \mathrm{open}| \leq \alpha |\mathrm{high} - \mathrm{low}|
\end{equation}

利用属于集体一致交易的 K 线成交量，计算一致交易类因子：
\begin{equation}
    \mathrm{TCV}_t = \frac{1}{d} \sum_{i=1}^{d} \left( \frac{\mathrm{ConsistentVolume}}{\mathrm{Volume}} \right)_{t-i}
\end{equation}

\subsection{变量定义}
\begin{itemize}
    \item ① \( \alpha \) 为一致参数，\( \alpha \) 越大，K 线一致性越强（上下引线越短，K 线越实体）；
    \item ② \( \mathrm{ConsistentVolume} \) 为 5 分钟实体 K 线的总成交量；
    \item ③ \( \mathrm{Volume} \) 为当日总成交量；
    \item ④ \( d \) 表示移动平均的周期参数。
\end{itemize}

\subsection{说明}
集体一致交易行为预示交易机会，如果一支股票在一段时间内大部分的交易属于集体一致交易行为，趋势性显著，则很可能这支股票在这段时间内有新信息进入，正在经历市场消化过程（price-in），未来股价很可能产生大幅变化，为交易者创造交易机会。

\subsection{参考文献}
\begin{enumerate}
    \item 刘均伟. 2018. 一致交易因子:挖掘集体行为背后的收益. 光大证券.
\end{enumerate}

\section{已实现亏损比率(行为金融因子-处置效应)}

\subsection{因子定义}

1. 提取股票分钟级前复权收盘价、成交量，以及日频流通换手率数据，计算股票每日分钟级量价分布情况，其中 \( t \) 日股票 \( i \) 在每个价位的筹码峰计算公式如下：
\begin{equation}
    \mathrm{Chip}_{i,t,p} = \mathrm{Vol}_{i,t,p} \times T_{i,t} + \mathrm{Chip}_{i,t-1,p} \times (1 - T_{i,t})
\end{equation}

2. 基于 \( t \) 日换手和 \( t-1 \) 日筹码峰，计算当天换手的筹码的亏损总金额作为已实现亏损：
\begin{equation}
    \text{已实现亏损}_{i,t} = \sum_p \mathrm{abs}(\mathrm{Close}_{i,t} - \mathrm{Price}_{i,p}) \times (\mathrm{Chip}_{i,t-1,p} \times T_{i,t}) \times \mathrm{I}_{\mathrm{Close}_{i,t} < \mathrm{Price}_{i,p}}
\end{equation}

3. 基于 \( t \) 日筹码峰，计算当天收盘时所有亏损筹码的账面亏损金额作为账面亏损总金额：
\begin{equation}
    \text{账面亏损}_{i,t} = \sum_p \mathrm{abs}(\mathrm{Close}_{i,t} - \mathrm{Price}_{i,p}) \times \mathrm{Chip}_{i,t,p} \times \mathrm{I}_{\mathrm{Close}_{i,t} < \mathrm{Price}_{i,p}}
\end{equation}

4. 计算当日兑现的亏损占筹码总亏损的比例，即为已实现亏损比率 \( \mathrm{CPLR} \)：
\begin{equation}
    \mathrm{CPLR}_{i,t} = \frac{\text{已实现亏损}_{i,t}}{\text{已实现亏损}_{i,t} + \text{账面亏损}_{i,t}}
\end{equation}

\subsection{说明}
已实现亏损比率衡量了投资者亏损时的卖出倾向；因子值越大，投资者越倾向于卖出亏损的股票；因子值越小，投资者越倾向于持有亏损的股票。

\subsection{参考文献}
\begin{enumerate}
    \item 陈升锐, 姚紫薇. 2025. 处置效应与 V 型处置效应在量化选股中的应用. 中信建投.
    \item Odean, T. (1998). \textit{Are investors reluctant to realize their losses?}. Journal of Finance, 53(5), 1775-1798.
\end{enumerate}

\section{量谷相对加权价格(高频因子-量价相关性类)}

\subsection{因子定义}

1. 对过去 20 日同时点成交量按照 1 倍标准差的标准划分，高于 1 倍标准差划分为“喷发成交量”，低于 1 倍标准差划分为“温和成交量”；

2. 对“喷发成交量”进一步按照连续与否划分：若当前分钟为“喷发成交量”，且前后 1 分钟均为“温和成交量”，则划为“孤立喷发成交量”，称为“量峰”；若当前分钟为“喷发成交量”，且前后 1 分钟也存在“喷发成交量”，则划为“连续喷发成交量”，称为“量岭”。将“温和成交量”称为“量谷”。

3. 计算如下因子的 20 日均值即为量谷相对加权价格因子：
\begin{equation}
    \text{量谷相对加权价格比} = \frac{\text{每日量谷的成交量加权价格}}{\text{同日总成交量加权价格}}
\end{equation}

\subsection{说明}
“量谷”时点的交易情绪较为低迷，其价格过度反应可能性较低，因此其价格相对水平越高，个股未来表现越好。

\subsection{参考文献}
\begin{enumerate}
    \item 魏建榕, 王志豪. 2025. 高频成交量的峰、岭、谷信息. 开源证券.
\end{enumerate}

\section{开盘后净主买占比(高频因子-资金流类)}

\subsection{因子定义}


\begin{equation}
    \text{开盘后净主买占比} = \frac{1}{T} \sum_{n=t}^{t-T+1} \frac{\sum_{j=1}^{N} \text{净主买成交额}_{i,j,n}}{\sum_{j=1}^{N} \text{成交额}_{i,j,n}}
\end{equation}

\begin{equation}
    \text{净主买成交额}_{i,j,n} = \text{主动买入成交额}_{i,j,n} - \text{主动卖出成交额}_{i,j,n}
\end{equation}

\subsection{变量定义}
\begin{itemize}
    \item ① 基于逐笔成交数据中的 BS 标志将逐笔成交数据合成为分钟级别的主动买卖金额，B 为主动买入（卖出方先挂单，买入方主动触碰卖单并成交）；S 为主动卖出（买入方先挂单，卖出方主动触碰买单并成交）；
    \item ② 剔除了处于涨跌停分钟上的主动买入金额和主动卖出金额数据；
    \item ③ \( i, j, n \) 分别表示第 \( i \) 只股票在第 \( n \) 个交易日内的第 \( j \) 分钟的成交额；
    \item ④ 开盘后的净主买占比用的是 9:30-10:00 范围内的数据；
    \item ⑤ 月度选股下，\( T=20 \) 个交易日；周度选股下，\( T=5 \) 个交易日。
\end{itemize}

\subsection{说明}
刻画了投资者在开盘后 30 分钟内净买入行为的强度，开盘后净主买占比越高，投资者的主动买入行为越强烈。

\subsection{参考文献}
\begin{enumerate}
    \item 冯佳睿, 袁林青. 2020. 基于主动买入行为的选股因子. 海通证券.
    \item 冯佳睿等. 2020. 高频因子的现实与幻想. 海通证券.
\end{enumerate}

\section{筹码分布变异系数(行为金融因子-筹码分布)}

\subsection{因子定义}
\begin{equation}
    \mathrm{coeff\_of\_var}_t = \frac{\mathrm{std}_t}{\mathrm{mean}_t}
\end{equation}

\subsection{变量定义}
\begin{itemize}
    \item ① 假设某只股票在 \( t \) 日的筹码分布共有 \( n \) 个价位（\( i = 1, 2, \ldots, n \)），\( \mathrm{vol}_{i,t} \) 为第 \( t \) 日 \( i \) 价位上的筹码峰高度，\( \mathrm{price}_{i,t} \) 为该筹码峰对应的价格；
    \item ② 筹码加权权重为 \( p_{i,t} = \dfrac{\mathrm{vol}_{i,t}}{\sum \mathrm{vol}_{i,t}} \)；
    \item ③ \( \mathrm{mean}_t = \sum (\mathrm{price}_{i,t} \times p_{i,t}) \) 为筹码分布加权平均成本；
    \item ④ \( \mathrm{std}_t = \sqrt{\sum (\mathrm{price}_{i,t} - \mathrm{mean}_t)^2 \times p_{i,t}} \) 为筹码成本加权标准差。
\end{itemize}

\subsection{说明}
筹码分布是基于历史成交量与成交价格推算出的“流通股票持仓成本分布”，描述了不同时点、不同价位上市场投资者的持仓情况，是对投资者实际持仓成本情况的近似估计；筹码分布变异系数衡量了筹码分布的离散程度；经测试，筹码分布变异系数与未来收益正相关。

\subsection{参考文献}
\begin{enumerate}
    \item 陈升锐, 姚紫薇. 2025. 筹码分布因子系统构建. 中信建投.
\end{enumerate}

\section{平均日内正跳跌收益(高频因子-波动跳跃类)}

\subsection{因子定义}

① 利用股票日内 5 分钟行情数据，计算 Jiang and Oomen (2008) 中的跳跃统计量 \( \mathrm{JS} \)，识别出股价发生跳跃的分钟时刻，具体识别细节见研报：
\begin{equation}
    \mathrm{JS} = N \frac{\hat{V}_{(0,1)}}{\sqrt{\hat{\Omega}_{\mathrm{SWV}}}} \left( 1 - \frac{\mathrm{RV}_N}{\mathrm{SwV}_N} \right) \sim N(0,1)
\end{equation}
\begin{equation}
    \mathrm{RV}_N = \sum_{k=1}^{N} r_k^2
\end{equation}
\begin{equation}
    \mathrm{SwV}_N = 2 \sum_{k=1}^{N} (R_k - r_k)
\end{equation}
\begin{equation}
    \hat{V}_{(0,1)} = \frac{1}{\mu_1^2} \sum_{k=1}^{N-1} |r_{k+1}| |r_k|
\end{equation}
\begin{equation}
    \hat{\Omega}_{\mathrm{SWV}} = \frac{\mu_6}{9} \cdot \frac{N^3 \mu_1^{-6}}{N - 5} \sum_{k=0}^{N-6} \prod_{l=1}^{6} |r_{k+l}|
\end{equation}
\begin{equation}
    \mu_p = 2^{p/2} \Gamma \left( \frac{p+1}{2} \right) / \sqrt{\pi}
\end{equation}

② 将所有发生价格跳跃且收益为正的时段的对数收益率累加即为当日的正跳跌收益因子值；

③ 以 10 个交易日为半衰期，将过去 20 个交易日的日内正跳跌收益指数加权平均，即为最终的平均日内正跳跌收益因子。

\subsection{说明}
股价跳跃体现了投资者对信息冲击的集中反应；股价跳跃收益类因子通常呈反转特性，与股票未来收益负相关。即过去跳跃收益较高的股票未来收益较小。若股票当日出现正跳跌又出现负跳跌时，全部累积会导致正跳跌与负跳跌相互抵消，低估股票每日的跳跃，可分别累加正跳跌时刻收益作为正跳跌收益因子，同理也可得负跳跌收益因子。此外还可只统计隔夜跳跃收益得到隔夜跳跃收益因子。

\subsection{参考文献}
\begin{enumerate}
    \item 任建, 杨帆. 2024. 如何识别股价跳跃. 招商证券.
    \item Jiang, G. J., \& Oomen, R. C. A. (2008). \textit{Testing for jumps when asset prices are observed with noise—a “swap variance” approach}. Journal of Econometrics, 144(2), 352-370.
    \item Jiang, G. J., \& Zhu, K. X. (2017). \textit{Information Shocks and Short-Term Market Underreaction}. Journal of Financial Economics, 124(1), 43-64.
\end{enumerate}

\section{累计跳跃绝对收益(高频因子-动量反转类)}

\subsection{因子定义}
\begin{equation}
    \mathrm{JAR} = e^{\sum_{t=1}^{D} \{|r_t| \cdot I_{\mathrm{jump\_SwV}, t}}\} - 1
\end{equation}
\begin{equation}
    I_{\mathrm{jump\_SwV}} = I(T_{\mathrm{SwV}} > \Phi^{-1}_{1-\alpha})
\end{equation}

\subsection{变量定义}
\begin{itemize}
    \item ① \( I_{\mathrm{jump\_SwV}} \) 为表示日内价格是否存在跳跃的示性函数；
    \item ② \( T_{\mathrm{SwV}} \) 为检验是否存在跳跃的检验统计量，具体计算逻辑见相关日历页；
    \item ③ \( r_t \) 为日度对数收益率；
    \item ④ 可以分开累计算发生跳跃时的正收益和负收益，也可以不对收益取绝对值。
\end{itemize}

\subsection{说明}
累计跳跃收益类因子是对存在显著跳跃交易日的收益的累加值，衡量了股价对信息冲击的反应。投资者对信息冲击的反应程度，短期内通常反应不足，长期内通常反应过度，对利好信息冲击容易反应过度，对坏信息冲击容易反应不足。

\subsection{参考文献}
\begin{enumerate}
    \item Jiang, G. J., \& Zhu, K. X. (2017). \textit{Information shocks and short-term market underreaction}. Journal of Financial Economics, 124(1), 43–64.
    \item 周飞鹏, 罗军, 安守宁. 2022. 再谈股价跳跃因子. 广发证券.
\end{enumerate}

\section{加权上引线频率(量价因子改进)}

\subsection{因子定义}

\begin{equation}
    \text{加权上影线频率}_{i,t} = \frac{\sum_{j=1}^{M_{i,t}} w_j \cdot \mathbb{1}_{\{ \text{上影线}_{i,j} > u \}}}{M_{i,t}}
\end{equation}

\begin{equation}
    \text{上影线}_{i,j} = \frac{\mathrm{High}_{i,j} - \max(\mathrm{Open}_{i,j}, \mathrm{Close}_{i,j})}{\mathrm{PrevClose}_{i,j}}
\end{equation}

\begin{equation}
    w_{\mathrm{decay},j} = 0.5^{\frac{t-j}{\lambda}}
\end{equation}

\subsection{变量定义}
\begin{itemize}
    \item ① \( M_{i,t} \) 为股票 \( i \) 在第 \( t \) 个回望期的天数，回望期可取 40 天；
    \item ② \( u \) 为上影线长度阈值，可取 1\%，只关注长上影线；
    \item ③ \( w_{\mathrm{decay},j} \) 为时间衰退指数加权权重，最近计数获得最大的权重，最远计数获得最小的权重；
    \item ④ \( \mathrm{High}_{i,j}, \mathrm{Open}_{i,j}, \mathrm{Close}_{i,j}, \mathrm{PrevClose}_{i,j} \) 分别为股票 \( i \) 在第 \( j \) 日的最高价、开盘价、收盘价和前收盘价。
\end{itemize}

\subsection{说明}
加权上影线频率为股票过去一段时间上影线长度大于一定阈值的天数加权之和，与股票未来收益负相关。

\subsection{参考文献}
\begin{enumerate}
    \item 郑琳琳, 盛宝丹. 2025. 加权影线频率与 K 线形态因子. 西南证券.
\end{enumerate}

\section{量峰分钟数(高频因子-成交分布类)}

\subsection{因子定义}

① 对过去 20 日同时点成交量按照 1 倍标准差的标准划分，高于 1 倍标准差划为“喷发成交量”，低于 1 倍标准差划为“温和成交量”；

② 对“喷发成交量”进一步按照连续与否划分：若当前分钟为“喷发成交量”，且前后 1 分钟均为“温和成交量”，则划为“孤立喷发成交量”，称为“量峰”；若当前分钟为“喷发成交量”，且前后 1 分钟也存在“喷发成交量”，则划为“连续喷发成交量”，称为“量岭”；将“温和成交量”称为“量谷”；

③ 统计过去 20 日量峰的分钟数，即为量峰分钟数因子。

\subsection{说明}
“量峰”对应的大额交易发生在情绪低迷处，更符合知情交易者的交易特征，可以通过统计发生“量峰”的分钟数量来衡量日内知情交易者参与的频率，因子值越高，知情交易参与度越高，未来收益也越高。

\subsection{参考文献}
\begin{enumerate}
    \item 魏建榕, 王志强. 2025. 高频成交量的峰、岭、谷信息. 开源证券.
\end{enumerate}

\section{每日平均异常成交(行为金融因子-投资者注意力)}

\subsection{因子定义}
\begin{equation}
    \mathrm{ABNVOLAVG} = \frac{1}{m} \sum_{t=1}^{m} \left( \frac{\mathrm{VOL}_{i,t}}{\overline{\mathrm{VOL}}_{i,t}} \right)
\end{equation}
其中，\(\overline{\mathrm{VOL}}_{i,t}\) 为股票 \( i \) 在交易日 \( t \) 滚动 1 年的平均成交量。

\subsection{变量定义}
\begin{itemize}
    \item ① \( \mathrm{VOL}_{i,t} \) 为股票 \( i \) 在交易日 \( t \) 的成交量；
    \item ② \( \overline{\mathrm{VOL}}_{i,t} \) 为股票 \( i \) 在交易日 \( t \) 滚动 1 年的平均成交量；
    \item ③ \( m \) 为交易日窗口，一般取为月度。
\end{itemize}

\subsection{说明}
每日平均异常成交为日频成交量除以过去一年平均成交量的均值，属于极端交易量类因子；异常成交量也是投资者注意力的常用代理指标，因子值越高，投资者对该股票的关注度越高；而有限注意力导致投资者更倾向于交易引起他们关注的股票，投资者关注度越高的股票通常意味着越多的投资购买；但 A 股市场的做多与做空不对称，使得投资者非理性买入高于非理性卖出，进而导致关注度高的股票由于投资者净买入而存在溢价，未来也更可能出现反转。

\subsection{参考文献}
\begin{enumerate}
    \item Cosemans, M., \& Frehen, R. (2021). \textit{Salience theory and stock prices: Empirical evidence}. Journal of Financial Economics, 140(2), 460-483.
    \item 陈升锐. 2024. 投资者有限关注及注意力捕捉与溢出. 中信建投.
\end{enumerate}

\section{主动买卖因子(高频因子-资金流类)}

\subsection{因子定义}

1. 逐日计算大单和中单总的主动买卖因子 \( \mathrm{ACT}_{\text{正向t}} \)，以及小单的主动买卖因子 \( \mathrm{ACT}_{\text{负向t}} \)：
\begin{equation}
    \mathrm{ACT}_{\text{正向t}} = \frac{\text{主动买入金额（大单+中单）} - \text{主动卖出金额（大单+中单）}}{\text{主动买入金额（大单+中单）} + \text{主动卖出金额（大单+中单）}}
\end{equation}
\begin{equation}
    \mathrm{ACT}_{\text{负向t}} = \frac{\text{主动买入金额（小单）} - \text{主动卖出金额（小单）}}{\text{主动买入金额（小单）} + \text{主动卖出金额（小单）}}
\end{equation}

2. 回溯过去 20 个交易日，取收益率最高 \( \lambda \) 比例的交易日，称为高收益日；取收益率最低 \( \lambda \) 比例的交易日，称为低收益日；

3. 对高收益日的 \( \mathrm{ACT}_{\text{正向t}} \) 因子取平均，记为 \( \mathrm{ACT}_{\text{正向}} \)；对低收益日的 \( \mathrm{ACT}_{\text{负向t}} \) 因子取平均，记为 \( \mathrm{ACT}_{\text{负向}} \)。

\subsection{说明}
切割后的主动买卖因子，细致的刻画了在市场上涨和下跌的不同情境下，不同交易者主动买卖意愿的差异；测试发现以大户和中户投资者为主的主动买卖因子，在高收益端呈现正向选股效应；而以小户为代表的主动买卖因子，在低收益端呈现负向选股效应。

\subsection{参考文献}
\begin{enumerate}
    \item 魏建榕. 2020. 主动买卖因子的正确用法. 开源证券.
\end{enumerate}

\section{基于 PB 的分析师锚效益偏差(行为金融因子-锚定效应)}

\subsection{因子定义}

\begin{equation}
    \mathrm{CAF\_BP}_{i,t} = \frac{\mathrm{FBP}_{i,t} - \mathrm{IFBP}_t}{\mathrm{IFBP}_t}
\end{equation}
\begin{equation}
    \mathrm{FBP}_{i,t} = \frac{1}{\mathrm{FPB}_{i,t}}
\end{equation}

\subsection{变量定义}
\begin{itemize}
    \item ① \( \mathrm{FPB}_{i,t} \) 为分析师在时间 \( t \) 对股票 \( i \) 的 PB 指标的预测值；
    \item ② \( \mathrm{IFBP}_t \) 为股票 \( i \) 所属行业中所有成分股的 \( \mathrm{FPB}_{i,t} \) 中位数。
\end{itemize}

\subsection{说明}
锚定效应指人们在作出判断时会将看到或听到的第一印象的影响作为锚点，并根据锚点对判断做调整，使得调整不充分偏向于该锚点而造成判断偏差的现象；假设分析师预测也会有锚定效应，即分析师以公司所在行业作为定价的“锚”，当预测的财务指标超出（低于）行业中位数时，分析师会做出悲观（乐观）预测的指标，以接近行业中位数，财报公布后，被预测的公司的财务指标将高于（低于）预测值，股价表现将更优（更差）。

\subsection{参考文献}
\begin{enumerate}
    \item Ashoura, S., \& Hao, Q. (2019). \textit{Do analysts really anchor? Evidence from credit risk and suppressed negative information}. Journal of Banking \& Finance, 98, 183-197.
    \item 陈升锐. 2023. 行为金融学在量化选股中的应用. 中信出版社.
\end{enumerate}

\section{成交量波峰计数因子(高频因子-成交分布类)}

\subsection{因子定义}

① 计算全部 240 根 K 线数据的均值与标准差，筛选出大于（均值 + 1 倍标准差）的分钟数据；

② 计算上一步筛选出的每一分钟与前一分钟之间的时间差，只有时间差超过 1 分钟（即相隔超过 1 根 K 线）时，该分钟数据才会被保留；

③ 经过上述两步筛选后，剩余的数据全部被视作局部峰值，峰值个数即为当日成交量波峰计数因子的取值，而最终调仓日的因子值仍为过去 20 个交易日因子值的均值。

\subsection{说明}
出现日内“放量”现象次数更多的股票可能有更多的趋势或知情交易者参与交易，进而在未来一段时间内有更高的收益率，因子与未来收益正相关。

\subsection{参考文献}
\begin{enumerate}
    \item 覃川桃, 郑起. 2020. 高频因子（九）：高频波动中的时间序列信息. 长江证券.
\end{enumerate}

\section{大单的动态知情交易概率(高频因子-资金流类)}

\subsection{因子定义}

构建下列自回归模型，计算股票 \( i \) 在日内区间 \( j \) 内的非预期收益率 \( \varepsilon_{i,j} \)：
\begin{equation}
    R_{i,j} = \gamma_0 + \sum_{k=1}^{4} \gamma_{1i,k} D_k^{\mathrm{Day}} + \sum_{k=1}^{48} \gamma_{2i,k} D_k^{\mathrm{Int}} + \sum_{k=1}^{12} \gamma_{3i,k} R_{i,j-k} + \varepsilon_{i,j}
\end{equation}

计算未预期收益率为正（负）时卖出（买入）交易量占比即为动态知情交易概率：
\begin{equation}
    \mathrm{DPIN}_{\mathrm{SIZE}}^{i,j} = \left[ \frac{\mathrm{NB}_{i,j}}{\mathrm{NT}_{i,j}} \cdot \mathbb{I}_{(\varepsilon_{i,j} < 0)} + \frac{\mathrm{NS}_{i,j}}{\mathrm{NT}_{i,j}} \cdot \mathbb{I}_{(\varepsilon_{i,j} > 0)} \right] \cdot \mathrm{LT}_{i,j}
\end{equation}

\subsection{变量定义}
\begin{itemize}
    \item ① \( R_{i,j} \) 为股票 \( i \) 在交易日 \( t \) 内区间 \( j \) 的收益率，\( R_{i,j-k} \) 为自回归的滞后收益率，滞后阶数 \( k=12 \)；
    \item ② \( D_k^{\mathrm{Day}} \) 为周内效应虚拟变量，当前交易日属于周 \( k \) 就取 1，其余取 0；同理，\( D_k^{\mathrm{Int}} \) 为日内效应虚拟变量，日内区间 \( j \) 的 \( D_j^{\mathrm{Int}} \) 取 1，其余取 0；
    \item ③ \( \mathrm{NB}_{i,j}, \mathrm{NS}_{i,j}, \mathrm{NT}_{i,j} \) 分别为股票 \( i \) 在日内区间 \( j \) 的主买成交笔数、主卖成交笔数、总成交笔数；
    \item ④ \( \mathbb{I}_{(\varepsilon_{i,j} < 0)} \) 为虚拟变量，未预期收益为负，取值为 1，否则取值为 0，反之亦然；
    \item ⑤ \( \mathrm{LT}_{i,j} \) 为大单的虚拟变量，若交易日 \( t \) 内，股票 \( i \) 在区间 \( j \) 内的总交易量超过当日各个区间内总交易量的中位数时，赋值为 1；
    \item ⑥ 自回归的时间窗口可以为从交易日 \( t \) 的区间 \( j \) 开始回溯过去 \( T \) 个交易日，如 5 分钟频率，则回溯过去 \( T \times 48 \) 个区间；
    \item ⑦ 可以进一步计算日度、周度内大单 DPIN 序列的均值、标准差等统计指标，刻画大单 DPIN 的时序特征。
\end{itemize}

\subsection{说明}
通常知情交易者更有可能进行大单交易，也就是说，大单交易区间内的反转交易更有可能来自知情交易者，所以可以进一步筛选出大单下的知情交易概率。

\subsection{参考文献}
\begin{enumerate}
    \item Chang, S. S., Chang, L. V., \& Wang, F. A. (2014). \textit{A dynamic intraday measure of the probability of informed trading and firm-specific return variation}. Journal of Empirical Finance, 29, 80-94.
\end{enumerate}

\section{衰减加权月频成交量(行为金融因子-投资者注意力)}

\subsection{因子定义}
\begin{equation}
    \mathrm{ATTN} = \frac{(12 - m) \cdot \mathrm{VOL}_{i, t}}{\sum_{m=1}^{11} m}
\end{equation}

\subsection{变量定义}
\begin{itemize}
    \item ① \( \mathrm{VOL}_{i, t} \) 为股票 \( i \) 在 \( t \) 日的成交量；
    \item ② \( m \) 为交易日窗口，一般取为月度。
\end{itemize}

\subsection{说明}
衰减加权月频成交量因子由股票成交量按时间线性衰减加权得到，属于极端交易量类因子；异常成交量也是投资者注意力的常用代理指标，因子值越高，投资者对该股票的关注度越高；而有限注意力导致投资者更倾向于交易引起他们关注的股票，投资者关注度越高的股票通常意味着越多的投资购买；但 A 股市场的做多与做空不对称，使得投资者非理性买入高于非理性卖出，进而导致关注度高的股票由于投资者净买入而存在溢价，未来也更可能出现反转。

\subsection{参考文献}
\begin{enumerate}
    \item Dong, D., Wu, K., Fang, J., Gozgor, G., \& Yan, C. (2021). \textit{Investor Attention Factors and Stock Returns: Evidence from China}. Journal of International Financial Markets, Institutions and Money, 77(4).
    \item 陈升锐. 2024. 投资者有限关注及注意力捕捉与溢出. 中信出版社.
\end{enumerate}

\section{绝对收益与调整后滞后成交量相关性(高频因子-量价相关性类)}

\subsection{因子定义}
\begin{equation}
    \mathrm{adj\_CORA\_A} = \mathrm{corr}(|\mathrm{Ret}_t|, \mathrm{adj\_Amount}_{t-1}), \quad \mathrm{Ret}_t \neq 0
\end{equation}
\begin{equation}
    \mathrm{adj\_Amount}_t = \frac{\mathrm{Amount}_t - \mu_t}{\sigma_t}
\end{equation}

\subsection{变量定义}
\begin{itemize}
    \item ① \( |\mathrm{Ret}_t| \) 为个股日内第 \( t \) 分钟的对数收益率的绝对值；
    \item ② \( \mathrm{adj\_Amount}_t \) 为个股日内第 \( t \) 分钟的调整后的成交额，\( \mu_t \) 和 \( \sigma_t \) 分别为前 20 个交易日相同第 \( t \) 分钟成交额的均值和标准差；
    \item ③ 月度因子可使用月末 5-15 个交易日因子的平均值来合成。
\end{itemize}

\subsection{说明}
分钟成交额在日内呈 U 型或 W 型分布，与收益率分布的不对等会影响相关性的计算。对成交额做标准化可以捕捉成交额的相对波动；若存在较多收益为 0 的样本也会使系数偏高，需将其剔除；修正后的因子能捕捉更真实纯粹的量价异动，因子表现大幅提升。

\subsection{参考文献}
\begin{enumerate}
    \item 朱定豪, 严佳炜. 2020. 量价关系的高频乐章. 方正证券.
\end{enumerate}

\section{非线性高频波动率(高频因子-波动跳跃类)}

\subsection{因子定义}

1. 在 Fama 三因子基础上加入流动性和反转因子，对过去 21 天个股收益率做线性时序回归，得到新特异率 \( = 1 - \) 拟合优度 \( R^2 \)：
\begin{equation}
    r_t = \alpha + \beta_{\mathrm{mkt}} \mathrm{MKT}_t + \beta_{\mathrm{smb}} \mathrm{SMB}_t + \beta_{\mathrm{hml}} \mathrm{HML}_t + \beta_{\mathrm{ret}} \mathrm{RET}_t + \beta_{\mathrm{liq}} \mathrm{LIQ}_t + \varepsilon_t
\end{equation}

2. 计算个股过去 21 天的日内分钟收益率序列（频率控制在 30 分钟内较好）的标准差，得到高频波动率 \( \mathrm{std}_i \)；再对该波动率进行横截面归一化和非线性变换：
\begin{equation}
    \mathrm{norm}_i = \frac{\mathrm{std}_i - \min(\{\mathrm{std}_i\})}{\max(\{\mathrm{std}_i\}) - \min(\{\mathrm{std}_i\})}
\end{equation}
\begin{equation}
    \mathrm{std}_i^{\mathrm{non\_linear}} = e^{\mathrm{norm}}
\end{equation}

3. 将上述非线性化后的波动率与新特异率开方相乘得到非线性波动率类因子：
\begin{equation}
    \text{特异波动率} = \sqrt{\text{新特异率}} \times \mathrm{std}_i^{\mathrm{non\_linear}}
\end{equation}

\subsection{说明}
原始的高频特异波动率因子因个股波动的头部区分能力较弱导致选股性较差，以非线性的方式改进个股波动部分乘数（加强头部波动率区分，淡化尾部波动率区分），可以保留其空头端的区分能力，同时淡化头部组分的分组信息，改进特异波动率因子的表现。

\subsection{参考文献}
\begin{enumerate}
    \item 覃川桃, 郑超. 2019. 高频因子(六): 特异视角下的波动率. 长江证券.
\end{enumerate}

\section{分析师共同覆盖间接关联动量(图谱网络-网络结构)}

\subsection{因子定义}

\begin{equation}
    \mathrm{CF2\_RET}_i = \frac{\sum_{j=1}^{N} m_{ij} \times \mathrm{Ret}_j}{\sum_{j=1}^{N} m_{ij}}
\end{equation}
\begin{equation}
    m_{ij} = \sum_{k=1}^{N} \log(n_{ik} + 1) \times \log(n_{kj} + 1)
\end{equation}

\subsection{变量定义}
\begin{itemize}
    \item ① \( m_{ij} \) 为股票 \( i \) 和股票 \( j \) 的间接关联强度；
    \item ② \( n_{ik}, n_{kj} \) 分别为连接股票 \( i \) 和股票 \( j \) 的中间股票 \( k \) 与双方的直接关联强度（共同覆盖的分析师数量）；
    \item ③ \( \mathrm{Ret}_j \) 为股票 \( j \) 过去 1 个月（20 个交易日）的收益率；
    \item ④ \( N \) 为股票池中的股票数量。
\end{itemize}

\subsection{说明}
间接关联较为隐蔽，不易被投资者察觉，也能带来领先滞后效应；对于直接覆盖的分析师数量较少的公司，间接关联带来的股价领先滞后效应可能更显著。

\subsection{参考文献}
\begin{enumerate}
    \item Ali, U., \& Hirshleifer, D. A. (2020). \textit{Shared analyst coverage: Unifying momentum spillover effects}. Journal of Financial Economics, 136(3), 649–675.
    \item 林晓明, 李子钰, 何康. 2022. 深挖分析师共同覆盖中的关联因子. 华泰证券.
\end{enumerate}

\section{最大涨幅(高频因子-动量反转类)}

\subsection{因子定义}
\begin{equation}
    \mathrm{MAX}_t = \prod_{i \in \mathrm{maxprod}} (1 + r_i)
\end{equation}

\subsection{变量定义}
\begin{itemize}
    \item ① \( \mathrm{maxprod} \) 为 \( t \) 日涨幅最大的前百分之十的分钟涨幅集合；
    \item ② \( r_i \) 为分钟涨跌幅。
\end{itemize}

\subsection{说明}
MAX 为当天前百分之十涨幅推动的股票具体涨幅。短期内，较大涨幅通常由大资金推动，这类“彩票型股票”也更得到投机性投资者的青睐，短期回报较高；长期内，随着投机者离场，股价又会下跌。

\subsection{参考文献}
\begin{enumerate}
    \item 陈升锐. 2023. 高频选股因子分类体系. 中信建投证券.
\end{enumerate}

\section{非流动性因子 Gammar(高频因子-流动性类)}

\subsection{因子定义}
在每月末，对股票 \( i \) 过去一个月的数据进行如下最小二乘回归：
\begin{equation}
    r^e_{i,d+1,t} = \theta_{i,t} + \phi_{i,t} r_{i,d,t} + \gamma_{i,t} \cdot \mathrm{sign}(r_{i,d,t}^e) \cdot v_{i,d,t} + \epsilon_{i,d+1,t}, \quad d = 1, \ldots, D
\end{equation}
\begin{equation}
    \mathrm{Gam}_{i,t} = -|\gamma_{i,t}|
\end{equation}

\subsection{变量定义}
\begin{itemize}
    \item ① \( r_{i,d,t} \) 为股票 \( i \) 在 \( t \) 月内 \( d \) 日收益率；
    \item ② \( r_{i,d,t}^e \) 为股票 \( i \) 在 \( t \) 月内 \( d \) 日的相对市场的超额收益率；
    \item ③ \( v_{i,d,t} \) 为股票 \( i \) 在 \( t \) 月内 \( d \) 日成交额；
    \item ④ 可将日度数据替换成分钟数据，计算得到高频的非流动性因子。
\end{itemize}

\subsection{说明}
Gam 更精细地刻画了成交额对股票收益的冲击；指标越小（绝对值越高），成交额对股票收益的冲击越大，股票流动性越差，股票未来预期收益越高。

\subsection{参考文献}
\begin{enumerate}
    \item Pastor, L., \& Stambaugh, R. F. (2003). \textit{Liquidity Risk and Expected Stock Returns}. Journal of Political Economy, 111(3), 642-685.
\end{enumerate}

\section{筹码分布成本宽带(行为金融因子-筹码分布)}

\subsection{因子定义}

\begin{equation}
    \mathrm{CBW}_t = \frac{\mathrm{Chip\_Highest}_t - \mathrm{Chip\_Lowest}_t}{\mathrm{Chip\_Lowest}_t}
\end{equation}

\subsection{变量定义}
\begin{itemize}
    \item ① 假设某只股票在 \( t \) 日的筹码分布共有 \( n \) 个价位（\( i=1,2,\ldots,n \)），\( \mathrm{vol}_{i,t} \) 为第 \( t \) 日 \( i \) 价位上的筹码峰高度，\( \mathrm{price}_{i,t} \) 为该筹码峰对应的价格，\( \mathrm{Chip\_Highest}_t \) 对应筹码最高价，\( \mathrm{Chip\_Lowest}_t \) 对应筹码最低价。
\end{itemize}

\subsection{说明}
筹码成本带宽表示当前市场上大多数流通股的买入成本区间宽度，即从最低持仓成本到最高持仓成本之间的价格范围，反映了投资者持股成本的集中或分散程度；因子值高，说明成本区间宽，筹码分散，投资者分歧较大；因子值低，说明成本区间窄，筹码集中；经测试，筹码分布成本宽带与未来收益负相关。

\subsection{参考文献}
\begin{enumerate}
    \item 陈升锐, 姚紫薇. 2025. 筹码分布因子系统构建. 中信建投.
\end{enumerate}

\section{成交量占比标准差(高频因子-成交分布类)}

\subsection{因子定义}
\begin{equation}
    \mathrm{std}\left(\frac{\mathrm{volume}_i}{\sum \mathrm{volume}_i}\right)
\end{equation}

\subsection{变量定义}
\begin{itemize}
    \item ① \( \mathrm{volume}_i \) 为日内分钟成交量，如60分钟等；成交量占比是成交量在个股内纵向可比指标，主要从成交的分布上给出衡量；
    \item ② 此外也可将成交量占比替换为对数成交量、换手率（全市场个股间横向可比指标，从流动性上给出衡量），计算成交量标准差类因子。
\end{itemize}

\subsection{说明}
成交量标准差类因子均衡量了个股过去一段时间交易活跃度的离散情况，与未来收益负相关；个股成交波动较大时，多空博弈激烈，异常交易显著，个股收益不确定性较强。

\subsection{参考文献}
\begin{enumerate}
    \item 覃川桃, 郑起. 2019. 高频因子(四): 高阶矩高频因子. 长江证券.
\end{enumerate}

\section{置信正态分布主动占比(高频因子-资金流类)}

\subsection{因子定义}


\begin{equation}
    \text{置信正态分布主动买入金额}_i = \mathrm{Amount}_i \times N\left(\frac{\mathrm{ret}_i}{0.1} \times 1.96\right)
\end{equation}

\begin{equation}
    \text{置信正态分布主动占比因子} = \frac{\sum_{i=1}^{N} \text{置信正态分布主动买入金额}_i}{\sum_{i=1}^{N} \mathrm{Amount}_i}
\end{equation}

\subsection{变量定义}
\begin{itemize}
    \item ① \( N(\cdot) \) 为标准正态分布累计函数；
    \item ② \( \mathrm{ret}_i \) 为每个时间段收益率，频率可为 1 分钟、5 分钟等。
\end{itemize}

\subsection{说明}
绝大部分时间段的成交几乎均有买卖双方各自驱动成交的部分，在计算主动买入金额时采用连续性值作为成交额所乘系数可以提高主买卖额估计准确度；此处收益率到主动买入占比的映射函数采用正态分布，且考虑了 A 股存在涨跌停限制，对收益率做了线性变换。

\subsection{参考文献}
\begin{enumerate}
    \item 覃川桃, 郑起. 2020. 高频因子(七): 分布估计下的主动成交占比. 长江证券.
\end{enumerate}

\section{极端跟随行为\_相关性因子(量价因子改进)}

\subsection{因子定义}

① 每个交易日 \( t \)，回看过去 \( k \) 个交易日 \( t-k, t-(k-1), \dots, t-1 \)，计算过去 \( k \) 日分钟成交量的 \( m\% \) 分位数，定义为成交量阈值，\( k=5, m=90 \)；将交易日 \( t \) 的每分钟成交量与上述阈值进行对比，若某一分钟的成交量大于上述阈值，则该分钟的交易被定义为趋势资金的交易；

② 对于趋势资金的每一分钟，取随后 \( n \) 分钟里成交额最大的一分钟，将其视为跟随趋势资金的极端行为，记录这一分钟的成交额，\( n=5 \)；

③ 每日将趋势资金的分钟成交额，与对应的极端跟随分钟成交额对齐，求这两个序列的相关系数；

④ 月底回看过去 20 个交易日，求 20 日相关系数的平均值，再做横截面市值、行业中性化处理。

\subsection{说明}
A股市场中“追涨”和“杀跌”的两股力量是非对称的，若某个股票上出现了强烈的羊群效应，通常对利好消息的反应过度现象会更加严重，股价未来下跌的概率也就更高，所以羊群效应通常会对股票未来收益产生负面影响。极端跟随行为\_相关性因子值越大，意味着极端跟随行为越剧烈，羊群效应越明显，从而导致股票未来收益越低。

\subsection{参考文献}
\begin{enumerate}
    \item 沈芷琦, 刘富兵, 阮俊烨. 2024. 基于趋势资金日内交易行为的选股因子. 国盛证券.
\end{enumerate}

\section{成交量“潮汐”的价格变动速率(高频因子-量价相关性类)}

\subsection{因子定义}

1. 对股票 \( i \)，计算 \( T \) 日内每分钟的领域成交量，即第 \( t \) 分钟及其前后 4 分钟（共计 9 分钟）的成交量总和；

2. 将领域成交量最高点所在分钟 \( t \) 记为“顶峰时刻”；将该时刻之前（第 5 到 \( t-1 \) 分钟）领域成交量最低点所在分钟 \( m \) 记为“涨潮时刻”，该分钟的领域成交量为 \( V_m \)，收盘价为 \( C_m \)；将该时刻之后（第 \( t+1 \) 到 233 分钟）领域成交量最低点所在分钟 \( n \) 记为“退潮时刻”，该分钟的领域成交量为 \( V_n \)，收盘价为 \( C_n \)；

3. 上述“涨潮时刻”～“退潮时刻”的全过程记为一次“潮汐”，计算全潮汐过程的价格变动速率：
\begin{equation}
    \text{全潮汐价格变动速率} = \frac{C_n - C_m}{C_m / (n-m)}
\end{equation}
对其计算最近 20 个交易日的均值即得到“全潮汐”因子。

\subsection{变量定义}
\begin{itemize}
    \item ① 计算因子需剔除开盘和收盘数据，仅考虑日内分钟频数据。
    \item ② \( V_m \) 和 \( V_n \) 大小可将“潮汐”过程进行强弱拆分，进一步构建“强势半潮汐”因子和“弱势半潮汐”因子。
\end{itemize}

\subsection{说明}
“全潮汐”因子衡量了投资者出售或购买股票意愿的强烈程度，与未来收益呈反比关系；在潮汐过程中，若股价快速下跌，表明部分投资者急于抛售，易导致过度反应，未来易发生补涨；若股价快速上涨，表明部分投资者急于建仓买入，同样易导致过度反应，未来也易发生补跌。

\subsection{参考文献}
\begin{enumerate}
    \item 曹春晓. 2022. 个股成交量的潮汐变化及“潮汐”因子构建. 方正证券.
\end{enumerate}

\section{已实现跳跃波动率(高频因子-波动跳跃类)}

\subsection{因子定义}

\begin{equation}
    \mathrm{RVJ}_t = \max\left(\mathrm{RV}_t - \hat{\mathrm{IV}}_t, 0\right)
\end{equation}

\begin{equation}
    \mathrm{RV}_t = \sum_{i=1}^{n} r_{i,n}^2 \rightarrow QV_t = \int_0^t \sigma_s^2 ds + \sum_{s \leq t} \Delta X_s^2 = \mathrm{IV}_t + \mathrm{QJ}_t
\end{equation}

\begin{equation}
    \hat{\mathrm{IV}}_t = \mu_{\frac{3}{2}}^{-3} \sum_{i=k}^{n} |r_{t_i}|^{\frac{2}{3}} |r_{t_{i-1}}|^{\frac{2}{3}} \cdots |r_{t_{i-k+1}}|^{\frac{2}{3}}
\end{equation}

\subsection{变量定义}
\begin{itemize}
    \item ① \( r_{t,i} \) 为某股票交易日 \( t \) 内的分钟对数收益率序列，如 5 分钟对数收益率序列；
    \item ② \( k=3 \)，\( n \) 表示该股票在交易日 \( t \) 中特定频率的分钟级别数据个数，如 5 分钟频率通常有 48 个数据点；
    \item ③ \( \mu_q \) 为标准正态分布的第 \( q \) 阶绝对矩；
    \item ④ 已实现跳跃波动率 \( \mathrm{RVJ}_t \) 整体取值为正，所以需要用 \( \max \) 对负差值做修正。
\end{itemize}

\subsection{说明}
已实现跳跃波动率刻画了股价波动率中股价跳跃部分的波动水平，可通过已实现波动剔除部分波动率估计得到；跳跃波动率衡量了股价在短时间内发生非连续性、剧烈变动（股价跳跃行为）的波动程度，跳跃行为能较好地动态反映市场情绪和信息冲击。

\subsection{参考文献}
\begin{enumerate}
    \item Barndorff-Nielsen, O. E., \& Shephard, N. (2004). \textit{Power and bipower variation with stochastic volatility and jumps}. Journal of Financial Econometrics, 2, 1-48.
    \item Andersen, T. G., Bollerslev, T., \& Diebold, F. X. (2007). \textit{Roughing it up: Including jump components in the measurement, modeling, and forecasting of return volatility}. Review of Economics and Statistics, 89, 701-720.
    \item Yu, B., Mizrach, B., \& Swanson, N. R. (2020). \textit{New Evidence of the Marginal Predictive Content of Small and Large Jumps in the Cross-Section}. Econometrics, MDPI, 8(2), 1-52.
\end{enumerate}

\section{未实现盈利量(行为金融因子-前景理论)}

\subsection{因子定义}

\begin{equation}
    \mathrm{CGO}_t = \frac{P_t - \mathrm{RP}_t}{\mathrm{RP}_t}
\end{equation}
\begin{equation}
    \mathrm{RP}_t = \frac{1}{k} \sum_{n=1}^{N} V_{t-n} \prod_{\tau=1}^{n-1} (1 - V_{t-\tau}) P_{t-n}
\end{equation}

\subsection{变量定义}
\begin{itemize}
    \item ① \( P_t \) 为 \( t \) 时刻的价格；
    \item ② \( V_t \) 为 \( t \) 时刻的换手率；
    \item ③ \( \dfrac{1}{k} \) 为权重归一化参数；
    \item ④ \( N \) 为截断时间长度，可以取 \( N=40, 60 \) 日。
\end{itemize}

\subsection{说明}
未实现盈利量衡量了当前投资者的盈亏水平，通过比较当前价格和锚定价格来确定。未实现盈利量通常与股票未来收益正相关，当投资者在盈利状态时，由于厌恶风险，更倾向于卖出股票，导致股价低估；在亏损状态下，由于偏好风险，更倾向于持有股票，导致股价高估。长期来看，上述现象在 A 股市场是成立的；但短期内，浮盈股票卖盘压力可能更大于被低估的幅度，而出现浮盈股票未来表现更差的情况。

\subsection{参考文献}
\begin{enumerate}
    \item Grinblatt, M., \& Han, B. (2005). \textit{Prospect theory, mental accounting, and momentum}. Journal of Financial Economics, 78(2), 311-339.
\end{enumerate}

\section{隔夜收益率(高频因子-动量反转类)}

\subsection{因子定义}
\begin{equation}
    \mathrm{OvernightRet}_{t,i} = \frac{\mathrm{open}_{t,i}}{\mathrm{close}_{t-1,i}} - 1
\end{equation}

\subsection{变量定义}
\begin{itemize}
    \item ① \( \mathrm{open}_{t,i} \) 和 \( \mathrm{close}_{t-1,i} \) 分别为 \( t \) 日开盘价和上一日 \( t-1 \) 收盘价；
    \item ② 若要计算周度、月度等频率的隔夜收益率，只需对区间内的日度隔夜涨跌幅求和。
\end{itemize}

\subsection{说明}
传统日度收益率可以划分为隔夜收益和日内收益，其中隔夜收益通常有一定程度的动量效应（对应日度收益率的动量）；隔夜收益与日内收益两者存在显著的负相关关系，对未来收益的预测方向也往往相反，许多研究认为正是双方的反转关系削弱了总收益的表现。

\subsection{参考文献}
\begin{enumerate}
    \item Branch, B., \& Ma, A. (2012). \textit{Overnight Return, the Invisible Hand Behind Intraday Returns}. Journal of Financial Markets, 2, 90-100.
    \item Lou, D., Polk, C., \& Skouras, S. (2019). \textit{A tug of war: Overnight versus intraday expected returns}. Journal of Financial Economics, 134(1), 192-213.
\end{enumerate}

\section{中间价变化率最大值(高频因子-流动性类)}

\subsection{因子定义}

\begin{equation}
    \mathrm{MPCmax}_{d,k} = \max \{ \mathrm{MPC}_{d,t,k}, \; t = k+1, \ldots, N_d \}
\end{equation}
\begin{equation}
    \mathrm{MPC}_{d,t,k} = \frac{M_{d,t} - M_{d,t-k}}{M_{d,t-k}}, \quad M_{d,t} = \frac{P_{d,t}^B + P_{d,t}^A}{2}
\end{equation}

\subsection{变量定义}
\begin{itemize}
    \item ① \( \mathrm{MPCmax}_{d,k} \) 为 \( d \) 日 \( \mathrm{MPC}_{t,k} \) 序列的最大值；
    \item ② \( M_{d,t} \) 为 \( d \) 日 \( t \) 时刻的市场中间价；
    \item ③ \( N_d \) 为 \( d \) 日的交易分钟数据。
\end{itemize}

\subsection{说明}
\( \mathrm{MPCmax} \) 描述了中间价的日内最大变动幅度，可用于捕捉主力市场操纵行为；长期看市场中间价的极端变动与股票收益显著负相关。

\subsection{参考文献}
\begin{enumerate}
    \item 丁鲁明, 陈升锐. 2021. 高频订单失衡及价差因子. 中信建投证券.
\end{enumerate}

\section{瓶颈公司(图谱网络-网络结构)}

\subsection{因子定义}
通过计算公司 \( i \) 在供应链网络中的中介中心性 Betweenness Centrality 来衡量公司在网络中的瓶颈（桥梁）地位：
\begin{equation}
    C_B(i) = \sum_{s \neq t \neq i \in V} \frac{\sigma_{st}(i)}{\sigma_{st}}
\end{equation}

\subsection{变量定义}
\begin{itemize}
    \item ① \( \sigma_{st} \) 为从供应商 \( s \) 到客户 \( t \) 的所有最短路径的数量；
    \item ② \( \sigma_{st}(i) \) 是从供应商 \( s \) 到客户 \( t \) 的所有最短路径中经过公司 \( i \) 的路径数量；
    \item ③ \( V \) 为图中节点的总数量。
\end{itemize}

\subsection{说明}
中介中心性衡量了节点在网络中作为最短路径“桥梁”的频率；做多中介中心性高的公司，做空中介中心性低的公司，能取得一定的超额收益。

\subsection{参考文献}
\begin{enumerate}
    \item Jussa et al. (2015). \textit{The Logistics of Supply Chain Alpha}. Deutsche Bank Markets Research.
\end{enumerate}

\section{买入浮亏偏离因子(高频因子-收益分布类)}

\subsection{因子定义}
\begin{equation}
    \mathrm{HCP} = \frac{\sum_{i=1}^{N} \overline{\mathrm{price}}_{\mathrm{buy},i} \cdot \mathbf{I}_{\{P_{\mathrm{buy},i} > \mathrm{close}\}}}{close} - 1
\end{equation}

\subsection{变量定义}

\begin{itemize}
    \item ① \( \mathrm{price}_{\mathrm{buy},i} \)、\( \overline{\mathrm{price}}_{\mathrm{buy}} \)、\( \mathrm{close} \) 分别为逐笔成交数据中买单的成交价、所有买单平均成交价、当日收盘价；
    \item ② 可以对该因子进行小单改进和尾盘改进，即只统计小单的占比和只统计 \([14:30,14:57)\) 尾盘区间的占比。
\end{itemize}

\subsection{说明}
买入浮亏偏离因子统计了高于收盘价买入的价格相较于收盘价的偏离幅度，与未来收益正相关；因子取值越高，投资浮亏越严重，根据遗憾规避理论，为了避免产生遗憾和后悔情绪，投资者存在惜售心理，未来抛压较低，带来更高的预期收益；遗憾规避理论认为非理性的投资者在做决策时，会倾向于避免产生后悔情绪并追求自豪感，避免承认之前的决策失误。

\subsection{参考文献}
\begin{enumerate}
    \item 高智威. 2023. 基于逐笔成交数据的遗憾规避因子. 国金证券.
\end{enumerate}

\section{一致买入交易(高频因子-成交分布类)}

\subsection{因子定义}

通过定义实体 K 线，将交易日中的每根 5 分钟 K 线划分为集体一致交易与非集体一致交易：
\begin{equation}
    |\mathrm{close} - \mathrm{open}| \leq \alpha |\mathrm{high} - \mathrm{low}|
\end{equation}

利用上涨的实体 K 线构造一致买入交易因子：
\begin{equation}
    \mathrm{PCV}_t = \frac{1}{d} \sum_{i=1}^{d} \left( \frac{\mathrm{ConsistentVolume}_{\mathrm{rise}}}{\mathrm{Volume}} \right)_{t-i}
\end{equation}

\subsection{变量定义}
\begin{itemize}
    \item ① \( \alpha \) 为一致参数，\( \alpha \) 越大，K 线一致性越强（上下引线越短，K 线越实体）；
    \item ② \( \mathrm{ConsistentVolume}_{\mathrm{rise}} \) 为上涨的 5 分钟实体 K 线的总成交量；
    \item ③ \( \mathrm{Volume} \) 为当日总成交量；
    \item ④ \( d \) 表示移动平均的周期参数。
\end{itemize}

\subsection{说明}
集体一致交易行为预示交易机会，如果一支股票在一段时间内大部分的交易属于集体一致交易行为，趋势性显著，则很可能这支股票在这段时间内有新信息进入，正在经历市场消化过程（price-in），未来股价很可能产生大幅变化，为交易者创造交易机会。

\subsection{参考文献}
\begin{enumerate}
    \item 刘均伟. 2018. 一致交易因子:挖掘集体行为背后的收益. 光大证券.
\end{enumerate}

\section{跳跃关联相对动量(图谱网络-动量溢出)}

\subsection{因子定义}

1. 利用股票日内 5 分钟行情数据，计算跳跃统计量 \( \mathrm{JS} \)，识别出股价发生跳跃的分钟时刻，具体计算细节见相关日历页：
\begin{equation}
    \mathrm{JS} = N \frac{\hat{V}_{(0,1)}}{\sqrt{\hat{\Omega}_{\mathrm{SWV}}}} \left( 1 - \frac{\mathrm{RV}_N}{\mathrm{SwV}_N} \right) \sim N(0,1)
\end{equation}

2. 计算关联跳跃次数：如果股票 \( i \) 在 \( t \) 日发生跳跃，股票 \( j \) 在 \( t-1 \) 日至 \( t+1 \) 日内也发生了一次同向跳跃，则记为一次关联跳跃；

3. 计算两只股票间的跳跃关联度，即在股票 \( i \) 所有的股价跳跃 \( \mathrm{jump\_num}_i \) 中关联跳跃 \( \mathrm{cojump\_num}_{i,j} \) 所占的比例：
\begin{equation}
    \mathrm{Corr}_{i,j} = \frac{\mathrm{cojump\_num}_{i,j}}{\mathrm{jump\_num}_i}
\end{equation}

4. 计算股票 \( i \) 在 \( t \) 时刻与之具有跳跃关联的其他股票过去一段时间的关联度加权收益（剔除关联度最低的 50\% 的关联关系），并对过去 20 天收益率 \( \mathrm{Ret}_{i,t} \) 做回归取残差即为跳跃关联相对动量：
\begin{equation}
    \mathrm{Peer\_Ret}_{i,t} = \frac{\sum_{j=1}^{N} \mathrm{Corr}_{i,j} \cdot \mathrm{Ret}_{j,t}}{\sum_{j=1}^{N} \mathrm{Corr}_{i,j}}, \quad i \neq j
\end{equation}


\subsection{说明}
具有相似属性的公司往往会受到类似的信息冲击，进而导致股价的同步跳跃。公司间股价跳跃的同步程度也可用于识别公司间的潜在关联程度；以股价跳跃关联度为权重加权关联股票的收益率可有效捕捉股票间的“领先-滞后”联动效应。

\subsection{参考文献}
\begin{enumerate}
    \item 任鹏, 刘凯, 杨航. 2025. 基于股价跳跃关联性的选股策略. 招商证券.
    \item 任鹏, 杨航. 2024. 如何识别股价跳跃. 招商证券.
\end{enumerate}

\section{量岭相对加权价格(高频因子-量价相关性类)}

\subsection{因子定义}


① 对过去 20 日同时点成交量按照 1 倍标准差的标准划分，高于 1 倍标准差划分为“喷发成交量”，低于 1 倍标准差划分为“温和成交量”；

② 对“喷发成交量”进一步按照连续与否划分：若当前分钟为“喷发成交量”，且前后 1 分钟均为“温和成交量”，则划为“孤立喷发成交量”，称为“量峰”；若当前分钟为“喷发成交量”，且前后 1 分钟也存在“喷发成交量”，则划为“连续喷发成交量”，称为“量岭”；将“温和成交量”称为“量谷”；

③ 计算如下因子的 20 日均值即为量岭相对加权价格因子：
\begin{equation}
    \text{量岭相对加权价格} = \frac{\text{每日量岭的成交量加权价格}}{\text{同日总成交量加权价格}}
\end{equation}

\subsection{说明}
量岭相对加权价格因子衡量了个人投资者交易的价格相对水平，因子值越高，越可能出现个人投资者交易的过度反应，未来收益越低。

\subsection{参考文献}
\begin{enumerate}
    \item 魏建榕, 王志豪. 2025. 高频成交量的峰、岭、谷信息. 开源证券.
\end{enumerate}

\section{大单买入强度(高频因子-资金流类)}

\subsection{因子定义}

\begin{equation}
    \frac{1}{T} \sum_{n=t}^{t-T+1} \frac{\mathrm{mean}(\text{大买单成交额}_{i,j,n})}{\mathrm{std}(\text{大买单成交额}_{i,j,n})}
\end{equation}

\subsection{变量定义}
\begin{itemize}
    \item ① 基于逐笔成交数据中的叫买与叫卖单号将逐笔成交数据合成为买卖单数据；
    \item ② 根据买卖单分布特征界定大小单，比如可将多日买卖单成交单数据对数调整后的“均值+1倍标准差”作为大单筛选阈值；
    \item ③ \( i, j, n \) 分别表示第 \( i \) 只股票在第 \( n \) 个交易日内的第 \( j \) 分钟的成交额；
    \item ④ 月度选股下，\( T = 20 \) 个交易日；周度选股下，\( T = 5 \) 个交易日。
\end{itemize}

\subsection{说明}
大单类因子刻画了大资金的交易行为，大资金往往具有信息优势，一般受到大资金关注的股票，未来通常具有更好的表现；大单买入强度因子同时衡量了大单买入强度和大单买入稳健性两方面的信息。

\subsection{参考文献}
\begin{enumerate}
    \item 冯佳睿, 袁林青. 2021. 大单的精细化处理与大单因子重构. 海通证券.
\end{enumerate}

\section{积极灵活筹码(行为金融因子-筹码分布)}

\subsection{因子定义}

\begin{equation}
    \frac{\sum_{i=T-20}^{T} \left( \text{主动买入且金额小于5万元的成交量} \right)_{i}}{ \text{当日成交量}_{i}}
\end{equation}


\subsection{说明}
积极灵活筹码占比越高，个股筹码中在各种市场表现下被持续持有的概率将高于卖出的概率，这在一定程度上影响了股票的多空力量，减小股票的卖出压力进而影响股票价格。

\subsection{参考文献}
\begin{enumerate}
    \item 任瞳, 许继宏. 2024. 基于逐笔数据的 AFH 改进因子. 招商证券.
\end{enumerate}

\section{平均日内正跳跌标准差(高频因子-波动跳跃类)}

\subsection{因子定义}

① 利用股票日内 5 分钟行情数据，计算 Jiang and Oomen (2008) 中的跳跃统计量 \( \mathrm{JS} \)，识别出股价发生跳跃的分钟时刻，具体识别细节见研报：
\begin{equation}
    \mathrm{JS} = N \frac{\hat{V}_{(0,1)}}{\sqrt{\hat{\Omega}_{\mathrm{SWV}}}} \left( 1 - \frac{\mathrm{RV}_N}{\mathrm{SwV}_N} \right) \sim N(0,1)
\end{equation}
\begin{equation}
    \mathrm{RV}_N = \sum_{k=1}^{N} r_k^2
\end{equation}
\begin{equation}
    \mathrm{SwV}_N = 2 \sum_{k=1}^{N} (R_k - r_k)
\end{equation}
\begin{equation}
    \hat{V}_{(0,1)} = \frac{1}{\mu_1^2} \sum_{k=1}^{N-1} |r_{k+1}| |r_k|
\end{equation}
\begin{equation}
    \hat{\Omega}_{\mathrm{SWV}} = \frac{\mu_6}{9} \cdot \frac{N^3 \mu_1^{-6}}{N - 5} \sum_{k=0}^{N-6} \prod_{l=1}^{6} |r_{k+l}|
\end{equation}
\begin{equation}
    \mu_p = 2^{p/2} \Gamma \left( \frac{p+1}{2} \right) / \sqrt{\pi}
\end{equation}

② 将所有发生价格跳跃且收益为正的时段的对数收益率累加即为当日的正跳跌收益因子值；

③ 以 10 个交易日为半衰期，将过去 20 个交易日的日内正跳跌收益求标准差即为最终的日内正跳跌标准差因子。

\subsection{说明}
股价跳跃体现了投资者对信息冲击的集中反应；股价跳跃收益类因子通常呈反转特性，与股票未来收益负相关，即过去跳跃收益较高的股票未来收益较小。

\subsection{参考文献}
\begin{enumerate}
    \item 任继敏, 杨帆. 2024. 如何识别股价跳跃. 招商证券.
    \item Jiang, G. J., \& Oomen, R. C. A. (2008). \textit{Testing for jumps when asset prices are observed with noise—a “swap variance” approach}. Journal of Econometrics, 144(2), 352-370.
    \item Jiang, G. J., \& Zhu, K. X. (2017). \textit{Information Shocks and Short-Term Market Underreaction}. Journal of Financial Economics, 124(1), 43-64.
\end{enumerate}

\section{理想反转(高频因子-动量反转类)}

\subsection{因子定义}


① 对单只股票，回溯取其过去 20 日的数据，计算每日日内逐笔成交金额分布的 13/16 分位值；

② 13/16 分位值高的 10 个交易日，涨跌幅加总，记作 \( M_{\text{high}} \)；
13/16 分位值低的 10 个交易日，涨跌幅加总，记作 \( M_{\text{low}} \)；

③ 计算理想反转因子：
\begin{equation}
    M = M_{\text{high}} - M_{\text{low}}
\end{equation}

\subsection{说明}
早期版本的理想反转因子是以“每日的平均单笔成交金额（成交金额/成交笔数）”作为切割标准；作者分析得出“反转之力的微观来源，是大单成交”，由高分位值较高（即大单成交较多）的那些交易日得到的反转因子，具有更强的反转特性。

\subsection{参考文献}
\begin{enumerate}
    \item 魏建榕. 2019. A 股反转之力的微观来源. 开源证券.
\end{enumerate}

\section{加权下引线频率(量价因子改进)}

\subsection{因子定义}

加权下影线频率：
\begin{equation}
    \mathrm{WeightedLowerShadowFreq}_{i,t} = \frac{\sum_{j=1}^{M_{i,t}} w_j \cdot \mathbb{1}_{\{ \text{下影线}_{i,j} > u \}}}{M_{i,t}}
\end{equation}

下影线：
\begin{equation}
    \text{下影线}_{i,j} = \frac{\min(\mathrm{Open}_{i,j}, \mathrm{Close}_{i,j}) - \mathrm{Low}_{i,j}}{\mathrm{PrevClose}_{i,j}}
\end{equation}

时间衰退权重：
\begin{equation}
    w_{\mathrm{decay},j} = 0.5^{\frac{t-j}{\lambda}}
\end{equation}

\subsection{变量定义}
\begin{itemize}
    \item ① \( M_{i,t} \) 为股票 \( i \) 在第 \( t \) 个回望期的天数，回望期可取 40 天；
    \item ② \( u \) 为下影线长度阈值，可取 1\%，只关注长下影线；
    \item ③ \( w_{\mathrm{decay},j} \) 为时间衰退指数加权权重，最近计数获得最大的权重，最远计数获得最小的权重；
    \item ④ \( \mathrm{Open}_{i,j}, \mathrm{Close}_{i,j}, \mathrm{Low}_{i,j}, \mathrm{PrevClose}_{i,j} \) 分别为股票 \( i \) 在第 \( j \) 日的开盘价、收盘价、最低价和前收盘价。
\end{itemize}

\subsection{说明}
加权下影线频率为股票过去一段时间下影线长度大于一定阈值的天数加权之和，与股票未来收益负相关。

\subsection{参考文献}
\begin{enumerate}
    \item 郑琳琳, 盛宝丹. 2025. 加权影线频率与 K 线形态因子. 西南证券.
\end{enumerate}

\section{量岭分钟数(高频因子-成交分布类)}

\subsection{因子定义}

1. 对过去 20 日同时点成交量按照 1 倍标准差的标准划分，高于 1 倍标准差划为“喷发成交量”，低于 1 倍标准差划为“温和成交量”；

2. 对“喷发成交量”进一步按照连续与否划分：若当前分钟为“喷发成交量”，且前后 1 分钟均为“温和成交量”，则划为“孤立喷发成交量”，称为“量峰”；若当前分钟为“喷发成交量”，且前后 1 分钟也存在“喷发成交量”，则划为“连续喷发成交量”，称为“量岭”；将“温和成交量”称为“量谷”；

3. 统计过去 20 日“量岭”的分钟数，即为量岭分钟数因子。

\subsection{说明}
“量岭”对应的连续性大额交易，更符合个人投资者的跟随交易特征，可以通过统计发生“量岭”的分钟数量来衡量个人投资者参与度，因子值越高，个人投资者参与度越高，未来收益越低。

\subsection{参考文献}
\begin{enumerate}
    \item 魏建榕, 王志豪. 2025. 高频成交量的峰、岭、谷信息. 开源证券.
\end{enumerate}

\section{基于跳跃幅度关联的点度中心性(图谱网络-网络结构)}

\subsection{因子定义}


① 利用股票日内 5 分钟行情数据，计算跳跃统计量 \( \mathrm{JS} \)，识别出股价发生跳跃的分钟时刻，具体计算细节见相关日历页：
\begin{equation}
    \mathrm{JS} = N \frac{\hat{V}_{(0,1)}}{\sqrt{\hat{\Omega}_{\mathrm{SWV}}}} \left( 1 - \frac{\mathrm{RV}_N}{\mathrm{SwV}_N} \right) \sim N(0,1)
\end{equation}

② 计算关联跳跃次数：如果股票 \( i \) 在 \( t \) 日发生跳跃，股票 \( j \) 在 \( t+1 \) 日内也发生了一次同向跳跃，则记为一次关联跳跃；

③ 计算两只股票间的跳跃关联度，\( \mathrm{cojump\_size}_{i,j} \) 为过去一段时间内股票 \( i \) 的所有跳跃中与股票 \( j \) 相关联的跳跃收益绝对值之和，\( \mathrm{jump\_size}_i \) 为过去一段时间股票 \( i \) 发生跳跃的交易日跳跃收益绝对值的和：
\begin{equation}
    \mathrm{Corr}_{i,j} = \frac{\mathrm{cojump\_size}_{i,j}}{\mathrm{jump\_size}_i}
\end{equation}

④ 对网络做稀释处理，剔除整个网络中关联度最低的 50\% 的关联关系，统计每只股票与之直接相连的股票数量即为该股票的点度中心性值。

\subsection{说明}
具有相似属性的公司往往会受到类似的信息冲击，进而导致股价的同步跳跃。公司间股价跳跃的同步程度也可用于识别公司间的潜在关联程度；跳跃关联网络的点度中心性因子衡量了股票在整个网络中的重要性或影响程度，与未来收益正相关。

\subsection{参考文献}
\begin{enumerate}
    \item 任晓, 刘凯, 杨航. 2025. 基于股价跳跃关联性的选股策略. 招商证券.
    \item 任晓, 杨航. 2024. 如何识别股价跳跃. 招商证券.
\end{enumerate}

\section{残差资金流强度因子(高频因子-资金流类)}

\subsection{因子定义}

计算小单资金流强度：分子为（小单买额-小单卖额）之和，分母为其绝对值；计算大单资金流强度：分子为（大单买额-大单卖额）之和，分母为其绝对值：
\begin{equation}
    S_t = \frac{\sum_{k=t-T}^{t} (\mathrm{buy}_k - \mathrm{sell}_k)}{\sum_{k=t-T}^{t} |\mathrm{buy}_k - \mathrm{sell}_k|}
\end{equation}

分别将大（小）单资金流强度与过去20日涨跌幅做回归，得到残差，即为大（小）单残差资金流因子：
\begin{equation}
    S_t = a + b \cdot \mathrm{Ret20}_t + \varepsilon_t
\end{equation}
\begin{equation}
    \text{残差资金流因子} = \varepsilon_t
\end{equation}

\subsection{变量定义}
\begin{itemize}
    \item ① 将主动大单净流入和主动中单净流入相加，计算主动大中单残差，再和非主动大单强度排序相加，可以改进原始大单残差资金流因子的表现；
    \item ② 使用非主动小单净流入计算非主动小单残差，可以提升原始小单残差资金流因子表现。
\end{itemize}

\subsection{说明}
残差资金流因子剥离了涨跌幅对资金流强度的影响，衡量的是在同等涨跌幅（\(\mathrm{Ret20}\)）的情况下资金流强度的选股能力。

\subsection{参考文献}
\begin{enumerate}
    \item 魏建梅, 高鹏. 2021. 大单与小单资金流的 alpha 能力. 开源证券.
    \item 魏建梅, 盛少成. 2022. 大单资金流的 alpha 探究. 交易精选与高频测算. 开源证券.
\end{enumerate}

\subsection{因子定义}
\section{成交量占比偏度()}

\subsection{因子定义}
\begin{equation}
    \mathrm{skew}\left(\frac{\mathrm{volume}_i}{\sum \mathrm{volume}_i}\right)
\end{equation}

\subsection{变量定义}
\begin{itemize}
    \item ① \( \mathrm{volume}_i \) 为日内分钟成交量，如 10 分钟等；成交量占比是成交量在个股内纵向可比指标，主要从成交的分布上给出衡量；
    \item ② 此外也可将成交量占比替换为对数成交量、换手率（全市场个股间横向可比指标，从流动性上给出衡量），计算成交量偏度类因子。
\end{itemize}

\subsection{说明}
成交量偏度类因子均衡量了个股过去一段时间的交易异常行为，与未来收益负相关；成交量呈右偏时，大部分时间段个股交易并不活跃，但存在特定区间成交量骤升的情况，而这类交易异常放量的个股，多头方存在短期筹码调整行为，价格被高估，之后回落至正常水平。

\subsection{参考文献}
\begin{enumerate}
    \item 覃川桃, 郑起. 2019. 高频因子(四): 高阶矩高频因子. 长江证券.
\end{enumerate}

\section{交易量变异系数(高频因子-资金流类)}

\subsection{因子定义}
\begin{equation}
    \mathrm{VCV}_{i,\gamma} = \frac{\hat{\sigma}_{V\{i,t \in \gamma\}}}{\hat{\mu}_{V\{i,t \in \gamma\}}}
\end{equation}

\subsection{变量定义}
\begin{itemize}
    \item ① \( \hat{\mu}_V \) 和 \( \hat{\sigma}_V \) 分别为过去 \( \gamma \) 区间内（如 1 周）的 5 分钟或 1 分钟成交额序列的均值和标准差；
    \item ② 也可用成交量、换手率等指标来计算成交量变异系数。
\end{itemize}

\subsection{说明}
VCV 衡量了市场的信息不对称性，市场总会存在知情交易者利用私有信息获取超额收益。VCV 与知情交易比例严格正相关；VCV 为负向因子，取值越大，知情交易者占比越高，股价会偏离真实价值，也更不会被大资金看好，未来更易利空。

\subsection{参考文献}
\begin{enumerate}
    \item Lof, M., \& van Bommel, J. (2022). \textit{Asymmetric information and the distribution of trading volume}. Available at: https://ssrn.com/abstract=4081052.
    \item 任建, 崔浩瀚. 2023. 再观知情交易者的踪迹. 招商证券.
\end{enumerate}

\section{筹码盈亏系数修正后的积极灵活筹码(行为金融因子-筹码分布)}

\subsection{因子定义}
基于最新收盘价计算过去买入筹码的涨跌幅作为筹码盈亏系数：
\begin{equation}
    \alpha_i = \frac{\mathrm{Close}_t}{\mathrm{TradePrice}_i}
\end{equation}

将筹码盈亏系数与筹码买入量相乘即可得到修正后的积极灵活筹码值：
\begin{equation}
    \mathrm{AFH\_Return} = \frac{\sum_{i=1}^{N} \alpha_i \times \mathrm{Vol\_AFH}_i}{\sum_{k=T-20}^{T} \mathrm{Volume}_k}
\end{equation}

\subsection{变量定义}
\begin{itemize}
    \item ① \( i \) 为筹码购入时的订单号，统计对象为过去 20 日的订单号；
    \item ② \( \mathrm{Vol\_AFH}_i \) 为主动买入时订单金额小于 5 万元的订单成交量；
    \item ③ \( \mathrm{TradePrice}_i \) 为该订单的买入价格；
    \item ④ \( \mathrm{Close}_t \) 为截止日的收盘价；
    \item ⑤ \( \mathrm{Volume}_k \) 为 \( k \) 日当天的总成交量。
\end{itemize}

\subsection{说明}
在处置效应的影响下，当投资者持有的股票盈利时，投资者出于对风险的厌恶，投资行为会更加谨慎，追求“落袋为安”，表现为卖出盈利的股票；而当投资者亏损时，投资者的行为更加激进，表现为将亏损的股票继续持有；可以根据最新收盘价计算过去买入筹码的盈亏状态对筹码量进行调整，提升原始积极灵活筹码的因子表现。

\subsection{参考文献}
\begin{enumerate}
    \item 任瞳, 许继宏. 2024. 基于逐笔数据的 AFH 改进因子. 招商证券.
\end{enumerate}

\section{滞后绝对收益与调整后成交量相关性(高频因子-量价相关性类)}

\subsection{因子定义}


\begin{equation}
    \mathrm{adj\_CORA\_R} = \mathrm{corr}(|\mathrm{Ret}_{t-1}|, \mathrm{adjAmount}_t), \quad \mathrm{Ret}_{t-1} \neq 0
\end{equation}
\begin{equation}
    \mathrm{adjAmount}_t = \frac{\mathrm{Amount}_t - \mu_t}{\sigma_t}
\end{equation}

\subsection{变量定义}
\begin{itemize}
    \item ① \( |\mathrm{Ret}_{t-1}| \) 为个股日内第 \( t-1 \) 分钟的对数收益率的绝对值，相比成交额滞后一期；
    \item ② \( \mathrm{adjAmount}_t \) 为个股日内第 \( t \) 分钟的调整后的成交额，\( \mu_t \) 和 \( \sigma_t \) 分别为前 20 个交易日相同第 \( t \) 分钟成交额的均值和标准差；
    \item ③ 月度因子可使用月末 5-15 个交易日因子的平均值来合成。
\end{itemize}

\subsection{说明}
分钟成交额在日内呈 U 型或 W 型分布，与收益率分布的不对等会影响相关性的计算，对成交额做标准化可以捕捉成交额的相对波动；若存在较多收益为 0 的样本也会使系数偏高，需将其剔除；修正后的因子能捕捉更真实纯粹的量价异动，因子表现大幅提升。

\subsection{参考文献}
\begin{enumerate}
    \item 朱定豪, 严佳炜. 2020. 量价关系的高频乐章. 方正证券.
\end{enumerate}

\section{高频成交量波动(高频因子-波动跳跃类)}

\subsection{因子定义}
\begin{equation}
    \sigma_{\mathrm{day}} = \mathrm{std}(\{\mathrm{vol}_i\})
\end{equation}
\begin{equation}
    \sigma_{\mathrm{vol}} = \frac{\mathrm{std}(\{\sigma_{\mathrm{day}}\})}{\mathrm{mean}(\{\sigma_{\mathrm{day}}\})}
\end{equation}

\subsection{变量定义}
\begin{itemize}
    \item ① \( \mathrm{vol}_i \) 为个股日内 240 根 1 分钟 K 线的成交量，\( \sigma_{\mathrm{day}} \) 取过去 20 个交易日数据；
    \item ② 除以均值是为了剔除量纲的影响；
    \item ③ 可以将分钟成交量替换成成交笔数、高频振幅等计算相应的高频波动因子。
\end{itemize}

\subsection{说明}
上述高频波动因子衡量了“波动的波动”，其优势在于能够同时刻画日内波动与日间波动；传统波动因子多以日度收益率为基础构建，除股价外，也可用成交量、成交笔数、振幅等指标来度量波动，因为在非理性驱动下这些指标也会呈现与股价相对一致的波动特点；样本区间内，这三种高频波动因子与传统日度波动因子相比，多头组收益更强，信息比和多空夏普比更高。

\subsection{参考文献}
\begin{enumerate}
    \item 覃川桃, 郑起. 2020. 高频因子（九）：高频波动中的时间序列信息. 长江证券.
\end{enumerate}

\section{客户重要性(图谱网络-网络结构)}

\subsection{因子定义}
利用 PageRank 算法计算客户重要性：
\begin{equation}
    p(i) = \frac{q}{V} + (1 - q) \sum_{j \in V: j \to i} \frac{p(j)}{k_j^{\mathrm{out}}}
\end{equation}

\subsection{变量定义}
\begin{itemize}
    \item ① \( p(i) \) 为客户 \( i \) 的 PageRank 得分；
    \item ② \( k_j^{\mathrm{out}} \) 是供应商 \( j \) 的出度，即从节点 \( j \) 发出的链接数量；
    \item ③ \( q \) 为阻尼因子，通常取值为 0.85；
    \item ④ \( j \in V: j \to i \) 为指向公司 \( i \) 的所有供应商节点的集合；
    \item ⑤ \( V \) 为图中节点的总数量。
\end{itemize}

\subsection{说明}
衡量客户重要性的 PageRank 得分相比客户节点度（入度），除了考虑指向客户的供应商数量外，还考虑了相连供应商本身的重要性，即使两个客户的入度相同，拥有更多“重要”供应商的客户也会有更高的 PageRank 得分。

\subsection{参考文献}
\begin{enumerate}
    \item Jussa et al. (2015). \textit{The Logistics of Supply Chain Alpha}. Deutsche Bank Markets Research.
\end{enumerate}

\section{耀眼收益率(高频因子-动量反转类)}

\subsection{因子定义}

1. 对股票 \( i \)，计算 t 日内 1 分钟成交量序列的环比差值，将成交量差值大于“该日差值序列均值+1倍标准差”的时刻定义为“成交量激增时刻 \( n \)”；

2. 利用分钟收盘价，计算“成交量激增时刻 \( n \)”的分钟收益率，即为“耀眼收益率”；

3. 将 t 日内所有耀眼收益率求均值即为“日耀眼收益率”；

4. 将各股票的“日耀眼收益率”减去该指标的横截面均值再取绝对值，即为“适度日耀眼收益率”。

\subsection{变量定义}
\begin{itemize}
    \item ① 计算因子需剔除开盘和收盘数据，仅考虑日内分钟频数据；
    \item ② 可对最近 20 个交易日的日度“适度日耀眼收益率”求均值和标准差，并等权合成，即可得“月耀眼收益率”因子。
\end{itemize}

\subsection{说明}
“耀眼收益率”衡量了日内成交量激增所引起的价格变动幅度；“耀眼收益率”与未来收益负相关，若成交量激增后的价格变动较小，说明利好消息未被市场广泛认可，适度承担该风险能获得风险补偿；若引起的价格变动较大，可能代表投资者已反应过度，未来会补跌。

\subsection{参考文献}
\begin{enumerate}
    \item 曹春晓. 2022. 成交量激增时刻蕴含的 alpha 信息. 方正证券.
\end{enumerate}

\section{盘口价差(高频因子-流动性类)}

\subsection{因子定义}
\begin{equation}
    \mathrm{spread} = \frac{2 \times (a_1 - b_1)}{a_1 + b_1}
\end{equation}

\subsection{变量定义}
\begin{itemize}
    \item ① \( a_1 \) 和 \( b_1 \) 分别为盘口委托快照数据的卖一价和买一价；
    \item ② 可对上述因子日内序列求均值得到日度因子值，再进行 20 个交易日指数移动平均得到月度因子值。
\end{itemize}

\subsection{说明}
盘口价差反映了卖一价和买一价的距离，与未来收益正相关，盘口价差越大，交易成本越高，市场宽度越大，流动性越差，而流动性较低的个股未来收益表现相对较好。也可对日内盘口价差求标准差，衡量价差的日内均匀程度。

\subsection{参考文献}
\begin{enumerate}
    \item 胡骏聪等. 2021. 高频视角下的微观流动性与波动性. 中金公司.
\end{enumerate}

\section{改进的分析师共同覆盖间接关联动量(图谱网络-动量溢出)}

\subsection{因子定义}
\begin{equation}
    \mathrm{ADJ\_CF2\_RET}_i = \frac{\sum_{j=1}^{N} m_{ij} \times \mathrm{Ret}_j}{\sum_{j=1}^{N} m_{ij}}
\end{equation}
\begin{equation}
    m_{ij} = \frac{\sum_{k=1}^{N} \log(n_{ik} + 1) \times \log(n_{kj} + 1)}{\log(\mathrm{num}_i + 1)}
\end{equation}

\subsection{变量定义}
\begin{itemize}
    \item ① \( m_{ij} \) 为股票 \( i \) 和股票 \( j \) 的间接关联强度；
    \item ② \( n_{ik}, n_{kj} \) 分别为连接股票 \( i \) 和股票 \( j \) 的中间股票 \( k \) 与双方的直接关联强度（共同覆盖的分析师数量）；
    \item ③ \( \mathrm{Ret}_j \) 为股票 \( j \) 过去 1 个月（20 个交易日）的收益率；
    \item ④ \( \mathrm{num}_i \) 为与股票 \( i \) 直接关联的公司数量；
    \item ⑤ \( N \) 为股票池中的股票数量。
\end{itemize}

\subsection{说明}
若一家公司有很多直接关联公司，市场难以迅速整合所有关联公司的信息，直接关联相比间接关联具有更显著的领先滞后效应；若一家公司的直接关联公司很少，信息传播效率较高，直接关联公司的信息会快速反映到股价上，此时进一步挖掘间接关联信息可能更有效；可以使用直接关联数量对间接关联做改进。

\subsection{参考文献}
\begin{enumerate}
    \item Ali, U., \& Hirshleifer, D. A. (2020). \textit{Shared analyst coverage: Unifying momentum spillover effects}. Journal of Financial Economics, 136(3), 649-675.
    \item 林晓明, 李子钰, 何康. 2022. 深挖分析师共同覆盖中的关联因子. 华泰证券.
\end{enumerate}

\section{均匀分布主动占比(高频因子-资金流类)}

\subsection{因子定义}

\begin{equation}
    \text{均匀分布主动买入金额}_i = \mathrm{Amount}_i \times \frac{\mathrm{ret}_i - 0.1}{0.2}
\end{equation}

\begin{equation}
    \text{均匀分布主动占比因子} = \frac{\sum_{i=1}^{N} \text{均匀分布主动买入金额}_i}{\sum_{i=1}^{N} \mathrm{Amount}_i}
\end{equation}

\subsection{变量定义}
\begin{itemize}
    \item ① \( \mathrm{ret}_i \) 为每个时间段收益率，频率可为 1 分钟、5 分钟等。
\end{itemize}

\subsection{说明}
绝大部分时间段的成交几乎均有买卖双方各自驱动成交的部分，在计算主动买入金额时采用连续性值作为成交额所乘系数可以提高主买卖额估计准确度；此处收益率到主动买入占比的映射函数采用均匀分布（对应价格变动对主动买卖力量均匀影响的情况），考虑到 A 股存在涨跌停限制，个股价格变动幅度一般在 -0.1 到 0.1 之间，对原收益率做了从 -0.1 至 0.1 到 0 至 1 之间的线性变换。

\subsection{参考文献}
\begin{enumerate}
    \item 覃川桃, 郑起. 2020. 高频因子(七): 分布估计下的主动成交占比. 长江证券.
\end{enumerate}

\section{一致卖出交易(高频因子-成交分布类)}

\subsection{因子定义}

通过定义实体 K 线，将交易日中的每根 5 分钟 K 线划分为集体一致交易与非集体一致交易：
\begin{equation}
    |\mathrm{close} - \mathrm{open}| \leq \alpha |\mathrm{high} - \mathrm{low}|
\end{equation}

利用下跌的实体 K 线构造一致卖出交易因子：
\begin{equation}
    \mathrm{NCV}_t = \frac{1}{d} \sum_{i=1}^{d} \left( \frac{\mathrm{ConsistentVolume}_{\mathrm{fall}}}{\mathrm{Volume}} \right)_{t-i}
\end{equation}

\subsection{变量定义}
\begin{itemize}
    \item ① \( \alpha \) 为一致参数，\( \alpha \) 越大，K 线一致性越强（上下引线越短，K 线越实体）；
    \item ② \( \mathrm{ConsistentVolume}_{\mathrm{fall}} \) 为下跌的 5 分钟实体 K 线的总成交量；
    \item ③ \( \mathrm{Volume} \) 为当日总成交量；
    \item ④ \( d \) 表示移动平均的周期参数。
\end{itemize}

\subsection{说明}
集体一致交易行为预示交易机会，如果一支股票在一段时间内大部分的交易属于集体一致交易行为，趋势性显著，则很可能这支股票在这段时间内有新信息进入，正在经历市场消化过程（price-in），未来股价很可能产生大幅变化，为交易者创造交易机会。

\subsection{参考文献}
\begin{enumerate}
    \item 刘均伟. 2018. 一致交易因子:挖掘集体行为背后的收益. 光大证券.
\end{enumerate}

\section{K 线形态(量价因子改进)}

\subsection{因子定义}

① 用过去 3 年（720 个交易日）所有股票的 2K 形态（每两个连续交易日 K 线构成一个 2K 形态）和该形态未来 20 个交易日的收益率作为 K 线形态表现打分的样本，形态 \( p \)（\( p \) 属于 768 种形态集合）的分值为：
\begin{equation}
    wr_p = \frac{W(r_p > 0)}{N(r_p)}
\end{equation}
\begin{equation}
    r_p = \frac{\mathrm{close}_{p+20}}{\mathrm{open}_{p+1}} - 1
\end{equation}

② 用形态胜率 \( wr_p \) 对股票最近 40 个交易日 K 线形态进行打分加总，采用衰退指数进行加权，得到形态因子 \( \mathrm{KP} \)：
\begin{equation}
    \mathrm{KP}_{i,T} = \sum_{t=T-40}^{T} w_t \cdot wr_{i,t}
\end{equation}
\begin{equation}
    wr_{i,t} = wr_p(\mathrm{pattern}_{i,t} = \mathrm{pattern}_p \mid p \in \{1,2,3,\ldots,768\})
\end{equation}
\begin{equation}
    w_t = 0.5^{\frac{T-t}{\lambda}}
\end{equation}

\subsection{变量定义}
\begin{itemize}
    \item ① \( r_p \) 为形态 \( p \) 出现后未来 20 个交易日的收益率；
    \item ② \( wr_p \) 为窗口期内形态 \( p \) 的胜率，\( N \) 为形态 \( p \) 出现的次数，\( W \) 为形态 \( p \) 样本中未来收益大于 0 的次数；
    \item ③ \( w_t \) 为时间衰退指数加权权重；
    \item ④ 768 种 2K 形态组合方式详见参考文献。
\end{itemize}

\subsection{说明}
K 线形态因子为股票最近一段时间 K 线形态得分加权之和，窗口期内测算各种 K 线形态的胜率，用以度量股票当前 K 线形态表现，与股票未来收益正相关。

\subsection{参考文献}
\begin{enumerate}
    \item 郑琳琳,盛宝丹.2025.加权影线与k线形态因子.西南证券
\end{enumerate}

\section{模糊关联度(高频因子-量价相关性类)}

\subsection{因子定义}

1. 对股票 \( i \) 计算 \( n \) 日内的 1 分钟收益率；计算第 \( t \) 分钟的波动率，即 \([t-4, t]\) 区间内分钟收益率的标准差；计算第 \( t \) 分钟的模糊性，即 \([t-4, t]\) 区间内上述分钟波动率的标准差；

2. 每日计算上述分钟模糊性序列与分钟成交额序列的相关系数，即为“模糊关联度”因子：
\begin{equation}
    \mathrm{FuzzyCorr}_d = \mathrm{corr}(\mathrm{Ambiguity}_t, \mathrm{Amount}_t)
\end{equation}

\subsection{变量定义}
\begin{itemize}
    \item ① 计算上述因子时，剔除开盘和收盘数据，仅保留日内交易数据，所以开盘前 9 分钟的模糊性为空值；
    \item ② 每月末可对最近 20 天的上述因子求均值和标准差并将两者等权合并，即得到月度“模糊关联度”因子。
\end{itemize}

\subsection{说明}
模糊关联度因子衡量了投资者成交金额随波动率模糊性变化而变化的程度，即从成交额维度衡量了投资者对模糊性的厌恶程度，与未来收益负相关；投资者对波动率模糊性的厌恶心理会使得投资者急于卖出而产生过度反应，未来很可能会补涨。

\subsection{参考文献}
\begin{enumerate}
    \item 曹春晓. 2022. 波动率的波动率与投资者模糊性厌恶. 方正证券.
    \item Kostopoulos, D., Meyer, S., \& Uhr, C. (2021). \textit{Ambiguity about volatility and investor behavior}. Journal of Financial Economics.
\end{enumerate}

\section{上行跳跃波动率(高频因子-波动跳跃类)}

\subsection{因子定义}
\begin{equation}
    \mathrm{RVJP}_t = \max\left(\mathrm{RS}_t^+ - \frac{1}{2} \hat{IV}_t, 0\right)
\end{equation}
\begin{equation}
    \mathrm{RS}_t^+ = \sum_{i=1}^{n} r_{t,i}^2 \mathrm{I}_{r_{t,i} > 0} \rightarrow_u \frac{1}{2} \int_0^t \sigma_s^2 ds + \sum_{s \leq t} \Delta X_s^2 \mathrm{I}_{\Delta X_s > 0}
\end{equation}
\begin{equation}
    \hat{IV}_t = \mu_{\frac{3}{2}}^{-3} \sum_{i=k}^{n} |r_{t,i}|^{\frac{2}{3}} |r_{t,i-1}|^{\frac{2}{3}} \cdots |r_{t,i-k+1}|^{\frac{2}{3}}
\end{equation}

\subsection{变量定义}
\begin{itemize}
    \item ① \( r_{t,i} \) 为某股票交易日 \( t \) 内的分钟对数收益率序列，如 5 分钟对数收益率序列；
    \item ② \( k=3 \)，\( n \) 表示该股票在交易日 \( t \) 中特定频率的分钟级别数据个数，如 5 分钟频率通常有 48 个数据点；
    \item ③ \( \mu_q \) 为标准正态分布的第 \( q \) 阶绝对矩；
    \item ④ 上行跳跃波动率 \( \mathrm{RVJP}_t \) 整体取值为正，所以需要用 \( \max \) 对负差值做修正。
\end{itemize}

\subsection{说明}
按收益涨跌情况，股价波动可分为上行波动和下行波动。跳跃波动也可据此做划分，上行跳跃波动就是上行波动中的跳跃部分；短期内，上行跳跃波动通常有负的风险溢价。

\subsection{参考文献}
\begin{enumerate}
    \item Yu, B., Mizrach, B., \& Swanson, N. R. (2020). \textit{New Evidence of the Marginal Predictive Content of Small and Large Jumps in the Cross-Section}. Econometrics, MDPI, 8(2), 1–52.
\end{enumerate}

\section{前景价值 TK(行为金融因子-前景理论)}

\subsection{因子定义}

① 以个股过去 \( N \) 日收益率序列作为未来收益分布的预测，其中负收益率从小到大为 \( r_{-m}, \ldots, r_{-1} \)，正收益率从小到大为 \( r_1, \ldots, r_n \)，每个收益出现的概率均为 \( \frac{1}{N} = m + n + 1 \)：
\begin{equation}
    (r_{-m}, \frac{1}{N}; \ldots; r_{-1}, \frac{1}{N}; r_0, \frac{1}{N}; r_1, \frac{1}{N}; \ldots; r_m, \frac{1}{N})
\end{equation}

② 利用价值函数、权重函数计算前景价值 \( \mathrm{TK} \)：
\begin{equation}
    \mathrm{TK} = \sum_{i=-m}^{-1} v(r_i)\left[ w^-\left(\frac{i + m + 1}{N}\right) - w^-\left(\frac{i + m}{N}\right) \right] + \sum_{i=1}^{n} v(r_i)\left[ w^+\left(\frac{n - i + 1}{N}\right) - w^+\left(\frac{n - i}{N}\right) \right]
\end{equation}

③ 价值函数、权重函数如下所示，函数参数取值为 \( \alpha = 0.88, \lambda = 2.25; \gamma = 0.61, \delta = 0.69 \)：
\begin{equation}
    v(x) =
    \begin{cases} 
        x^\alpha, & x \geq 0 \\ 
        -\lambda(-x)^\alpha, & x < 0 
    \end{cases}
\end{equation}
\begin{equation}
    \pi_l =
    \begin{cases} 
        w^+(p_i + \cdots + p_n) - w^+(p_{i+1} + \cdots + p_n), & 0 \leq l \leq n \\ 
        w^-(p_m + \cdots + p_n) - w^-(p_{i+1} + \cdots + p_n), & 0 \leq i \leq n \\ 
        w^-(p_m + \cdots + p_i) - w^-(p_m + \cdots + p_{i-1}), & -m \leq i < 0 
    \end{cases}
\end{equation}
\begin{equation}
    w^+(P) = \frac{P^\gamma}{(P^\gamma + (1 - P)^\gamma)^{\frac{1}{\gamma}}}
\end{equation}
\begin{equation}
    w^-(P) = \frac{P^\delta}{(P^\delta + (1 - P)^\delta)^{\frac{1}{\delta}}}
\end{equation}

\subsection{说明}
前景价值代表了投资者给每支个股的“打分”，基于前景理论，由于价值量 \( \mathrm{TK} \) 高的股票更具吸引力，投资者总会去买那些价值量 \( \mathrm{TK} \) 高的股票（选项），导致价值高的股票被超买，而价值低的股票由于没人关注而被低估，股票未来的预期收益率和当前的前景价值 \( \mathrm{TK} \) 呈负相关。

\subsection{参考文献}
\begin{enumerate}
    \item Barberis, N., Mukherjee, A., \& Wang, B. (2016). \textit{Prospect theory and stock returns: An empirical test}. Review of Financial Studies, 29(11), 3068–3107.
\end{enumerate}

\section{日内收益率(高频因子-动量反转类)}

\subsection{因子定义}
\begin{equation}
    \mathrm{intraday\_ret}_{t,i} = \frac{\mathrm{close}_{t,i}}{\mathrm{open}_{t,i}} - 1
\end{equation}

\subsection{变量定义}
\begin{itemize}
    \item ① \( \mathrm{open}_{t,i} \) 和 \( \mathrm{close}_{t,i} \) 分别为 \( t \) 日开盘价和 \( t \) 日收盘价；
    \item ② 若要计算周度、月度等频率的日内收益率，只需对区间内的日度日内涨跌幅求和。
\end{itemize}

\subsection{说明}
传统日度收益率可以划分为隔夜收益和日内收益，其中日内收益通常有一定程度的反转效应（对应日度收益率的反转）；隔夜收益与日内收益两者存在显著的负相关关系，对未来收益的预测方向也往往相反，许多研究认为正是双方的反转关系削弱了总收益的表现。

\subsection{参考文献}
\begin{enumerate}
    \item Branch, B., \& Ma, A. (2012). \textit{Overnight Return, the Invisible Hand Behind Intraday Returns}. Journal of Financial Markets, 2, 90-100.
    \item Lou, D., Polk, C., \& Skouras, S. (2019). \textit{A tug of war: Overnight versus intraday expected returns}. Journal of Financial Economics, 134(1), 192-213.
\end{enumerate}

\section{中间价变化率偏度(高频因子-流动性类)}

\subsection{因子定义}
高频因子-流动性类

\begin{equation}
    \mathrm{MPCskew}_{d,k} = \frac{1}{N_d - k - 1} \sum_{t=k+1}^{N_d} \left( \frac{\mathrm{MPC}_{d,t,k} - \overline{\mathrm{MPC}}_{d,t,k}}{\sigma_{d,t,k}} \right)^3
\end{equation}
\begin{equation}
    \mathrm{MPC}_{d,t,k} = \frac{M_{d,t} - M_{d,t-k}}{M_{d,t-k}}, \quad M_{d,t} = \frac{P_{d,t}^B + P_{d,t}^A}{2}
\end{equation}

\subsection{变量定义}
\begin{itemize}
    \item ① \( \mathrm{MPCskew}_{d,k} \) 为 \( d \) 日 \( \mathrm{MPC}_{t,k} \) 序列的偏度；
    \item ② \( M_{d,t} \) 为 \( d \) 日 \( t \) 时刻的市场中间价；
    \item ③ \( N_d \) 为 \( d \) 日的交易分钟数据；
    \item ④ \( \overline{\mathrm{MPC}}_{d,t,k} \) 和 \( \sigma_{d,t,k} \) 分别为 \( \mathrm{MPC}_{d,t,k} \) 序列的均值和标准差。
\end{itemize}

\subsection{说明}
MPCskew 描述了中间价极端变动偏离平均值的幅度，可用于捕捉主力市场操纵行为；长期看市场中间价的极端变动与股票收益显著负相关。

\subsection{参考文献}
\begin{enumerate}
    \item 丁鲁明, 陈升锐. 2021. 高频订单失衡及价差因子. 中信建投证券.
\end{enumerate}

\section{跳跃关联非正跳跃相对动量(图谱网络-动量溢出)}

\subsection{因子定义}

① 利用股票日内 5 分钟行情数据，计算跳跃统计量 \( \mathrm{JS} \)，识别出股价发生跳跃的分钟时刻并计算当天的负跳跃收益 \( \mathrm{NegJump\_Ret}_{j,t} \) 和非跳跃收益 \( \mathrm{NoJump\_Ret}_{j,t} \)，具体计算细节见相关日历页：
\begin{equation}
    \mathrm{JS} = N \frac{\hat{V}_{(0,1)}}{\sqrt{\hat{\Omega}_{\mathrm{swv}}}} \left( 1 - \frac{\mathrm{RV}_N}{\mathrm{SwV}_N} \right) \sim N(0,1)
\end{equation}

② 计算关联跳跃次数：如果股票 \( i \) 在 \( t \) 日发生跳跃，股票 \( j \) 在 \( t-1 \) 日至 \( t+1 \) 日内也发生了一次同向跳跃，则记为一次关联跳跃；

③ 计算两只股票间的跳跃关联度，即在股票 \( i \) 所有的股价跳跃 \( \mathrm{jump\_num}_i \) 中关联跳跃 \( \mathrm{cojump\_num}_{i,j} \) 所占的比例：
\begin{equation}
    \mathrm{Corr}_{i,j} = \frac{\mathrm{cojump\_num}_{i,j}}{\mathrm{jump\_num}_i}
\end{equation}

④ 在 \( t \) 时刻，以跳跃关联度为权重，分别加权股票 \( i \) 的那些关联股票的过去 20 天的负跳跃收益和非跳跃收益（剔除关联度最低的 50\% 的关联关系），并对过去 20 天收益率做回归取残差，将两个残差相加即为跳跃关联非正跳跃相对动量：
\begin{equation}
    \mathrm{Peer\_Negjump\_Ret}_{i,t} = \frac{\sum_{j=1}^{N} \mathrm{Corr}_{i,j} \cdot \mathrm{NegJump\_Ret}_{j,t}}{\sum_{j=1}^{N} \mathrm{Corr}_{i,j}}, \quad i \neq j
\end{equation}
\begin{equation}
    \mathrm{Peer\_Nojump\_Ret}_{i,t} = \frac{\sum_{j=1}^{N} \mathrm{Corr}_{i,j} \cdot \mathrm{NoJump\_Ret}_{j,t}}{\sum_{j=1}^{N} \mathrm{Corr}_{i,j}}, \quad i \neq j
\end{equation}
\begin{equation}
    \mathrm{Ret}_{i,t} = \alpha_1 + \beta_1 \cdot \mathrm{Peer\_Negjump\_Ret}_{i,t} + \varepsilon_{i,t}^{(1)}
\end{equation}
\begin{equation}
    \mathrm{Ret}_{i,t} = \alpha_2 + \beta_2 \cdot \mathrm{Peer\_Nojump\_Ret}_{i,t} + \varepsilon_{i,t}^{(2)}
\end{equation}
\begin{equation}
    \mathrm{NonPositiveJumpRelativeMomentum}_{i,t} = \varepsilon_{i,t}^{(1)} + \varepsilon_{i,t}^{(2)}
\end{equation}

\subsection{说明}
具有相似属性的公司往往会受到类似的信息冲击，进而导致股价的同步跳跃。公司间股价跳跃的同步程度也可用于识别公司间的潜在关联程度；以股价跳跃关联度为权重加权关联股票的收益率可有效捕捉股票间的“领先-滞后”联动效应。由于个股正向跳跃收益与自身未来收益显著负相关，在加权关联公司的跳跃收益时剔除掉正跳跃收益部分（该部分收益会抵消关联股票的动量溢出效应），有助于提升因子表现。

\subsection{参考文献}
\begin{enumerate}
    \item 任晓, 刘凯, 杨航. 2025. 基于股价跳跃关联性的选股策略. 招商证券.
    \item 任晓, 杨航. 2024. 如何识别股价跳跃. 招商证券.
\end{enumerate}

\section{卖出反弹偏离因子(高频因子-收益分布类)}

\subsection{因子定义}
\begin{equation}
    \mathrm{LCP} = \frac{\sum_{i=1}^{N} \overline{\mathrm{price}}_{\mathrm{sell},i} \cdot \mathbf{I}_{\{P_{\mathrm{sell},i} < \mathrm{close}\}}}{close} - 1
\end{equation}

\subsection{变量定义}
\begin{itemize}
    \item ① \( \mathrm{price}_{\mathrm{sell},i} \) 为逐笔成交数据中卖单的成交价，\( \mathrm{close} \) 为当日收盘价；
    \item ② 可以对该因子进行小单改进和尾盘改进，即只统计小单的占比和只统计 \([14:30,14:57)\) 尾盘区间的占比。
\end{itemize}

\subsection{说明}
卖出反弹偏离因子统计了低于收盘价卖出的价格相较于收盘价的偏离幅度，与未来收益正相关；因子取值越负，卖出后价格反弹越大，根据遗憾规避理论，为了避免产生遗憾和后悔情绪，投资者卖出股票后不愿再次买回的意愿较高，未来买入动力较低，有更低的预期收益；遗憾规避理论认为非理性的投资者在做决策时，会倾向于避免产生后悔情绪并追求自豪感，避免承认之前的决策失误。

\subsection{参考文献}
\begin{enumerate}
    \item 高智威. 2023. 基于逐笔成交数据的遗憾规避因子. 国金证券.
\end{enumerate}

\section{量峰间隔峰度(高频因子-成交分布类)}

\subsection{因子定义}

① 对过去 20 日同时点成交量按照 1 倍标准差的标准划分，高于 1 倍标准差划分为“喷发成交量”，低于 1 倍标准差划分为“温和成交量”；

② 对“喷发成交量”进一步按照连续与否划分：若当前分钟为“喷发成交量”，且前后 1 分钟均为“温和成交量”，则划为“孤立喷发成交量”，称为“量峰”；若当前分钟为“喷发成交量”，且前后 1 分钟也存在“喷发成交量”，则划为“连续喷发成交量”，称为“量岭”；将“温和成交量”称为“量谷”；

③ 对过去 20 日同日前后两个量峰之间的时间间隔分布，计算其峰度，即为量峰间隔峰度因子。

\subsection{变量定义}
\begin{itemize}
    \item 此外，还可统计两个量峰之间的时间间隔分布的标准差、偏度等指标，即可得到量峰间隔标准差因子、量峰间隔偏度因子。
\end{itemize}

\subsection{说明}
量峰间隔峰度因子衡量了知情交易者参与交易的时间间隔的分布情况，知情交易者通常是在特定时间点集中交易，交易频率不定，而不是持续不断地买卖。所以时间间隔分布更倾向于尖峰厚尾的特点，与未来收益正相关。（报告中测试发现，量峰间隔偏度因子也为正向因子，但量峰间隔标准差因子为负向因子。）

\subsection{参考文献}
\begin{enumerate}
    \item 魏建榕, 王志豪. 2025. 高频成交量的峰、岭、谷信息. 开源证券.
\end{enumerate}

\section{筹码分布相对价位(行为金融因子-筹码分布)}

\subsection{因子定义}

\begin{equation}
    \mathrm{CKDP}_t = \frac{\mathrm{close}_t - \mathrm{mean}_t}{\mathrm{Chip\_Highest}_t - \mathrm{Chip\_Lowest}_t}
\end{equation}

\subsection{变量定义}
\begin{itemize}
    \item ① 假设某只股票在 \( t \) 日的筹码分布共有 \( n \) 个价位（\( i=1,2,\ldots,n \)），\( \mathrm{vol}_{i,t} \) 为第 \( t \) 日 \( i \) 价位上的筹码峰高度，\( \mathrm{price}_{i,t} \) 为该筹码峰对应的价格，\( \mathrm{Chip\_Highest}_t \) 对应筹码最高价，\( \mathrm{Chip\_Lowest}_t \) 对应筹码最低价；
    \item ② \( \mathrm{mean}_t = \sum (\mathrm{price}_{i,t} \times p_{i,t}) \) 为筹码分布加权平均成本；
    \item ③ \( \mathrm{close}_t \) 为 \( t \) 日收盘价。
\end{itemize}

\subsection{说明}
筹码分布相对价位衡量了当前价格相对平均成本的偏离程度，可用于衡量投资者盈利或亏损程度。因子值小于零，说明多数投资者亏损；因子值大于零，说明多数投资者盈利；因子值接近零，说明股价接近平均成本，多空博弈相对激烈。经测试，筹码分布成本宽带与未来收益负相关。

\subsection{参考文献}
\begin{enumerate}
    \item 陈升锐, 姚紫薇. 2025. 筹码分布因子系统构建. 中信建投.
\end{enumerate}

\section{量谷加权价格分位点(量谷价券价格分位点)}

\subsection{因子定义}

① 对过去 20 日同时点成交量按照 1 倍标准差的标准划分，高于 1 倍标准差划为“喷发成交量”，低于 1 倍标准差划为“温和成交量”；

② 对“喷发成交量”进一步按照连续与否划分：若当前分钟为“喷发成交量”，且前后 1 分钟均为“温和成交量”，则划为“孤立喷发成交量”，称为“量峰”；若当前分钟为“喷发成交量”，且前后 1 分钟也存在“喷发成交量”，则划为“连续喷发成交量”，称为“量岭”；将“温和成交量”称为“量谷”；

③ 将“日内最高价、日内最低价、昨日收盘价”三者的最高值与最低值作为[昨收,今收]区间的价格高低位，并计算每日量谷成交量加权价格的相对分位点，20 日的分位点均值即为量谷加权价格分位点因子；
\begin{equation}
    \mathrm{ValleyPriceQuantile}_t = \frac{1}{20} \sum_{d=1}^{20} \frac{\mathrm{VWAP}_{\mathrm{valley},d} - \min(\mathrm{Low}_d, \mathrm{PrevClose}_d)}{\max(\mathrm{High}_d, \mathrm{PrevClose}_d) - \min(\mathrm{Low}_d, \mathrm{PrevClose}_d)}
\end{equation}

④ 需对量谷加权价格分位点因子做 20 日反转因子的中性化处理。

\subsection{说明}
量谷加权价格分位点因子衡量了交易低迷时点的价格水平，在不太可能存在股价过度反应的情况下，股价越高，且可持续，未来收益也可能越高。

\subsection{参考文献}
\begin{enumerate}
    \item 魏建榕, 王志豪. 2025. 高频成交量的峰、岭、谷信息. 开源证券.
\end{enumerate}

\section{主买强度(高频因子-资金流类)}

\subsection{因子定义}

\begin{equation}
    \frac{1}{T} \sum_{n=t}^{t-T+1} \frac{\mathrm{mean}(\text{主动买入成交额}_{i,j,n})}{\mathrm{std}(\text{主动买入成交额}_{i,j,n})}
\end{equation}

\subsection{变量定义}
\begin{itemize}
    \item ① 基于逐笔成交数据中的 BS 标志将逐笔成交数据合成为分钟级别的主动买卖金额，B 为主动买入（卖出方先挂单，买入方主动触碰卖单并成交）；S 为主动卖出（买入方先挂单，卖出方主动触碰买单并成交）；
    \item ② 剔除了处于涨跌停分钟上的主动买入金额和主动卖出金额数据；
    \item ③ \( i, j, n \) 分别表示第 \( i \) 只股票在第 \( n \) 个交易日内的第 \( j \) 分钟的成交额；
    \item ④ 月度选股下，\( T = 20 \) 个交易日；周度选股下，\( T = 5 \) 个交易日。
\end{itemize}

\subsection{说明}
同时刻画了投资者主动买入行为的强度和稳健性，但日内主买强度与未来收益负相关（使用日度数据计算的日间主买强度与未来收益正相关），该因子与市值存在较强的相关性。

\subsection{参考文献}
\begin{enumerate}
    \item 冯佳睿, 袁林青. 2020. 基于主动买入行为的选股因子. 海通证券.
\end{enumerate}

\section{平均换手率(行为金融因子-投资者注意力)}

\subsection{因子定义}
\begin{equation}
    \mathrm{TURNAVG} = \frac{1}{m} \sum_{t=1}^{t-m+1} \mathrm{TURN}_{i,t}
\end{equation}

\subsection{变量定义}
\begin{itemize}
    \item ① \( \mathrm{TURN}_{i,t} \) 为股票 \( i \) 在交易日 \( t \) 的换手率；
    \item ② \( m \) 为交易日窗口，一般取为月度。
\end{itemize}

\subsection{说明}
平均换手率属于股票交易活跃度类因子，通过衡量股票交易活跃度来反映投资者关注度，股票交易越活跃，投资者对该股票的关注度越高；而投资者有限注意力会推动关注度高的股票更易出现非理性净买入溢价，进而导致股价后续反转。

\subsection{参考文献}
\begin{enumerate}
    \item 陈升锐. 2024. 投资者有限关注及注意力捕捉与溢出. 中信建设.
\end{enumerate}

\section{平均换手日内正跳跌收益(高频因子-波动跳跃类)}

\subsection{因子定义}

1. 利用股票日内 5 分钟行情数据，计算 Jiang and Oomen (2008) 中的跳跃统计量 \( \mathrm{JS} \)，识别出股价发生跳跃的分钟时刻，具体识别细节见研报：
\begin{equation}
    \mathrm{JS} = N \frac{\hat{V}_{(0,1)}}{\sqrt{\hat{\Omega}_{\mathrm{SWV}}}} \left( 1 - \frac{\mathrm{RV}_N}{\mathrm{SwV}_N} \right) \sim N(0,1)
\end{equation}
\begin{equation}
    \mathrm{RV}_N = \sum_{k=1}^{N} r_k^2
\end{equation}
\begin{equation}
    \mathrm{SwV}_N = 2 \sum_{k=1}^{N} (R_k - r_k)
\end{equation}
\begin{equation}
    \hat{V}_{(0,1)} = \frac{1}{\mu_1^2} \sum_{k=1}^{N-1} |r_{k+1}| |r_k|
\end{equation}
\begin{equation}
    \Omega_{\mathrm{SWV}} = \frac{\mu_6}{9} \cdot \frac{N^3 \mu_1^{-6}}{N - 5} \sum_{k=0}^{N-6} \prod_{l=1}^{6} |r_{k+l}|
\end{equation}
\begin{equation}
    \mu_p = 2^{p/2} \Gamma \left( \frac{p+1}{2} \right) / \sqrt{\pi}
\end{equation}

2. 将所有发生价格跳跃且收益为正的时段的对数收益率累加即为当日的正向跳跃收益；

3. 以 10 个交易日为半衰期，将每日的日内正向跳跃收益乘上当日的换手率，再取过去 20 个交易日的指数加权平均收益即为最终的平均换手日内正跳跌收益因子：
\begin{equation}
    \mathrm{AvgTurnoverPosJumpRet}_{i,t} = \frac{\sum_{d=1}^{20} w_d \cdot (\mathrm{PosJumpRet}_{i,t-d+1} \cdot \mathrm{Turnover}_{i,t-d+1})}{\sum_{d=1}^{20} w_d}
\end{equation}
\begin{equation}
    w_d = 0.5^{\frac{d-1}{10}}
\end{equation}

\subsection{说明}
股价跳跃体现了投资者对信息冲击的集中反应；股价跳跃收益类因子通常呈反转特性，与股票未来收益负相关，即过去跳跃收益较高的股票未来收益较小。

\subsection{参考文献}
\begin{enumerate}
    \item 任建稳, 杨帆. 2024. 如何识别股价跳跃. 招商证券.
    \item Jiang, G. J., \& Oomen, R. C. A. (2008). \textit{Testing for jumps when asset prices are observed with noise—a “swap variance” approach}. Journal of Econometrics, 144(2), 352-370.
    \item Jiang, G. J., \& Zhu, K. X. (2017). \textit{Information Shocks and Short-Term Market Underreaction}. Journal of Financial Economics, 124(1), 43-64.
\end{enumerate}

\section{APM 因子(高频因子-动量反转类)}

\subsection{因子定义}

对选定股票，回溯取其过去20日数据，记录上午的股票收益率为 \( r_t^{\mathrm{am}} \)，指数收益率为 \( R_t^{\mathrm{am}} \)；逐日下午的股票收益率为 \( r_t^{\mathrm{pm}} \)，指数收益率为 \( R_t^{\mathrm{pm}} \)。

将得到的40组上午与下午 \((r, R)\) 的收益率数据进行回归：
\begin{equation}
    r_i = \alpha + \beta R_i + \varepsilon_i
\end{equation}
得到残差项 \( \varepsilon_i \)；

以上得到的40个残差 \( \varepsilon_i \)中，上午残差记为 \( \varepsilon_i^{\mathrm{am}} \)，下午残差记为 \( \varepsilon_i^{\mathrm{pm}} \)，进一步计算每日上午与下午残差的差：
\begin{equation}
    \delta_t = \varepsilon_t^{\mathrm{am}} - \varepsilon_t^{\mathrm{pm}}
\end{equation}

构造统计量 \( \mathrm{stat} \) 来衡量上午与下午残差的差异程度，计算公式如下（\( \mu \) 为均值，\( \sigma \) 为标准差）：
\begin{equation}
    \mathrm{stat} = \frac{\mu(\delta_t)}{\sigma(\delta_t) / \sqrt{N}}
\end{equation}

为了消除动量因子影响，将统计量 \( \mathrm{stat} \) 对动量因子进行横截面回归：
\begin{equation}
    \mathrm{stat}_j = b \cdot \mathrm{Ret20}_j + \varepsilon_j
\end{equation}
其中 \( \mathrm{Ret20} \) 为股票过去20日的收益率，代表动量因子；回归得到的残差值 \( \varepsilon \) 即为APM因子。

\subsection{说明}
基于“知情交易者更加倾向于在每日上午进行交易，上午的价格行为蕴藏了更多可用于选股的信息量”的想法，构建了刻画了股票上午与下午的价格行为差异程度的APM因子；改进的APM因子是将个股和指数的上午收益 \( r_t^{\mathrm{am}} \) 替换成了隔夜收益 \( r_t^{\mathrm{overnight}} \)。

\subsection{参考文献}
\begin{enumerate}
    \item 魏建榕. 2020. APM因子模型的进阶版. 开源证券.
\end{enumerate}

\section{新闻网络节点度(图谱网络-网络结构)}

\subsection{因子定义}
① 以标题中股票为 lead，以同时出现在新闻正文中的其他股票为 follower，确定领边 \( l_{ij,T} \)，并构建股票共现图的邻接矩阵 \( W_T \)：
\begin{equation}
    l_{ij,T} = \bigcup_{m_d \in D_T} \{(i,j)_{m_d} \mid j \text{ in } m_d \text{ title, } i \text{ in } m_d \text{ headline, } i \neq j\}
\end{equation}
\begin{equation}
    W_T = [\omega_{ij,T}]_{n \times n} = \left[ \frac{l_{ij,T}}{\sum_{j=1}^n l_{ij,T}} \right]_{n \times n}
\end{equation}

② 基于股票共现图邻接矩阵 \( W_T \) 计算节点度 degree：
\begin{equation}
    A_{i,t}(\omega_{ij,t-l:T}) = \# \{j : \omega_{ij,t-l:T} \neq 0\} + \# \{j : \omega_{ji,t-l:T} \neq 0\}, \quad j = 1, 2, \cdots, n
\end{equation}

\subsection{说明}
新闻节点度为新闻共现网络中与节点相连的邻居节点的数量，节点度较高，说明该公司在新闻报道中频繁与其他多个公司共同提及，有较高的媒体关注度；媒体在报道时更愿意报道积极的、乐观的消息，有较高关注度的公司表现会优于那些受到较少媒体关注的公司。

\subsection{参考文献}
\begin{enumerate}
    \item Hu, J., \& Härdle, W. (2021). \textit{Networks of news and cross-sectional returns}. IRTG 1792 Discussion Papers 2021-023, Humboldt University of Berlin, International Research Training Group 1792 “High Dimensional Nonstationary Time Series”.
\end{enumerate}

\section{时间加权平均的股票相对价格位置(高频因子-收益分布类)}

\subsection{因子定义}
计算股票相对价格位置，即股票当期价格相对区间最高最低价的分数：
\begin{equation}
    \mathrm{RPP}_{i,t} = \frac{P_{i,t} - L_i}{H_i - L_i}
\end{equation}

将股票相对价格位置对时间进行积分，得到时间加权平均的相对价格位置 \( \mathrm{ARPP} \)：
\begin{equation}
    \mathrm{ARPP}_i = \int_0^T \mathrm{RPP}_{i,t} \, dt \quad \Rightarrow \quad \mathrm{ARPP}_{i,t} = \frac{\int_0^T P_{i,t} \, dt - L_i}{H_i - L_i}
\end{equation}

\subsection{变量定义}
\begin{itemize}
    \item ① \( \mathrm{RPP}_{i,t} \) 表示股票 \( i \) 在某一时间区间内的时刻 \( t \) 的相对价格位置；
    \item ② \( H_i \) 和 \( L_i \) 为该时间区间内股票 \( i \) 的最高价和最低价；
    \item ③ \( \displaystyle \int_0^T \mathrm{RPP}_{i,t} \, dt \) 表示股票 \( i \) 在时间区间 \( 0 \) 到 \( T \) 内的时间加权平均价格（即 TWAP），可以用分钟线高开低收均值在时序上的均值作为区间的 TWAP。
\end{itemize}

\subsection{说明}
衡量股票在价格相对高位停留的时间长短，股票在价格相对高位停留的时间越长，因子取值越大。

\subsection{参考文献}
\begin{enumerate}
    \item 朱剑涛.2020.基于时间尺度度量的日内买卖压力.东方证券
\end{enumerate}

\section{反向交易系数修正后的积极灵活筹码(行为金融因子-筹码分布)}

\subsection{因子定义}
① 计算筹码买入价格相对当日开盘价的涨跌幅作为该筹码的反向交易系数：
\begin{equation}
    \alpha_i = \frac{\mathrm{TradePrice}_i}{\mathrm{Open}_t}
\end{equation}

② 将反向修正系数与筹码买入量相乘即可得到修正后的积极灵活筹码值：
\begin{equation}
    \mathrm{AFH\_Open} = \frac{\sum_{i=1}^{N} \alpha_i \times \mathrm{Vol\_AFH}_i}{\sum_{t=T-20}^{T} \mathrm{Volume}_t}
\end{equation}

\subsection{变量定义}
\begin{itemize}
    \item ① \( i \) 为筹码购入时的订单号，统计对象为过去 20 日的订单号；
    \item ② \( \mathrm{Vol\_AFH}_i \) 为主动买入时订单金额小于 5 万元的订单成交量；
    \item ③ \( \mathrm{TradePrice}_i \) 为该订单的买入价格；
    \item ④ \( \mathrm{Open}_t \)、\( \mathrm{Volume}_t \) 分别为订单买入当日的开盘价和当日总成交量。
\end{itemize}

\subsection{说明}
处置效应下的“卖盈持亏”行为是一种反向交易行为，而在日内下跌过程中买入筹码的投资者会更具有反向交易的思维，此时的买入筹码也更具处置效应，因此可以计算筹码买入价格相对当日开盘价的涨跌幅作为该筹码的反向交易系数，用于修正原始的积极灵活筹码，提升因子表现。

\subsection{参考文献}
\begin{enumerate}
    \item 任瞳, 许继宏. 2024. 基于逐笔数据的 AFH 改进因子. 招商证券.
\end{enumerate}

\section{下跌时刻笔均成交金额(高频因子-成交分布类)}

\subsection{因子定义}

① 将日内 240 个 1 分钟数据集合标记为 \( T = \{t_1, t_2, \dots, t_{240}\} \)，判断股票 \( i \) 的每分钟涨跌幅，将下跌时刻标记为集合 \( P = \{t_{i1}, t_{i2}, \dots, t_{ik}\} \)；

② 对下跌时刻进行汇总，根据其成交额 \( \mathrm{Amt}_t \) 之和除以成交笔数 \( \mathrm{DealNum}_t \) 之和，构造下跌时刻的笔均成交金额 \( \mathrm{ATD}_{\mathrm{Down}} \)：
\begin{equation}
    \mathrm{ATD}_{\mathrm{Down}} = \frac{\sum_{t \in P} \mathrm{Amt}_t}{\sum_{t \in P} \mathrm{DealNum}_t}
\end{equation}

③ 同理，将全天的成交金额除以全天的成交笔数，得到全天的笔均成交金额 \( \mathrm{ATD}_T \)：
\begin{equation}
    \mathrm{ATD}_T = \frac{\sum_{t \in T} \mathrm{Amt}_t}{\sum_{t \in T} \mathrm{DealNum}_t}
\end{equation}

④ 将下跌时刻的笔均成交金额除以全天的笔均成交金额，即得到去量纲化后的下跌时刻标准化笔均成交金额因子 \( \mathrm{SATD}_{\mathrm{Down}} \)：
\begin{equation}
    \mathrm{SATD}_{\mathrm{Down}} = \frac{\mathrm{ATD}_{\mathrm{Down}}}{\mathrm{ATD}_T}
\end{equation}

\subsection{变量定义}
\begin{itemize}
    \item ① 特殊时刻 \( P = \{t_{i1}, t_{i2}, \dots, t_{ik}\} \) 除了下跌时刻外，还可以是上涨时刻、横盘时刻等。由此可进一步得到上涨时刻标准化笔均成交金额因子、横盘时刻标准化笔均成交额因子等。
\end{itemize}

\subsection{说明}
笔均成交额类因子通过汇总“特殊的、具有较多信息含量”时刻的成交金额情况来刻画主力资金的行为动向；笔均成交金额越大，说明该特殊时间内资金优势越明显，属于主力资金的可能性越大。上述因子是从“股价涨跌”维度来寻找特殊时刻，测试发现股价“下跌时刻”的笔均成交金额越高，股价未来的表现越好，可能刻画了下跌时刻主力资金的“抄底行为”。

\subsection{参考文献}
\begin{enumerate}
    \item 张欣悦, 张宇. 2025. 日内特殊时刻蕴含的主力资金 Alpha 信息. 国信证券.
\end{enumerate}

\section{价格冲击偏差(高频因子-资金流类)}

\subsection{因子定义}

\begin{equation}
    \mathrm{return}_i = \gamma^{\mathrm{up}} \cdot I_i \cdot \mathrm{MF}_i + \gamma^{\mathrm{down}} \cdot (1 - I_i) \cdot \mathrm{MF}_i
\end{equation}
\begin{equation}
    \mathrm{MF}_i = \frac{\mathrm{MoneyFlow}_i}{\mathrm{Amount}_i}
\end{equation}
\begin{equation}
    \gamma_{\mathrm{bias}} = \frac{\gamma^{\mathrm{up}} - \gamma^{\mathrm{down}}}{(\gamma^{\mathrm{up}} + \gamma^{\mathrm{down}})/2}
\end{equation}

\subsection{变量定义}
\begin{itemize}
    \item ① \( \mathrm{return}_i \) 为第 \( i \) 个 5 分钟线的收益率；
    \item ② \( \mathrm{MoneyFlow}_i \)、\( \mathrm{Amount}_i \) 分别为第 \( i \) 个 5 分钟内的主动净流入金额和成交额；
    \item ③ \( \mathrm{MF}_i \) 为主动净买入占比；
    \item ④ \( I_i \) 为示性函数，当 \( \mathrm{MF}_i \) 大于 0 时取 1，否则取 0。
\end{itemize}

\subsection{说明}
价格冲击偏差主要表征了股票在某一时间区间上涨或者下跌的难易程度，价格冲击偏差较大（正值）的股票，相同比例的主动订单对其股价向下的冲击小于向上的冲击，股票容易上涨，反之，股票容易下跌。

\subsection{参考文献}
\begin{enumerate}
    \item 朱剑涛. 2016. 非对称价格冲击带来的超额收益. 东方证券.
\end{enumerate}

\section{基于日收益率的注意力因子}

\subsection{因子定义}
本文的因子定义包含原始因子和经过处理的因子两种形式。原始因子的核心计算如下：
\begin{equation}
\begin{aligned}
\text{因子原始值}_{i, t} &= \frac{1}{T} \sum_{\tau = t-T+1}^{t} (r_{i, \tau} - R_{m, \tau})^2
\end{aligned}
\end{equation}
\vspace{6pt}
其中，$T$为一个自然年的交易日总数。在此基础上，采用“当前自然年因子值减去历史因子值”的处理方法得到最终因子：
\begin{equation}
\begin{aligned}
\text{处理后的因子}_{i, t} &= \text{因子原始值}_{i, t}^{\text{(当前自然年)}} - \text{因子原始值}_{i}^{\text{(历史)}}
\end{aligned}
\end{equation}

\subsection{变量定义}
\begin{itemize}
    \item $r_{i, \tau}$ 表示在交易日 $\tau$ 基金 $i$ 的日收益率。
    \item $R_{m, \tau}$ 表示在交易日 $\tau$ 的市场基准收益率。文中明确指出：“市场收益率，笔者直接用了全市场所有同类型基金收益率的均值”。
    \item $T$ 表示计算窗口期的长度，文中明确为“一个自然年的数据”，即过去一个自然年所包含的交易日总数。
    \item $\text{因子原始值}_{i}^{\text{(历史)}}$ 表示基金 $i$ 在“当前自然年”之前（即历史数据）的因子原始值的长期均值。文中未给出其具体计算窗口，仅提及“历史均值”。
\end{itemize}

\subsection{说明}
这个因子，笔者是在研报当中看到的，然后在大A全市场所有股票中进行了第一次复现，发现效果很不错。于是，昨天在转债市场上用了一下这个因子。虽然，结果谈不上惊艳，但也还算可以。既然股票和转债上都尝试了这个因子了，而且它的计算只与收益率有关，那为啥不在基金上面也做一下尝试呢？

基于之前的研究，笔者发现选基因子用当前自然年的均值减去历史均值会有奇效（在《量化选基系列（十八）：我似乎找到了量化选基因子的正确打开方式！》这篇文章有过相关介绍），因此在实证分析的时候，也加入了通过这种方式处理之后的结果。

\subsection{参考文献}
本文为独立研究，文中引用的参考文献如下：
\begin{enumerate}
    \item 量化选基系列（十八）：我似乎找到了量化选基因子的正确打开方式！
    \item 量化选基系列（十四）：精心设计了一个均值回归因子，结果却变成了强者恒强！
    \item 量化选基系列（十五）：一个对均值回归走火入魔的博主，整了个新活！结果在FOF基金上效果炸裂！
\end{enumerate}










\end{document}